\documentclass[11pt,a4paper]{article}
\usepackage[T1]{fontenc}
\usepackage[utf8]{inputenc}
\usepackage{lmodern,amsmath,amssymb,amsthm,mathtools,microtype,mathrsfs}
\usepackage[margin=25mm,top=23mm,bottom=25mm]{geometry}
\usepackage{booktabs,array,longtable,enumitem,xcolor,fancyhdr}
\usepackage[colorlinks=true,linkcolor=blue!45!black,urlcolor=blue!45!black,citecolor=blue!45!black]{hyperref}
\hypersetup{pdftitle={Reading the Alpoge--Claude RH and Dirichlet-GRH proof: working version 0.42},pdfauthor={ChatGPT},pdfsubject={Partial human-readable reconstruction with explicit arithmetic inputs},pdfkeywords={Riemann hypothesis, Lean, theta lifts, prime sums, Arakelov slopes}}
\setlength{\parindent}{0pt}
\setlength{\parskip}{5pt plus 1pt}
\setlist{topsep=3pt,itemsep=3pt,leftmargin=*}
\setlength{\headheight}{14pt}
\pagestyle{fancy}\fancyhf{}
\fancyhead[L]{\small RH and Dirichlet GRH: a reconstruction}
\fancyhead[R]{\small Working version 0.42}
\fancyfoot[C]{\thepage}
\renewcommand{\headrulewidth}{0.3pt}
\numberwithin{equation}{section}
\newtheorem{theorem}{Theorem}[section]
\newtheorem{proposition}[theorem]{Proposition}
\newtheorem{lemma}[theorem]{Lemma}
\newtheorem{corollary}[theorem]{Corollary}
\theoremstyle{remark}\newtheorem{remark}[theorem]{Remark}
\newcommand{\Rea}{\operatorname{Re}}
\newcommand{\HS}{\mathrm{HS}}
\newcommand{\tr}{\operatorname{tr}}
\newcommand{\codim}{\operatorname{codim}}
\newcommand{\covol}{\operatorname{covol}}
\newcommand{\OK}{\mathcal O_K}
\newcommand{\src}[2]{\href{https://github.com/CocoSgt/riemann-hypothesis-lean/blob/7cbf7a0c32fca1badf05931879a37c3f8a2b03b4/src/#1}{\texttt{#2}}}
\newcommand{\status}[1]{\textbf{Status:} #1}
\begin{document}
\begin{center}
{\LARGE\bfseries RH and Dirichlet GRH: new reconstruction chapters}\\[7pt]
{\large Version 0.42 --- 23 November 2026}\\[5pt]
{\large Ordinary structure, rigidity, and quantitative reductions}
\end{center}
\textbf{Scope.} This reader's edition contains the current audit and
the new chapters only. It is a partial reconstruction, not a completed
human proof of RH or GRH. References to earlier named claims refer to
the complete cumulative manuscript and the preserved component files
in the accompanying archive. The underlying pinned formal proof is
accepted as correct. The exact current boundary is stated below.
\setcounter{tocdepth}{1}
\tableofcontents
\clearpage
\begingroup\emergencystretch=3em
\section{Current dependency audit: ordinary structure, rigidity, and quantitative reductions}
\label{sec:dependency-audit-v42}

\subsection{Current status and reading order}

This version gives a complete conditional terminal argument and completes
the roadmap's finite-level, ordinary-structure and base-change milestone.
It also proves the original classical control and local Hecke-algebra
theorems on their stated arithmetic inputs. Both previously unresolved
theorem-matching obligations are now reconstructed: the full original
\texttt{Ad\allowbreak Vanish\allowbreak K} and the full original \texttt{Host\allowbreak Rigid\_\allowbreak s}.
The finite-field row calculations and the global quantitative reductions
give further proved inputs to the selected route.

The underlying Alpoge--Claude Lean proof is accepted as fully checked and
correct. This manuscript remains a \textbf{partial human reconstruction}:
the common arithmetic and geometric objects, their compatible arithmetic
specializations, and the five original global quantitative producers are
not all reconstructed. No Lean build or kernel replay is claimed.

There are 154 new and 2363 cumulative proof claims. The selected graph
retains 121 nodes and 161 logical edges. Its open register has 31
whole obligations, 22 of class C and 9 of class B, with no remaining
class-D matching obligation. Of these, 24 are shared by RH and
Dirichlet GRH and 7 are additional for Dirichlet GRH; none is RH-only.
The count is an inventory of scopes, not a percentage of completion.
Every predecessor claim, source array and dependency edge is preserved.

Read CMP1--CMP6 first for the complete conditional final deduction.
The ordinary entrance is FNL, OST, OCT and OHA; the two rigidity chapters
are ADK and HSR. RRT and GQT supply the finite row arguments and sufficient
quantitative replacements. The source pin and exact claim-to-source
records are in the accompanying provenance ledger. The separate file
\texttt{RH\_\allowbreak new\_\allowbreak proofs\_\allowbreak v0\_\allowbreak 42.\allowbreak tex} contains this audit and only the
current increment. Historical chapters below retain their original
hypotheses; this audit governs their older frontier descriptions.

\subsection{The completed ordinary entrance}

FNL1--FNL16 prove the actual finite adelic factorization in its required
order, and all eighteen original \texttt{HSet.\allowbreak Ax} fields. They verify the
local split sections, auxiliary congruence condition, Iwahori topology,
level joins, original Artin maps and inverse-Frobenius density on the
same objects used by the ordinary theory.

OST1--OST28 construct one factorial-limit projector onto the actual
ordinary module. They prove its exact value filtration, finite freeness,
invertibility of every positive-monoid operator, and base change by the
original coefficientwise map for every allowed local map of complete
Noetherian local rings. Neither nilpotents nor nonflat local maps are
excluded. The local matrix and product-group axioms are proved before
the theorem is installed on each original level.

OCT1--OCT19 prove the original control theorem by identifying an integral
ordinary classical lattice with the actual ordinary module. The proof
retains the original coefficient domain, injective embedding and all ten
\texttt{Wt\allowbreak Rep} fields. The zeroth-coordinate control map, normalized
positive operators, depth comparisons and finite away sums are checked
on the same functions. OHA1--OHA12 prove spherical commutativity and the
full original \texttt{hi\allowbreak Comm}/\texttt{hi\allowbreak Alg}: the literal residual part
is the appropriate idempotent summand of the finite commuting operator
algebra, its Hecke algebra is finite complete local with the prescribed
residue field, and the positive generators are units in that algebra.

The additional weight-representation chapter constructs the original algebraic representation's integral lattice and all ten \texttt{Wt\allowbreak Rep} assertion fields, with the specified complex embedding retained. It supplies the algebraic representation required by control; arithmetic specialization and separating fibers remain separate.\par SMP1--SMP10 prove the full original \texttt{Sp.\allowbreak abs\allowbreak Spec}: one map on all ambient functions, fixed before every level, depth, basis and operator. Finite modules of untwisted values prove both eigenvalue-integrality assertions. The full finite away-sum clause uses only its original fixed-space hypothesis, without requiring its coefficient-valued image to be ordinary.\par The arithmetic-specialization chapter installs this same map in the original \texttt{Spec\allowbreak Int\allowbreak Iff.\allowbreak spec}, with the actual coefficient point, finite-part character, prescribed complex embedding and all original \texttt{Is\allowbreak Spec} fields. Arithmetic nonvanishing and separation of fibers remain separate producers.

The full integral type and separating arithmetic fibers,
the four original Hida endpoints and all-level automorphy on one ideal
still have to be supplied. Both original rank-four Links, the fixed
complex embedding, residue characteristic three and one common
coefficient threshold remain in scope. The selected Raw implication
proved in the predecessor is reused, not reopened.

\subsection{Two whole matching obligations now closed}

ADK1--ADK6 prove the exact application of Thorne's adjoint
Bloch--Kato vanishing theorem. The comparison includes normalized
Satake parameters, the infinite-place algebraicity dictionary,
conjugate self-duality, finite coefficient descent, irreducibility on
the full infinite cyclotomic kernel, and the literal local finite
conditions. It uses the full matrix adjoint, including the scalar line
and negative-transpose complex-conjugation action. It holds throughout
the source contract, in every rank at least two and every allowed prime.

HSR1--HSR9 prove the original host-rigidity implication. The dyadic
conductor gives an actual separating local element. An allowed
dual-number deformation yields a Bloch--Kato cocycle by the written
ordinary crystalline argument and the actual local comparison maps.
The original automorphic representation is twisted by
$\epsilon^{-1}$, its pairing is the actual antidiagonal matrix $J$,
and its full adjoint is the representation to which ADK6 applies.
Surjectivity of the original quotient map gives one global primitive
with the required sign. No new host-existence or automorphy theorem is
inferred, and the distinct \texttt{RIG\_\allowbreak II} producer remains open.

\subsection{The conditional main theorem and quantitative progress}

CMP1--CMP6 give the final derivation from two precisely stated common
arithmetic input contracts. Exact character orthogonality, including
ramified prime readings, identifies the centered prime sum. Fourier
averaging precedes estimation. The arithmetic comparison gives the
one-sided square-root bound, and the fully written pole argument
excludes zeros to its right. Functional equations and the elementary
imprimitive Euler factors finish RH and unrestricted Dirichlet GRH.
The conclusion remains conditional until the two input contracts are
produced. CMP1, CMP3, CMP4 and CMP5 reorganize inherited proofs; they
are not counted as new arithmetic input theorems.

GQT1--GQT7 show that only three counting tests are needed for the
retained-row defect, and give polynomial error conditions sufficient
for the final spectral consumer. All common spaces, metrics, measures,
integral lattices and normalization factors remain fixed as required.
The exact original stronger spectral producer is not claimed.

RRT1--RRT16 prove positivity and normalization of the actual finite
weights, probability transfer, squarefree and divisor bounds, the
$7/8$ invariant-subset estimate, an eventual long-cycle estimate and
the full original root-to-twist collision estimate with its constants.
They also prove the finite tensor identity and the actual repeated-label
enumeration on supplied multiplicity identities. Four of the five
displayed terms of $W_g$ now have local producers. The remaining local
inputs are two exact character laws (with the original label matching)
and the root-injectivity collision estimate. The global row identities,
adelic inflation comparison and three arithmetic counts remain to be
constructed. Thus \texttt{GLOBAL\_\allowbreak RETAINED\_\allowbreak DEFECT} is still open.

\subsection{The remaining route}

The next large arithmetic endpoint is one same-ideal Hida/support
package on both original rank-four Links. A second endpoint is one
compatible geometric and local family carrying all required readings,
tests, point identities, lattices and base rows. Separate witnesses
for those properties do not suffice. The one-character, congruence,
host-counting and additional Dirichlet-GRH branches retain their original
quantifiers. The denominator, retained-defect, period/height, local-gain
and spectral estimates must all be installed on that common family.

Classes A, B, C and D retain the roadmap's meaning: checked published
application, independent elementary/moderate reconstruction, substantial
reconstruction, and unresolved source matching respectively.
Classification and completion are separate. The two D questions have
been resolved by proofs; this does not reclassify the remaining C tasks
as easy. All 245 historical compound targets remain traceable.

\subsection{Every remaining selected obligation}

The following contracts retain the exact selected outputs. Progress
inside a compound contract is recorded immediately after its statement.
Completed subtheorems are available inputs; a compound contract remains
open until all its actual input producers and compatibility conditions
have been reconstructed.
\paragraph{1. GLOBAL\_\allowbreak{}DENOMINATOR — C — Actual LinkN local table and binding producer}

Closed(Des $\mathcal C$ 3) LinkN: the actual etaSum is bounded with exact dZ0 and same-host vvTerm, subject only to the literal ShaFinOdd and summability guards.


\paragraph{2. GLOBAL\_\allowbreak{}RETAINED\_\allowbreak{}DEFECT — C — Original retained-row defect bound LinkIRR}

Rows.DelpNonneg and delp$\le$exp(-x\textasciicircum{}(ez/\allowbreak{}7)) uniformly for every late actual row datum.


\emph{Current subfrontier.} Two exact local character laws and normalized label comparison; root-injectivity collision estimate; original HidaSetting, reduction/inflation EtOne, RowsAdd/RowsGood, common AllRows object and three arithmetic cylinder counts. Four of five W\_\allowbreak g terms have local proofs.


\paragraph{3. GLOBAL\_\allowbreak{}PERIOD\_\allowbreak{}HEIGHT — C — Full actual period and height comparison LinkDIVt}

Closed conditional TuSize -> DGamSize -> (LinkEta -> DIVinN -> LinkDIV), with every original LinkDIV\_\allowbreak{}at clause.


\paragraph{4. GLOBAL\_\allowbreak{}LOCAL\_\allowbreak{}GAIN — C — Original global LinkG installation}

One uniform LinkG with actual local lattice lengths, served-orbit coverage and MeanKappaBound; delta(x)->0 before varying actual indices.


\paragraph{5. GLOBAL\_\allowbreak{}SPECTRAL — C — Original normalized kernel/\allowbreak{}Gram estimates LinkS}

Closed TuSize -> LinkS with exact Gamma normalization, off-diagonal Hilbert-Schmidt bound, diagonal error and Dall/\allowbreak{}Dstar estimate, errors exp(-x\textasciicircum{}(epsilon/\allowbreak{}8)).


\emph{Current subfrontier.} Actual normalized Hilbert–Schmidt and dimension estimates; sufficient weaker polynomial rate proved only for the selected final consumer.


\paragraph{6. TASK\_\allowbreak{}PROP\_\allowbreak{}GK\_\allowbreak{}INSTALL — B — Install OneChar and the congruence integer into PropGK}

Complete the finite orbit/\allowbreak{}Frobenius/\allowbreak{}valuation argument producing one nonzero integer before every later coefficient prime, with the original e,\allowbreak{} N,\allowbreak{} kernel,\allowbreak{} p and fixed-K constants.


\paragraph{7. TASK\_\allowbreak{}ONE\_\allowbreak{}CHAR\_\allowbreak{}PROOF — C — Reconstruct the original one-character argument}

From exact InpZ and AutomorphyUnramifiedCoefficients, prove the original OneChar trace/\allowbreak{}length estimate on the supplied coefficient datum and GD input.


\paragraph{8. TASK\_\allowbreak{}INP\_\allowbreak{}Z\_\allowbreak{}ASSEMBLY — B — Install the exact Kato input record}

Build the original CohInp/\allowbreak{}InpZ record using exactly lemM,\allowbreak{} lemMprime,\allowbreak{} lemP,\allowbreak{} RigII,\allowbreak{} AdVanishK and number on the same arithmetic objects. No rank-four replacement is permitted.


\paragraph{9. RIG\_\allowbreak{}II — C — Type-II polarized adjoint rigidity}

Marked original datum, actual type-II point and modularity of all source extensions imply vanishing for its exact adRep.


\paragraph{10. AUT\_\allowbreak{}UNRAM\_\allowbreak{}COEF — C — Actual unramified-coefficient automorphy}

For each actual Dat with coefficient ramification one and fixed iota, choose one marked g1 supplying FinA, AutB and LemW together.


\paragraph{11. CNR3 — C — Uniform rank-three and rank-two congruence integer}

Fixed quadratic K -> c,\allowbreak{} A -> e,\allowbreak{} N,\allowbreak{} p,\allowbreak{} w,\allowbreak{} finite Galois-stable roots-of-unity orbit -> nonzero integer NN -> every later coefficient prime, closure, embedding, Frobenius and finite admissible collection.


\paragraph{12. COUNT\_\allowbreak{}II\_\allowbreak{}SUPPORT — C — Two actual rank-four Links with common threshold and same-ideal Hida/\allowbreak{}P1}

For each of the TWO original QX rank-four Links, at fixed iota choose ONE common coefficient threshold before every later coefficient DVR and actual universal B,\allowbreak{} Lambda; produce one ideal J carrying ring,\allowbreak{} dense,\allowbreak{} pt,\allowbreak{} end1,\allowbreak{} end2,\allowbreak{} twist AND P1core. Inherited JSP/\allowbreak{}CIP/\allowbreak{}CAP/\allowbreak{}CCH/\allowbreak{}CFO supply the same-J consumer, every-map reciprocity/\allowbreak{}universal NiceI Norm, ArPl under unchanged CohP, supplied-type IsRes, and the full selected h=0,\allowbreak{} a=n=4,\allowbreak{} g=0 Choice with coupled B1/\allowbreak{}B2, D=4[L+:Q]+4T.q and B2 constant zero. The dualN bound retains one C before every m and an eventual N threshold uniform over each original admissible rep-prime under its finiteness premise, on actual TWPl. LIM constructs the full original S4. PRG now constructs the full original RRObj/\allowbreak{}RMaps, every finite-level map and actual ordered splitting target; PMS constructs the actual all-level PMod and full S5 with exact support, dimension, fiber and move fields, plus piKer/\allowbreak{}piSurj and the same-J annihilator equivalence. SCI constructs the actual quotient determinant, OoPair/\allowbreak{}ooMaps and point/\allowbreak{}OoPt on one splitting prime. The selected CRW full-block cotangent bound forces the support ideal to vanish, proving selected normalized Core. BCN supplies full Core=>CoreBC and its selected normalization/\allowbreak{}Unit consumer; PQL supplies full original TWQuot and matched pro-ell descent. Together SCI13 proves selected original Raw at h=0,\allowbreak{} a=n=4,\allowbreak{} k.g=0, with final B,\allowbreak{} J,\allowbreak{} K,\allowbreak{} f and every Selmer quantifier unchanged. NTL17 now proves full original neatU for arbitrary integral tau0 and every level. OTW12 proves one original theta before every K with complete-coefficient IsTheta/\allowbreak{}IsWt; OTW13 and CAP17 assemble full UnitaryHeckeData on supplied exact TypeForAlt. FNL, OST, OCT and OHA now supply the full original finite-level and HSet axioms, ordinary projector, freeness, base change, control on its weight representation, and local Hecke algebra. The algebraic weight representation is also constructed in the new weight chapter. The remaining arithmetic producers are the full original integral TypeForAlt; separating automorphic arithmetic fibers; the four Hida endpoints PointTheoremBaseObject, ArithmeticPointDecomposition, IrrArithPointsAutomorphicAway and TwistOk on this same ideal; and all-level same-ideal LevelAlt, TWAut and HasAut on both original Links. The selected Raw patching, normalization and descent are no longer open producers. Preserve full CohP, LocCFT, N4, exact determinant/\allowbreak{}pro-ell conditions, nontrivial residual cyclotomic character, ell=3, fixed iota and one common coefficient threshold. No arbitrary-Kap B1/\allowbreak{}B2 theorem, Sel1=0 assertion, arbitrary-chain Core or complete same-J arithmetic support producer is claimed.


\emph{Current subfrontier.} Full integral TypeForAlt; nonvanishing and separating arithmetic fibers on the original Hecke blocks; four same-J Hida outputs and all-level automorphy on both Links. Full WtRep, abstract Sp.absSpec and arithmetic SpecIntIff.spec are reconstructed with the actual coefficient point, fixed iota and one map at every original level.


\paragraph{13. TASK\_\allowbreak{}CASE3\_\allowbreak{}AUT\_\allowbreak{}JOIN — B — Complete the original Case-III automorphy/\allowbreak{}image join}

Install the actual Case-III ProModAlt, SeedIVAlt and shared PMExAlt with the inherited CohomologyHyp/\allowbreak{}Ferrero-Washington/\allowbreak{}HFit/\allowbreak{}SeedII facts into the original whole Case-III automorphy/\allowbreak{}image conclusion.


\paragraph{14. TASK\_\allowbreak{}CASE3\_\allowbreak{}PROMOD\_\allowbreak{}JOIN — C — Reconstruct the actual Case-III chain-to-promodularity proof}

Use the original tensor-factor tower, ChainRankFour, ChainDesc4, seeded rank-four chain reduction, actual PMExAlt and SeedIVAlt to construct ProModAlt on every original Ctx3 star point at the same fixed iota, preserving all finite-coefficient, inertia and support maps.


\paragraph{15. SEED\_\allowbreak{}IV\_\allowbreak{}ALT — C — Actual Case-III seed with prechosen residual datum}

Base Ctx3/\allowbreak{}Four tower, iota and Amb -> w,\allowbreak{} s before every later Dies tower -> compatible finite-coefficient seed with exact residual support/\allowbreak{}root entries.


\paragraph{16. TENSOR\_\allowbreak{}FACTOR\_\allowbreak{}TOWER — B — Actual tensor-factor tower and characters}

Original finite-field, tame, factor and character constraints; source-level construction is sizeable and has not been reconstructed here.


\paragraph{17. CHAIN\_\allowbreak{}RANK\_\allowbreak{}FOUR — C — Actual SomeIn-to-AllIn propagation}

Same source Space and PM ideal, residual point Good4, original bound 30(C+1)\textasciicircum{}2 <= d; geometry must preserve all identities.


\paragraph{18. CHAIN\_\allowbreak{}DESC4 — B — Original residual chain descent}

Finite residual descent retaining source Chain4 compatibility.


\paragraph{19. RANK4\_\allowbreak{}CHAIN\_\allowbreak{}REDUCTION — C — Seeded reduction to the original residual chain}

Select T1 before every later Deep tower; retain one compatible seed and Chain4 for every original residual lift over Q.


\paragraph{20. TASK\_\allowbreak{}COUNT\_\allowbreak{}V\_\allowbreak{}INSTALL — B — Complete the original CountV finite assembly}

From exactly the listed source host/\allowbreak{}cohomology inputs, complete countV\_\allowbreak{}of\_\allowbreak{}finiteness\_\allowbreak{}duality\_\allowbreak{}euler\_\allowbreak{}host, including original host/\allowbreak{}frame formation, normalized length, the CaseV guard and R1<=2\textasciicircum{}c2.


\paragraph{21. THM\_\allowbreak{}S6\_\allowbreak{}S — C — Original unitary-three ordinary support and automorphy}

Original host/\allowbreak{}frame and fixed iota -> one marked g1 -> every universal R,\allowbreak{} u -> weight-zero finiteness and every Type-I BPrime\_\allowbreak{}s point.


\paragraph{22. HOST\_\allowbreak{}FRAME\_\allowbreak{}LENGTH — B — Original Case-V host/\allowbreak{}frame and normalized-length installation}

Actual HostData, frame, local-condition equality, abstract count and normalized original Selmer length must all refer to the same objects.


\paragraph{23. HOST\_\allowbreak{}NUMBER\_\allowbreak{}S — C — Uniform original-host rank-three catalogue}

One exponent before all closures, coefficient rings, levels, characters, frames, embeddings and marked points; one separating trace function controls all later finite pairwise nonisomorphic collections.


\paragraph{24. TASK\_\allowbreak{}HOST\_\allowbreak{}VAN\_\allowbreak{}E\_\allowbreak{}INSTALL — B — Transport K/\allowbreak{}duality to literal HostVanE}

Use the already supplied LemmaK/\allowbreak{}local-global duality and finiteness on the actual CountV host/\allowbreak{}frame and exact original coefficient modules.


\paragraph{25. TASK\_\allowbreak{}HOST\_\allowbreak{}XPP\_\allowbreak{}INSTALL — B — Transport CorX to literal HostXpp}

Preserve Psi\textasciicircum{}-2 eta, the prime-to-ell quotient, actual local conditions and the original CountV host while installing BR7.


\paragraph{26. DESIGN\_\allowbreak{}CHOICES\_\allowbreak{}COMMON — C — Construct the original common design family and bundle}

For the one target HostExW, construct the original compatible DesignChoices on Cdes=WithTw(WithSm Cprime), where Cprime=WithChN(A intersect target class), A is the one Built construction class, and cT(eps,\allowbreak{} y)=cTmin(y). Retain every original field, tuple, word, metric and common bundle duty. DSZ10 and RMP11 give size/\allowbreak{}rank properties of class members, not member existence. The original RecExP/\allowbreak{}RecAtG and common object/\allowbreak{}row route remains coupled. Construct ONE original common design family R and its recursive/\allowbreak{}geometric half with original Inputs/\allowbreak{}Pins/\allowbreak{}Holds/\allowbreak{}ExAt/\allowbreak{}InGeo and required CoreP. PKT now reconstructs full exists\_\allowbreak{}pkT, rowPk\_\allowbreak{}nonempty and exists\_\allowbreak{}moved with Copy/\allowbreak{}Coh on those original given inputs and a supplied fitting later stage: one catalogue, both ChT, complete LfT, actual integral lattices and all metric/\allowbreak{}measure/\allowbreak{}period constants are retained; one numerical record precedes every eta and V/\allowbreak{}S/\allowbreak{}N/\allowbreak{}G. RDR supplies one Side0 before all readers, ReadPk/\allowbreak{}pins and full ReadL with exact o.eta=q.eta. The initial actual geometry, Fits/\allowbreak{}division/\allowbreak{}central conditions, primitive rational tests and eta agreement must still be produced. Inherited CSR supplies finite columns, one ReadC before later twists, period/\allowbreak{}evaluation and scalar/\allowbreak{}error transports, RowSideDetermined and same-R Bundle/\allowbreak{}BTw/\allowbreak{}BWit/\allowbreak{}BunU consumers. RPI supplies exact RowPairInt for actual pairs and finite translates with original mu\_\allowbreak{}b/\allowbreak{}fd0; first literal DivDef3 at the same updated R,\allowbreak{} F,\allowbreak{} Y,\allowbreak{} Z,\allowbreak{} C,\allowbreak{} sEE follows from its retained original readings, Pins/\allowbreak{}Holds, stored nonzero IsItW chain, five real-group predicates and group tie. Remaining common fields include host/\allowbreak{}frame, MultEq nonzero-part support, that original nonzero chain/\allowbreak{}group tie, eigensystem/\allowbreak{}carrier/\allowbreak{}parts and all-embedding logLam, HD/\allowbreak{}d3b, NEb, actual rational tests and base rows. Retain every simultaneous family hypothesis and eventual threshold. Stronger second DivRows3 remains separate; universal quotient-height CuspRapidDecay and actual common member existence are not inferred from the comparison constructors. Version 0.41 supplies the finite phase-correct rational theta route: exact coefficient extension, split representatives, square-law/\allowbreak{}PhiIs, Pins from CoversAt, class-sum/\allowbreak{}point-identity transfer to OneMapOverL, actual-good-projector RatG and both-level LiftTestB. Its original higher covers, all-tuple class-sum uniqueness including phase rescaling, one PointIdentCore before O, base ClassSumsSpan and geometric LatFull remain explicit producers. MLat is available from the original LatFull input. This sharpens the rational-test work without claiming actual common geometry, Fits or eta agreement.


\paragraph{27. LOCAL\_\allowbreak{}NBUILT\_\allowbreak{}COMMON — C — Actual local specialization on the common selected family}

Construct, with the required closed-class property on the same Des(C,3), one eventual family of actual NBuilt certificates for the original Fm.D(x,\allowbreak{} i), retaining the certificate before forgetting to TSel. HLP now reconstructs full HeckeLift0, all-index commutation, HeckeStd0 and HeckeIntTw from the original MtInput and complex Hecke facts, in the full ambient GenAt range k<=nV. Every maximal ideal is an actual supported residual character class; literal MtSub is the O-span of projections of the original MtLat generators, is finite free, and localizes onto the exact NCov.part. Literal Psi has its exact range and off-space projection error, both original span assertions hold, and ActEx/\allowbreak{}ActHK are constructed on the same lattice under original semisimplicity. Original MjLat relation quotients, universal coefficient-map relation fields, selected tame generators, coefficient embeddings and all identity transports give the complete nonempty\_\allowbreak{}spec constructor on its full printed input list. A single compositum combines the four original coefficient requirements with big fixed beforehand. Whole ndSpec\_\allowbreak{}exists still needs actual common geometric/\allowbreak{}choice MtInput, complex HeckeStdPlain and the no-Jordan producer, original F3 bridges with valuation/\allowbreak{}GoodCoef/\allowbreak{}iq, original F4 common-overfield reading of the entire coefficient field, and exact arithmetic character/\allowbreak{}reciprocity/\allowbreak{}admissibility suppliers. Original AutSpec, Galois specialization and the remaining NBuilt fields on that same installation also remain. Do not replace the actual projected lattice or maps by abstract objects with the same dimensions.


\paragraph{28. SECOND\_\allowbreak{}LAYER\_\allowbreak{}COMPATIBLE\_\allowbreak{}SELECTION — C — Original second-layer selector and class duties}

For the same A, supply Duty on A and a Closed common original TSel family whose selected records are the original true-column rec7 records (or have the entire dependent tuple transported), with their real witness and original coefficient-field root data. Existence/\allowbreak{}compatibility remains arithmetic; merely abstractly isomorphic fields or independently selected columns do not suffice. CSR21–22 now construct the same-R Bundle/\allowbreak{}BTw/\allowbreak{}BWit/\allowbreak{}BunU outputs from the exact selected realization/\allowbreak{}row-side packet and NEb, with the original EvP/\allowbreak{}ParEv quantifier order. Actual common selection, the five stated realization packets and NEb remain upstream; none is replaced by an independently chosen compatible record. Version 0.38 closes the selected analytic extension and reconstructs the first updated DivDef3 consumer; it leaves the original common selection, geometric readings, rational tests, base rows, host/\allowbreak{}parts/\allowbreak{}NEb and exact retained DI premises upstream, with unchanged EvP/\allowbreak{}ParEv thresholds. Version 0.39 supplies the original given-data reader constructors and exact RowPairInt/\allowbreak{}first DivDef3 on original premises. It leaves one actual common geometric/\allowbreak{}comparison/\allowbreak{}test member, eta agreement, compatible selection, rational tests, base rows, host/\allowbreak{}parts/\allowbreak{}NEb and original EvP/\allowbreak{}ParEv thresholds upstream. Version 0.40 also supplies original comparison/\allowbreak{}RowPk/\allowbreak{}fitting-stage Moved algorithms and exact common data transport. The actual initial geometric member, fitting tests, eta agreement and compatible second-layer selector with original eventual thresholds remain upstream. Version 0.41 supplies the finite phase-correct rational theta route: exact coefficient extension, split representatives, square-law/\allowbreak{}PhiIs, Pins from CoversAt, class-sum/\allowbreak{}point-identity transfer to OneMapOverL, actual-good-projector RatG and both-level LiftTestB. Its original higher covers, all-tuple class-sum uniqueness including phase rescaling, one PointIdentCore before O, base ClassSumsSpan and geometric LatFull remain explicit producers. MLat is available from the original LatFull input. This sharpens the rational-test work without claiming actual common geometry, Fits or eta agreement.


\paragraph{29. BASE\_\allowbreak{}THETA\_\allowbreak{}READINGS — C — Actual base reading on the common bundle}

For every family in the original class, eventually on the SAME selected Bundle and for each i<m, supply the actual R,\allowbreak{} F,\allowbreak{} Read,\allowbreak{} ReadL, weak models/\allowbreak{}Weil data and phase-one pins. Identify their LiftTie.theta, mtRead, mjRead and ThDef3 basis with that Bundle.P.A3, P.Zs, P.S0 and its stored maps. Exact field, embeddings, metrics, measures, row family and weak choices are retained. This remains one part of the coupled BTH12 producer. PKT now supplies the complete original comparison packet, RowPk and fitting-stage Moved with Copy/\allowbreak{}Coh on original given geometry, Inputs/\allowbreak{}Pins/\allowbreak{}Holds/\allowbreak{}ExAt/\allowbreak{}InGeo and Fits. Its ChT/\allowbreak{}LfT/\allowbreak{}Reps, actual lattice carriers, original Haar and norm constants, period maps, representatives and all-place RrT are constructed together; its one numerical record precedes every eta/\allowbreak{}V/\allowbreak{}S/\allowbreak{}N/\allowbreak{}G. RDR gives the original level-zero side, moved-Read and ReadPk/\allowbreak{}pin constructors on those objects, with one Side0 before every reader and all twelve Read ties. Full ReadL retains the exact o.eta(j,\allowbreak{} v)=q.eta(j,\allowbreak{} v) guard for every 1<=j<=m and place v; arbitrary Over does not imply it. Inherited CSR fixes one ReadC before all TwArgs, transports both metrics and all thirteen reading maps, and gives actual period/\allowbreak{}evaluation with kappa=(c*nP)\textasciicircum{}(-1/\allowbreak{}2). RPI supplies exact row-pair integrability, so first updated DivDef3 follows from the retained original Pins/\allowbreak{}Holds, nonzero chain, five real-group predicates and group tie. Initial theta extension remains complete. Remaining work is actual initial common geometry and its original predicates, Fits/\allowbreak{}primitive rational tests, eta agreement and outer family/\allowbreak{}eventual/\allowbreak{}same-Bundle duties, reFace compatibility, MultEq support, original chain/\allowbreak{}group data, actual base rows and any stronger DivRows3 theta square. No new compatible reading, family threshold or Haar normalization is substituted. Version 0.41 supplies the finite phase-correct rational theta route: exact coefficient extension, split representatives, square-law/\allowbreak{}PhiIs, Pins from CoversAt, class-sum/\allowbreak{}point-identity transfer to OneMapOverL, actual-good-projector RatG and both-level LiftTestB. Its original higher covers, all-tuple class-sum uniqueness including phase rescaling, one PointIdentCore before O, base ClassSumsSpan and geometric LatFull remain explicit producers. MLat is available from the original LatFull input. This sharpens the rational-test work without claiming actual common geometry, Fits or eta agreement.


\paragraph{30. BASE\_\allowbreak{}THETA\_\allowbreak{}RATIONAL\_\allowbreak{}TESTS — C — Simultaneous original base rational tests}

On the SAME original Bundle/\allowbreak{}reading/\allowbreak{}theta map, for every i<m supply one finite separating complex-linear test family whose values lie in the fixed embedded coefficient field on theta of EVERY original rational source vector and on EVERY stored target-coordinate basis vector. BTH8/\allowbreak{}BTH11 use this same LiftTestB for rational adjoint image and matrix coordinates. RTG now reconstructs prescribed-field coefficient extension, original split class representatives/\allowbreak{}RepsOK, the self-dual split-lattice square law and finite residue PhiIs; every Pins field follows from the actual CoversAt. With the base coefficient record literally fixed, exact class-sum existence/\allowbreak{}uniqueness including phase-rescaled O, PointIdentCore with common lambda,\allowbreak{} z before every O, and base ClassSumsSpan imply all-tuple finite lift span and OneMapOverL. The phase is exactly cj(twAt embLL A3.A.v0)(R.phs i). RatG uses the actual good projector eOp(y), and full LiftTestB follows at BOTH original levels. These finite implication proofs are complete. Remaining original producers are the prescribed higher covers, those actual class sums/\allowbreak{}uniqueness, PointIdentCore, base ClassSumsSpan and the base/\allowbreak{}common geometric datum; exact MLat is supplied only from original LatFull, whose geometric production remains separate. Finite dimension or an unrelated rational theta model cannot replace these same-object identities. Original outer eventual-family/\allowbreak{}Bundle quantifiers remain.


\paragraph{31. SLR\_\allowbreak{}BASE\_\allowbreak{}ROWS\_\allowbreak{}PRODUCER — C — Actual base row identities on that same bundle}

For the SAME selected Bundle, common original fields and readings used by the other three BTH12 inputs, supply actual RowsHold0\_\allowbreak{}3 uniformly eventually for every family. Its original primitive row realization, spectrum/\allowbreak{}type and basis identifications must be retained. BTH supplies a theta consumer only; it does not construct the base row identity. No separate witness may be chosen merely to satisfy this row requirement.

VII. PUBLISHED INPUTS, VALIDATION AND DELIVERY

No new published theorem import is needed for these 87 claims. The original reciprocity application from CIP is reused with its already checked split-CM-place, normalization and coefficient hypotheses. The complete prior publication registries, including their exact theorem statements and checked hypotheses, remain in the archive. Independently proved finite algebra in the current chapters is not mislabeled as a new literature application.

Every new claim has exact source declarations, inclusive physical line ranges, full-file hashes and proof-block hashes. All predecessor claim fields, source arrays and mathematical dependency edges remain. Independent mathematical review and source binding are recorded separately from draft TeX syntax checks; neither is represented as a fresh Lean kernel verification. The source correctness remains the standing assumption.

RH\_\allowbreak{}reconstruction\_\allowbreak{}v0\_\allowbreak{}41.zip is the cumulative continuation snapshot. RH\_\allowbreak{}readable\_\allowbreak{}reconstruction\_\allowbreak{}v0\_\allowbreak{}41.tex contains all historical and current human proofs. RH\_\allowbreak{}dependency\_\allowbreak{}audit\_\allowbreak{}v0\_\allowbreak{}41.txt is this self-contained audit, including every remaining selected contract. Component chapters, source/\allowbreak{}proof ledgers, current and historical graphs, target dispositions, independent reviews, validation and source-anchored continuation notes are included. No PDF is generated.

The next-step notes are continuation\_\allowbreak{}notes\_\allowbreak{}v0\_\allowbreak{}41/\allowbreak{}next\_\allowbreak{}count\_\allowbreak{}ii.txt, next\_\allowbreak{}analytic\_\allowbreak{}and\_\allowbreak{}common.txt and next\_\allowbreak{}division\_\allowbreak{}and\_\allowbreak{}readings.txt. They distinguish work plans from proved results. Historical RF16 remains governed by DEX11. The complete human RH and unrestricted Dirichlet-GRH reconstruction remains unfinished.

\clearpage
\section{A human-readable conditional main proof}
\label{sec:conditional-main-v42}

The argument below reaches RH and GRH for Dirichlet $L$-functions
from two explicitly stated arithmetic inputs.  It gives the whole
terminal proof in ordinary mathematical notation, including the
normalizations, the dependence of the constants and the passage from
a one-sided prime-sum estimate to zeros.  The two arithmetic inputs
are not asserted here.  The current reconstruction has proved many
of their component implications, but their common construction and
the five quantitative estimates are still incomplete.  The final
table locates every v0.41 obligation in this main argument.  The
adjoint theorem in the next section closes one reference-matching
obligation; it does not supply these two arithmetic inputs.

\subsection{The fixed analytic target}

Let $\Lambda$ be the von Mangoldt function and let $\gamma$ be
Euler's constant.  For $y>0$ put
\begin{align*}
 G(y)&=\sum_{r\le y}\Lambda(r)(y/r-1)
                         -y\log y+(1+\gamma)y,\\
 G_\chi(y)&=\sum_{r\le y}\Lambda(r)\chi(r)(y/r-1)
                         +yL'/L(1,\chi)\qquad(\chi\ne1).
\end{align*}
Here $\chi$ is primitive, and $G_1$ means $G$; the sum is empty
if $y<1$.  The plus sign before $yL'/L(1,\chi)$ cancels the
unwanted linear Mellin pole.  Fix a primitive target $\psi$ of
order $o$.  Let $\mathcal C_\psi$ be the set of distinct primitive
inducers of its powers.  Its positive signal is
\[
 H_\psi(x)=G(x)+\frac1{3o+1}
             \sum_{\chi\in\mathcal C_\psi\setminus\{1\}}
                                      \operatorname{Re}G_\chi(x).
 \tag{CMP.1}
\]
This family and these coefficients are fixed throughout the proof.
For $\psi=1$ it is simply $G$.

The standard analytic facts used below are continuation and the
functional equations of primitive Dirichlet $L$-functions, their
Euler products, and their nonvanishing on $\operatorname{Re}s=1$.
The reconstruction already supplies the needed nonvanishing and
primitive-character transitions in QW14 and GH6.  Complex analysis
then gives CMP4 below.  No unproved zero-free half-plane inside the
critical strip is among the analytic inputs.

\subsection{First arithmetic input: one compatible finite construction}

At a scale $x\ge3$ and an accuracy parameter $\varepsilon>0$,
the following objects and assertions constitute a
\emph{compatible finite datum}.  They are deliberately stated on
one set of objects.

Take a finite abelian group $A$ of order $n$, a subgroup $A_0$ of
index $o$, and write
\[
 G_0=(\mathbb Z/2\mathbb Z)^2\times A,\quad
 \Phi=\{(0,\allowbreak{} b,\allowbreak{} h):b\in\mathbb Z/2\mathbb Z,\ h\in A\},
 \quad f=|\Phi|=2n.
\]
Set $c_0=(1,0,0)$, $c_\infty=(1,1,0)$,
$\chi_d(a,\allowbreak{} b,\allowbreak{} h)=(-1)^a$ and $\chi_q(a,\allowbreak{} b,\allowbreak{} h)=(-1)^b$.
Supply a positive modulus $N_0$ and an onto homomorphism
\[
 r:(\mathbb Z/N_0\mathbb Z)^\times\longrightarrow G_0.
\]
For every character $\xi$ of $G_0$, let $\xi^*$ be the primitive
inducer of $\xi\circ r$, and set
$\mathcal G_\xi=G_{\xi^*}$ and
$\ell_\xi=L'/L(1,\xi^*)$ for $\xi\ne1$.
These are actual equalities, not freely chosen approximations to
the logarithmic derivatives.

For $p\mid N_0$, define $I_p$ to be the image of the $p$-primary
CRT unit factor.  Define $F_p$ modulo $I_p$ by the class of $p$
in the prime-to-$p$ factor.  For $p\nmid N_0$, put $I_p=0$ and
$F_p=r(p)$.  The choice of a representative of $F_p$ will not
matter.  As in the source host, supply a finite prime set $S$ and
a prime $d\notin S$ such that
\[
 I_p=0\ (p\notin S\cup\{d\}),\qquad
 I_p\subset\Phi, F_p\in\Phi\ (p\in S),\qquad
 I_d=\{0,\allowbreak{} c_0\}.                                    \tag{CMP.2}
\]

Put $\mathcal B=\{\xi:\xi(0,0,\allowbreak{} A_0)=1\}$; it has $4o$
elements.  Supply a subset $K\subset\mathcal B$ containing $1$
such that $\xi(c_0)=1$ for $\xi\in K$, $\chi_q\notin K$, and
$\xi\mapsto\xi^*$ identifies $K$ with
$\mathcal C_\psi$.  Supply nonempty $T\subset A$ and
$T_2\subset A_0$, and set
\[
 n_T=|T|,\quad n_2=|T_2|,\quad d_1=n_T/n,\quad d_2=n_2/n.
\]
The required bias in every coset is
\[
 \left|\#(T\cap(h+A_0))-
       \frac{3+\mathbf1_{h\in A_0}}{3o+1}\,\allowbreak{} n_T\right|
                   \le x^{-1+\varepsilon}n_T.
 \tag{CMP.3}
\]
For constants $K_0,\allowbreak{} C_E$ that stay fixed as the windows move,
require
\[
 |\ell_\xi|\le K_0\ (\xi\in K\setminus\{1\}),\qquad
 |\ell_\xi|\le C_E\log x\ (\xi\in\mathcal B\setminus K),
 \quad
 |\mathcal G_\xi(x)|\le x^{1/2+\varepsilon}
                 \ (\xi\in\mathcal B\setminus K).
 \tag{CMP.4}
\]

Choose $m\ge1$ integral, $b,\allowbreak{} b',\allowbreak{} X',\allowbreak{} k>0$, $\eta,\theta_0\ge0$,
$0\le a_E\le2xb$, numbers $a_t>0$ for $t\in T$, and positive
weights $w_{u,\allowbreak{} j}:\Phi\to\mathbb R_{>0}$ for
$u\in A_0$, $1\le j\le m$.  The thick and thin weights satisfy
\begin{gather*}
 |a_t/b-x|\le1,\qquad b'X'\le xb,\\
 b\le w_{u,\allowbreak{} j}(0,1,\allowbreak{} s+u)\le b+\eta+1\quad(s\in T_2),\\
 w_{u,\allowbreak{} j}(\rho)\le2b'
   \quad\bigl(\rho\notin\{(0,1,\allowbreak{} s+u):s\in T_2\}\bigr),\\
 |\mathcal G_\xi(y)|\le y\theta_0
           \quad(\xi\in\widehat G_0, y\ge X'/3),\\
 2f(x+1)\theta_0\le x^{-1},\qquad
 6f(x+1)(\eta+1)\le b/x.                         \tag{CMP.5}
\end{gather*}
The scalar sizes are
\[
 d_1\ge\tfrac12x^{-\varepsilon/8},\quad
 x^{-\varepsilon}\le d_2\le1,\quad
 x^{1/2-\varepsilon}\le m\le x^{1/2+\varepsilon/2},\quad
 xm\le k\le n^2x^{1/2-\varepsilon}.                \tag{CMP.6}
\]
Define the positive normalizations
\[
 Q=n_Tmb,\qquad D=Qd_2/2,
\]
and the error budget
\[
 B=n_Ta_E(1+m^2/k)
       +m\left(2nmbx^\varepsilon+
            \frac{n_2kbx^\varepsilon}{n(nx^{-\varepsilon})}\right).
 \tag{CMP.7}
\]
The symbols $D$ and $B$ have this meaning only in this main proof;
in particular $D$ is a scalar, not a replacement geometric object.

\subsection{Second arithmetic input: a uniform upper bound for one mean gap}

For $t\in A$, the bijection
$\rho_t(g)=c_\infty+g+(0,0,\allowbreak{} t)$ maps $G_0\setminus\Phi$ to
$\Phi$.  Put $(z)_+=\max(z,0)$ and define
\begin{align*}
 \mathcal A(a,\allowbreak{} t,\allowbreak{} w)
 &=\sum_{p}\sum_{j\ge1}\log p\,
       \frac2{|I_p|}
       \sum_{\substack{g\notin\Phi\\g-jF_p\in I_p}}
                    (a/p^j-w(\rho_t(g)))_+,\\
 \mathcal M(a,\allowbreak{} w)
 &=a\left(\log a-1-\gamma+\ell_{\chi_d}
                         -\frac1f\sum_{\rho\in\Phi}\log w(\rho)\right),\\
 \Gamma_u
 &=\sum_{t\in T}\sum_{j=1}^m
       \{\mathcal A(a_t,\allowbreak{} t,\allowbreak{} w_{u,\allowbreak{} j})
                         -\operatorname{Re}\mathcal M(a_t,\allowbreak{} w_{u,\allowbreak{} j})\},\\
 \overline\Gamma&=\frac1{|A_0|}\sum_{u\in A_0}\Gamma_u.
                                                               \tag{CMP.8}
\end{align*}
All prime-power sums here are finite: use a cutoff above all ratios
$a_t/w_{u,\allowbreak{} j}(\rho)$.  Conditions (CMP.2) make this precisely the
source's asked prime sum: primes in $S$ give zero, each unramified
term outside $\Phi$ has coefficient $2$, and the exceptional prime
$d$ has coefficient $1$ on its inertia coset outside $\Phi$.
The $d$ term is retained even when a particular cutoff makes it zero.

The quantitative input is the single inequality
\[
                    \overline\Gamma\le x^{C\varepsilon}B.
 \tag{CMP.9}
\]
The source's pointwise inequalities $\Gamma_u\le x^{C\varepsilon}B$
for every $u$ imply (CMP.9).  The next proof needs only (CMP.9);
this is a genuine weakening of the sufficient arithmetic input.
It allows cancellation between different translates before any
upper estimate is applied.

The complete order of choices is as follows:
\begin{quote}
For each fixed primitive $\psi$, choose $C\ge0$ and
$\varepsilon_0>0$.  For each $0<\varepsilon<\varepsilon_0$,
choose $K_0,\allowbreak{} C_E$ and any further constants of the construction.
For every $Y_0$, construct $Y\ge\max(3,\allowbreak{} Y_0)$ such that, for
every real $x\in[Y,\allowbreak{} Y^{3/2}]$, one compatible datum satisfies
(CMP.2)--(CMP.6) and its own inequality (CMP.9).
\end{quote}
Equivalently, the first input produces a common family on
arbitrarily late windows, and the second is uniform on that same
family after one further threshold.  The family may depend on
$\varepsilon$; the datum may depend on $x$.  The target,
coefficient $C$, and positive character weights may not be chosen
anew for each window.  Existence of separate data satisfying
separate clauses is insufficient.

% BEGIN CMP1
\begin{lemma}[CMP1: primitive values and the exact balance identity]
\label{claim:CMP1}
Every compatible datum satisfies
\[
 \mathcal A(a,\allowbreak{} t,\allowbreak{} w)-\mathcal M(a,\allowbreak{} w)
   =\frac1f\sum_{g\notin\Phi}\sum_{\xi\in\widehat G_0}
       \overline{\xi(g)}w(\rho_t(g))
                      \mathcal G_\xi(a/w(\rho_t(g))).
 \tag{CMP.10}
\]
\end{lemma}
\begin{proof}
For every prime power,
\[
 \xi^*(p^j)=\frac1{|I_p|}\sum_{h\in I_p}\xi(jF_p+h).
\]
If $\xi$ is nontrivial on inertia, both sides are zero: its
primitive conductor is divisible by $p$.  Otherwise it factors
through the prime-to-$p$ modulus, and both sides are $\xi(jF_p)$.
Character orthogonality consequently gives
\[
 \frac1f\sum_\xi\overline{\xi(g)}\xi^*(p^j)
          =\frac2{|I_p|}\mathbf1_{g-jF_p\in I_p}.
\]
Substitute into (CMP.8) and interchange finite sums.  For a
nonprincipal character, replacing the uncentered prime sum by
$w\mathcal G_\xi(a/w)$ subtracts $a\ell_\xi$.
For the principal character the extra terms are
$a\log a-a\log w-(1+\gamma)a$.
The character sum over $G_0\setminus\Phi$ equals $f$ at $1$,
$-f$ at $\chi_d$, and zero elsewhere.  These terms give exactly
$\mathcal M$, including its plus sign before $a\ell_{\chi_d}$.
This proves (CMP.10) at ramified as well as unramified primes.
\end{proof}
% END CMP1

% BEGIN CMP2
\begin{lemma}[CMP2: average first, then keep the positive target]
\label{claim:CMP2}
For all sufficiently large $x$, with threshold depending only on
$o,\allowbreak{} K_0,\allowbreak{} C_E,\varepsilon$, the same datum satisfies
\[
 H_\psi(x)\le\frac{\overline\Gamma}{D}
        +(4+4o)x^{1/2+\varepsilon}+\frac4{d_2x}.
 \tag{CMP.11}
\]
This conclusion uses the structural and cutoff conditions but not
the gap bound (CMP.9).
\end{lemma}
\begin{proof}
We recall the elementary uniform estimates established in WB3.
There is an absolute $C_0$ such that, for $y,\allowbreak{} z\ge1$,
\begin{align*}
 |G(y)-G(z)|&\le C_0|y-z|,\\
 |\mathcal G_\xi(y)-\mathcal G_\xi(z)|
 &\le |y-z|\{\log\max(y,\allowbreak{} z)+C_0+|\ell_\xi|\}\quad(\xi\ne1),\\
 |w\mathcal G_\xi(a/w)-w'\mathcal G_\xi(a/w')|
 &\le3a|w-w'|/\min(w,\allowbreak{} w').
                                                               \tag{CMP.12}
\end{align*}
For completeness the only arithmetic estimates in them are
$\sum_{r\le y}\Lambda(r)\le3y$ and
$\log k-1\le\sum_{r\le k}\Lambda(r)/r\le\log k+3$.
The first follows from central binomial coefficients and summing
prime powers; the second follows from
$\log(k!)=\sum_{r\le k}\Lambda(r)\lfloor k/r\rfloor$.
Differentiate the ramp sums between integers for the first two
inequalities in (CMP.12).  For the last use the Lipschitz property
of $(a/r-w)_+$; the nonprincipal linear term cancels exactly.
For the principal character the additional logarithmic difference
is bounded by $a|w-w'|/\min(w,\allowbreak{} w')$, and the difference of the two
nonnegative terms is still bounded by the displayed constant $3$.

Replace each thick weight by $b$ in (CMP.10).  The error per
character is at most $3(x+1)(\eta+1)$ by (CMP.12).
At a thin point,
$a_t/w\ge(x-1)b/(2b')\ge X'/3$, so its whole term has modulus
at most $(x+1)b\theta_0$.  There are $f$ points and $2f$
characters.  After division by $f$, (CMP.5) bounds the total error
for each $t,\allowbreak{} j$ by $2b/x$.
The ideal thick term is
\[
 \mathcal T_u(t,\allowbreak{} y)=\frac b f\sum_{s\in T_2}\sum_\xi
             \xi(-(1,0,\allowbreak{} s+u-t))\mathcal G_\xi(y).
\]
Average this term over $u\in A_0$.  The characters outside
$\mathcal B$ vanish, whereas those in $\mathcal B$ are trivial
on both $u$ and $s$.  Since $f=2n$, the result is
\[
 \frac1{|A_0|}\sum_u\mathcal T_u(t,\allowbreak{} y)
       =\frac{bd_2}{2}\sum_{\xi\in\mathcal B}
                                   \xi(1,0,\allowbreak{} t)\mathcal G_\xi(y).
\]
Thus (CMP.10), before any upper bound on $\Gamma_u$, gives
\[
 \frac{mbd_2}{2}\sum_{t\in T}
       \operatorname{Re}\sum_{\xi\in\mathcal B}
               \xi(1,0,\allowbreak{} t)\mathcal G_\xi(a_t/b)
                  \le\overline\Gamma+2Q/x.          \tag{CMP.13}
\]

For a nontrivial character of $A/A_0$, (CMP.3) makes its mean on
$T$ differ from $c=(3o+1)^{-1}$ by at most
$ox^{-1+\varepsilon}$.  For a character trivial on $A$, that
mean is exactly $1$.  The factors $\xi(c_0)$ are $1$ for
retained characters.  Among these, only $1$ is trivial on $A$:
the only other possibility would be the excluded $\chi_q$.
The retained mean therefore gives exactly (CMP.1).  There are
only $4o$ characters to estimate.  Moving $a_t/b$ to $x$ costs
$O(\log x)$ per character by (CMP.12), with the constants in
(CMP.4).  The Fourier bias costs
$O(x^\varepsilon(\log x+1))$, since
$|\mathcal G_\xi(x)|\le x(\log x+C_0+K_0)$ for a retained
nonprincipal character.  All these constants are fixed as the
host grows.  Their sum is at most $4x^{1/2+\varepsilon}$ beyond
one common threshold.  Omitted characters contribute at worst
$-\sum_{\mathcal B\setminus K}|\mathcal G_\xi(x)|$.
With this sum denoted by $\Sigma(x)$, (CMP.13) becomes
\[
 D\{H_\psi(x)-\Sigma(x)-4x^{1/2+\varepsilon}\}
                       \le\overline\Gamma+2Q/x.
\]
By (CMP.4), $\Sigma(x)\le4o x^{1/2+\varepsilon}$.
Divide by $D=Qd_2/2>0$ to obtain (CMP.11).
\end{proof}
% END CMP2

% BEGIN CMP3
\begin{proposition}[CMP3: the mean gap has square-root scale]
\label{claim:CMP3}
If (CMP.9) holds, then for all sufficiently large $x$
\[
                     H_\psi(x)\le x^{1/2+(C+4)\varepsilon}.
 \tag{CMP.14}
\]
\end{proposition}
\begin{proof}
Put $Z=x^{1/2+13\varepsilon/8}$.  Dividing (CMP.7) by $Q$ gives
\[
 \frac BQ=\frac{a_E}{mb}+\frac{a_Em}{bk}
           +\frac{2mx^\varepsilon}{d_1}
           +\frac{d_2kx^{2\varepsilon}}{d_1n^2}.
\]
The four terms are bounded respectively by $2Z,2Z,4Z,2Z$.
Indeed the first is at most $2x/m$, the second at most
$2xm/k\le2$, and the last two use precisely the upper bounds
for $m,\allowbreak{} k$ and the lower bound for $d_1$ in (CMP.6).
Hence $B\le10QZ$.
Retain $\Sigma$ in the penultimate inequality of CMP2 and divide
by $Q/2$.  By $d_2\ge x^{-\varepsilon}$ and $C\ge0$ this yields
\[
 H_\psi(x)-\Sigma(x)
          \le26x^{1/2+(C+21/8)\varepsilon}
          \le x^{1/2+(C+3)\varepsilon},
\]
the last step holding once $26\le x^{3\varepsilon/8}$.
Now use $\Sigma(x)\le4o x^{1/2+\varepsilon}$ and
$x^\varepsilon\ge1+4o$.  This gives (CMP.14).
The source values $e_D=2,\allowbreak{} e_L=0$ give
$C=e_D+e_L+3=5$, and hence the exponent $1/2+9\varepsilon$.
This identifies the selected route's constants; their arithmetic
production is still part of the two inputs.
\end{proof}
% END CMP3

% BEGIN CMP4
\begin{theorem}[CMP4: the one-sided analytic criterion]
\label{claim:CMP4}
Fix $0<\theta<1$.  Suppose that for every $\alpha>0$ and every
$Y_0$ there is $Y\ge\max(3,\allowbreak{} Y_0)$ such that
\[
       \int_Y^{Y^{3/2}}H_\psi(x)_+x^{-\theta-1}\,\allowbreak{} dx\le Y^\alpha.
 \tag{CMP.15}
\]
Then every constituent of $\mathcal C_\psi$ has no zero in
$\operatorname{Re}s>\theta$, apart from the principal pole.
In particular it is sufficient to have
$H_\psi(x)\le C_\theta x^\theta$ throughout arbitrarily late
windows, with $C_\theta$ fixed.
\end{theorem}
\begin{proof}
Here is the complete pole argument, consolidating AN2--AN3 and AQ3.
For $\operatorname{Re}s>1$ the Mellin transforms are
\begin{align*}
 M_1(s)&=-\frac{\zeta'/\zeta(s)}{s(s-1)}
             -\frac1{(s-1)^2}+\frac{1+\gamma}{s-1},\\
 M_\chi(s)&=-\frac{L'/L(s,\chi)-sL'/L(1,\chi)}{s(s-1)}
                                                    \quad(\chi\ne1).
\end{align*}
The apparent singularities at $1$ are removable.  For the principal
term use $\zeta'/\zeta(s)=-(s-1)^{-1}+\gamma+O(s-1)$.
Set $a_1=1$ and $a_\chi=(3o+1)^{-1}$ otherwise.  The transform
of the real function $H_\psi$ is
\[
 \Phi(s)=\tfrac12\sum_{\chi\in\mathcal C_\psi}
                    a_\chi\{M_\chi(s)+M_{\bar\chi}(s)\}.
\]
It extends meromorphically to $\operatorname{Re}s>0$.  Every
constituent zero $\rho$ there produces a simple pole with residue
\[
 \frac{b_\rho}{\rho(1-\rho)},\qquad
 b_\rho=\tfrac12\sum_\chi a_\chi
                    \{m_\chi(\rho)+m_{\bar\chi}(\rho)\}>0,
                                                               \tag{CMP.16}
\]
where $m_\chi(\rho)$ is its multiplicity, zero if it is not a
zero of that function.  Thus coincident zeros cannot cancel.  A
real pole in $(0,1)$ has positive residue.  The Euler product and
nonvanishing on the line $1$ put all poles in $0<\operatorname{Re}s<1$.

Put $g(u)=H_\psi(e^u)$.  The fixed finite family gives the crude
bound $|g(u)|\le M e^{2u}$ for $u\ge0$.  At
$\sigma=336$, define
\[
 A_k(z)=\frac{(-1)^k}{k!}\Phi^{(k)}(z)
       =\int_0^\infty g(u)\frac{u^k}{k!}e^{-zu}\,\allowbreak{} du.
\]
For $U=\log Y\ge1$ and $k=\lceil420U\rceil$, the tails obey
\[
 \int_{(0,\infty)\setminus[U,3U/2]}
          \frac{u^k}{k!}e^{-334u}\,\allowbreak{} du
 \le\frac{e^{-21U}}{320^{k+1}}+
                         \frac{e^{14U}}{348^{k+1}}
 \ll336^{-k}.                                      \tag{CMP.17}
\]
On the upper tail insert $e^{14(u-3U/2)}\ge1$; on the lower
tail insert $e^{14(U-u)}\ge1$, and extend the two integrals to
the positive half-line.  The last inequality uses
$\log(336/320)\le1/20$, $\log(348/336)\ge1/29$,
$420U\le k\le420U+1$.  The constants do not depend on the window.

Reality of $g$ gives, for every real $t$,
\[
 |A_k(\sigma+it)|+A_k(\sigma)
       \le2\int_0^\infty g(u)_+\frac{u^k}{k!}e^{-\sigma u}\,\allowbreak{} du.
\]
On the good interval, use the elementary bound
$u^ke^{-(\sigma-\theta)u}/k!\le(\sigma-\theta)^{-k}$.
Condition (CMP.15), after $x=e^u$, and (CMP.17) imply
\[
 |A_k(\sigma+it)|+A_k(\sigma)
       \le2e^{\alpha U}(\sigma-\theta)^{-k}+O(\sigma^{-k})
       \le2R_\alpha^{-k}+O(\sigma^{-k}),
 \quad R_\alpha=(\sigma-\theta)e^{-\alpha/420}.
                                                               \tag{CMP.18}
\]
The extra real derivative is what permits a one-sided bound.

Suppose a pole lies to the right of $\theta$, and put
$D_0=\min\{\sigma-\theta,\inf_\rho|\sigma-\rho|\}$.
If some pole has horizontal distance $<D_0$ from
$\operatorname{Re}s=\sigma$, choose $t$ close to its imaginary
part so that $z=\sigma+it$ has a unique nearest pole $\rho$
at distance $a<D_0$.  Only finitely many poles compete in the
relevant compact disks.  Each pair with different imaginary parts
excludes at most one $t$; pairs with the same imaginary part have
different horizontal distances.  Hence such a choice is possible.
Choose $a<r<D_0$ with no other pole in the closed $r$-disk about
$z$.  The $r$-disk about $\sigma$ is pole-free.  Subtraction of
the principal part and Cauchy's derivative estimate give
\[
 A_k(z)=\frac{c}{(z-\rho)^{k+1}}+O(r^{-k}),\qquad
 A_k(\sigma)=O(r^{-k}),\qquad c\ne0.
\]
Choose $\alpha>0$ so small that $R_\alpha>r$, before choosing
arbitrarily late windows.  Equation (CMP.18) would force
$|c|a^{-k-1}\le O(r^{-k})+2R_\alpha^{-k}$ as $k\to\infty$,
which is impossible since $a<r<R_\alpha$.

In the remaining case every pole has horizontal distance at least
$D_0$.  The hypothesized right-hand pole forces
$D_0<\sigma-\theta$.  A nearest pole to $\sigma$ exists in a
compact disk.  It must be the unique real pole
$\beta=\sigma-D_0$, because total distance $D_0$ cannot equal a
larger horizontal distance.  Its residue $c$ is positive by
(CMP.16).  For $D_0<r<\sigma-\theta$ sufficiently close to
$D_0$, Cauchy's estimate gives
$A_k(\sigma)=cD_0^{-k-1}+O(r^{-k})$.
Choose $\alpha$ with $R_\alpha>r$ and apply (CMP.18) at $t=0$.
Its left side is $2A_k(\sigma)>0$ for late $k$, again a
contradiction.  There can therefore be no right-hand pole.

Finally a pointwise bound in the statement makes the integral in
(CMP.15) at most $(C_\theta/2)\log Y$, which is at most
$Y^\alpha$ for sufficiently large $Y$.  This proves that
special case too.
\end{proof}
% END CMP4

% BEGIN CMP5
\begin{theorem}[CMP5: the conditional main conclusion]
\label{claim:CMP5}
If the two arithmetic inputs above hold for every fixed primitive
target with the stated order of choices, then RH and GRH for all
Dirichlet $L$-functions hold.  For RH alone it suffices to supply
the construction for $\psi=1$.
\end{theorem}
\begin{proof}
Fix a proposed constituent zero with real part $\beta>1/2$.
Choose $0<\delta<\min(1/2,\beta-1/2)$ and then
$0<\varepsilon<\min(\varepsilon_0,\delta/(C+4))$.
All thresholds in CMP2--CMP3 are now fixed.  The arithmetic input
supplies arbitrarily late windows beyond them.  At every real
point of each such window, CMP3 gives
$H_\psi(x)\le x^{1/2+\delta}$.
CMP4 at $\theta=1/2+\delta$ contradicts the proposed zero.
The same applies to the principal constituent, whose pole at $1$
has been explicitly removed from the argument.

The primitive functional equation reflects a zero in the open
critical strip to the left of $1/2$ into a zero to the right of
$1/2$ of the conjugate primitive character.  Its gamma, conductor
and root-number factors are nonzero in that strip.  For a character
induced by $\chi^*$,
\[
 L(s,\chi)=L(s,\chi^*)
       \prod_{\substack{p\mid N\\p\nmid\operatorname{cond}\chi^*}}
                                      (1-\chi^*(p)p^{-s}).
\]
The additional factors are nonzero for $\operatorname{Re}s>0$,
since $|\chi^*(p)p^{-s}|<1$.  Thus the assertion passes to all
Dirichlet characters.  The ordinary functional equation and the
known trivial zeros give the usual RH formulation as well.
The scope here is Dirichlet GRH; no assertion about arbitrary
automorphic $L$-functions is inferred.
\end{proof}
% END CMP5

% BEGIN CMP6
\begin{corollary}[CMP6: an even smaller sufficient averaged input]
\label{claim:CMP6}
Fix $1/2<\theta<1$ and choose $0<\varepsilon<\theta-1/2$.
For the same compatible data, suppose that on the selected windows
the nonnegative function
\[
 \mathcal E(x)=\frac{(\overline\Gamma(x))_+}{D(x)}
                                      +\frac4{d_2(x)x}
 \tag{CMP.19}
\]
is measurable and, for every $\alpha>0$, has arbitrarily late
windows on which
$\int_Y^{Y^{3/2}}\mathcal E(x)x^{-\theta-1}\,\allowbreak{} dx\le Y^\alpha$.
Then the same constituent zero-free half-plane
$\operatorname{Re}s>\theta$ follows, without (CMP.9).
\end{corollary}
\begin{proof}
CMP2 gives
$H_\psi(x)_+\le\mathcal E(x)+(4+4o)x^{1/2+\varepsilon}$.
The weighted integral of the second term is bounded by
$(4+4o)Y^{1/2+\varepsilon-\theta}/
             (\theta-1/2-\varepsilon)$.
Use the hypothesis with $\alpha/2$ and choose the window late
enough that adding this bounded term still leaves at most
$Y^\alpha$.  CMP4 applies.  Measurability is stated because an
arbitrary pointwise choice of auxiliary data need not provide it.
\end{proof}
% END CMP6

\subsection{Where the remaining arithmetic work enters}

CMP1, the cutoff/Fourier part of CMP2, CMP3, CMP4 and CMP5 are
consolidations of existing complete arguments WB1--WB8, SC1,
AN2--AN3 and AQ1--AQ4.  The explicit averaging-before-estimation
in CMP2 and the resulting input (CMP.19) are additional sufficient
weakening results.  Neither constructs the missing arithmetic data.
The new GQT1--GQT5 packet similarly isolates three sufficient row
counts and two bad-set weight estimates for the retained-defect
input; GQT6--GQT7 show that a uniform polynomial spectral error
suffices for the selected global comparison.  These substitutions
may make reconstruction easier, but their actual suppliers remain
to be written on the same family.

For orientation, the source obtains (CMP.9) by comparing integral
period maps with their archimedean norms.  The denominator and
local-gain estimates force divisibility of the same integral
matrices.  The period/height estimate and spectral estimate control
their norms and the size of their common target.  The retained-row
estimate ensures that discarding exceptional rows leaves the
required positive rank.  The arithmetic slope comparison then
produces the budget (CMP.7).  The geometry and deformation theory
construct those matrices, lattices and rows together, rather than
independently.  This explains why the remaining frontier includes
both construction and quantitative statements.

The following table maps every one of the 33 v0.41 leaf identifiers.
The category letters concern the reconstruction: B denotes a
remaining installation, C a remaining human proof of a producer,
and D a remaining exact match to a published theorem.  They are
not claims that the corresponding mathematics is unproved in the
literature.  ``Construction'' refers to the first arithmetic input;
``gap'' to the second on that construction.

\begin{longtable}{p{0.44\linewidth}p{0.06\linewidth}p{0.41\linewidth}}
\textbf{v0.41 leaf} & \textbf{Kind} & \textbf{Role in the conditional proof}\\\hline
\endhead
\texttt{GLOBAL\_\allowbreak DENOMINATOR} & C & Gap: actual local denominator table and bounds.\\
\texttt{GLOBAL\_\allowbreak RETAINED\_\allowbreak DEFECT} & C & Gap: enough actual rows survive with the original ranks.\\
\texttt{GLOBAL\_\allowbreak PERIOD\_\allowbreak HEIGHT} & C & Gap: norms, eigenvalues, metrics and height on the same period maps.\\
\texttt{GLOBAL\_\allowbreak LOCAL\_\allowbreak GAIN} & C & Gap: local divisibility gain with the required uniform loss.\\
\texttt{GLOBAL\_\allowbreak SPECTRAL} & C & Gap: normalized kernel and target-dimension estimates.\\
\texttt{TASK\_\allowbreak PROP\_\allowbreak GK\_\allowbreak INSTALL} & B & Construction: install one-character control and its congruence integer.\\
\texttt{TASK\_\allowbreak ONE\_\allowbreak CHAR\_\allowbreak PROOF} & C & Construction: the actual one-character argument.\\
\texttt{TASK\_\allowbreak INP\_\allowbreak Z\_\allowbreak ASSEMBLY} & B & Construction: the exact Kato input record.\\
\texttt{RIG\_\allowbreak II} & C & Construction: type-II polarized tangent rigidity.\\
\texttt{AD\_\allowbreak VANISH\_\allowbreak K\_\allowbreak MATCH} & D & Construction: full polarized adjoint vanishing; supplied by ADK1--ADK6 here.\\
\texttt{AUT\_\allowbreak UNRAM\_\allowbreak COEF} & C & Construction: automorphy over the actual unramified coefficient field.\\
\texttt{CNR3} & C & Construction: uniform rank-three/rank-two congruence integer.\\
\texttt{COUNT\_\allowbreak II\_\allowbreak SUPPORT} & C & Construction: two compatible rank-four links and same-ideal support.\\
\texttt{TASK\_\allowbreak CASE3\_\allowbreak AUT\_\allowbreak JOIN} & B & Dirichlet construction: exact Case-III automorphy and image join.\\
\texttt{TASK\_\allowbreak CASE3\_\allowbreak PROMOD\_\allowbreak JOIN} & C & Dirichlet construction: original chain-to-promodularity proof.\\
\texttt{SEED\_\allowbreak IV\_\allowbreak ALT} & C & Dirichlet construction: seed with the prechosen residual datum.\\
\texttt{TENSOR\_\allowbreak FACTOR\_\allowbreak TOWER} & B & Dirichlet construction: actual tensor factors and characters.\\
\texttt{CHAIN\_\allowbreak RANK\_\allowbreak FOUR} & C & Dirichlet construction: propagation along the actual chain.\\
\texttt{CHAIN\_\allowbreak DESC4} & B & Dirichlet construction: residual chain descent.\\
\texttt{RANK4\_\allowbreak CHAIN\_\allowbreak REDUCTION} & C & Dirichlet construction: seeded reduction to that chain.\\
\texttt{TASK\_\allowbreak COUNT\_\allowbreak V\_\allowbreak INSTALL} & B & Construction: original finite CountV assembly.\\
\texttt{THM\_\allowbreak S6\_\allowbreak S} & C & Construction: unitary-three ordinary support and automorphy.\\
\texttt{HOST\_\allowbreak FRAME\_\allowbreak LENGTH} & B & Construction: actual host/frame with normalized lengths.\\
\texttt{HOST\_\allowbreak RIGID\_\allowbreak S} & D & Construction: original host's allowed tangents; retained open after ADK6.\\
\texttt{HOST\_\allowbreak NUMBER\_\allowbreak S} & C & Construction: uniform catalogue for the original rank-three hosts.\\
\texttt{TASK\_\allowbreak HOST\_\allowbreak VAN\_\allowbreak E\_\allowbreak INSTALL} & B & Construction: literal HostVanE from the supplied duality.\\
\texttt{TASK\_\allowbreak HOST\_\allowbreak XPP\_\allowbreak INSTALL} & B & Construction: literal HostXpp from the supplied CorX.\\
\texttt{DESIGN\_\allowbreak CHOICES\_\allowbreak COMMON} & C & Common geometry: design family and one original bundle.\\
\texttt{LOCAL\_\allowbreak NBUILT\_\allowbreak COMMON} & C & Common geometry: local specialization on that bundle.\\
\texttt{SECOND\_\allowbreak LAYER\_\allowbreak COMPATIBLE\_\allowbreak SELECTION} & C & Common geometry: compatible second-layer selector.\\
\texttt{BASE\_\allowbreak THETA\_\allowbreak READINGS} & C & Common geometry: actual base readings, maps, metrics and measures.\\
\texttt{BASE\_\allowbreak THETA\_\allowbreak RATIONAL\_\allowbreak TESTS} & C & Common geometry: simultaneous rationality tests for those maps.\\
\texttt{SLR\_\allowbreak BASE\_\allowbreak ROWS\_\allowbreak PRODUCER} & C & Common geometry: actual base row identities.\\
\end{longtable}

This is a map of a selected sufficient route, not a proof that the
route is logically minimal.  In particular the seven Case-III and
chain items are specific to the general Dirichlet extension on this
route; the ordinary target avoids that additional branch.  The
full exact scope of each leaf remains the contract in the v0.41
register, carried into the accompanying JSON without weakening.
ADK6 is the one whole-leaf discharge proposed by this packet.

\subsection{Source appendix for the conditional proof}

All formal-source references are pinned to
\texttt{7cbf7a0c32fca1badf05931879a37c3f8a2b03b4}.
The human proofs reused are
\texttt{window\_\allowbreak bridge/\allowbreak main.\allowbreak tex} (WB1--WB8),
\texttt{analytic\_\allowbreak envelope/\allowbreak analytic\_\allowbreak envelope.\allowbreak tex} (AQ1--AQ5),
and the AN2--AN3 and SC1 sections of
\texttt{RH\_\allowbreak readable\_\allowbreak reconstruction\_\allowbreak v0\_\allowbreak 41.\allowbreak tex}.
Their exact hashes, and the source definitions for the original
character, gap and window joins, are recorded in
\texttt{conditional\_\allowbreak main\_\allowbreak v42/\allowbreak contracts.\allowbreak json}.
The correspondence is
\texttt{Ch\allowbreak Balance} (CMP1),
\texttt{Ch\allowbreak Core}, \texttt{Gm\allowbreak AEq}, \texttt{Sig\allowbreak ALe} (CMP2),
\texttt{Core\allowbreak Gm} (CMP3), and
\texttt{Top\allowbreak Win}, \texttt{landau\_\allowbreak window}, \texttt{lemma\allowbreak L}
(CMP4--CMP5).  These references are provenance for the
reconstructed implications, not assumptions used to bypass either
unproved arithmetic input.


\section{The full polarized adjoint: a checked application of Thorne's theorem}
\label{sec:adk-match}

This section resolves the matching question called
\texttt{AD\_\allowbreak VANISH\_\allowbreak K\_\allowbreak MATCH} in the v0.41 register.  The conclusion is
the entire original \texttt{Ad\allowbreak Vanish\allowbreak K}, for every rank $n\ge2$; the
selected Kato application uses rank three.  The proof is a published
theorem application with explicit coefficient, normalization, action,
and local-condition comparisons.  It is independent of the unfinished
ordinary Hecke and host-construction arguments.  The separate
\texttt{HOST\_\allowbreak RIGID\_\allowbreak S} obligation is not closed by this result.

\subsection{The exact statement to be supplied}

Let $K$ be an imaginary quadratic field, let $\Omega$ be an algebraic
closure of $K$ containing the specified copy of $\mathbb Q$, let $\ell$
be any rational prime, and fix
$\iota:\overline{\mathbb Q}_{\ell}\simeq\mathbb C$.
Write $G_K=\operatorname{Gal}(\Omega/K)$,
$G_{\mathbb Q}=\operatorname{Gal}(\Omega/\mathbb Q)$, and
$\epsilon:G_K\to\overline{\mathbb Q}_{\ell}^{\times}$ for the actual
cyclotomic character, obtained by the specified inclusion
$\mathbb Z_\ell\hookrightarrow\overline{\mathbb Q}_\ell$.
Suppose $n\ge2$ and
\[
 r:G_K\longrightarrow\operatorname{GL}_n(\overline{\mathbb Q}_\ell)
 \tag{ADK.1}
\]
is continuous.  The original automorphy hypothesis says that there is
a regular algebraic cuspidal representation $\pi$ of
$\operatorname{GL}_n(\mathbb A_K)$, conjugate self-dual, such that $r$
is semisimple and has the following unramified characteristic
polynomials outside a finite set.  If $w$ has norm $q_w$, if $F_w$ is
\emph{arithmetic} Frobenius, and if
$\{\alpha_{w,1},\ldots,\alpha_{w,\allowbreak{} n}\}$ is its Satake multiset in the
source normalization, then
\[
 \det\bigl(T-\iota r(F_w^{-1})\bigr)
 =\prod_{i=1}^{n}
   \bigl(T-q_w^{(n-1)/2}\alpha_{w,\allowbreak{} i}\bigr).
 \tag{ADK.2}
\]
Thus (ADK.2) uses geometric Frobenius; it is not the formula obtained
by applying the same right side to $F_w$.
The remaining irreducibility hypothesis is
\[
                       r|_{\ker\epsilon}\text{ is irreducible}.
 \tag{ADK.3}
\]
It is a characteristic-zero assertion over the full infinite
cyclotomic kernel, not residual irreducibility or irreducibility after
adjoining just $\ell$th roots of unity.

Let $V=M_n(\overline{\mathbb Q}_\ell)$ and let $\rho$ be a
$\overline{\mathbb Q}_\ell$-linear representation of $G_{\mathbb Q}$
on $V$.  The original \texttt{Is\allowbreak Ad\allowbreak Of} hypothesis supplies a complex
conjugation $c\in G_{\mathbb Q}$, an embedding
$j:\Omega\hookrightarrow\mathbb C$ inducing that conjugation, and
$P\in\operatorname{GL}_n(\overline{\mathbb Q}_\ell)$, such that
\begin{align}
 r(s)^tP r(csc^{-1})&=\epsilon(s)^{1-n}P &&(s\in G_K),
                 \tag{ADK.4}\\
 \rho(s)X&=r(s)Xr(s)^{-1} &&(s\in G_K),
                 \tag{ADK.5}\\
 \rho(c)X&=-P^{-1}X^tP &&(X\in V).
                 \tag{ADK.6}
\end{align}
The intended conclusion is that every continuous $\rho$-cocycle
$a:G_{\mathbb Q}\to V$ satisfying the following two conditions is a
coboundary:
\begin{enumerate}
\item At every finite $p\ne\ell$, its restriction to inertia is
  $s\mapsto\rho(s)v_p-v_p$ for some $v_p\in V$.
\item At $\ell$, after identifying the local decomposition group
  using any permitted embedding into the fixed local algebraic
  closure, $a(s)\otimes1=\rho_{\mathrm{cris}}(s)b-b$ for some
  $b\in V\otimes_{\mathbb Q_\ell}B_{\mathrm{cris}}$.
\end{enumerate}
Here the action on the tensor product is the diagonal action, including
the Galois action on $B_{\mathrm{cris}}$.  These are the literal
\texttt{Is\allowbreak BK.\allowbreak away} and \texttt{Is\allowbreak BK.\allowbreak above} conditions, as verified
below.  In particular the conclusion is about all of $M_n$, including
its scalar matrices.

% BEGIN ADK1
\begin{lemma}[ADK1: one finite coefficient field for a representation and a cocycle]
\label{claim:ADK1}
A continuous finite-dimensional representation of a profinite group
over $\overline{\mathbb Q}_\ell$ has all its matrix coefficients in
one finite extension of $\mathbb Q_\ell$.  Once its coefficients are
in $E$, each continuous cocycle has all its values in a finite
extension $E'/E$.  Any finitely many fixed matrices may be included
in $E'$.
\end{lemma}
\begin{proof}
NR19 gives a countable cofinal family of finite subfields of
$\overline{\mathbb Q}_\ell$.  Briefly, approximate any algebraic
$\ell$-adic number by an element of a fixed countable dense copy of
$\overline{\mathbb Q}$ more closely than its distance to its other
conjugates.  The nonarchimedean triangle inequality makes every
automorphism fixing the approximant fix the original number.  The
fixed-field theorem then puts that number in the finite local field
generated by the approximant.  Taking finite subsets gives the claimed
cofinal family; taking successive composita makes it increasing.
Each finite local field is complete, hence closed in the algebraic
closure equipped with its $\ell$-adic topology.

For a continuous representation $u$, the inverse images
$u^{-1}(\operatorname{GL}_d(E_i))$ are closed subgroups covering the
compact group.  Baire's theorem makes one of them have nonempty
interior, hence be open.  It has finite index.  Adjoining the entries
of $u$ at finitely many coset representatives captures every matrix
coefficient in a finite field.  Adjoin the entries of any additional
fixed matrices as well.

If the action is already defined over $E$, enlarge every $E_i$ to
contain $E$.  For a cocycle $a$, the sets
$H_i=\{g:a(g)\in E_i^d\}$ are again closed subgroups: the identities
\[
 a(gh)=a(g)+u(g)a(h),\qquad
 a(g^{-1})=-u(g^{-1})a(g)
\]
show closure under product and inverse.  Apply the same Baire and
finite-coset argument.  This argument does not assume that an arbitrary
compact subset of the algebraic closure is contained in a finite
extension; it uses the group or cocycle laws essentially.
\end{proof}
% END ADK1

% BEGIN ADK2
\begin{lemma}[ADK2: the exact infinite cyclotomic subgroup]
\label{claim:ADK2}
With the original embeddings,
\[
 \ker\epsilon=G_{K(\mu_{\ell^\infty})}.
 \tag{ADK.7}
\]
Consequently (ADK.3) is exactly the irreducibility hypothesis in
Thorne's Theorem~6.1.
\end{lemma}
\begin{proof}
The $\ell$-adic cyclotomic character is the compatible system of
actions on $\mu_{\ell^m}$, for all $m\ge1$.  The inclusion of
$\mathbb Z_\ell$ into the coefficient algebraic closure is
injective.  Thus $\epsilon(s)=1$ if and only if $s$ fixes every such
root of unity, which is equivalent to fixing their compositum with
$K$.  The induced topology is the closed-subgroup topology on both
sides.  This identifies the actual infinite subgroup, without taking
a finite cyclotomic approximation.  Irreducibility on this subgroup
also implies irreducibility on $G_K$.
\end{proof}
% END ADK2

% BEGIN ADK3
\begin{lemma}[ADK3: the normalized attached representation is the given one]
\label{claim:ADK3}
The automorphic representation in (ADK.2) is polarizable, and its
standard normalized Galois representation $r_{\pi,\iota}$ is
isomorphic to the given $r$.  The polarization has multiplier
$\epsilon^{1-n}$ on $G_K$ and a symmetric pairing at the specified
complex conjugation.
\end{lemma}
\begin{proof}
We first unpack the automorphic dictionary; no conclusion is inferred
from a predicate's name.  In
\texttt{Automorphic\allowbreak Membership/\allowbreak Expo\allowbreak Conj\allowbreak Of\allowbreak Is\allowbreak Csd.\allowbreak lean},
\texttt{Cusp\allowbreak Sub} is a nonzero irreducible submodule of smooth
functions on $\operatorname{GL}_n(\mathbb A_K)$, stable under the
finite adelic group, the maximal compact group at infinity, and the
complexified Lie algebra.  Its vectors are left
$\operatorname{GL}_n(K)$-invariant, have an open finite-adelic
stabilizer, are finite under the maximal compact and the center of
the enveloping algebra, and have moderate growth.  The literal
\texttt{Is\allowbreak Cuspidal} tests the maximal standard parabolics indexed
by $0<r<n$.  Vanishing for every proper standard parabolic follows
by integrating first over a containing maximal parabolic's
unipotent radical and then over the remaining quotient; the
unipotent rational quotients are compact, so Fubini applies.
Thus the integral of each vector over every proper standard
parabolic unipotent radical is zero.  These are the usual cuspidal
automorphic $(\mathfrak g,\allowbreak{} K_\infty)\times G(\mathbb A_f)$ modules;
the standard smooth, admissible realization of an irreducible
cuspidal automorphic representation is precisely this realization.
\texttt{CG.\allowbreak mk} passes to its intertwining-isomorphism class and
changes none of these actions.

There are two normalization checks in this dictionary.
First, the source's central polynomial is
\[
 \operatorname{cap}_{\tau}(u)
    =\operatorname{qdet}(E_{\tau},\allowbreak{} u+(n-1)/2).
\]
On a highest-weight vector of weight $(\lambda_0,\ldots,
\lambda_{n-1})$ it has eigenvalue
$\prod_{i=0}^{n-1}(u+\lambda_i+(n-1)/2-i)$;
this is the literal \texttt{cap\_\allowbreak highest} formula in
\texttt{Automorphic\allowbreak Membership/\allowbreak Continuous\allowbreak Conj\allowbreak Fin.\allowbreak lean},
lines 147--157.  Thus its roots record the ordinary
Harish--Chandra parameters with the usual half-sum-of-positive-roots
shift.  \texttt{CG.\allowbreak Is\allowbreak Reg\allowbreak Alg}, through
\texttt{is\allowbreak Reg\allowbreak Alg\_\allowbreak iff}, requires for every embedding $\tau$ that
these $n$ parameters are distinct and belong to
$\mathbb Z+(n-1)/2$.  Sorting them as
$z_0>\cdots>z_{n-1}$ gives the integral dominant weight
$\lambda_i=z_i-(n-1)/2+i$: indeed
$\lambda_i-\lambda_{i+1}=z_i-z_{i+1}-1\ge0$.
Conversely every integral dominant weight produces such distinct
parameters.  This is exactly regular algebraicity in BLGGT,
\S2.1, for both complex embeddings.  The source defines a fallback
multiset consisting of $n$ zeros when an infinitesimal character
does not exist; since $n\ge2$, distinctness excludes that fallback.
Consequently the regularity hypothesis really does supply the
stated central eigencharacter.

Second, \texttt{CG.\allowbreak Has\allowbreak Sat} in
\texttt{Galois\allowbreak Cohomology2/\allowbreak Strong\allowbreak Multiplicity\allowbreak One.\allowbreak lean},
lines 61--66, descends the predicate
\texttt{Has\allowbreak Sat\allowbreak GL} of the actual local representation.  With
$\operatorname{vol}(\operatorname{GL}_n(\mathcal O_w))=1$, it
says that the spherical fixed space is nonzero and that, for every
nonzero fixed vector and every $0\le j\le n$, the characteristic
function of the double coset of
$\operatorname{diag}(\varpi I_j,\allowbreak{} I_{n-j})$ acts by
\[
             q_w^{j(n-j)/2} e_j(\alpha_{w,1},\ldots,
                                    \alpha_{w,\allowbreak{} n}).
 \tag{ADK.9}
\]
This is obtained by unfolding \texttt{Is\allowbreak Sat\allowbreak Vec\allowbreak GL} and
\texttt{sat\allowbreak Eig}, not by selecting a different local constituent.
The factor in (ADK.9) is exactly that of normalized Satake
parameters, or equivalently normalized unramified induction.
The elementary symmetric functions determine the multiset, so the
source parameters and the parameters used by normalized local
Langlands agree.  The source's further \texttt{alg\allowbreak Sat} multiplies
each of them by $q_w^{(n-1)/2}$, as displayed in (ADK.2).

Finally \texttt{Dual\allowbreak Is\allowbreak Conj} is a nondegenerate complex bilinear
pairing $B$ on the cusp module with
$B(\pi(g)v,\pi(g^c)w)=B(v,\allowbreak{} w)$ for finite adelic and compact
elements, and
$B(Xv,\allowbreak{} w)+B(v,\allowbreak{} X^cw)=0$ for Lie algebra elements.
Here \texttt{conj\allowbreak Op} is entrywise CM conjugation at finite places,
and the corresponding conjugate actions on compact and Lie algebra
elements.  The map $v\mapsto B(v,\cdot)$ intertwines $\pi$ with
the contragredient of $\pi^c$.  It lands in the smooth,
$K_\infty$-finite contragredient: a compact open subgroup fixing
$v$ fixes its functional by the first pairing identity, and the
$K_\infty$-span of that functional is the image of the finite
$K_\infty$-span of $v$.  Nondegeneracy and irreducibility make this
an isomorphism of admissible irreducible modules.  This checks
$\pi^c\simeq\pi^\vee$, with a bilinear pairing rather than an
unintended Hermitian one.

Let $\delta_K$ denote the quadratic Hecke character of
$K/\mathbb Q$.  Its pullback by the idele norm from $K$ is trivial,
and its infinity value at $-1$ is $-1$.  Since
$\pi^c\simeq\pi^\vee$, the pair
$(\pi,\delta_K^n)$ is polarized with the precise infinity sign
$(-1)^n$ required in BLGGT, \S2.1.  Regular algebraicity and
cuspidality are the original hypotheses, not new conclusions of a
finite coefficient argument.

BLGGT, Theorem~2.1.1, attaches a continuous semisimple
$r_{\pi,\iota}$ and identifies its local Weil--Deligne representation
away from $\ell$ with that of
$\pi_w\otimes|\det|_w^{(1-n)/2}$.  At an unramified place the
geometric Frobenius eigenvalues are therefore exactly
$q_w^{(n-1)/2}\alpha_{w,\allowbreak{} i}$.  This is (ADK.2), with the same
$\iota$ and the same direction of Frobenius.

By ADK1 put both continuous representations over one finite
coefficient field.  Chebotarev density makes equality of their
unramified Frobenius traces outside a finite set imply equality of
their continuous traces on $G_K$.  The characteristic-zero
semisimple character criterion then supplies an isomorphism.  One
elementary proof of that criterion uses the algebra generated by the
two actions on their direct sum: its radical kills a semisimple
module, and its simple matrix factors have independent trace
characters.  Equality of traces gives equal multiplicities of all
simple factors.  Thus we may transport $r_{\pi,\iota}$ to the
original space carrying $r$.

The Galois multiplier in BLGGT is
$\delta_K^n\epsilon^{1-n}$ on $G_{\mathbb Q}$.  Its restriction to
$G_K$ is $\epsilon^{1-n}$; its value at complex conjugation is
$(-1)^n(-1)^{1-n}=-1$.  The total oddness assertion in that theorem
and the polarization convention in \S2.1 therefore give a symmetric
nondegenerate matrix $P_0$ with
$r(s)^tP_0r(csc^{-1})=\epsilon(s)^{1-n}P_0$.
Every complex conjugation arising from an embedding of $\Omega$ in
$\mathbb C$ is allowed in this assertion, including the one in
(ADK.4).  The quadratic character used for the infinity sign is
trivial on $G_K$ and changes neither the attached representation nor
its Frobenius polynomials.
\end{proof}
% END ADK3

% BEGIN ADK4
\begin{lemma}[ADK4: equality of the full polarized adjoint actions]
\label{claim:ADK4}
The original $\rho$ in (ADK.5)--(ADK.6) is the full unitary adjoint
$\operatorname{ad}r_{\pi,\iota}$ in Thorne's theorem.
\end{lemma}
\begin{proof}
The matrices $P$ and $P_0$ satisfy the same intertwining equation.
Dividing their equations shows that $P^{-1}P_0$ commutes with
$r(cG_Kc^{-1})=r(G_K)$.  This representation is irreducible over
the algebraically closed coefficient field by ADK2, so Schur's lemma
makes $P_0$ a nonzero scalar multiple of $P$.  Hence
\[
              -P_0^{-1}X^tP_0=-P^{-1}X^tP.
 \tag{ADK.8}
\]
For clarity, the minus-transpose formula is the differential of the
disconnected component of the standard group
$\mathcal G_n=(\operatorname{GL}_n\times\operatorname{GL}_1)
\rtimes\{1,\allowbreak{} j\}$, where
$j(g,\allowbreak{} a)j^{-1}=(a g^{-t},\allowbreak{} a)$; see BLGGT, \S1.1.
The differential of $g\mapsto g^{-t}$ at the identity is
$X\mapsto-X^t$.  With the pairing $P_0$, conjugation by the
corresponding component gives $X\mapsto-P_0^{-1}X^tP_0$.
The connected subgroup acts by $X\mapsto r(s)Xr(s)^{-1}$.
Thus the two actions agree on $G_K$ and on $c$.  Since $K$ is
imaginary quadratic, those elements generate $G_{\mathbb Q}$,
proving equality on the entire original group.

Thorne's adjoint acts on $\operatorname{Lie}\operatorname{GL}_n=M_n$;
it is not restricted to trace-zero matrices.  In particular $c$
acts by $-1$ on the scalar line.  No trace-zero substitution, omitted
quadratic line, or division by $n$ is used.  The action is continuous:
on each of the two open cosets of $G_K$ it is expressed by the
continuous matrices of $r$ and the fixed matrix $P$.
\end{proof}
% END ADK4

% BEGIN ADK5
\begin{lemma}[ADK5: the original Bloch--Kato condition, including coefficients]
\label{claim:ADK5}
For the representation $\rho$ of ADK4, the source's
\texttt{Bloch\allowbreak Kato\allowbreak Selmer\allowbreak Vanishes} is precisely vanishing of the usual
Bloch--Kato Selmer group over $\mathbb Q$, with its literal
algebraic-closure coefficient and cocycle scope.
\end{lemma}
\begin{proof}
The original \texttt{Z1} consists of continuous cocycles.  At
$p\ne\ell$, the assertion that a cocycle is a coboundary on inertia
is exactly membership in the unramified local subgroup
$\ker(H^1(G_{\mathbb Q_p},\allowbreak{} V)\to H^1(I_p,\allowbreak{} V))$.
Changing the stored valuation above $p$ conjugates the decomposition
and inertia groups.  For a global cocycle the change of restriction
is conjugation followed by a coboundary, so this kernel condition is
unchanged.  Thus the existential choice of a place in
\texttt{Is\allowbreak BK.\allowbreak away} is the usual local condition.

At $\ell$, \texttt{dec\allowbreak J} is the graph of the restriction map from
the local absolute Galois group along the chosen embedding
$\Omega\hookrightarrow\overline{\mathbb Q}_\ell$.  The second
projection identifies this graph with the local group.  Indeed the
image of $\Omega$ is the subfield of elements algebraic over the
specified global field, which is stable under local automorphisms;
restriction determines the first component uniquely.  This
identification is topological, too.  A basic open subgroup of the
global Galois group fixes a finite set $S\subset\Omega$.
Its inverse image under local restriction fixes the finite extension
$\mathbb Q_\ell(j(S))$, and is therefore open in the local Galois
group.  Restriction is continuous; hence the inverse to the second
projection, $s\mapsto(\operatorname{res}_j(s),\allowbreak{} s)$, is continuous.
The second projection is itself continuous as a restricted product
projection.  In particular a continuous global cocycle pulls back
to the usual continuous local cocycle.  On the graph,
\texttt{act\allowbreak C} is precisely the diagonal action on
$V\otimes_{\mathbb Q_\ell}B_{\mathrm{cris}}$.
The original \texttt{finite\allowbreak Cocycles} requires a tensor primitive,
which says exactly that the local cohomology class dies after this
tensor extension.  BX1 identifies the original tied period-ring
model with the usual $B_{\mathrm{cris}}$, equivariantly and with
the stated cyclotomic character.  This proves the local comparison
without replacing finite cohomology by a larger de Rham condition.
At the real place positive-degree cohomology is zero by averaging
over a group of order two in characteristic zero.

Finally, ADK1 gives a finite field $E$ containing all coefficients
of the action and all values of an individual continuous cocycle.
Every local algebraic-closure primitive uses finitely many vector
coefficients, so lies over some finite $E'/E$.  An $E$-linear
retraction $q:E'\to E$ with $q(1)=1$ commutes with the $E$-defined
action.  Apply it to an inertia primitive, or apply
$q\otimes1$ to a crystalline tensor primitive, to return it to
$E$.  Explicitly, for the coordinatewise retraction $q_V$, a pure
tensor $v\otimes b$ and a local element $s$ satisfy
\[
 (q_V\otimes1)\bigl(\rho(s)v\otimes s(b)\bigr)
     =\rho(s)q_V(v)\otimes s(b)
     =s\bigl((q_V\otimes1)(v\otimes b)\bigr).
\]
Linearity extends this equality to every tensor; it uses the
$E$-linearity of $q_V$ and the identity on the period-ring factor,
so includes the nontrivial $B_{\mathrm{cris}}$ action.  The same map
fixes $a(s)\otimes1$ since $a(s)\in V_E$.
This is done separately at each place; it does not posit one
finite field containing infinitely many selected local primitives.
Therefore the source cocycle belongs to $H^1_f(\mathbb Q,\allowbreak{} V_E)$.
Conversely a finite-coefficient Selmer cocycle remains one after
scalar extension.  A coboundary in either presentation has the
literal form $g\mapsto\rho(g)v-v$, as required by the source.
\end{proof}
% END ADK5

% BEGIN ADK6
\begin{theorem}[ADK6: the full original \texttt{Ad\allowbreak Vanish\allowbreak K}]
\label{claim:ADK6}
Under (ADK.1)--(ADK.6), every source Bloch--Kato cocycle is a
coboundary.  This proves \texttt{Ad\allowbreak Vanish\allowbreak K} with all its original
quantifiers, in particular for rank three and for every prime
allowed there.
\end{theorem}
\begin{proof}
We use Thorne, \emph{On the vanishing of adjoint Bloch--Kato Selmer
groups of irreducible automorphic Galois representations},
Theorem~6.1 (arXiv:2207.04925v1, printed p.~24): a regular algebraic
cuspidal polarized representation over a CM field has vanishing
unitary adjoint Selmer group when its attached representation
remains irreducible over the infinite cyclotomic extension.
Apply it to $F=K$, $F^+=\mathbb Q$, the same $n$, prime $\ell$ and
embedding $\iota$, and the polarized pair of ADK3.  Its
irreducibility hypothesis is exactly ADK2.  Its attached
representation and full adjoint are the original ones by ADK3--ADK4.
Thus $H^1_f(\mathbb Q,\rho)=0$.  ADK5 transfers this statement to
every literal source cocycle, proving the claimed endpoint.

There is no residual-image assumption, bound $\ell>n$, exclusion of
$\ell=3$, ordinary hypothesis, or hypothesis that every place splits
in $K$ in the cited theorem.  The split/unramified/Iwahori conditions
introduced later inside Thorne's proof are reductions in that proof,
not extra assumptions of Theorem~6.1.  This is why the application
does not import the much larger ordinary-lifting package used by a
different source route.
\end{proof}
% END ADK6

\subsection{The separate host rigidity question}

The remaining host endpoint is now a located reconstruction task,
not an unidentified reference.  The exact source is
\texttt{EHost\allowbreak Rigid\allowbreak U.\allowbreak host\allowbreak Rigid\allowbreak UE\_\allowbreak alt}; its final assembly occupies
lines 390--440.  Its conclusion concerns the original quotient
$\Gamma=G_{\mathbb Q,\allowbreak{} S}$, the original marked element $g_1$, a point
which becomes integral after the specified torus conjugation, and
every dual-number deformation satisfying its prescribed local flags.
The source representation is the full adjoint of
$r:G_{\mathbb Q,\allowbreak{} S}\to\operatorname{GL}_3\rtimes\{\pm1\}$, pulled
back by the actual quotient map.

Its \texttt{BPrime\_\allowbreak s.\allowbreak cyc} already states irreducibility on
$\ker(\eta_4)\cap\ker(\epsilon)$, and
$\ker(\eta_4)=G_{\mathbb Q(i)}$.  Its automorphy field supplies the
cyclotomic twist $r|_{G_{\mathbb Q(i)}}\otimes\epsilon^{-1}$.
The twist has trivial conjugate-pair product after its prescribed
multiplier adjustment, and scalar twists do not change the connected
adjoint action.  A full matching proof must still write the
disconnected action on the same pairing.  More importantly, global
Selmer vanishing yields the host's rigidity only after showing that
each allowed tangent cocycle satisfies the original local finite
conditions.

The away-from-$\ell$ part is finite-index averaging and the literal
diagonal-conjugacy condition.  The substantive local step is
\texttt{rig\_\allowbreak above}, lines 145--191 of \texttt{EHost\allowbreak Rigid\allowbreak U.\allowbreak lean}:
it uses \texttt{rig\_\allowbreak sub} and \texttt{rig\_\allowbreak eig} to obtain an actual
open subgroup, fixed inertia weights $(e,0,-e)$, and the required
non-cyclotomic eigenvalue ratios before applying
\texttt{Rig\allowbreak Loc.\allowbreak rig\_\allowbreak loc}.  Those hypotheses are not automatic from
the phrase ``ordinary''.  \texttt{Rig\allowbreak Fin.\allowbreak rig\_\allowbreak fin} and
\texttt{Rig\allowbreak Desc.\allowbreak rig\_\allowbreak desc} give the finite-coefficient local
crystalline-extension and descent arguments.  Once they are written
for these same objects, the quotient map's surjectivity descends the
global coboundary back to the original tangent cocycle; its sign must
agree with the source's \texttt{cob} convention.

Accordingly ADK6 discharges \texttt{AD\_\allowbreak VANISH\_\allowbreak K\_\allowbreak MATCH}, subject to
integration review.  It does not discharge \texttt{HOST\_\allowbreak RIGID\_\allowbreak S},
\texttt{RIG\_\allowbreak II}, ordinary support, or the existence of any
automorphic host.  The suggested next host packet is the complete
\texttt{rig\_\allowbreak sub}--\texttt{rig\_\allowbreak eig}--\texttt{rig\_\allowbreak loc}
comparison plus the explicit twist/pairing identification.

\subsection{References and literal source correspondence}

The source pin throughout is
\texttt{7cbf7a0c32fca1badf05931879a37c3f8a2b03b4}.
\texttt{Automorphic\allowbreak Forms/\allowbreak Congruence\allowbreak Number\allowbreak Rank\allowbreak Three.\allowbreak lean},
lines 158--175, gives \texttt{Is\allowbreak Ad\allowbreak Of} and \texttt{Ad\allowbreak Vanish\allowbreak K}.
\texttt{Automorphic\allowbreak Forms/\allowbreak Bloch\allowbreak Kato\allowbreak Selmer\allowbreak Vanishes.\allowbreak lean},
lines 69--86 and 117--144, gives the normalized attachment,
automorphy and exact local Selmer definitions.
\texttt{Padic\allowbreak Hodge\allowbreak Theory/\allowbreak Finite\allowbreak Cohomology.\allowbreak lean},
lines 122--151, gives the tensor action and finite cocycles.
These paths are relative to
\texttt{src/\allowbreak Development/\allowbreak Automorphic\allowbreak Arithmetic/\allowbreak Automorphic\allowbreak Varieties/\allowbreak }.
The accompanying JSON records the full paths, ranges and hashes.

The published statements used are:
\begin{enumerate}
\item J.~A. Thorne, \emph{On the vanishing of adjoint Bloch--Kato
  Selmer groups of irreducible automorphic Galois representations},
  Pure Appl. Math. Q. \textbf{18} (2022), no.~5, Theorem~6.1;
  \url{https://arxiv.org/html/2207.04925v1} and
  \url{https://arxiv.org/pdf/2207.04925}.  The full adjoint is
  specified on printed pp.~3--4; the theorem is on p.~24.
\item T.~Barnet-Lamb, T.~Gee, D.~Geraghty and R.~Taylor,
  \emph{Potential automorphy and change of weight},
  Ann. of Math. \textbf{179} (2014), 501--609.
  Section~1.1, pp.~512--513, defines $\mathcal G_n$;
  Section~2.1, p.~535, specifies the pairing and sign;
  Theorem~2.1.1, p.~538, supplies normalized attachment and
  total oddness.
  \url{https://annals.math.princeton.edu/wp-content/uploads/annals-v179-n2-p03-p.pdf}.
\item Chebotarev density and the characteristic-zero semisimple
  character criterion, used with both representations in one finite
  coefficient field.  The Chebotarev application is the same standard
  one retained in RX5; the character argument is recalled in ADK3.
\end{enumerate}


\section{Host rigidity from the literal adjoint and local deformation conditions}
\label{sec:host-rigidity-v42}

The remaining question called \texttt{HOST\_\allowbreak RIGID\_\allowbreak S} is a
reconstruction matching question: do the particular deformations in
the source give Bloch--Kato classes for the particular polarized
adjoint to which ADK6 applies?  This section answers that question.
The pairing is the source's antidiagonal matrix, the scalar twist is
exactly $\epsilon^{-1}$, and the local argument includes arbitrary
dual-number changes of basis and unramified diagonal variations.
The conclusion does not construct the automorphic points or prove
the separate quantitative host-number bound.

\subsection{The original hypotheses, in ordinary notation}

Fix an algebraic closure $\Omega$ of $\mathbb Q$, an odd prime
$\ell$, and a finite $\mathbb Z_\ell$-subalgebra
$O\subset\overline{\mathbb Q}_\ell$ which is a discrete valuation
ring.  Let $N\ne0$ be even, let
$\psi:(\mathbb Z/N\mathbb Z)^\times\to O^\times$ be a Dirichlet
character, and let $e,\allowbreak{} o\ge1$.  For a prime $p$, write $\psi_p$ for
the CRT component of $\psi$ supported at the $p$-part of the modulus.
This is the source's \texttt{comp p}; it is different from the
decomposition according to the order of the character.  Write
$\psi=\theta\kappa$, where $\theta$ has order prime to $\ell$
and $\kappa$ has $\ell$-power order.  These are
\texttt{th0} and \texttt{kap}.  Put
\[
 c_2=v_2(\operatorname{cond}\psi_2),\qquad
 S_b=\{p\mid N:p\ne\ell,\ \kappa_p\ne1\},\qquad
 T_b=\{\ell\}\cup S_b.
 \tag{HSR.1}
\]
The complete \texttt{Host\allowbreak Data} assumptions are
\[
 \psi(-1)=(-1)^e,\quad c_2\ge4,\quad
 \theta_p^{\,2o}=1\ (p\ne2),\quad
 8op^2<2^{c_2}\ (p\in T_b).
 \tag{HSR.2}
\]
The prime-to-$\ell$ and $\ell$-power order parts exist because
$\psi$ is a character of a finite group.  The explicit exponent
in the source's definition of \texttt{l\allowbreak Part} is the CRT projector
on the $\ell$-primary subgroup of the finite cyclic group generated
by $\psi$.

Let $S$ be exactly the finite primes dividing $\ell N$, let
$\Gamma=G_{\mathbb Q,\allowbreak{} S}$, and let
$\pi:G_{\mathbb Q}\twoheadrightarrow\Gamma$ be the specified
quotient.  For each $p\in S$ use the stored valuation above $p$,
with decomposition and inertia groups $D_p,\allowbreak{} I_p\subset G_{\mathbb Q}$.
The character $\eta:\Gamma\to\{\pm1\}$ and the character
$\Psi:\Gamma\to O^\times$ satisfy
\[
 \eta\pi=\eta_4,\qquad
 \Psi\pi=\epsilon^e\psi_{\mathrm{Gal}}.
 \tag{HSR.3}
\]
Here $\eta_4$ is the quadratic character of $\mathbb Q(i)$,
and $\epsilon$ is the actual $\ell$-adic cyclotomic character.
A stored element $c\in\Gamma$ is the image of an actual complex
conjugation $c_\infty\in G_{\mathbb Q}$.  This is precisely the
content of \texttt{Is\allowbreak Host\allowbreak Frame}, including its choices of valuations,
quotient map, distinguished $\ell$-place, and subset $S_b$.
An actual complex conjugation sends $i$ to $-i$, so
$\eta(c)=\eta_4(c_\infty)=-1$.

Let
\[
 J=\begin{pmatrix}0&0&1\\0&1&0\\1&0&0\end{pmatrix},\qquad
 \vartheta(a)=Ja^{-t}J,
 \qquad \mathcal G_3=\operatorname{GL}_3\rtimes\{\pm1\},
 \tag{HSR.4}
\]
where $-1$ acts by $\vartheta$.  The source's group
\texttt{GG 3} is exactly $\mathcal G_3$.  If
$R(g)=(a(g),\eta(g))$, its adjoint action on the full space
$V=M_3(\overline{\mathbb Q}_\ell)$ is
\[
 \operatorname{Ad}_R(g)X=
 \begin{cases}
 a(g)Xa(g)^{-1},&\eta(g)=1,\\
 -a(g)JX^tJa(g)^{-1},&\eta(g)=-1.
 \end{cases}
 \tag{HSR.5}
\]

Fix the marked element $g_1$ required by the source and a point
$r:\Gamma\to\mathcal G_3(\overline{\mathbb Q}_\ell)$ satisfying
\texttt{Is\allowbreak Char\allowbreak Zero\allowbreak Pt}.  Thus there are a finite $O$-subalgebra
$A\subset\overline{\mathbb Q}_\ell$ which is a DVR, a unit
$t\in\overline{\mathbb Q}_\ell^\times$, and an integral point
$r_A:\Gamma\to\mathcal G_3(A)$ such that
\[
 r_A=\operatorname{conj}_{T_t}r,
 \qquad T_t=\operatorname{diag}(t,1,\allowbreak{} t^{-1}).
 \tag{HSR.6}
\]
Scalar inclusions in (HSR.6) are understood.  The integral point has
multiplier $\eta$, has continuous matrix entries modulo every power
of $\mathfrak m_A$, has $r_A(c)^{\rm mat}=1$, and satisfies
\[
 r_A(g)^{\rm mat}\equiv
 \operatorname{diag}(\Psi(g),1,\Psi(g)^{-1})\pmod{\mathfrak m_A}
 \qquad(g\in\Gamma).
 \tag{HSR.7}
\]
The remaining normalization and local conditions on this integral
point are retained from the source, although the proof below does
not need to use each of them separately.

For clarity, here is the full conclusion \texttt{Rigid r} to be
proved.  Let $z:\Gamma\to V$ be any continuous cocycle for (HSR.5),
and let $r_\delta$ be any deformation over
$\overline{\mathbb Q}_\ell[\delta]/(\delta^2)$ with the same
multiplier and
\[
 r_\delta(g)^{\rm mat}=(1+\delta z(g))r(g)^{\rm mat}.
 \tag{HSR.8}
\]
For $p\notin T_b$ assume that $z$ vanishes on
\[
 N_p=I_p\cap\ker\eta_4\cap\ker\overline{\Psi\pi}.
 \tag{HSR.9}
\]
At $\ell$, assume that some dual-number basis $h$ makes $r_\delta$
upper triangular on $D_\ell\cap\ker\eta_4$, with inertia diagonal
\[
                  (\Psi,1,\Psi^{-1}).
 \tag{HSR.10}
\]
The right side of (HSR.10) has no $\delta$ term.  At every
$p\in S_b$, assume that such a basis makes the restriction diagonal
on $D_p\cap\ker\eta_4$, again with (HSR.10) on inertia.  These
are literally \texttt{Tri\allowbreak In} and \texttt{Diag\allowbreak In}; their bases are
not required to be constant or close to the identity.  We must find
one $X\in V$ with
\[
                  z(g)=X-\operatorname{Ad}_r(g)X
                  \qquad(g\in\Gamma).
 \tag{HSR.11}
\]
The sign in (HSR.11) is the source's \texttt{cob} convention.

\subsection{Integral points and the local separating element}

% BEGIN HSR1
\begin{lemma}[HSR1: continuity, normalization, and finite coefficients]
\label{claim:HSR1}
The point $r$ has continuous matrix entries, has multiplier $\eta$,
and satisfies $r(c)=(1,-1)$.  Its adjoint action is continuous.
For every cocycle and deformation in (HSR.8), all matrix coefficients
of $r,\allowbreak{} z,\allowbreak{} r_\delta$, together with any fixed finite set of dual-number
bases, lie in one finite extension of $\mathbb Q_\ell$.
\end{lemma}
\begin{proof}
Every element of $A$ is integral over $\mathbb Z_\ell$, hence has
absolute value at most one: otherwise the leading term of a monic
integral equation would have strictly larger absolute value than
all its other terms.  Every unit of $A$ has absolute value one,
by applying the same assertion to it and its inverse.  If $\varpi$
is a uniformizer, write $\ell=u\varpi^b$.  Here $b\ge1$, since
$\ell$ has absolute value less than one and cannot be a unit of
$A$.  Thus $q=|\varpi|=\ell^{-1/b}<1$, and an element of
$\mathfrak m_A^n$ has absolute value at most $q^n$.  Local constancy
modulo $\mathfrak m_A^n$ for every $n$ now proves norm continuity
of every entry of $r_A$, hence of $r$ under (HSR.6).

Conjugation does not change the multiplier.  Also
$\vartheta(T_t)=T_t$.  Since $\eta(c)=-1$, the matrix part of
$T_t r(c)T_t^{-1}$ is consequently
$T_t r(c)^{\rm mat}T_t^{-1}$.  Its equality to $1$ proves
$r(c)^{\rm mat}=1$.  The character $\eta$ is locally constant:
its composition with the quotient map is the continuous finite
character $\eta_4$, and the quotient topology gives continuity of
$\eta$.  Formula (HSR.5), on its two open cosets, proves continuity
of the adjoint action.

The fraction field of $A$ is finite over $\mathbb Q_\ell$ because
$A$ is finite over $O$ and $O$ is finite over $\mathbb Z_\ell$.
Adjoining $t$ gives a finite field containing every entry of $r$ and
its adjoint action.  ADK1's cocycle argument then puts all values of
the chosen $z$ in a finite extension.  Formula (HSR.8) includes the
deformation, and adjoining the finitely many entries of the chosen
bases includes those bases and their inverses.

For the more general \texttt{Rig\allowbreak Desc} interface one can start
directly with a continuous dual-number representation on any compact
group.  Take an increasing countable cofinal family $E_i$ of finite
local coefficient fields, as in NR19 and ADK1.  The sets where all
constant and $\delta$ matrix coefficients belong to $E_i$ are
closed subgroups and cover the group.  They are subgroups because
the semidirect-product operations are defined over $E_i[\delta]$;
a matrix over this ring which is invertible over the larger ring
is already invertible over it, since its constant determinant is
nonzero.  Baire's theorem makes one subgroup open.  Adjoining the
coefficients at its finitely many coset representatives and of the
specified basis gives the asserted finite field.  The topology on
each finite field is its subspace topology, so all descended matrix
maps remain continuous.  This proves exactly the coefficient
descent used in \texttt{Rig\allowbreak Desc}, rather than assuming that arbitrary
compact subsets of the algebraic closure have finite coefficients.
\end{proof}
% END HSR1

% BEGIN HSR2
\begin{lemma}[HSR2: the dyadic conductor forces separation after reduction]
\label{claim:HSR2}
Suppose $u\in(\mathbb Z/N\mathbb Z)^\times$ has $2$-primary
coordinate equal to $\ell$.  Under (HSR.2),
\[
                         \overline{\psi(u)}^{\,4}\ne1.
 \tag{HSR.12}
\]
\end{lemma}
\begin{proof}
Put $s=v_2(\ell^2-1)\ge3$, write $o=2^jm$ with $m$ odd, and
put $a=v_2(N)$ and $d=s+j+1$.  Squaring $1+2^sb$ with $b$ odd
raises the exact $2$-adic valuation of its difference from $1$ by
one when $s\ge2$; raising to an odd power does not change it.
Induction gives
\[
 v_2(\ell^{4o}-1)=d,\qquad
 2^d\le2o(\ell^2-1)<8o\ell^2<2^{c_2}.
 \tag{HSR.13}
\]
Thus $d<c_2\le a$.

For completeness, an element $v=1+2^d b$ with $b$ odd and $d\ge2$
generates the group $1+2^d\mathbb Z/2^a\mathbb Z$.
The same squaring calculation shows
$v_2(v^{2^k}-1)=d+k$, so its order modulo $2^a$ is exactly
$2^{a-d}$, the size of that group.  If
$\psi_2(u)^{4o}=1$, then the character on
$(\mathbb Z/2^a\mathbb Z)^\times$ kills the generator
$\ell^{4o}$ of this group.  It therefore factors through
$(\mathbb Z/2^d\mathbb Z)^\times$.  Its conductor divides $2^d$,
contradicting the definition of $c_2$.  We have proved
$\psi_2(u)^{4o}\ne1$.

Reduction to the residue field kills all $\ell$-power order roots
of unity.  It is injective on prime-to-$\ell$ roots: if $\zeta^b=1$
with $\ell\nmid b$ and $\bar\zeta=1$, then
$1+\zeta+\cdots+\zeta^{b-1}$ is a unit, and
$(\zeta-1)(1+\cdots+\zeta^{b-1})=0$ implies $\zeta=1$.
Since the $2$-primary unit group has $2$-power order and $\ell$ is
odd, $\kappa_2=1$ and $\theta_2=\psi_2$.  The CRT product
decomposition and (HSR.2) give
\[
                   \theta(u)^{4o}=\psi_2(u)^{4o}\ne1.
 \tag{HSR.14}
\]
If $\overline{\psi(u)}^4=1$, reduction of $\kappa(u)$ gives
$\overline{\theta(u)^4}=1$, and injectivity on prime-to-$\ell$
roots gives $\theta(u)^4=1$.  This contradicts (HSR.14).
Notice also that $\kappa_2=1$ proves $2\notin S_b$.
\end{proof}
% END HSR2

% BEGIN HSR3
\begin{lemma}[HSR3: the exact local subgroup, with cyclotomic value one]
\label{claim:HSR3}
For the stored decomposition group at $\ell$, there is an open
normal subgroup
\[
 H=D_\ell\cap\ker\eta_4\cap\ker(\psi_\ell)_{\mathrm{Gal}}
 \tag{HSR.15}
\]
and an element $g\in H$ with
\[
 \Psi\pi(s)=\epsilon(s)^e\quad(s\in H\cap I_\ell),\qquad
 \epsilon(g)=1,\qquad (\Psi\pi(g))^2-1\in O^\times.
 \tag{HSR.16}
\]
In particular this proves the full \texttt{rig\_\allowbreak sub} output, with
its final inequality $|\epsilon(g)-1|<1$ strengthened to equality
$\epsilon(g)=1$.
\end{lemma}
\begin{proof}
Choose an embedding $j:\Omega\hookrightarrow
\overline{\mathbb Q}_\ell$ inducing the stored valuation.  One
construction embeds the completions of finite subextensions at
that valuation into the local algebraic closure and takes a
compatible system of these finite embeddings.  Each finite
compatibility problem has a solution by extension of embeddings;
compactness of the product of the finite nonempty sets gives the
system.  Restriction identifies $G_{\mathbb Q_\ell}$
topologically with $D_\ell$ and identifies their inertia groups.
Here are the details relevant to the chosen model.  The image
$j\Omega$ is the field of elements algebraic over $\mathbb Q$
in the local closure and is dense there, by the rational-polynomial
approximation and conjugate-separation argument used in NR19.
An element of $D_\ell$ preserves the normalized absolute value,
so extends isometrically to the completion.  It fixes the closure
$\mathbb Q_\ell$ of $\mathbb Q$, and consequently preserves
the elements algebraic over $\mathbb Q_\ell$.  This gives the
inverse to restriction.  Restriction is continuous: fixing a finite
set $T\subset\Omega$ pulls back to the open subgroup fixing
$\mathbb Q_\ell(j(T))$.  The continuous bijection from the compact
local Galois group to the Hausdorff group $D_\ell$ is a
homeomorphism.  Its valuation-preserving construction identifies
inertia and preserves the action on every root of unity.

The image of local inertia under $\epsilon$ is all of
$\mathbb Z_\ell^\times$.  Indeed
$\Phi_{\ell^m}(1+T)$ is Eisenstein at $\ell$: its constant term
is $\ell$, its reduction modulo $\ell$ is
$T^{\ell^{m-1}(\ell-1)}$, and it is monic.  Thus
$\mathbb Q_\ell(\mu_{\ell^m})/\mathbb Q_\ell$ is totally
ramified of degree $\ell^{m-1}(\ell-1)$, and inertia maps onto
$(\mathbb Z/\ell^m\mathbb Z)^\times$.  Compactness gives the
surjection on the inverse limit as well.

Take any lift $\sigma_0$ of arithmetic Frobenius in the local
unramified quotient.  Choose $\tau\in I_\ell$ with
$\epsilon(\tau)=\epsilon(\sigma_0)^{-1}$ and set
$\sigma=\sigma_0\tau$.  Then $\epsilon(\sigma)=1$.
Roots of unity of order prime to $\ell$ lie in the maximal
unramified extension: they are the unique Hensel lifts of roots
of the same order in finite residue fields.  Hence $\sigma$
still acts on them by raising to the $\ell$th power.  Its action
on $\mu_N$ therefore has $2$-primary coordinate $\ell$.
HSR2 gives $\overline{\psi_{\mathrm{Gal}}(\sigma)}^4\ne1$.

Since $\epsilon(\sigma)=1$, the element $\sigma$ fixes all
$\ell$-power roots of unity and kills $(\psi_\ell)_{\mathrm{Gal}}$.
Set $g=\sigma^2$.  This kills $\eta_4$, so $g\in H$, and
\[
 \overline{\Psi\pi(g)}^2
   =\overline{\psi_{\mathrm{Gal}}(\sigma)}^4\ne1.
\]
An element of the local ring $O$ is a unit exactly when its residue
is nonzero, proving the last assertion in (HSR.16).
Both characters cutting out $H$ have finite image and open kernel,
so $H$ is open and normal in $D_\ell$.  Finally inertia at $\ell$
fixes all prime-to-$\ell$ roots of unity, so on it
$\psi_{\mathrm{Gal}}=(\psi_\ell)_{\mathrm{Gal}}$.
Equations (HSR.3) and (HSR.15) give the first assertion of
(HSR.16).
\end{proof}
% END HSR3

% BEGIN HSR4
\begin{lemma}[HSR4: separation in every triangular basis]
\label{claim:HSR4}
Suppose $g\in\Gamma$ has $\eta(g)=1$ and
$\Psi(g)^2-1\in O^\times$.  In any basis in which $r(g)^{\rm mat}$
is upper triangular, its diagonal entries $\lambda_1,\lambda_2,
\lambda_3$ satisfy
\[
 |\lambda_i|=1,\qquad |\lambda_i-\lambda_j|=1\ (i\ne j),
 \qquad \lambda_i\ne u\lambda_j\quad(|u-1|<1, i\ne j).
 \tag{HSR.17}
\]
This is the exact \texttt{rig\_\allowbreak eig} conclusion.
\end{lemma}
\begin{proof}
The three entries $d=(\Psi(g),1,\Psi(g)^{-1})$ are units with
pairwise distinct residues.  Indeed $\Psi(g)^2-1$ being a unit
makes both $\Psi(g)-1$ and $\Psi(g)+1$ units, and the other
differences follow after multiplication by $\Psi(g)^{-1}$.
By HSR1 and (HSR.7), the integral model $M=r_A(g)^{\rm mat}$
has entries of absolute value at most one and reduces to
$\operatorname{diag}(\bar d_1,\bar d_2,\bar d_3)$.
Its characteristic polynomial is monic integral.  Every root has
absolute value at most one, by the same leading-term argument as
in HSR1.  Reducing its factorization into linear factors gives
\[
 \prod_i(T-\bar\lambda_i)=\prod_i(T-\bar d_i).
\]
The roots on the right are distinct and nonzero.  Unique
factorization over the residue algebraic closure makes the
$\bar\lambda_i$ a permutation of them.  This proves the first
two assertions of (HSR.17).  The triangular diagonals are these
roots because $\eta(g)=1$ makes (HSR.6) ordinary matrix
conjugation.  Finally $|\lambda_j(u-1)|<1$ and
$|\lambda_i-\lambda_j|=1$ imply
$|\lambda_i-u\lambda_j|=1$ by the ultrametric inequality.
\end{proof}
% END HSR4

\subsection{The local deformation is a finite adjoint class}

% BEGIN HSR5
\begin{proposition}[HSR5: the precise RigFin and RigLoc implication]
\label{claim:HSR5}
Let $D$ be an open subgroup of a local absolute Galois group at
$\ell$, let $r:D\to\mathcal G_3(\overline{\mathbb Q}_\ell)$
have locally constant multiplier and continuous matrix entries,
let $z$ be a continuous adjoint cocycle, and let $r_\delta$ be a
deformation with the same multiplier satisfying (HSR.8).
Equivalently, $D$ may be an abstract topological group equipped
with the specified open embedding into $G_{\mathbb Q_\ell}$;
transport along that embedding gives the present formulation.
Suppose an open normal
$H\subset D\cap\ker\eta$ and a dual-number basis give an upper
triangular representation on $H$ whose diagonal characters
$\kappa_1,\kappa_2,\kappa_3$ satisfy
\[
 \kappa_i|_{H\cap I}=\epsilon^{w_i},\qquad
 w_1>w_2>w_3,
 \qquad \bar\kappa_i\ne\bar\kappa_j\epsilon\quad(i<j).
 \tag{HSR.18}
\]
The bar here means reduction modulo $\delta$, not modulo $\ell$.
Then the original adjoint cocycle has a primitive in
$M_3(\overline{\mathbb Q}_\ell)\otimes_{\mathbb Q_\ell}
B_{\mathrm{cris}}$ on all of $D$.  This conclusion is exactly
the source's \texttt{finite\allowbreak Cocycles}, including its tensor order
and its diagonal Galois action.
\end{proposition}
\begin{proof}
The coefficient descent in HSR1 gives a finite field $E$ containing
the representation, cocycle, and basis.  It applies here because
open local Galois subgroups are compact.  All polynomial equations,
the upper-triangular conditions, and the non-equalities in
(HSR.18) descend along the injective inclusion
$E[\delta]\hookrightarrow\overline{\mathbb Q}_\ell[\delta]$.
On $H$, both $r$ and $r_\delta$ are ordinary linear
representations, because their multiplier is one.  Under the
identification $H=G_L$ for a finite extension of the local base,
$H\cap I$ is precisely $I_L$.

We recall the actual calculation from RX2 to record why all these
hypotheses, including the dual-number ones, suffice.  For
$A=E$ or $E[\delta]$, an upper triangular representation with
(HSR.18) is crystalline as an underlying $E$-representation.
Each $\kappa_i\epsilon^{-w_i}$ is unramified.  Even when it
contains a nontrivial unramified $\delta$-variation, its underlying
$E$-representation is unramified and hence crystalline; twisting
by $\epsilon^{w_i}$ preserves this property.  For the induction
step, let $V_0$ be the first two (or first $n-1$) coordinates and
put $W=V_0\otimes_A A(\kappa_n^{-1})$.  By induction and tensor
and quotient stability, $W$ is crystalline.  Its $E$-character
constituents are $\bar\kappa_i/\bar\kappa_n$, with repetitions
when $A$ has dual numbers.  All inertia exponents are positive,
so $\operatorname{Fil}^0D_{\mathrm{dR}}(W)=0$ in the convention
where $\epsilon$ has Hodge--Tate weight $-1$.
The constituents of $W^\vee(1)$ are nontrivial by (HSR.18),
whence $H^0(L,\allowbreak{} W^\vee(1))=0$.

The local Euler characteristic and Tate duality give
\[
 H^2(L,\allowbreak{} W)=0,\qquad
 \dim_E H^1(L,\allowbreak{} W)=\dim_E H^0(L,\allowbreak{} W)+[L:\mathbb Q_\ell]\dim_EW.
\]
The Bloch--Kato finite dimension formula gives the same expression
for $\dim_EH^1_f(L,\allowbreak{} W)$, since $\operatorname{Fil}^0=0$.
Thus $H^1=H^1_f$ for this $W$.  RX2 supplies the lattice limit
argument deriving these rational statements from the already
installed finite local duality and Euler identities, and BX1--BX4
supply the finite dimension formula for the actual period model.
The last column of the upper triangular representation, divided
by its last diagonal character, is a cocycle in $W$.  Its
$B_{\mathrm{cris}}$ primitive splits the corresponding underlying
$E$-extension after tensoring.  The crystalline comparison
isomorphisms for its submodule and quotient therefore give one
for the extension.  This proves the induction, with all
unramified dual-number diagonal variations included.

Consequently the underlying $E$-self-extension of $r|_H$ defined
by $r_\delta|_H$ is crystalline.  Exactness of $D_{\mathrm{cris}}$
and the comparison isomorphism give an equivariant splitting after
tensoring with $B_{\mathrm{cris}}$.  Choose it compatibly with $E$
by splitting first over the semisimple algebra
$L_0\otimes_{\mathbb Q_\ell}E$.  In the vector-space splitting
already used in (HSR.8), this is precisely an element $b$ with
\[
 z(h)\otimes1=\operatorname{Ad}_{r,\mathrm{cris}}(h)b-b
                  \qquad(h\in H).
 \tag{HSR.19}
\]

Here is the descent to $D$, also valid if $\ell$ divides $[D:H]$.
Subtract the coboundary in (HSR.19) from $z\otimes1$, obtaining
a cocycle $f$ which vanishes on $H$.  Normality and the cocycle
identity show that it factors through $Q=D/H$ with values in
the $H$-invariants.  For $S_f=\sum_{q\in Q}f(q)$, permutation
of the finite sum gives
\[
                  gS_f-S_f=-|Q|f(g).
\]
Thus $b-|Q|^{-1}S_f$ is a primitive on $D$.  The action here is
the full disconnected adjoint on the first tensor factor and
the actual Galois action on $B_{\mathrm{cris}}$ on the second;
no connectedness assumption on $D$ was made.  Extending
coefficients to the algebraic closure carries this primitive to
the required one.  BX1--BX2 identify the period ring and this
literal tensor-primitive condition with \texttt{finite\allowbreak Cocycles}.
Equivariance of scalar extension is immediate on pure tensors
$X\otimes b$ and hence on their sums.  This proves both the
finite-coefficient \texttt{Rig\allowbreak Fin} and the algebraic-closure
\texttt{Rig\allowbreak Loc} interfaces.
\end{proof}
% END HSR5

% BEGIN HSR6
\begin{lemma}[HSR6: the exact inertia primitives away from ell]
\label{claim:HSR6}
Every deformation satisfying (HSR.8)--(HSR.10) has an adjoint
coboundary on inertia at each finite $p\ne\ell$.
\end{lemma}
\begin{proof}
If $p\notin S$, the quotient $\pi$ kills inertia, and
$z(1)=0$, so the assertion holds with primitive zero.
Suppose $p\in S_b$.  By HSR2, $p\ne2$; hence $\eta_4$ is
trivial on $I_p$, since the fourth roots of unity are unramified
at odd primes.  Write the supplied diagonalizing basis as
$h=h_0+\delta h_1$, with $h_0$ invertible, and write
$C=h^{-1}r_\delta(s)h$ for $s\in I_p$.  Its matrix is diagonal
with the fixed inertia entries (HSR.10), so has no $\delta$
coefficient.  Put $R=r(s)^{\rm mat}$.  The constant and
$\delta$ terms of $hC=(1+\delta z(s))Rh$ give
\[
 h_0C_0=Rh_0,\qquad
 h_1C_0=Rh_1+z(s)Rh_0.
\]
With $U=h_1h_0^{-1}$, this says
\[
 z(s)=U-RUR^{-1}
     =\operatorname{Ad}_r(s)(-U)-(-U).
 \tag{HSR.20}
\]
The single matrix $-U$ works for all inertia elements.  This is
the source's \texttt{rig\_\allowbreak diag}, with its sign explicitly fixed.

If $p\in S\setminus T_b$, the cocycle vanishes on $N_p$ from
(HSR.9).  This is a normal subgroup of finite index in $I_p$:
it is the intersection of the kernels of two finite characters,
$\eta_4$ and $\bar\Psi$, with inertia.  The same finite quotient
sum as in HSR5, now without a period-ring tensor, proves that
$z|_{I_p}$ is a coboundary.  In detail $z$ factors through
$I_p/N_p$ with values invariant under $N_p$, and
$-\sum_{q\in I_p/N_p}z(q)/[I_p:N_p]$ is a primitive.  This
also handles $p=2$, without asserting that $\eta_4$ is
unramified there.  These cases exhaust all $p\ne\ell$.
\end{proof}
% END HSR6

% BEGIN HSR7
\begin{proposition}[HSR7: all source tangent classes satisfy literal IsBK]
\label{claim:HSR7}
Let $R=r\pi$ and $Z=z\pi$.  For every allowed tangent in the
source's \texttt{Rigid} definition,
\[
       Z\in Z^1(G_{\mathbb Q},\operatorname{Ad}_R),
       \qquad \texttt{Is\allowbreak BK}_{\ell}(\operatorname{Ad}_R,\allowbreak{} Z).
 \tag{HSR.21}
\]
\end{proposition}
\begin{proof}
Continuity follows from HSR1 and continuity of the quotient map;
the cocycle law follows by composition.  The away condition is
HSR6, using the stored valuation when $p\in S$ and any valuation
when $p\notin S$.

At $\ell$, use the subgroup $H$ and element $g$ of HSR3 and the
dual-number upper-triangular basis supplied in (HSR.10).  On $H$
the diagonal characters have inertia exponents
\[
                         (w_1,\allowbreak{} w_2,\allowbreak{} w_3)=(e,0,-e),
 \tag{HSR.22}
\]
which are strictly decreasing because $e\ge1$.  Reduction of
the basis modulo $\delta$ is an ordinary invertible matrix and
triangularizes $r(g)^{\rm mat}$.  HSR4 gives pairwise distinct
diagonal entries at $g$.  Since $\epsilon(g)=1$, this proves
$\bar\kappa_i(g)\ne\bar\kappa_j(g)\epsilon(g)$ for every
$i<j$, hence the character non-equalities in (HSR.18).
HSR5 gives the required tensor primitive on the entire stored
$D_\ell$.

To check the exact \texttt{Is\allowbreak BK.\allowbreak above} interface, take the embedding
$j$ of HSR3.  A field homomorphism
$\tau:\mathbb Q\to\overline{\mathbb Q}_\ell$ is the unique
rational inclusion, so $j|_{\mathbb Q}=\tau$ as required.
The source's \texttt{dec\allowbreak J} consists of pairs $(s,\allowbreak{} u)$ such that
$u j=j s$.  It is the graph of restriction from
$G_{\mathbb Q_\ell}$; its first projection has image exactly
the stored $D_\ell$ and its second projection is a topological
isomorphism.  HSR3 supplies the identification with that stored
group, and ADK5 gives the graph topology directly by open
stabilizers of finite global sets.  Transport (HSR.19) along these
maps.  On a pure tensor the transported action is
\[
            X\otimes b\longmapsto
            \operatorname{Ad}_R(s)X\otimes u(b),
\]
which is exactly the source's \texttt{act\allowbreak C}.  Thus the transported
primitive proves its \texttt{finite\allowbreak Cocycles} assertion, not just
a condition on an unspecified local representation.  This
verifies both clauses of the literal \texttt{Is\allowbreak BK}.
\end{proof}
% END HSR7

\subsection{The exact global polarization and the endpoint}

% BEGIN HSR8
\begin{proposition}[HSR8: the automorphic twist and full disconnected adjoint]
\label{claim:HSR8}
Suppose the original \texttt{BPrime\_\allowbreak s} holds for
$a(s)=R(s)^{\rm mat}$ and the specified $K$ with
\texttt{Is\allowbreak Qi K}.  Then ADK6 applies to
$\operatorname{Ad}_R$ on the full $M_3$, with its original
coefficient and local-condition scope.
\end{proposition}
\begin{proof}
By definition \texttt{Is\allowbreak Qi K} says $[K:\mathbb Q]=2$ and
$K$ contains a square root of $-1$.  The algebra structure
$K\hookrightarrow\Omega$ provided by
\texttt{BPrime\_\allowbreak s.\allowbreak aut} therefore has image $\mathbb Q(i)$.
The image of restriction $G_K\to G_{\mathbb Q}$ is exactly
$\ker\eta_4$.  For example, the action on $i$ is its action on
$\mu_4$, and the nontrivial element of
$\operatorname{Gal}(\mathbb Q(i)/\mathbb Q)$ has character
value $-1$.

Unfolding \texttt{Af.\allowbreak Aut\allowbreak On\allowbreak Alt K iota (-1) a} gives a
representation $r_0:G_K\to\operatorname{GL}_3
(\overline{\mathbb Q}_\ell)$ whose matrix is literally
$a(s)$ on $G_K$, and says that
\[
                         r_K=\epsilon^{-1}r_0
 \tag{HSR.23}
\]
satisfies \texttt{Is\allowbreak Aut\allowbreak Csd\allowbreak A}.  Its first conjunct is
\texttt{Aut\allowbreak Csd0}.  Inspecting that definition gives the very same
existential cuspidal representation, regular algebraicity,
conjugate self-duality, and \texttt{Att\allowbreak C} condition as
\texttt{Is\allowbreak Aut\allowbreak Csd}.  The normalized Satake and infinite-place
dictionary, including the geometric Frobenius convention, is
already checked in ADK3.  Thus (HSR.23) satisfies the literal
automorphy hypothesis of ADK6.  HSR1 gives continuity of $r_0$
and hence of $r_K$; the entries of the inverse are continuous as
well, so this is continuity into $\operatorname{GL}_3$.

The field $\Omega$ is an algebraic closure of this specified
copy of $K$: it is algebraically closed and algebraic over
$\mathbb Q\subset K$.  Under the above restriction map,
$\ker\epsilon|_{G_K}$ identifies with
$\ker\eta_4\cap\ker\epsilon$ in $G_{\mathbb Q}$.
The source uses $\texttt{cyc\allowbreak Z}$ for the latter character and
$\texttt{cyc\allowbreak Bar}$ for its inclusion in
$\overline{\mathbb Q}_\ell^\times$; the inclusion is injective,
so their kernels coincide.  The \texttt{BPrime\_\allowbreak s.\allowbreak cyc}
assertion says that $r_0$ is irreducible on this exact subgroup.
The twist in (HSR.23) is one there.  Therefore $r_K$ is
irreducible on its full infinite cyclotomic kernel, exactly the
hypothesis of ADK6.  No residual or finite-level version is being
substituted.  The additional \texttt{BPrime\_\allowbreak s.\allowbreak opens} hypothesis
is retained but is unnecessary for this implication.

It remains to check the pairing and the disconnected action.
HSR1 and \texttt{Is\allowbreak Host\allowbreak Frame.\allowbreak cc} give
$R(c_\infty)=(1,-1)$, with $c_\infty$ induced by an actual
complex embedding.  The semidirect-product law implies, for
$s\in G_K$,
\[
 r_0(c_\infty s c_\infty^{-1})
      =Jr_0(s)^{-t}J,\qquad
 r_0(s)^tJr_0(c_\infty s c_\infty^{-1})=J.
\]
The cyclotomic character is invariant under conjugation, so
(HSR.23) gives
\[
 r_K(s)^tJr_K(c_\infty s c_\infty^{-1})
                       =\epsilon(s)^{-2}J
                       =\epsilon(s)^{1-3}J.
 \tag{HSR.24}
\]
Also $J^t=J=J^{-1}$ and, for every $X\in M_3$,
\[
 \operatorname{Ad}_R(s)X=r_K(s)Xr_K(s)^{-1},\qquad
 \operatorname{Ad}_R(c_\infty)X=-J^{-1}X^tJ.
 \tag{HSR.25}
\]
Equations (HSR.24)--(HSR.25), together with the actual complex
embedding, are every clause of \texttt{Is\allowbreak Ad\allowbreak Of}, with $P=J$.
In particular the scalar line is present and complex conjugation
acts on it by $-1$.  There is no passage to trace zero and no
undetermined scalar in the pairing.  ADK6 now proves
\[
 H^1_f(\mathbb Q,\operatorname{Ad}_R)=0
\]
in the literal source sense \texttt{Bloch\allowbreak Kato\allowbreak Selmer\allowbreak Vanishes}.
\end{proof}
% END HSR8

% BEGIN HSR9
\begin{theorem}[HSR9: the original HostRigid\_\allowbreak s statement]
\label{claim:HSR9}
For every choice of data quantified in the source's
\texttt{Host\allowbreak Rigid\_\allowbreak s}, an \texttt{Is\allowbreak Char\allowbreak Zero\allowbreak Pt} satisfying
\texttt{BPrime\_\allowbreak s} is rigid in the exact sense (HSR.11).
Consequently \texttt{HOST\_\allowbreak RIGID\_\allowbreak S} is supplied by the
human proofs ADK1--ADK6, HSR1--HSR8, and the previously installed
local inputs BX1--BX4, DX1, NR19, and RX2.
\end{theorem}
\begin{proof}
Take an arbitrary continuous cocycle and arbitrary allowed
deformation from the definition of rigidity.  HSR7 makes
$Z=z\pi$ a literal Bloch--Kato cocycle for
$\operatorname{Ad}_{r\pi}$.  HSR8 and ADK6 give a single
$v\in M_3(\overline{\mathbb Q}_\ell)$ such that
\[
 z(\pi(s))=\operatorname{Ad}_r(\pi(s))v-v
                       \qquad(s\in G_{\mathbb Q}).
\]
The map $\pi$ is onto.  Set $X=-v$ and choose a preimage of
each $g\in\Gamma$ to obtain (HSR.11).  This proves every
quantifier of the original rigidity conclusion.  In particular
the primitive is one global matrix, not a family of unrelated
local primitives.
\end{proof}
% END HSR9

\subsection{Scope and source comparisons}

The proof has supplied the entire original
\texttt{Host\allowbreak Rigid\_\allowbreak s} implication.  It retains the original
marked-element, parity, coefficient, valuation, and
\texttt{BPrime\_\allowbreak s} hypotheses, and it works at every odd
$\ell$ permitted by \texttt{Host\allowbreak Data}.  It uses only the
$p=\ell$ instance of the numerical inequality in (HSR.2), and
only the \texttt{cyc} and \texttt{aut} parts of
\texttt{BPrime\_\allowbreak s}.  These observations simplify this proof;
they do not manufacture those hypotheses for a host or remove
them from any downstream source statement.

The local calculation is an explicit application of the already
proved RX2, with its finite local duality and exact crystalline
period model.  The global calculation is an explicit application
of the published theorem matched in ADK1--ADK6.  There is no new
automorphy theorem or new vanishing conjecture hidden in this
step.  Conversely this section does not supply
\texttt{Host\allowbreak Number\_\allowbreak s}, the existence or automorphy of its points,
the quantitative congruence estimates, or the other remaining
arithmetic producers.  It closes a proof-reconstruction leaf;
it does not by itself complete the reconstruction of RH or GRH.

\begin{center}
\begin{tabular}{p{0.31\textwidth}p{0.61\textwidth}}
Original interface & Human proof and exact comparison \\
\hline
\texttt{Is\allowbreak Char\allowbreak Zero\allowbreak Pt}, \texttt{Rig\allowbreak Desc}
 & HSR1: norm continuity from the actual finite DVR and Baire descent.\\
\texttt{Dyadic\allowbreak Character\allowbreak Power\allowbreak Ne\allowbreak One}, \texttt{rig\_\allowbreak sub}
 & HSR2--HSR3: dyadic conductor, finite characters, stored local group,
   and the stronger witness $\epsilon(g)=1$.\\
\texttt{Rig\allowbreak Sep}, \texttt{rig\_\allowbreak eig}
 & HSR4: reduction of the actual characteristic polynomial in every
   triangular basis.\\
\texttt{rig\_\allowbreak fin}, \texttt{rig\_\allowbreak loc}
 & HSR5 and RX2: finite coefficients, dual-number ordinary
   crystalline calculation, and descent with the diagonal period action.\\
\texttt{rig\_\allowbreak diag}, \texttt{rig\_\allowbreak away}, \texttt{rig\_\allowbreak above}
 & HSR6--HSR7: explicit inertia primitives and the literal
   \texttt{dec\allowbreak J}/\texttt{finite\allowbreak Cocycles} match.\\
\texttt{BPrime\_\allowbreak s}, \texttt{Aut\allowbreak On\allowbreak Alt}, \texttt{Is\allowbreak Ad\allowbreak Of}
 & HSR8: exact twist $\epsilon^{-1}$, full infinite cyclotomic
   kernel, and actual matrix $P=J$.\\
\texttt{Host\allowbreak Rigid\_\allowbreak s}
 & HSR9: ADK6 applied to the literal global cocycle, quotient
   surjectivity, and the required sign of \texttt{cob}.\\
\end{tabular}
\end{center}

The accompanying \texttt{provenance.\allowbreak json} gives pinned source
paths, declaration ranges, range hashes, inherited human-proof
locations, and claim dependencies for this table.


\section{Finite representatives for the original ordinary-form levels}
\label{sec:finite-levels-v42}

The ordinary projector is constructed on functions on a definite unitary
adelic quotient.  Before applying any argument about finite coefficient
rings, one needs a finite set of adelic representatives for the exact level
of those functions.  This chapter supplies that input.  It proves the full
original \texttt{Fu\allowbreak Cmp\allowbreak Mem.\allowbreak fin\allowbreak U} for arbitrary integral coefficient
representation and every permitted compact-open auxiliary level.

The standing source pin is
\texttt{7cbf7a0c32fca1badf05931879a37c3f8a2b03b4}.
All definitions below are read at that pin.  The only non-elementary input
is the already reconstructed finite-class theorem FV9, applied through an
explicit dictionary.  No integral type, ordinary structure, automorphic
Galois representation, or new published theorem is assumed.

\subsection{The objects and endpoint}

Let \(L\) be the original CM field, \(F=L^+\), and \(\sigma\) its complex
conjugation.  Let \(\ell,\allowbreak{} s,\allowbreak{} n,\allowbreak{} V,\allowbreak{} Q,\allowbreak{} c,\allowbreak{} m,\tau _0,\allowbreak{} K\) be arbitrary original
\texttt{fin\allowbreak U} data: \(s\) is a Hida setting, \(n\ge1\),
\(V:s.\texttt{Lev}\), \(Q:\texttt{QData}\,\allowbreak{} V\),
\(\tau _0:\prod_{u\in s.S_t}\mathrm{GL}_n(\mathcal O_{L_{w(u)}})
\to\mathrm{GL}_m(c.\mathcal O)\), and \(K_q\) is a compact open
subgroup of \(\mathrm{GL}_n(L_{w_q})\) for each \(q\in Q\).
Write \(v(w)=w\cap\mathcal O_F\), \(v_q=v(w_q)\), and
\[
 S=\{v:v\mid\ell\}\cup s.S_t\cup\{s.v_0\},\qquad
 V.S'=S\cup\{v_q:q\in Q\}.
\]
The \(v_q\)'s are distinct and outside \(S\).  Put
\[
 U_{\mathbb A}=\{g\in\mathrm{GL}_n(\mathbb A_L):
                  \sigma(g)^{\mathsf t}g=1\},\quad
 U_f=\{g\in U_{\mathbb A}:g_\infty=1\},\quad
 U_\infty=\{g\in U_{\mathbb A}:g_f=1\}.
\]
Here conjugation on the full adele ring is the original
\texttt{conj\allowbreak Ad}; the infinite coordinates are read through the original
\texttt{inf\allowbreak Eq}.  The subgroup \(\Gamma\) is the rational matrix image
\(U(F)\), literally \texttt{rat\allowbreak U}.  These groups are respectively
\texttt{UA}, \texttt{fin\allowbreak U}, and \texttt{inf\allowbreak U} in
\texttt{Base\allowbreak Change\allowbreak Definite\allowbreak Unitary.\allowbreak lean}; the name \texttt{fin\allowbreak U} there
is a subgroup and is distinct from the endpoint theorem.

Let \(a:U_{\mathbb A}\to G^a\) be the range-restricted projection onto
all finite places of \(L\) away from \(\ell\).  For each place \(w\),
\(\operatorname{cmp}_w\) is evaluation at \(w\) when \(w\nmid\ell\)
and the trivial homomorphism otherwise.  Let \(U_b\subset G^a\) impose
integrality at the stored \(s.S_t\) places and the stored principal
congruence condition at \(w(s.v_0)\).  On this subgroup the original
\(\operatorname{sec}_w\) is the inverse of its own-place projection
on the subgroup supported above \(v(w)\), when that projection is
bijective, and is the trivial homomorphism otherwise.  Set
\[
 U_Q=\{u\in U_b:\operatorname{cmp}_w(u)\in\mathrm{GL}_n(\mathcal O_{L_w})
      \ (v(w)\notin V.S'),\quad
      \operatorname{cmp}_{w_q}(u)=1\ (q\in Q)\},
\]
\[
 U(K)=U_Q\vee\bigvee_{q\in Q}\operatorname{sec}_{w_q}(K_q),\qquad
 G^\ell=\prod_{v\mid\ell}\mathrm{GL}_n(L_{w(v)}).
\]
All joins here are subgroup joins in the source's original \(U_b\).
With \(X=(\texttt{hset}\,\allowbreak{} Q\,\tau_0).X\), the assertion to prove is
\begin{equation}\label{eq:fnl-target}
 \exists\mathcal S\subset G^a\times G^\ell\text{ finite},\quad
 \forall g\ \exists\gamma\in\Gamma,\ x\in\mathcal S,\ u\in U(K),\
 b\in X.I.\operatorname{Iw}(0):\quad
             g=\operatorname{glob}(\gamma)\,\allowbreak{} x\,(u,\allowbreak{} b).
\end{equation}
The finite set may depend on the displayed initial datum and \(K\), but
is chosen before \(g\).  It does not require a bound on its size.

% BEGIN FNL1
\paragraph{FNL1. Finite coordinate patching in an arbitrary group.}
\label{claim:FNL1}
Suppose \(I\) is finite, \(p_i:G\to H_i\) are homomorphisms and
\(P\subseteq G\) is a subgroup.  If for each \(i\) and prescribed
\(b_i\in H_i\) one has \(y_i\in P\) with
\(p_i(y_i)=b_i\) and \(p_j(y_i)=1\) for every \(j\ne i\), then there
is \(y\in P\) with \(p_i(y)=b_i\) for all \(i\).
Choose any ordering of \(I\) and multiply the \(y_i\)'s in that order.
Every coordinate of the product has exactly one possibly nonidentity
factor, so gives the required value.  Commutation of the \(y_i\)'s is
unnecessary.  For \(I=\varnothing\), use \(y=1\).
This is the original \texttt{fu\_\allowbreak patch} argument.
% END FNL1

% BEGIN FNL2
\paragraph{FNL2. The identity-form adelic dictionary, including topology.}
\label{claim:FNL2}
The finite unitary group for the algebraic \(F\)-group
\[
 \mathcal U(A)=\{g\in\mathrm{GL}_n(A\otimes_F L):
                     (1\otimes\sigma)(g)^{\mathsf t}g=1\}
\]
is topologically isomorphic to the literal \(U_f\).  The isomorphism
commutes with every finite-place evaluation and sends the finite image of
each rational matrix to its finite adelic extension with infinite
component \(1\).

Indeed the algebra map
\[
 \alpha:\mathbb A_{F,\allowbreak{} f}\otimes_F L\longrightarrow\mathbb A_{L,\allowbreak{} f},
 \qquad a\otimes x\longmapsto \bigl(a_{v(w)}x\bigr)_w
\]
is the original \texttt{fin\allowbreak Ad\allowbreak Equiv}.  Locally it is
\(F_v\otimes_F L\simeq\prod_{w\mid v}L_w\).
To check topology, choose an \(F\)-basis of \(L\).  After inverting
one nonzero integer it is a basis of an order containing and contained
in fixed multiples of \(\mathcal O_L\).  Outside finitely many places
its integral lattice agrees with the product of the local integer
rings.  Each local map and its inverse are continuous linear maps
between finite dimensional local vector spaces.  They therefore induce
inverse continuous maps of the restricted products.  Equivalently this
is the tensor topology of FV4 in these coordinates.

For a simple tensor,
\(\alpha((1\otimes\sigma)(a\otimes x))=\alpha(a\otimes\sigma x)\)
is \(\operatorname{conjFin}(\alpha(a\otimes x))\); here conjugation
fixes the embedded \(F\)-adeles and carries the \(w\) coordinate to
the conjugate place.  Linearity gives the equality on every tensor.
Applying it entrywise identifies the two unitary equations, including
the identity Gram matrix.  Extending a finite matrix by the identity
at infinity then gives \(U_f\); projection to the finite components
is its inverse.  Both maps are continuous in the displayed coordinates.
For \(x\in L\), \(\alpha(1\otimes x)\) is its diagonal finite
adelic image.  Thus rational matrices and all local projections agree
literally.  Finally, every \(g\in U_{\mathbb A}\) has a unique
factorization \(g=g_f g_\infty=g_\infty g_f\) with the indicated
finite and infinite components: extend each component by the identity
in the other factor.  This also proves that \(U_f\) and \(U_\infty\)
commute.

This dictionary checks the specialization \(J=1\), \(E=L\),
\(F=L^+\), and \(\sigma=c_L\) of FV1--FV4 and FV9.  The CM condition,
Hermitian symmetry, and nonzero determinant are automatic from these
data.  In particular no change of rational form or rational subgroup
has entered the finite-class application.
% END FNL2

% BEGIN FNL3
\paragraph{FNL3. Finite classes with the source's infinite factor in its exact position.}
\label{claim:FNL3}
For every open subgroup \(K'\subset U_f\), there is a finite
\(\mathcal T\subset U_f\subset U_{\mathbb A}\) such that
\begin{equation}\label{eq:fnl-class}
 \forall g\in U_{\mathbb A}\ \exists\gamma\in\Gamma,\ t\in\mathcal T,
 z\in U_\infty,\ k\in K':\qquad g=\gamma t z k.
\end{equation}

Apply FV9 through FNL2 to get
\(g_f=\gamma_f t k\).  For the actual rational adelic matrix
\(\gamma\), define \(z=(\gamma t)^{-1}gk^{-1}\).
Its finite component is
\(t^{-1}\gamma_f^{-1}g_f k^{-1}=1\), so \(z\in U_\infty\).
Its definition gives \(g=\gamma t z k\) with exactly that order.
The selected \(t\)'s depend only on \(K'\), not on \(g\).
This proves the full \texttt{Def\allowbreak U.\allowbreak Class\allowbreak Fin.\allowbreak class\allowbreak Set\_\allowbreak finite} conclusion
used here.  It uses the inherited checked Borel application in FV9;
it does not infer finiteness merely from finite volume.  Compactness
of \(K'\) is not required.
% END FNL3

% BEGIN FNL4
\paragraph{FNL4. All the places at which sections are needed are actually split.}
\label{claim:FNL4}
For every \(v\in V.S'\), the extension \(L/F\) is split at \(v\).
Furthermore \(v(w(v))=v\), and
\(\ell\in w\) if and only if \(\ell\in v(w)\).

The latter assertions follow by pulling back the maximal ideal of the
same stored valuation subring along
\(\mathcal O_F\to\mathcal O_L\to\Omega\); this is also NTL15.
For the split assertion, the hypothesis \(V.\texttt{split}\) says
that the absolute decomposition group at the stored place of \(\Omega\)
is contained in \(\operatorname{Gal}(\Omega/L)\).
Its image in \(\operatorname{Gal}(L/F)\) is the finite decomposition
group at \(w(v)\), hence that group is trivial.  A place in a quadratic
Galois extension has trivial decomposition group precisely when there
are two distinct places above it, each with completion \(F_v\).

Here the asserted surjectivity onto the finite decomposition group has
no unmentioned splitting assumption.  If an automorphism stabilizes
\(w(v)\), lift it to \(\Omega\); its image of the stored valuation
and the stored valuation both lie over \(w(v)\).  An element of
\(\operatorname{Gal}(\Omega/L)\) carries the former to the latter,
and correcting the lift gives an element of the absolute decomposition
group.  The conjugacy of extensions of a place follows first for finite
normal extensions from transitivity of primes, and then for \(\Omega\)
by compactness of the compatible finite Galois groups.  For completeness,
finite transitivity follows by prime avoidance and the norm: if a prime
\(\mathfrak q\) were outside the orbit of \(\mathfrak p\), choose an
integer in \(\mathfrak q\) outside every conjugate of \(\mathfrak p\).
Its norm belongs to the common base prime and is outside
\(\mathfrak p\), a contradiction.

In particular every \(w(v)\) for \(v\mid\ell\), and every added
\(w_q\), is a split place.  The proof retains the given \(V\); no
replacement set of places is chosen.
% END FNL4

% BEGIN FNL5
\paragraph{FNL5. The actual split sections and their trivial-map fallback.}
\label{claim:FNL5}
At a split place \(w\), put \(\bar w=\sigma w\).  Conjugation gives
an isomorphism \(c_w:L_w\to L_{\bar w}\).  For
\(h\in\mathrm{GL}_n(L_w)\), form the adelic unitary matrix \(j_w(h)\)
whose \(w\) coordinate is \(h\), whose \(\bar w\) coordinate is
\(c_w(h)^{-\mathsf t}\), and whose other finite and all infinite
coordinates are \(1\).  This is a continuous homomorphism.  The two
unitary equations at the split pair force exactly this relation, so
\(j_w\) is inverse to the \(w\) projection on the source subgroup
\(\operatorname{locU}_w\).  In particular it is trivial at each
place whose underlying \(F\)-place differs from \(v(w)\).

If \(v(w)\notin S\), its away projection belongs to \(U_b\): at all
the stored \(S_t\) and auxiliary-congruence coordinates it is the
identity.  Also \(w,\bar w\nmid\ell\).  Projection at \(w\) on the
source \(\operatorname{locA}_w\subset U_b\) is bijective.  Surjectivity
is provided by \(a(j_w(h))\).  For injectivity, an element in its
kernel has all coordinates away from \(v(w)\) equal to \(1\), its
\(w\) coordinate equal to \(1\), and therefore its \(\bar w\)
coordinate equal to \(1\) by the unitary relation; the away-range
coordinates separate its elements.  Thus
\[
       \operatorname{sec}_w(h)=a(j_w(h))
          \qquad(v(w)\notin S).
\]
This is an equality with the definition of the original section:
the bijective alternative of its definition has been verified and
its trivial-map fallback cannot occur at any \(w_q\).  At \(w\mid\ell\)
we use \(j_w\) in the full adelic group, not the away-section fallback.
The opposite-place inverse transpose and the distinction between these
two types of sections are essential to these identities.
% END FNL5

% BEGIN FNL6
\paragraph{FNL6. Surjectivity onto the actual frame and annihilation of the infinite group.}
\label{claim:FNL6}
The homomorphism
\[
 \Phi:U_{\mathbb A}\longrightarrow G^a\times G^\ell,
 \qquad g\longmapsto\bigl(a(g),(g_{w(v)})_{v\mid\ell}\bigr)
\]
is surjective, restricts to the original \(\operatorname{glob}\) on
\(\Gamma\), and kills \(U_\infty\).

For a prescribed \((a_0,\allowbreak{} b)\), choose \(g\) with \(a(g)=a_0\), using
that \(G^a\) is the image, not an independently specified restricted
product.  For each \(v\mid\ell\) apply FNL5 with
\(h_v=g_{w(v)}^{-1}b_v\).  Its matrix \(j_{w(v)}(h_v)\) belongs to
\(\ker a\), because both places above \(v\) divide \(\ell\).
It changes no selected \(w(v')\) coordinate with \(v'\ne v\).
The index set is finite, and FNL1 gives \(y\in\ker a\) with
\(y_{w(v)}=h_v\) for all \(v\).  Consequently \(\Phi(gy)=(a_0,\allowbreak{} b)\).
On \(\Gamma\) the displayed expression is the definition of
\texttt{hset.\allowbreak glob}.  An infinite-only matrix has every finite
coordinate equal to \(1\), hence \(\Phi(U_\infty)=1\).
% END FNL6

% BEGIN FNL7
\paragraph{FNL7. The exact first Iwahori subgroup is open.}
\label{claim:FNL7}
For a nonarchimedean completion \(E=L_w\), the source subgroup
\(\mathrm{I}_1(0)\) is the inverse image of the upper unipotent
subgroup of \(\mathrm{GL}_n(k_E)\) under reduction on
\(\mathrm{GL}_n(\mathcal O_E)\).  In particular it is open.  If
\(b_v\in\mathrm{I}_{1,\allowbreak{} v}(0)\) for each \(v\mid\ell\), then
\((b_v)_v\in X.I.\operatorname{Iw}(0)\).

Indeed the three defining factors of \(\mathrm{I}_1(0)\) are the
integral upper unipotents \(N\), the diagonal matrices
\(T_c(0)\) congruent to \(1\) modulo the maximal ideal, and the
lower unipotents \(N_m(0)\) congruent to \(1\).  They all reduce
to upper unipotents.  Conversely, suppose \(g\) has such reduction.
The following upper-diagonal-lower elimination proves that
\(g=x t y\) with these three types of factors.  Write
\[
 g=\begin{pmatrix}A&b\\c&d\end{pmatrix}
 =\begin{pmatrix}1&bd^{-1}\\0&1\end{pmatrix}
  \begin{pmatrix}A-bd^{-1}c&0\\0&d\end{pmatrix}
  \begin{pmatrix}1&0\\d^{-1}c&1\end{pmatrix}.
\]
Here \(d\equiv1\) is a unit, \(c\equiv0\), and the Schur
complement reduces to an upper unipotent matrix of size \(n-1\).
Induction gives the required factors; the zero-size case is trivial.
The factorization lies in the subgroup join defining \(\mathrm I_1(0)\).
Thus the claimed equality holds, and reduction is continuous to a
finite discrete matrix group.  Since the integral general linear
group is open, so is this inverse image.

Finally choose such factorizations \(b_v=x_vt_vy_v\) in each factor.
Their finite coordinatewise products satisfy
\(x\in X.I.N\), \(t\in X.I.T_c(0)\subset X.I.T\), and
\(y\in X.I.N_m(0)\).  Hence \(b=xty\in X.I.\operatorname{Iw}(0)\)
by its original subgroup-join definition.  This reconstructs the
\texttt{fu\_\allowbreak Iw} use of Iwahori factorization without importing the
representation-theoretic parts of \texttt{Gl\allowbreak Bridge}.
% END FNL7

% BEGIN FNL8
\paragraph{FNL8. The source's auxiliary subgroup is open on the same adelic group.}
\label{claim:FNL8}
In \(U_{\mathbb A}\) define the original subgroup
\[
 \begin{split}
 K_0=\{k:\;& k_w\in\mathrm{GL}_n(\mathcal O_{L_w})
                          \quad\hbox{for every finite }w;\\
 &k_{w(s.v_0)}\equiv1\pmod{\mathfrak m_{w(s.v_0)}};\quad
 k_{w_q}\in K_q\ (q\in Q);\quad
 k_{w(v)}\in\mathrm I_{1,\allowbreak{} v}(0)\ (v\mid\ell)\}.
 \end{split}
\]
It is open, and \(K_0\cap U_f\) is open in \(U_f\).
The all-integral condition is open: the product
\(\widehat{\mathcal O}_L=\prod_w\mathcal O_{L_w}\) is an open
subring in the finite adele ring, and membership of a matrix and its
inverse in \(M_n(\widehat{\mathcal O}_L)\) is an open condition in
\(\mathrm{GL}_n(\mathbb A_{L,\allowbreak{} f})\).  It is exactly coordinatewise
membership in the integral general linear groups.  Finite projection
and inclusion of \(U_{\mathbb A}\) are continuous by FNL2.
The principal congruence condition is the inverse image of the
identity under continuous reduction.  Each \(K_q\) is open by the
original datum and each \(\mathrm I_{1,\allowbreak{} v}(0)\) is open by FNL7.
There are only finitely many of these additional conditions, so their
intersection with the all-integral subgroup is open.  Taking its
preimage along \(U_f\hookrightarrow U_{\mathbb A}\) proves the last
assertion.  Notice that compactness of \(K_0\) was not asserted or
needed.  The auxiliary prime and every supplied \(K_q\) are unchanged.
% END FNL8

% BEGIN FNL9
\paragraph{FNL9. Transport of the actual finite subgroup into the actual level.}
\label{claim:FNL9}
For every \(k\in K_0\), write \(\Phi(k)=(u,\allowbreak{} b)\).  Then
\(u\in U(K)\) and \(b\in X.I.\operatorname{Iw}(0)\).

The \(U_b\) conditions on \(u=a(k)\) follow from the integral and
auxiliary-congruence conditions on \(k\).  If a tested coordinate
lies over \(\ell\), the source \(\operatorname{cmp}\) is the
trivial map, whose value belongs to each tested subgroup; away from
\(\ell\) it is exactly \(k_w\).  For each \(q\), put
\(h_q=k_{w_q}\in K_q\).  FNL5 gives
\(y_q=\operatorname{sec}_{w_q}(h_q)\), with its own \(q\) coordinate
\(h_q\), all other added-place coordinates equal to \(1\), and
all coordinates over places outside \(V.S'\) equal to \(1\).
Applying FNL1 inside the intersection of the indicated kernel
subgroups and \(\bigvee_q\operatorname{sec}_{w_q}(K_q)\) gives
\(y\) with all these properties simultaneously.

Set \(u_0=uy^{-1}\).  At \(w_q\) its coordinate is
\(h_qh_q^{-1}=1\); here \(w_q\nmid\ell\) follows from
\(v_q\notin S\).  At a place over \(F\setminus V.S'\), the
coordinate of \(y\) is \(1\), so \(u_0\) remains integral.
Therefore \(u_0\in U_Q\), and
\(u=u_0y\in U_Q\vee\bigvee_q\operatorname{sec}_{w_q}(K_q)=U(K)\).
This order of factors uses no unproved commutation or normality.
The assertion about \(b\) follows from its coordinatewise
\(\mathrm I_1(0)\) membership and FNL7.
% END FNL9

% BEGIN FNL10
\paragraph{FNL10. Full original finite-level theorem.}
\label{claim:FNL10}
The assertion \eqref{eq:fnl-target} holds for every original datum
specified above.  In particular the full original
\texttt{Upper\allowbreak Theory.\allowbreak V.\allowbreak Hi.\allowbreak fin\allowbreak U} is reconstructed.

Apply FNL3 to the open subgroup \(K_0\cap U_f\), using FNL8,
and obtain a finite set \(\mathcal T\subset U_{\mathbb A}\).
Set \(\mathcal S=\Phi(\mathcal T)\).  Given any element of the
actual frame \(G^a\times G^\ell\), lift it to \(g\in U_{\mathbb A}\)
by FNL6 and write \(g=\gamma t z k\) by FNL3.  Applying \(\Phi\)
and using FNL6 and FNL9 gives
\[
 \Phi(g)=\operatorname{glob}(\gamma)\,\Phi(t)\,\Phi(z)\,\Phi(k)
         =\operatorname{glob}(\gamma)\,\Phi(t)\,(u,\allowbreak{} b),
 \quad u\in U(K),\ b\in X.I.\operatorname{Iw}(0).
\]
This is exactly the defining statement of \(X.\texttt{Fin\allowbreak Level}\).
The coefficient representation \(\tau _0\), its dimension \(m\),
and the coefficient ring have not been restricted at any step.
The given rank, auxiliary prime, actual section maps, Iwahori factors,
and order of multiplication have all been retained.  No condition
\(\ell>n\) is introduced; in particular rank four at \(\ell=3\)
remains in the original scope.
% END FNL10

\subsection{The remaining original Hecke-setting axioms}

The finite-level theorem also permits completion of the original
\texttt{HSet.\allowbreak Ax} package.  Its group-theoretic fields must refer to
the same sections as FNL5.  The arithmetic ingredients are inherited
CIP4 (the exact chosen local Artin map) and CAP6 (degree-one
Chebotarev with a prescribed Frobenius coset).  This removes an
input-producing task; it is distinct from the full ordinary/control
or Hida construction.

% BEGIN FNL11
\paragraph{FNL11. Coordinate equality, opposite places, and commuting supports.}
\label{claim:FNL11}
The family \(\operatorname{cmp}_w\) separates points of \(G^a\).
If \(v(w')=v(w)\), then
\[
 \operatorname{cmp}_w(u)=1\ \Longrightarrow\
 \operatorname{cmp}_{w'}(u)=1,\qquad
 \operatorname{cmp}_w(u)\in\mathrm{GL}_n(\mathcal O_{L_w})
 \ \Longrightarrow\
 \operatorname{cmp}_{w'}(u)\in\mathrm{GL}_n(\mathcal O_{L_{w'}}).
\]
The first assertion is the definition of the range inside a product.
For the latter assertions lift \(u\) to \(U_{\mathbb A}\).
Above a nonsplit place there is only one coordinate.  Above a split
place the other coordinate is the field conjugate inverse transpose,
which preserves the identity and integral invertible matrices.
At places above \(\ell\), both \(\operatorname{cmp}\)'s are the
trivial map, so the assertions still hold.

If \(u,\allowbreak{} x\in U_b\), \(\operatorname{cmp}_w(u)=1\), and \(x\) is
supported over \(v(w)\), then \(ux=xu\).  On the two possible
coordinates above \(v(w)\), the preceding assertion makes \(u=1\);
on every other coordinate \(x=1\).  Thus the products agree on
every coordinate.  This proves the equality in the literal subgroup,
not only in an abstract local quotient.  Two sections supported over
distinct \(F\)-places therefore commute, even before their inputs
are restricted to compact subgroups.
% END FNL11

% BEGIN FNL12
\paragraph{FNL12. Coefficient triviality and the three section commutation fields.}
\label{claim:FNL12}
Put
\(W'=\{w:v(w)\notin V.S',\ w\text{ split}\}\), the original
\texttt{HSet.\allowbreak Wp}.  For every added place \(w_q\) and every
\(w\in W'\), the original coefficient representation
\(\tau=\tau_0\circ\operatorname{toKT}\) is trivial on the image of
the corresponding section.  Each such section has trivial coordinates
at every place in \(s.S_t\), since its supporting \(F\)-place is
outside \(S\).  Hence its image under \(\operatorname{toKT}\) is
\(1\), and its image under \(\tau_0\) is \(1\).
This supplies both \texttt{tau\_\allowbreak Q} and
\texttt{coefficient\allowbreak Rep\_\allowbreak trivial\_\allowbreak on\_\allowbreak unramified}.

An element of \(U_Q\) has trivial \(w_q\) coordinate, so FNL11
makes it commute with \(\operatorname{sec}_{w_q}(h)\) for every
\(h\).  Distinct added places have different underlying places by
\(Q.\texttt{inj}\), so their sections commute by FNL11.  A place
in \(W'\) lies outside \(V.S'\), while each \(v_q\) lies inside
it, so those sections commute too.  These prove exactly
\texttt{Q\_\allowbreak UQ}, \texttt{Q\_\allowbreak Q}, and
\texttt{sections\allowbreak Commute\allowbreak At\allowbreak Split}.
% END FNL12

% BEGIN FNL13
\paragraph{FNL13. Integral unramified sections and the required commuting remainder.}
\label{claim:FNL13}
For every \(w\in W'\),
\[
 \operatorname{sec}_w(\mathrm{GL}_n(\mathcal O_{L_w}))\subset U_Q.
\]
At \(w\) this follows from the original inverse-section identity
FNL5.  At the conjugate place it follows from FNL11; away from the
pair the section is \(1\).  Each added-place coordinate is \(1\)
because its underlying place is in \(V.S'\).  These are all defining
conditions of \(U_Q\), proving the inclusion.

For \(u\in U_Q\) let \(h=\operatorname{cmp}_w(u)\), which is
integral, and put \(u'=\operatorname{sec}_w(h)^{-1}u\).
The inclusion just proved implies \(u'\in U_Q\), and FNL5 gives
\(\operatorname{cmp}_w(u')=h^{-1}h=1\).
Thus
\[
 u=\operatorname{sec}_w(h)u',\qquad
 [u',\operatorname{sec}_w(g)]=1\quad\hbox{for every }g,
\]
where the second equality is FNL11.  This is the original
\texttt{W\_\allowbreak dec} order, and the preceding inclusion is
\texttt{W\_\allowbreak le}.
% END FNL13

% BEGIN FNL14
\paragraph{FNL14. The same-place alternative preserves the actual sections and maximal compact groups.}
\label{claim:FNL14}
Let \(w,\allowbreak{} w'\in W'\).  If their underlying \(F\)-places differ,
all their section images commute.  If their underlying places agree,
there is a group isomorphism
\[
 e:\mathrm{GL}_n(L_{w'})\longrightarrow\mathrm{GL}_n(L_w)
\]
carrying the integral general linear group onto the integral general
linear group and satisfying
\(\operatorname{sec}_{w'}(g)=\operatorname{sec}_w(e(g))\) for every
\(g\).  For equal places take the identity.  For the two conjugate
places take the field-conjugate inverse-transpose isomorphism.
FNL5 identifies both sides coordinatewise, and FNL11 verifies the
integral equivalence in both directions.  This proves the exact
\texttt{Wm\_\allowbreak mul\_\allowbreak Wm\_\allowbreak comm} alternative; it never assumes that two
opposite-place sections commute.
% END FNL14

% BEGIN FNL15
\paragraph{FNL15. Density of inverse Frobenius on the original quotient.}
\label{claim:FNL15}
Let \(G_0=\ker\eta\subset G_{F,\allowbreak{} S'}\) be the original quotient group
in \(V.\texttt{frame}\), and
\(p=V.\operatorname{pr}:G_L\to G_0\) the original map.  It is a
continuous surjection.  Moreover
\begin{equation}\label{eq:fnl-dense}
 \{p(\sigma^{-1}):w\in W',\ A\text{ lies over }w,\
       \sigma\text{ is arithmetic Frobenius at }A\}
                    \quad\text{is dense in }G_0.
\end{equation}

The original quotient map \(q:G_F\to G_{F,\allowbreak{} S'}\) is continuous and
onto.  Its sign character has
\(\eta(q(g))=1\) exactly when \(g\) fixes \(L\), by the stored
\(V.\texttt{h}\eta\) identity.  A lift under \(q\) of any element
of \(G_0\) therefore belongs to \(G_L\), proving surjectivity of
\(p\).  Continuity follows from continuity of restriction to the
same algebraic closure, the quotient topology, and the subgroup
topology.  Consequently \(g\mapsto p(g^{-1})\) is continuous and
surjective as well.

Only finitely many places of \(L\) lie over \(V.S'\): there are
finitely many base places, each with finitely many factors in the
number field extension.  Add the finitely many places ramified in
\(L/F\) to this excluded set.  CAP6, applied with \(E=L\), says
that arithmetic Frobenius elements at the remaining degree-one
places meet every open coset of \(G_L\).  Such cosets form a basis
of its profinite topology, so those Frobenius elements are dense.
For one of these places \(w\), absolute degree one means
\(k_w=\mathbb F_p\).  Hence also \(k_{v(w)}=\mathbb F_p\) and the
relative residue degree is one.  The relative ramification index is
one by our exclusion.  The quadratic Galois degree formula
\(2=e f g\) then gives two primes above \(v(w)\), so \(w\) is split.
It is therefore in \(W'\).

The image of a dense subset under a continuous surjection is dense:
every nonempty open set has a nonempty open preimage.  Apply this
to \(g\mapsto p(g^{-1})\) and the Frobenius set just obtained.
This proves \eqref{eq:fnl-dense} with the source's exact inverse
and its arbitrary valuation-over-place witnesses.  CAP6 supplies
both those witnesses and the actual arithmetic Frobenius action;
no density assertion about an unrelated quotient is substituted.
% END FNL15

% BEGIN FNL16
\paragraph{FNL16. The complete original \texttt{HSet.\allowbreak Ax} on the constructed setting.}
\label{claim:FNL16}
For every original datum in FNL10, the setting
\(H=\texttt{hset}\,\allowbreak{} Q\,\tau_0\) satisfies all eighteen fields of
\(H.\texttt{Ax}\).

Here is the complete assembly, retaining the original field names.
\begin{center}\small
\begin{tabular}{p{0.43\linewidth}p{0.49\linewidth}}
\texttt{odd}, \texttt{pos} & The stored \(s.\texttt{odd}\) and \(n\ge1\).\\
\texttt{Sl\_\allowbreak over}, \texttt{Sl\_\allowbreak inj} & FNL4, by contraction of places.\\
\texttt{art\_\allowbreak spec} & CIP4 on the same stored valuation and \(\operatorname{artOf}\).\\
\texttt{Q\_\allowbreak away}, \texttt{Wp\_\allowbreak away} & The defining place exclusions and FNL4.\\
\texttt{fin} & FNL10 for every permitted \(K\).\\
\texttt{tau\_\allowbreak Q}, \texttt{coefficient\allowbreak Rep\_\allowbreak trivial\_\allowbreak on\_\allowbreak unramified} & FNL12.\\
\texttt{Q\_\allowbreak UQ}, \texttt{Q\_\allowbreak Q}, \texttt{sections\allowbreak Commute\allowbreak At\allowbreak Split} & FNL12.\\
\texttt{W\_\allowbreak le}, \texttt{W\_\allowbreak dec} & FNL13.\\
\texttt{Wm\_\allowbreak mul\_\allowbreak Wm\_\allowbreak comm} & FNL14.\\
\texttt{prj\_\allowbreak cont}, \texttt{dense} & FNL15.
\end{tabular}
\end{center}
In particular \(w(v)=w(v')\) implies \(v=v'\) by contraction,
which verifies \texttt{Sl\_\allowbreak inj} without a choice of new embeddings.
The local Artin field uses CIP4's validity of the original Hilbert-choice
map, so no further \texttt{Literature\allowbreak Rank\allowbreak NHypothesis} premise is needed
in this human proof beyond its already proved class-field-theory input.
This is the full original \texttt{ax\allowbreak U} conclusion and, in particular,
the conclusion of \texttt{CLoc\allowbreak Conj.\allowbreak ax\allowbreak U\_\allowbreak s}.  The latter source variant
carries a larger literature package because of its Artin-map route;
CIP4 is a proved replacement for that use.
% END FNL16

\subsection{What this supplies, and what still has to be constructed}

Together with NTL17 and OTW12, FNL10 supplies finite level and the actual
\texttt{Frame.\allowbreak Free} stabilizer condition needed at the entrance to
ordinary theory.  FNL16 supplies the entire original
\texttt{HSet.\allowbreak Ax}, including inverse-Frobenius density, so this package
is no longer an unconstructed antecedent.  The argument does not by
itself supply ordinary Hecke commutation, classical control,
arithmetic-point separation, the integral type, or same-ideal all-level
automorphy.  The outer \texttt{COUNT\_\allowbreak II\_\allowbreak SUPPORT} obligation remains
open until those producers have been proved.  The finite-class,
local-reciprocity and Chebotarev results are used as completed
mathematical inputs with their actual objects and maps checked.


\section{Ordinary projectors, finite freeness, and coefficient base change}
\label{sec:ordinary-structure-v42}

\status{This chapter reconstructs the original generic endpoints
\texttt{Ord\allowbreak Free.\allowbreak ord\allowbreak Struct}, \texttt{Ord\allowbreak Free.\allowbreak ord\allowbreak Free}, and
\texttt{Ord\allowbreak BCMap\allowbreak Ord.\allowbreak ord\allowbreak BC}, including their elementary input producers.
The hypotheses are precisely the original finite-level, stabilizer, and
continuous-character hypotheses on a supplied Frame.  The coefficient
rings are arbitrary complete Noetherian local rings with finite residue
fields; reducedness, flatness of a coefficient map, and characteristic
zero are not required.  The construction uses the actual function spaces,
double-coset operators, and coefficientwise maps.  Applying it to the
original arithmetic data still uses the separately constructed Frame and
its finite-level and weight inputs.  Classical control, arithmetic fibers,
and the same-ideal Hida endpoints are not conclusions of this chapter.}

\subsection{The actual spaces and their hypotheses}

Write $G=G_a\times G_p$ and let $\operatorname{glob}:\Gamma\to G$
and $U_b\leq G_a$ be the supplied Frame.  Its Iwahori datum consists of
subgroups $N\leq N_f$, $T\leq T_f$, decreasing subgroups
$T_c(c)\leq T$ and $N^-(c)$, a submonoid $P\subset T_f$, and
$z\in P$.  We retain every original Iwahori axiom.  The ones used
below have the following mathematical content:

\begin{enumerate}
\item $T_f$ is commutative and normalizes $N_f$; $T$ normalizes
both $N$ and every $N^-(c)$.
\item For $a\in P$, $aNa^{-1}\subset N$ has finite index, and
$a^{-1}N^-(c)a\subset N^-(c)$.  In addition,
$z^{-1}N^-(c)z\subset N^-(c+1)$.
\item For $y\in N^-(c)$ and $x\in N$, there are
$x'\in N$, $t\in T_c(c)$, and $y'\in N^-(c)$ with
$yx=x'ty'$.
\item If $xty=1$, with $x\in N_f$, $t\in T_f$, and
$y\in N^-(0)$, then $x=t=1$ (and hence $y=1$).
\item The indices $[N^-(0):N^-(c)]$ and $[T:T_c(c)]$ are
finite.  For every $a\in P$, some $h\geq0$ satisfies
$z^h a^{-1}\in P$.
\end{enumerate}

Put $I(c)=\langle N,\allowbreak{} T,\allowbreak{} N^-(c)\rangle$ and
$I_1(c)=\langle N,\allowbreak{} T_c(c),\allowbreak{} N^-(c)\rangle$.  Let $(A,\mathfrak m)$
be a complete Noetherian local ring with finite residue field, let
$d\geq0$, and fix
\[
 \tau:U_b\longrightarrow\operatorname{GL}_d(A),\qquad
 \theta:T\longrightarrow A^\times,\qquad U\leq U_b.
\]
The original continuous-character condition is
\begin{equation}\label{eq:ost-contchar}
 \forall k\geq0\ \exists c\geq0\quad
             \theta(t)-1\in\mathfrak m^k\quad(t\in T_c(c)).
\end{equation}
The finite-level condition is finiteness of
$\operatorname{glob}(\Gamma)\backslash G/(U\times I(0))$,
with the factorization order
$g=\operatorname{glob}(\gamma)x(u,\allowbreak{} b)$.  The stabilizer condition
is the original \texttt{Free}: whenever
\begin{equation}\label{eq:ost-stabilizer}
 \operatorname{glob}(\gamma)g=g(u,\allowbreak{} b),\quad u\in U,\ b\in I(0),
\end{equation}
one has $b\in Z(G_p)\cap T$ and $\theta(b)\tau(u)=1$.
This is a condition on stabilizers; it does not presume module freeness.

Let $\mathcal F_A$ be the module of all functions $G\to A^d$.
For $a\in G_p$ write $r_a f(g)=f(g(1,\allowbreak{} a))$.  The actual space
$M_A=\mathrm{Sp}_A(\tau,\theta,\allowbreak{} U)$ consists of functions satisfying
\begin{align}
 f(\operatorname{glob}(\gamma)g)&=f(g), &
 r_n f&=f\quad(n\in N),\\
 r_t f&=\theta(t)^{-1}f\quad(t\in T), &
 \tau(u)f(g(u,1))&=f(g)\quad(u\in U),\\
 \forall k\ \exists c\quad
 r_y f-f&\in\mathrm{valIn}_A(k)\quad(y\in N^-(c)).
 \label{eq:ost-space}
\end{align}
Here $\mathrm{valIn}_A(k)$ means that every coordinate of every
function value lies in $\mathfrak m^k$.  Write
$V_A(k)=M_A\cap\mathrm{valIn}_A(k)$.  In particular these
are value filtrations, before any identification with intrinsic
module filtrations.

For $a\in P$ choose representatives $R_a$ of the left cosets
$N/(aNa^{-1})$.  On $N$-fixed functions the original operator is
\begin{equation}\label{eq:ost-ua}
                 U_a f=\sum_{n\in R_a}r_{na}f.
\end{equation}
It agrees with the source's sum using quotient representatives.
The ordinary space is the literal source submodule
\begin{equation}\label{eq:ost-ord-definition}
 O_A=M_A\cap\bigcap_{k\geq0}
       \sum_{q\geq1}\ker\bigl(U_z^q-1:
                         \mathcal F_A\to\mathcal F_A/
                                      \mathrm{valIn}_A(k)\bigr).
\end{equation}
The sum notation in this definition is made concrete in OST10.

\subsection{Finite coordinates without changing the function module}

% BEGIN OST1
\paragraph{OST1. Iwahori factorization and finite representatives at every depth.}
\label{claim:OST1}
One has
\[
 I(c)=NTN^-(c)=N^-(c)TN,\qquad
 I_1(c)=NT_c(c)N^-(c),\qquad I(c)\cap N_f=N.
\]
Moreover $[I(0):I(c)]$ and $[I(0):I_1(c)]$ are finite.
Consequently the finite-level hypothesis gives, for every $c$, a
finite list $x_1,\ldots,\allowbreak{} x_s$ of distinct double-coset representatives
such that every $g\in G$ can be written
\begin{equation}\label{eq:ost-cosets}
             g=\operatorname{glob}(\gamma)x_i(u,\allowbreak{} b),
                       \quad u\in U,\quad b\in I(c).
\end{equation}
The index $i$ is unique.
\begin{proof}
More generally let $T_c(c)\leq S\leq T$.  The set $NSN^-(c)$
is closed under multiplication: in a product $xty\,\allowbreak{} x't'y'$, use
$yx'=x''s y''$ with $s\in T_c(c)$, move torus elements through
$N$ and $N^-(c)$ by their normalization, and collect the three
factors.  It contains the three subgroups, and it is closed under
inversion, since $(xty)^{-1}=y^{-1}t^{-1}x^{-1}$ is a product
of three elements of that same multiplicatively closed set.
It is therefore their generated subgroup.  This proves the forward
factorizations; inversion proves the reversed one.  If $g=xty$
also belongs to $N_f$, apply the uniqueness axiom to
$(g^{-1}x)ty=1$.  It gives $g=x\in N$.

The reversed factorization gives $I(0)=N^-(0)I(c)$, and
$N^-(c)\subset I(c)$.  Thus the finite coset set
$N^-(0)/N^-(c)$ surjects onto $I(0)/I(c)$.
Similarly $I(c)=T I_1(c)$ and $T_c(c)\subset I_1(c)$, so
$[I(c):I_1(c)]\leq[T:T_c(c)]$.
Refine a finite list at level $I(0)$ by representatives of
$I(0)/I(c)$ and discard duplicate double cosets.  This gives
\eqref{eq:ost-cosets} with the stated uniqueness.  No normality
of $I(c)$ or averaging denominator is used.
\end{proof}
% END OST1

% BEGIN OST2
\paragraph{OST2. The torus coordinate and its congruence variation.}
\label{claim:OST2}
There is a well-defined function $t_0:I(0)\to T$ characterized by
$t_0(y t n)=t$ for $y\in N^-(0)$, $t\in T$, $n\in N$.
It satisfies
\begin{gather}
 t_0(1)=1,\quad t_0(sb)=s t_0(b),\quad
 t_0(bs)=t_0(b)s,\quad t_0(bn)=t_0(b),\\
 y\in N^-(c)\quad\Longrightarrow\quad
 t_0(by)=t_0(b)s_c\quad\hbox{for some }s_c\in T_c(c).
 \label{eq:ost-torus}
\end{gather}
Here $s\in T$, $n\in N$, and $b\in I(0)$.
\begin{proof}
Existence of the reversed factorization is OST1.  If
$ytn=y't'n'$, rearranging gives an identity with successive factors
in $N_f,\allowbreak{} T_f,\allowbreak{} N^-(0)$.  Explicitly, from
$n=t^{-1}y^{-1}y't'n'$ one obtains
\[
 (nn'^{-1})(t'^{-1}t)
           \bigl(t^{-1}y'^{-1}yt\bigr)=1.
\]
The last factor belongs to $N^-(0)$ by torus normalization.
The uniqueness axiom gives $t'^{-1}t=1$.
The first four formulas now follow by moving a torus factor through
$N^-(0)$ or $N$, or multiplying two $N$ factors.

For the last formula write $b=y_0tn$.  Apply the swap axiom at
depth $c$ to $y^{-1}n^{-1}=n'sy'$ with $s\in T_c(c)$.
Then $ny=(y')^{-1}s^{-1}(n')^{-1}$.  Move $(y')^{-1}$
past $t$ and combine it with $y_0$; the resulting reversed
factorization of $by$ has torus coordinate $ts^{-1}$.
Since $T_c(c)$ is a subgroup, this is the required assertion.
\end{proof}
% END OST2

% BEGIN OST3
\paragraph{OST3. Integral sections for the actual coefficient functions.}
\label{claim:OST3}
Fix the representatives of OST1 at depth $c$.  There is an
$A$-linear map
\[
 s_c:(A^d)^s\longrightarrow M_A
\]
such that $(s_c v)(x_i)=v_i$.  If
$\theta(T_c(c))\subset1+\mathfrak m^k$, then every $s_cv$ is
$N^-(c)$-invariant modulo $\mathrm{valIn}_A(k)$.
\begin{proof}
For a factorization \eqref{eq:ost-cosets} define
\begin{equation}\label{eq:ost-section}
 (s_cv)(g)=\theta(t_0(b))^{-1}\tau(u^{-1})v_i.
\end{equation}
To verify independence of the factorization, compare
$\operatorname{glob}(\gamma)x_i(u,\allowbreak{} b)$ and
$\operatorname{glob}(\gamma_1)x_i(u_1,\allowbreak{} b_1)$.  The index is the
same by OST1.  The resulting stabilizer relation is
\[
 \operatorname{glob}(\gamma_1^{-1}\gamma)x_i
       =x_i(u_1u^{-1},\allowbreak{} b_1b^{-1}).
\]
By \eqref{eq:ost-stabilizer}, $s=b_1b^{-1}\in T$ and
$\theta(s)\tau(u_1u^{-1})=1$.  Hence
\[
 t_0(b_1)=s t_0(b),\qquad
 \tau(u_1^{-1})=\theta(s)\tau(u^{-1}).
\]
These factors cancel in \eqref{eq:ost-section}.  The formula is
linear in $v$ and gives $v_i$ at $x_i$.

Left $\Gamma$-invariance follows by changing $\gamma$.
Right $N$-invariance and the $T$-character law follow from OST2.
For $u'\in U$, replace $u$ by $uu'$ and use
$\tau(u')\tau((uu')^{-1})=\tau(u^{-1})$; this is the required
away-factor equivariance.  Finally fix $k$ and choose
$c'\geq c$ with $\theta(T_c(c'))\subset1+\mathfrak m^k$.
For $y\in N^-(c')$, the last formula of OST2 changes
\eqref{eq:ost-section} by multiplication by
$\theta(s_{c'})^{-1}$, which is congruent to one modulo
$\mathfrak m^k$.  Indeed $a^{-1}-1=-a^{-1}(a-1)$ for a unit
$a$.  All entries of $\tau(u^{-1})$ lie in $A$, so the
congruence holds for every $g$ and every coordinate.  This proves
the last condition of \eqref{eq:ost-space}.  If the character is
already trivial at depth $c$ modulo $\mathfrak m^k$, take $c'=c$.
\end{proof}
% END OST3

% BEGIN OST4
\paragraph{OST4. Finite values detect congruences at a fixed depth.}
\label{claim:OST4}
Suppose $f\in M_A$, $r_yf-f\in\mathrm{valIn}_A(k)$ for
$y\in N^-(c)$, and $f(x_i)\in(\mathfrak m^k)^d$ at all
representatives of OST1.  Then $f\in V_A(k)$.
\begin{proof}
Write $g=\operatorname{glob}(\gamma)x_i(u,\allowbreak{} b)$ and factor
$b=nty$ by OST1.  Modulo $\mathfrak m^k$, remove $y$ using
the assumed invariance, remove $t$ using its unit character,
and remove $n$ using exact $N$-invariance.  Away-factor
equivariance then expresses $f(g)$ modulo $\mathfrak m^k$ as
$\theta(t)^{-1}\tau(u^{-1})f(x_i)$.  Its entries belong to
$\mathfrak m^k$.  This proof uses precisely the integral matrix
$\tau(u^{-1})$, not an unspecified lattice in an isomorphic space.
\end{proof}
% END OST4

\subsection{The original operators and level lowering}

% BEGIN OST5
\paragraph{OST5. Double-coset algebra with no division.}
\label{claim:OST5}
On $N$-fixed vectors in any representation, the sum
\eqref{eq:ost-ua} is independent of its representatives, has
$N$-fixed image, and satisfies
\[
 U_aU_b=U_{ab},\qquad U_1=1,\qquad r_tU_a=U_a r_t
                    \quad(a,\allowbreak{} b\in P,\ t\in T).
\]
For function representations it commutes with all coefficientwise
ring maps, preserves every value filtration, and commutes with
the away action $f(g)\mapsto\tau(u)f(g(u,1))$.
\begin{proof}
Replacing $n\in R_a$ by $n a n_0a^{-1}$ changes
$r_{na}f$ into $r_{na}r_{n_0}f$, which is the same vector.
Left multiplication of $R_a$ by an element of $N$ permutes its
cosets, proving $N$-invariance.  If $R_a$ and $R_b$ are chosen,
then
\[
             \{n(a n'a^{-1}):n\in R_a,\ n'\in R_b\}
\]
is a representative set for $N/(abN(ab)^{-1})$: first reduce
modulo $aNa^{-1}$ and then conjugate the remaining factor by
$a^{-1}$ and reduce modulo $bNb^{-1}$.  Uniqueness follows
by reversing those two steps.  Expanding the two finite sums
proves the product identity.  For $a=1$ there is one coset.
Conjugation by $t\in T$ permutes the quotient because $t$
normalizes $N$ and commutes with $a$; this proves the torus
identity.  The final assertions follow term by term from the
finite sum.  No index is inverted, so these identities hold in
every characteristic and with nilpotent coefficients.
\end{proof}
% END OST5

% BEGIN OST6
\paragraph{OST6. Level preservation and strict lowering by $U_z$.}
\label{claim:OST6}
If a vector $v$ is fixed by $N$, by $T_c(c)$, and by every
$a^{-1}N^-(c)a$, then $U_av$ is fixed by $N^-(c)$.
Consequently $U_a$ preserves $I_1(c)$-invariants, and if
$d'\geq c$, then $U_z$ sends vectors fixed by
$N,\allowbreak{} T_c(c),\allowbreak{} N^-(d'+1)$ to vectors fixed by
$N,\allowbreak{} T_c(c),\allowbreak{} N^-(d')$.
\begin{proof}
Fix $y\in N^-(c)$.  For each representative $n\in R_a$,
the swap axiom gives $yn=x t y'$ with $x\in N$,
$t\in T_c(c)$, $y'\in N^-(c)$.  Choose $n'\in R_a$
with $x=n'a n_0a^{-1}$, $n_0\in N$.  Since $t$ commutes
with $a$,
\[
 yn a=n'a k_n,\qquad k_n=n_0t(a^{-1}y'a)\in I(c),
                       \qquad r_{k_n}v=v.
\]
The resulting map $n\mapsto n'$ is injective.  If two images
coincide, then
$a^{-1}n^{-1}\widetilde n a=k_n^{-1}k_{\widetilde n}$
belongs to $I(c)\cap N_f=N$ by OST1.  Thus $n$ and
$\widetilde n$ represent the same coset and are equal.
It is a permutation of the finite set, so summing gives
$r_yU_av=U_av$.

The remaining $N$ and torus invariances are OST5.
For level preservation use
$a^{-1}N^-(c)a\subset N^-(c)$.
For strict lowering apply the argument at depth $d'$,
using $T_c(d')\subset T_c(c)$ and
$z^{-1}N^-(d')z\subset N^-(d'+1)$.
\end{proof}
% END OST6

% BEGIN OST7
\paragraph{OST7. Stability of the actual automorphic space and eventual lowering.}
\label{claim:OST7}
Every $U_a$, $a\in P$, preserves $M_A$ and its value
filtration.  If $\theta(T_c(c))\subset1+\mathfrak m^k$,
then for every $f\in M_A$ there is $n_0$ such that
\begin{equation}\label{eq:ost-eventual-lowering}
 n\geq n_0,\ y\in N^-(c)
    \quad\Longrightarrow\quad
        r_yU_z^nf-U_z^nf\in\mathrm{valIn}_A(k).
\end{equation}
The difference is measured in the ambient function module;
right translation by $y$ is not assumed to preserve $M_A$.
\begin{proof}
Left $\Gamma$-invariance, the $N$ and $T$ laws, and the
away action are preserved by OST5.  For continuity modulo
$\mathfrak m^k$, choose a depth at which both the character
and $f$ are trivial on the respective congruence subgroups.
The reduction of $f$ in $(A/\mathfrak m^k)^d$ is then
$I_1(c')$-fixed.  OST6 makes its image under $U_a$
$I_1(c')$-fixed as well.  Finite sums commute with reduction,
so the required uniform continuity lifts back to $A$.

For the second assertion, choose $c_1$ such that $f$ is
$N^-(c_1)$-invariant modulo $\mathfrak m^k$.  For $n\geq c_1$,
its reduction is fixed by $N,\allowbreak{} T_c(c),\allowbreak{} N^-(c+n)$.
Apply the last assertion of OST6 $n$ times.  The reduction of
$U_z^nf$ is fixed by $N^-(c)$, proving the congruence.
The threshold depends on $f,\allowbreak{} k,\allowbreak{} c$, as permitted by the source.
\end{proof}
% END OST7

% BEGIN OST8
\paragraph{OST8. Separated completeness and the initial value-filtration rules.}
\label{claim:OST8}
The modules $V_A(k)$ decrease with $k$ and satisfy
\[
 \mathfrak m^kM_A\subset V_A(k),\qquad
 \mathfrak m^kV_A(1)\subset V_A(k+1),\qquad
 \bigcap_kV_A(k)=0.
\]
If $f_{k+1}-f_k\in V_A(k)$, there is a unique $f\in M_A$
with $f-f_k\in V_A(k)$ for all $k$.
\begin{proof}
The inclusions follow by multiplying coordinates in the relevant
ideal powers.  Separation follows from separation of $A$.
For existence, take the $\mathfrak m$-adic limit of each
coordinate $f_k(g)_i$.  For every fixed precision the convergence
is uniform over all $g,\allowbreak{} i$, because the difference hypothesis is
a value-filtration statement.  All exact linear transformation
laws in \eqref{eq:ost-space} pass to the limit; the matrices and
scalar characters preserve the ideal powers.  For the continuity
law at precision $k$, use the depth supplied by $f_k$ and write
\[
 r_yf-f=r_y(f-f_k)+(r_yf_k-f_k)-(f-f_k).
\]
Every term is zero modulo $\mathfrak m^k$ uniformly in $g,\allowbreak{} i$.
Thus the limit lies in $M_A$, and uniqueness again follows
from separation.
\end{proof}
% END OST8

% BEGIN OST9
\paragraph{OST9. Finite-coordinate approximation of late operator iterates.}
\label{claim:OST9}
For every $k$ there are a finite set $D_k$, an $A$-linear
evaluation map $v_k:M_A\to A^{D_k}$, and a section
$s_k:A^{D_k}\to M_A$, such that
\begin{align}
 f\in V_A(j)&\ \Longrightarrow\ v_k(f)\in(\mathfrak m^j)^{D_k},
 \label{eq:ost-evaluation-filtration}\\
 \forall f\ \exists n_0\ \forall n\geq n_0\quad
 U_z^nf-s_kv_k(U_z^nf)&\in V_A(k).
 \label{eq:ost-finite-approx}
\end{align}
\begin{proof}
Choose $c$ using \eqref{eq:ost-contchar}, choose $x_i$ using
OST1, and set $D_k=\{1,\ldots,\allowbreak{} s\}\times\{1,\ldots,\allowbreak{} d\}$.
Let $v_k$ evaluate at those coordinates and let $s_k$ be OST3.
The first inclusion is immediate.  For large $n$, OST7 and
OST3 say that the two terms in \eqref{eq:ost-finite-approx}
are invariant modulo $\mathfrak m^k$ under $N^-(c)$.
Their difference vanishes at every $x_i$.  OST4 gives the
desired value congruence.  Only these finitely many evaluations
have finite reductions; finiteness of the whole quotient
$M_A/V_A(k)$ is neither asserted nor needed.
\end{proof}
% END OST9

\subsection{A complete-ring ordinary projector}

% BEGIN OST10
\paragraph{OST10. The original ordinary module is periodic at each precision.}
\label{claim:OST10}
Writing $T_0=U_z|_{M_A}$, the original definition is equivalent to
\begin{equation}\label{eq:ost-periodic}
 f\in O_A\quad\Longleftrightarrow\quad
 f\in M_A\ \hbox{and}
 \forall k\ \exists q\geq1\quad T_0^qf-f\in V_A(k).
\end{equation}
Every $U_a$ with $a\in P$ preserves $O_A$.
\begin{proof}
If $T_0^qf-f\in V_A(k)$, the telescoping identity
\[
 T_0^{qr}f-f=\sum_{i=0}^{r-1}T_0^{qi}(T_0^qf-f)
\]
shows the same assertion for every positive multiple of $q$.
Thus the kernels in \eqref{eq:ost-ord-definition} form a
directed family: periods $q,\allowbreak{} q'$ are both replaced by $qq'$.
The sum of a directed family of submodules is its union,
which proves \eqref{eq:ost-periodic}.  OST5 gives
$U_aT_0=T_0U_a$ on $M_A$, and OST7 preserves the space and
its filtration.  Apply $U_a$ to each periodic congruence.
\end{proof}
% END OST10

% BEGIN OST11
\paragraph{OST11. Every vector is eventually periodic modulo every precision.}
\label{claim:OST11}
For $f\in M_A$ and $k\geq0$, there are $a\geq0$ and $q\geq1$
such that $T_0^{a+q}f-T_0^af\in V_A(k)$.
\begin{proof}
The ring $A/\mathfrak m^k$ is finite.  Indeed, finite generators
of $\mathfrak m$ show that each
$\mathfrak m^j/\mathfrak m^{j+1}$ is finite dimensional over
the finite residue field, and induction proves the assertion.
Use OST9 and examine $v_k(T_0^nf)$ modulo $\mathfrak m^k$
for $n\geq n_0$.  Two of these tuples are equal, at indices
$a<a+q$.  Their section values differ by an element of
$\mathfrak m^kM_A\subset V_A(k)$, since $s_k$ is linear
on a finite free module.  Subtracting the two approximation
congruences proves the assertion.  Subsequent iterates repeat
with period $q$ because $T_0$ preserves $V_A(k)$.
\end{proof}
% END OST11

% BEGIN OST12
\paragraph{OST12. One factorial-limit projector on the complete space.}
\label{claim:OST12}
There is a unique $A$-linear endomorphism $e_A$ of $M_A$
satisfying
\begin{equation}\label{eq:ost-factorial-limit}
 \forall f\in M_A\ \forall k\ \exists J\ \forall j\geq J
                 \quad T_0^{j!}f-e_Af\in V_A(k).
\end{equation}
It preserves every $V_A(k)$, has image exactly $O_A$, and
is the identity on $O_A$.  In particular $e_A^2=e_A$.
\begin{proof}
Fix $f,\allowbreak{} k$ and the eventual period $q$ and starting index $a$
from OST11.  Once $j!\geq a$ and $q\mid j!$, the residue
of $T_0^{j!}f$ modulo $V_A(k)$ is independent of $j$:
all those exponents lie in the same residue class modulo $q$
in the periodic tail.  Therefore the factorial sequence is
Cauchy in the value filtration.  To use exactly OST8, choose
increasing integers $J_k$ so that all factorial terms of
index at least $J_k$ agree modulo $V_A(k)$.  The sequence
$T_0^{J_k!}f$ then satisfies its successive-difference
condition.  Its limit has \eqref{eq:ost-factorial-limit}
for the entire factorial sequence, not merely this subsequence.
It is unique by separation.

For $f+g$ or $bf$, the finite sum or scalar multiple of the
limits has the same defining congruences, using the larger
of the finitely many thresholds involved.  Uniqueness proves
linearity.  If $f\in V_A(k)$, every iterate is in that
submodule, so the limit is also in it.  The eventual period
$q$ for $f$ at precision $k$ gives
$T_0^q e_Af-e_Af\in V_A(k)$ by comparing with one sufficiently
late factorial term.  Thus $e_Af\in O_A$ by OST10.
If $f\in O_A$, a period $q$ already holds starting at zero.
For $q\mid j!$, $T_0^{j!}f-f\in V_A(k)$; hence
$e_Af-f\in V_A(k)$ for every $k$, proving $e_Af=f$.
The map is chosen once before $f,\allowbreak{} k$; no uniform threshold
over all vectors is asserted.
\end{proof}
% END OST12

% BEGIN OST13
\paragraph{OST13. Finite-coordinate approximation of ordinary vectors.}
\label{claim:OST13}
The maps of OST9 also satisfy, for all $f\in O_A$,
\begin{equation}\label{eq:ost-ordinary-approx}
                   f-e_As_kv_k(f)\in V_A(k).
\end{equation}
\begin{proof}
Choose a factorial index large enough for both OST9 and OST12.
Put $g=T_0^{j!}f$.  Then $g-f\in V_A(k)$, and
$g-s_kv_k(g)\in V_A(k)$.  By
\eqref{eq:ost-evaluation-filtration},
$s_kv_k(g-f)\in\mathfrak m^kM_A\subset V_A(k)$.
The identity
\[
 f-s_kv_k(f)=(g-s_kv_k(g))-(g-f)+s_kv_k(g-f)
\]
puts its left side in $V_A(k)$.  Apply the filtration-preserving
projector and use $e_Af=f$.
\end{proof}
% END OST13

% BEGIN OST14
\paragraph{OST14. One fixed finite generating map makes every next correction.}
\label{claim:OST14}
Fix the data at precision one, let $D=D_1$, and define
$\sigma=e_As_1:A^D\to O_A$.
For every $j\geq0$ and $f\in O_A\cap V_A(j)$ there is
$b\in(\mathfrak m^j)^D$ such that
\begin{equation}\label{eq:ost-one-step}
                         f-\sigma b\in V_A(j+1).
\end{equation}
\begin{proof}
Use OST13 at precision $j+1$.  Writing the standard vectors
of $A^{D_{j+1}}$ as $\delta_p$, put
$w_p=e_As_{j+1}\delta_p$ and
$a_p=(v_{j+1}f)_p\in\mathfrak m^j$.
Then $f-\sum_pa_pw_p\in V_A(j+1)$.
At precision one, each $w_p$ satisfies
$w_p-\sigma v_1(w_p)\in V_A(1)$.
Set $b=\sum_pa_pv_1(w_p)$.  Its coordinates lie in
$\mathfrak m^j$, and
\[
 f-\sigma b=
 (f-\sum_pa_pw_p)+\sum_pa_p(w_p-\sigma v_1(w_p)).
\]
The last sum belongs to $\mathfrak m^jV_A(1)\subset V_A(j+1)$
by OST8.  This proves the claim without assuming the desired
reverse filtration inclusion in advance.
\end{proof}
% END OST14

% BEGIN OST15
\paragraph{OST15. Successive approximation over the whole coefficient ring.}
\label{claim:OST15}
For every $j$ and $f\in O_A\cap V_A(j)$ there is
$b\in(\mathfrak m^j)^D$ with $\sigma b=f$.
\begin{proof}
Start with $r_0=f$.  Apply OST14 repeatedly to choose
$b_n\in(\mathfrak m^{j+n})^D$ and
$r_{n+1}=r_n-\sigma b_n\in V_A(j+n+1)$.
The series $b=\sum_{n\geq0}b_n$ converges coordinatewise
in the finite free module $A^D$.  Its sum belongs to
$(\mathfrak m^j)^D$: this follows either from the definition
of the inverse-limit topology or by reducing every partial
sum modulo $\mathfrak m^j$.  For every $N$,
\[
 f-\sigma b=r_N+
          \sigma\Bigl(\sum_{n<N}b_n-b\Bigr).
\]
The first term is in $V_A(j+N)\subset V_A(N)$.
The second is in $\mathfrak m^NO_A\subset V_A(N)$,
since its coefficient vector belongs to $(\mathfrak m^N)^D$.
Separation proves $f=\sigma b$.  No completeness or flatness
of $O_A$ was presumed; only completeness of $A$ and the
already proved value-filtration separation were used.
\end{proof}
% END OST15

% BEGIN OST16
\paragraph{OST16. Finite generation and exact ordinary filtration.}
\label{claim:OST16}
The module $O_A$ is finite over $A$ and, for every $j\geq0$,
\begin{equation}\label{eq:ost-exact-filtration}
                     O_A\cap\mathrm{valIn}_A(j)=\mathfrak m^jO_A.
\end{equation}
\begin{proof}
The case $j=0$ of OST15 gives surjectivity of the fixed map
$\sigma:A^D\to O_A$, hence finite generation.  For general
$j$, the same result expresses every vector in the left side
of \eqref{eq:ost-exact-filtration} as
$\sum_{p\in D}b_p\sigma\delta_p$ with $b_p\in\mathfrak m^j$.
The reverse inclusion follows by multiplying function values.
This proves both directions on the literal ordinary functions.
\end{proof}
% END OST16

% BEGIN OST17
\paragraph{OST17. Finite freeness without a reducedness hypothesis.}
\label{claim:OST17}
The module $O_A$ is finite free over $A$.
\begin{proof}
Let $v=v_1|_{O_A}$ and $\Psi=\sigma v$.  By OST13 and OST16,
$f-\Psi f\in\mathfrak mO_A$ for every $f$.  Nakayama's
lemma gives surjectivity of $\Psi$.  For completeness, the
form of Nakayama used here follows by choosing finite generators
of a module $Q=\mathfrak mQ$, expressing them through a matrix
with entries in $\mathfrak m$, and multiplying by the adjugate
of $1$ minus that matrix; its determinant is a unit.  Apply
this to $Q=\operatorname{coker}\Psi$.

A surjective endomorphism of a finite module over a Noetherian
ring is injective.  Indeed, the increasing sequence of kernels
of its powers stabilizes.  If $x\in\ker\Psi$, choose $y$
with $\Psi^ny=x$ at a stabilization index $n$; then
$y\in\ker\Psi^{n+1}=\ker\Psi^n$, giving $x=0$.
Thus $\Psi$ is invertible.  The maps
\[
 O_A\xrightarrow{v}A^D
       \xrightarrow{\Psi^{-1}\sigma}O_A
\]
have composite the identity.  Hence $O_A$ is a finite
projective module.  A finite projective module over a local
ring is free: lift a basis of its residue space to obtain
a surjection $A^r\to O_A$, split it, and call the kernel
$K$.  Reducing the direct sum modulo $\mathfrak m$ gives
$K/\mathfrak mK=0$, and finite Nakayama gives $K=0$.
The argument also covers $O_A=0$ and $d=0$.
\end{proof}
% END OST17

% BEGIN OST18
\paragraph{OST18. The distinguished ordinary operator is invertible.}
\label{claim:OST18}
The restriction of $U_z$ to $O_A$ is bijective.
\begin{proof}
It preserves $O_A$ by OST10.  For $f\in O_A$, choose its
period $q\geq1$ at precision one.  Then
$f-U_z(U_z^{q-1}f)\in V_A(1)\cap O_A=\mathfrak mO_A$.
Thus $U_z$ is surjective modulo $\mathfrak m$.  Apply finite
Nakayama to its cokernel, then the Noetherian kernel argument
of OST17.  This proves bijectivity over $A$, not only over
its residue field.
\end{proof}
% END OST18

% BEGIN OST19
\paragraph{OST19. Full original ordinary-structure endpoint.}
\label{claim:OST19}
For every original Frame $X$, every $A$ of class
\texttt{Is\allowbreak CNL}, every $d\geq0$, every $\tau,\theta,\allowbreak{} U$ with
\texttt{Cont\allowbreak Char}, \texttt{Fin\allowbreak Level}, and \texttt{Free}, the
literal $X.\mathrm{Ord}(\tau,\theta,\allowbreak{} U)$ is finite free,
$X.U(z)$ is bijective on it, and its value filtration is its
intrinsic $\mathfrak m$-adic filtration.  There is one
$A$-linear map $e_A:X.\mathrm{Sp}\to X.\mathrm{Ord}$ which
fixes every ordinary function and satisfies
\eqref{eq:ost-factorial-limit} with the original
$X.U(z)$ and $X.\mathrm{valIn}$.
\begin{proof}
The spaces, maps, actions, and filtrations were identified
literally in \eqref{eq:ost-space}--\eqref{eq:ost-ord-definition}.
OST1--OST9 produce every input to the abstract construction
from the three stated hypotheses on that same Frame.
OST10--OST12 identify the projector image with its actual
ordinary submodule.  OST16--OST18 give the remaining
assertions.  Regarding $e_A$ as a map to its proved image
changes no function value.  These are exactly the fields
of \texttt{ord\allowbreak Struct}, including one projector before all
later vectors and precisions.
\end{proof}
% END OST19

% BEGIN OST20
\paragraph{OST20. Every positive-monoid operator is an ordinary unit.}
\label{claim:OST20}
For every $a\in P$, $U_a$ is bijective on $O_A$.  Together
with OST17 this is the full original \texttt{ord\allowbreak Free} endpoint.
\begin{proof}
Choose $h\geq0$ with $b=z^ha^{-1}\in P$ by the original
Iwahori axiom.  Since $a,\allowbreak{} b\in T_f$, they commute and
$ab=ba=z^h$.  OST5 gives, on $O_A$,
\[
                  U_aU_b=U_bU_a=U_z^h.
\]
The last map is invertible by OST18.  An explicit inverse of
$U_a$ is $U_bU_z^{-h}$: all these maps commute, and either
composite is the identity.  This includes $h=0$ and adds no
restriction on the residue characteristic.
\end{proof}
% END OST20

\subsection{The original coefficientwise base-change theorem}

% BEGIN OST21
\paragraph{OST21. Local coefficient maps preserve every original condition.}
\label{claim:OST21}
Let $\varphi:A\to B$ be any local homomorphism of complete
Noetherian local rings with finite residue fields.  Put
$\tau_B=\varphi(\tau)$ and $\theta_B=\varphi(\theta)$
entrywise.  The coefficientwise map
\[
 C_\varphi:\mathcal F_A\longrightarrow\mathcal F_B,
                 \qquad C_\varphi(f)(g)_i=\varphi(f(g)_i)
\]
restricts to $M_A\to M_B$ and $j:O_A\to O_B$.
The same finite-level and stabilizer hypotheses hold for $B$,
and $\theta_B$ is continuous in the required sense.
\begin{proof}
Locality gives $\varphi(\mathfrak m_A^k)\subset\mathfrak m_B^k$
for every $k$.  The exact transformation laws pass through
the ring map, and every value congruence maps to its counterpart
over $B$.  Finite sums in OST5 give
$C_\varphi U_a=U_a C_\varphi$.  Thus the periodic description
OST10 maps into the corresponding description over $B$.
The finite-level hypothesis contains no coefficient ring.
In the stabilizer predicate, the geometric assertions
$b\in Z(G_p)\cap T$ are unchanged, and applying $\varphi$
entrywise to $\theta(b)\tau(u)=1$ gives its exact mapped
version.  Applying $\varphi$ to
\eqref{eq:ost-contchar} proves continuity of $\theta_B$.
This proof does not assume that $\varphi$ is injective,
surjective, flat, or finite.
\end{proof}
% END OST21

% BEGIN OST22
\paragraph{OST22. One finite family gives simultaneous residual reconstruction.}
\label{claim:OST22}
Choose a depth $c$ at which $\theta$ is trivial modulo
$\mathfrak m_A$, and use the representatives $x_i$ and
the section $s_c$ of OST1 and OST3.  For
$p=(i,\allowbreak{} j)$ let $s_p=s_c\delta_p$ and $w_p=e_A s_p\in O_A$.
For $f\in O_A$ and $g\in O_B$ one has respectively
\begin{align}
 f-\sum_p f(x_i)_j w_p&\in\mathfrak m_AO_A,\label{eq:ost-bc-A}\\
 g-\sum_p g(x_i)_j j(w_p)&\in\mathfrak m_BO_B.
 \label{eq:ost-bc-B}
\end{align}
The evaluations in the two rings satisfy
$\operatorname{ev}_{B,\allowbreak{} p}(j(f))=\varphi(\operatorname{ev}_{A,\allowbreak{} p}(f))$.
\begin{proof}
Equation \eqref{eq:ost-bc-A} is OST13 at precision one
followed by OST16.  The evaluation compatibility is the
definition of $j$.  To prove \eqref{eq:ost-bc-B}, set
\[
 h=\sum_p g(x_i)_j C_\varphi(s_p),\qquad
 g'=\sum_p g(x_i)_j j(w_p)\in O_B.
\]
All $C_\varphi(s_p)$ lie in $M_B$ and are
$N^-(c)$-invariant modulo $\mathfrak m_B$ by OST3 and
OST21.  They have the same delta-function evaluations.
By OST12 over $B$, $U_z^{n!}g\equiv g$ modulo the
value filtration for sufficiently large $n$.  By OST7
the same late iterate is $N^-(c)$-invariant modulo
$\mathfrak m_B$.  Hence $g$ itself has that invariance.
Since $h$ agrees with $g$ at each $x_i$, OST4 gives
$g-h\in\mathrm{valIn}_B(1)$.

For each of the finitely many $p$, OST12 over $A$ gives
$U_z^{n!}s_p-w_p\in\mathrm{valIn}_A(1)$ for all
sufficiently large $n$.  Choose one $n$ beyond those
finitely many thresholds and the threshold for $g$.
Naturality of the finite sums and OST21 imply
$U_z^{n!}h-g'\in\mathrm{valIn}_B(1)$.
Applying $U_z^{n!}$ to $g-h$, and then subtracting the
two congruences, proves $g-g'\in\mathrm{valIn}_B(1)$.
Both terms are ordinary, so OST16 over $B$ converts this
into \eqref{eq:ost-bc-B}.  This argument uses the same
$A$-valued sections and ordinary vectors in both rings;
an unrelated basis of $O_B$ would not prove the claim.
\end{proof}
% END OST22

% BEGIN OST23
\paragraph{OST23. A finite-module base-change criterion with an explicit inverse.}
\label{claim:OST23}
Suppose $O$ and $O'$ are finite modules over local Noetherian
rings $A$ and $B$, $j:O\to O'$ is $A$-linear, and finite
families $w_p\in O$, $A$-linear maps $v_p:O\to A$, and
$B$-linear maps $v'_p:O'\to B$ satisfy the evaluation
compatibility and the two congruences
of OST22.  Then
\[
 L:B\otimes_A O\longrightarrow O',\qquad b\otimes f\longmapsto b\,\allowbreak{} j(f)
\]
is an isomorphism.
\begin{proof}
Let $\Psi(f)=\sum_p v_p(f)w_p$.  Its reduction modulo
the maximal ideal is the identity.  The finite Nakayama and
Noetherian kernel arguments in OST17 make $\Psi$ an
automorphism.  Define
\[
 R:O'\longrightarrow B\otimes_AO,\qquad
         R(g)=\sum_p v'_p(g)(1\otimes w_p).
\]
On a pure tensor the compatibility of evaluations gives
$RL=1\otimes\Psi$, an automorphism, and thus $L$ is
injective.  The second congruence says
$g-LR(g)\in\mathfrak m_BO'$.  The finite module
$\operatorname{coker}L$ is therefore equal to its maximal-ideal
multiple and vanishes by Nakayama.  Hence $L$ is surjective.
Its inverse is $(1\otimes\Psi)^{-1}R$.
No flatness of $B$ over $A$ was used; injectivity comes from
this actual left inverse.
\end{proof}
% END OST23

% BEGIN OST24
\paragraph{OST24. Full original coefficientwise ordinary base change.}
\label{claim:OST24}
Under the original hypotheses of OST19 and every local map
$A\to B$ between rings of class \texttt{Is\allowbreak CNL}, the
coefficientwise map of OST21 induces an isomorphism
\begin{equation}\label{eq:ost-full-base-change}
 B\otimes_A X.\mathrm{Ord}_A(\tau,\theta,\allowbreak{} U)
       \ \xrightarrow{\ \sim\ }
 X.\mathrm{Ord}_B(\varphi\tau,\varphi\theta,\allowbreak{} U),
 \qquad b\otimes f\longmapsto\bigl(g\mapsto b\varphi(f(g))\bigr).
\end{equation}
This is the exact \texttt{Is\allowbreak Base\allowbreak Change B j} assertion of
\texttt{Ord\allowbreak BCMap\allowbreak Ord.\allowbreak ord\allowbreak BC}, with its stipulated formula for $j$.
\begin{proof}
OST21 transports every input hypothesis to the same Frame
and level over $B$.  OST19 gives finite modules on both sides.
OST22 supplies the evaluation-compatible data of OST23, which
proves \eqref{eq:ost-full-base-change}.  Its map is the
tensor extension of the original coefficientwise map, by the
displayed pure-tensor formula.  This identifies the map, not
merely its source and target ranks.  The proof includes
nonreduced rings and maps with kernels.
\end{proof}
% END OST24

% BEGIN OST25
\paragraph{OST25. The projector commutes with coefficient change and positive operators.}
\label{claim:OST25}
For every coefficient map in OST21 and every $a\in P$,
\[
 C_\varphi e_A=e_B C_\varphi\quad\hbox{on }M_A,
                  \qquad U_ae_A=e_AU_a.
\]
\begin{proof}
The first equality follows by applying $C_\varphi$ to
\eqref{eq:ost-factorial-limit}.  It preserves precision $k$
by locality and commutes with $U_z^{j!}$ by OST5, so both
sides are the unique factorial limit of the same function
over $B$.  For the second equality, $U_a$ preserves the
value filtration and commutes with $U_z$ on $M_A$ by OST5
and OST7.  Apply it to the same limit and use uniqueness.
These equalities also show that the base-change isomorphism
of OST24 intertwines the original positive-monoid operators.
\end{proof}
% END OST25

\subsection{Installing the actual local matrix datum}

% BEGIN OST26
\paragraph{OST26. Every axiom of the original local general-linear Iwahori datum.}
\label{claim:OST26}
Let $K$ be any of the original nonarchimedean local fields, with
valuation ring $\mathcal O$, maximal ideal $\mathfrak p$, and the
source's chosen uniformizer $\pi=\texttt{unif}(K)$.  Normalize the
additive valuation by $v(\pi)=1$.  For rank $n\geq0$, the literal
\texttt{gl\allowbreak Data} is
\begin{align*}
 N_f&=\{\text{upper unipotent matrices over }K\},\qquad
 T_f=\{\text{invertible diagonal matrices over }K\},\\
 N&=N_f\cap\operatorname{GL}_n(\mathcal O),\qquad
 T=\{\operatorname{diag}(u_i):u_i\in\mathcal O^\times\},\\
 T_c(c)&=\{\operatorname{diag}(u_i):u_i\equiv1
                                  \pmod{\mathfrak p^{c+1}}\},\\
 N^-(c)&=\{\text{lower unipotent matrices congruent to }1
                                  \pmod{\mathfrak p^{c+1}}\},\\
 P&=\{\operatorname{diag}(\pi^{a_i}):a_0\geq\cdots\geq a_{n-1},
                                      \ a_i\in\mathbb Z\},\\
 z&=\operatorname{diag}(\pi^{n-1-i})_{0\leq i<n}.
\end{align*}
These actual groups satisfy all the original \texttt{Iwa\allowbreak Data.\allowbreak Ax}
fields, so they supply the exact \texttt{gl\allowbreak Iwa} used above.
\begin{proof}
The source's $\mathrm{Kc}(r)$ is the kernel of reduction from
$\operatorname{GL}_n(\mathcal O)$ modulo $\mathfrak p^r$, mapped
into $\operatorname{GL}_n(K)$.  The integral map is injective, so
this is precisely the displayed congruence condition.  Inverse
transpose carries upper unipotents to lower unipotents, proving
the stated interpretation of $N^-(c)$.  Diagonal units identify
the source's valuation-one torus with $T$.  The selected
$\texttt{unif}(K)$ has uniformizer valuation by its defining
choice; no different uniformizer or monoid is substituted.

The inclusions, decreasing congruence subgroups, commutativity
of $T_f$, and membership of $z$ in $P$ follow directly.
Conjugation by $\operatorname{diag}(t_i)$ multiplies the
$(i,\allowbreak{} j)$ entry by $t_i/t_j$.  It therefore preserves upper
unipotents, and when the $t_i$ are units it preserves the
integral upper unipotents and every lower congruence group.
For $a=\operatorname{diag}(\pi^{a_i})\in P$, the difference
$a_i-a_j$ is nonnegative when $i<j$, proving
$aNa^{-1}\subset N$.  For $i>j$, conjugation by $a^{-1}$
multiplies the lower entry by $\pi^{a_j-a_i}$ with
nonnegative exponent, proving
$a^{-1}N^-(c)a\subset N^-(c)$.  For $a=z$ this exponent
equals $i-j\geq1$, which gives the required strict inclusion
into $N^-(c+1)$.

For the swap axiom, $y\in N^-(c)$ and $x\in N$ imply
$yx\equiv x\pmod{\mathfrak p^{c+1}}$.  We give the integral
upper-diagonal-lower elimination that applies to any matrix
$g$ with upper-unipotent reduction modulo $\mathfrak p^r$,
$r\geq1$.  Write
\[
 g=\begin{pmatrix}A&b\\c_0&d_0\end{pmatrix}
   =\begin{pmatrix}1&bd_0^{-1}\\0&1\end{pmatrix}
    \begin{pmatrix}A-bd_0^{-1}c_0&0\\0&d_0\end{pmatrix}
    \begin{pmatrix}1&0\\d_0^{-1}c_0&1\end{pmatrix}.
\]
Here $d_0\equiv1$ is a unit, $c_0\equiv0$, and the Schur
complement again has upper-unipotent reduction.  Induction on
$n$ yields $g=x'ty'$ with $x'\in N$,
$t\in T_c(r-1)$, and $y'\in N^-(r-1)$.  Apply this with
$r=c+1$.  For uniqueness, if an upper-unipotent matrix times
a diagonal matrix times a lower-unipotent matrix is $1$,
then the first two factors give a matrix that is simultaneously
upper triangular and lower unipotent.  It must be diagonal
with diagonal $1$, so all three factors are $1$.  This proves
the exact stronger $N_f,\allowbreak{} T_f,\allowbreak{} N^-(0)$ uniqueness field.

All required indices are finite.  For $a\in P$, choose
$D\geq0$ at least every $a_i-a_j$ with $i<j$.
The kernel of reduction of $N$ modulo $\mathfrak p^D$
is contained in $aNa^{-1}$, by the entrywise conjugation
formula.  Since $\mathcal O/\mathfrak p^D$ is finite, this
gives $[N:aNa^{-1}]<\infty$.  The same reduction argument
gives $[N^-(0):N^-(c)]<\infty$, while reduction of diagonal
units gives $[T:T_c(c)]<\infty$.  When $D=0$ the quotient
has one element; no positive lower bound for $D$ is needed.

For the exhaustion field, conjugating an upper-unipotent
matrix $x$ by $z^h$ multiplies its $(i,\allowbreak{} j)$ entry, $i<j$,
by $\pi^{h(j-i)}$.  A large enough $h$ makes the finitely
many entries integral, and an integral unipotent matrix has
integral inverse by the finite nilpotent geometric series.
Thus $z^h xz^{-h}\in N$.  For the positive-monoid domination
field, choose
\[
          h\geq\max_{0\leq i<n-1}(a_i-a_{i+1}),\qquad h\geq0.
\]
Then $h(n-1-i)-a_i$ is decreasing, so $z^ha^{-1}\in P$.
For the final torus-generation field, write each diagonal
entry $t_i=u_i\pi^{d_i}$ with $u_i\in\mathcal O^\times$.
Choose $h\geq0$ so that $d_i+h(n-1-i)$ is decreasing.
Then
\[
 \operatorname{diag}(t_i)
 =\operatorname{diag}(u_i)\,
   \operatorname{diag}(\pi^{d_i+h(n-1-i)})\,(z^h)^{-1},
\]
which is exactly $s a b^{-1}$ with $s\in T$, $a,\allowbreak{} b\in P$.
For ranks zero and one all missing adjacent-index conditions
are empty; choose $h=0$.  These calculations account for every
field of the original record.
\end{proof}
% END OST26

% BEGIN OST27
\paragraph{OST27. The product datum and the original Frame are the same objects.}
\label{claim:OST27}
The finite product of original Iwahori data, with all subgroups
and $P$ taken coordinatewise and $z=(z_v)_v$, satisfies all
the original axioms.  In particular OST26 constructs the
literal $\mathrm{iwS}(n,\allowbreak{} S)$, hence the Iwahori field of the
original $H.X$ for $H=\mathrm{hset}(Q,\tau_0)$.
\begin{proof}
Inclusions, conjugation laws, decrease with $c$, uniqueness,
and the swap identity are coordinatewise.  Choose the factors
in each coordinate for the swap and assemble their product
vectors.  Conjugation of the product $N$ is the product of
the conjugated coordinate subgroups; a finite product of their
finite coset sets proves each finite-index condition.

The two fields with a natural exponent need a common exponent.
For exhaustion choose $h_v$ in every coordinate and take their
maximum $h$.  If an exponent works, every larger one works:
conjugate further by a power of $z_v\in P_v$, which preserves
$N_v$.  For domination use the same maximum, since
$z_v^h a_v^{-1}=z_v^{h-h_v}(z_v^{h_v}a_v^{-1})\in P_v$.
The torus-generation factors are chosen coordinatewise and
then assembled.  The empty product also satisfies these
statements, with exponent zero.

The definitions in \texttt{Tc\allowbreak Anti}, \texttt{Distinct}, and
\texttt{Coefficient\allowbreak Rep\allowbreak Trivial\allowbreak On\allowbreak Unramified} take exactly these
coordinatewise groups to define $\mathrm{piIwa}$ and
$\mathrm{iwS}$.  In \texttt{Is\allowbreak Theta.\allowbreak hset} the index set is
the original finite set $V.\mathrm{frame}.V_\ell$ and the
local field in coordinate $v$ is the actual completion
$L_{s.w(v)}$.  Thus every group, congruence depth, operator
representative quotient, and selected uniformizer used in
OST1--OST25 is that of the original Frame, with no replacement
of the group by an abstractly isomorphic one.
\end{proof}
% END OST27

% BEGIN OST28
\paragraph{OST28. Full ordinary structure at every original arithmetic level.}
\label{claim:OST28}
Fix any original $\ell,\allowbreak{} L,\Omega,\allowbreak{} s,\allowbreak{} n\geq1,\allowbreak{} V,\allowbreak{} Q,\allowbreak{} c,\allowbreak{} d,\tau_0,\allowbreak{} r,\Lambda$
in the scope of OTW12, set $H=\mathrm{hset}(Q,\tau_0)$,
$A=\Lambda.\Lambda$, and fix the single character $\theta$
provided there before choosing a level.  For every original
compact-open auxiliary level $K$, the actual module
\[
 H.X.\mathrm{Ord}\bigl(H.\tau_A(A),\theta,\allowbreak{} H.\mathrm{lev}(K)\bigr)
\]
is finite free; its original positive-monoid operators are
bijective; it has the exact value filtration and one
factorial-limit projector of OST19.  For every local map
$A\to B$ with $B$ of class \texttt{Is\allowbreak CNL}, its ordinary
base change is exactly the original coefficientwise map.
\begin{proof}
OTW1 verifies \texttt{Is\allowbreak CNL} on this exact complete ring.
OST26--OST27 supply the actual Iwahori datum, including all
local matrix and finite-product axioms.  FNL10 supplies
\texttt{Fin\allowbreak Level} at each $H.\mathrm{lev}(K)$.  OTW12 gives
\texttt{Cont\allowbreak Char} for its already chosen $\theta$, and its
full \texttt{Is\allowbreak Wt} field gives \texttt{Free} for
$H.\tau_A(A)$ at every such level, using NTL17.
Apply OST19, OST20, and OST24 on these very objects.
The projector may depend on $K$, as the source allows, but
within each level it is chosen before all functions and
precisions.  The character is still fixed before every $K$.
Neither \texttt{Type\allowbreak For\allowbreak Alt} nor any of the other
$H.\mathrm{Ax}$ fields is needed for these applications.
The proof does not introduce any coefficient threshold or
restriction on $\ell$, and retains the original rank-four
case $\ell=3$.
\end{proof}
% END OST28

\subsection{Source correspondence and the remaining arithmetic entrance}

The definitions above are those of
\texttt{Continuous\allowbreak Torus\allowbreak Character.\allowbreak lean}: the Iwahori record at
lines 29--67, $\mathrm{ContChar}$ at 136--137, the Frame and
its finite-level and stabilizer predicates at 151--175, and
$\mathrm{Sp}$, $\mathrm{valIn}$, $\mathrm{Ord}$, and the
coefficient map at 199--237.  The precise line intervals are
recorded with file hashes in this chapter's source manifest.
The source endpoints are
\texttt{Ord\allowbreak Free.\allowbreak lean}, lines 295--437 and 441--481, and
\texttt{Ord\allowbreak BCMap\allowbreak Ord.\allowbreak lean}, lines 198--232.  The auxiliary
source statements are expanded into OST1--OST18 rather than
invoked as unproved mathematical inputs.

This completes the ordinary-structure and ordinary-base-change
part of roadmap M1 on the original arithmetic Frames in OST28.
The finite-level theorem supplies \texttt{Fin\allowbreak Level}; NTL17
and OTW12 supply the stabilizer and complete-character inputs.
The chapter introduces no new coefficient threshold, no choice
of a replacement ideal, and no restriction excluding residue
characteristic three.  Its construction is independent of the
missing integral-type theorem.  The ordinary maps are the
Frame's actual double-coset maps, distinct from the separate
local theta-Hecke construction.

The following remain downstream arithmetic tasks: the
integral local type, classical weight/control and automorphic
fibers, arithmetic-point separation, the four Hida endpoints,
and all-level automorphy on the same ideal.  None follows merely
from finite freeness or base change.  Thus no whole
\texttt{COUNT\_\allowbreak II\_\allowbreak SUPPORT} obligation is retired by this
chapter alone.


\section{The actual algebraic weight representation with all integral fields}
\label{sec:weight-representation-v42}

\status{This chapter constructs the existential weight representation
used by the original \texttt{Prod\allowbreak Zpow\allowbreak Const.\allowbreak wt\allowbreak Rep}, on the specified
algebraic representation and with the specified field isomorphism.
It proves all ten assertion fields of \texttt{Wt\allowbreak Rep}.  The proof uses
the already reconstructed characteristic-zero polynomial results
HP1--HP2, and a saturated lattice over the actual coefficient DVR.
It therefore does not need a new general standard-monomial basis
theorem over arbitrary domains.  The source endpoint permits the
basis chosen here: its required identification is an explicit
intertwining isomorphism with the unchanged original
\texttt{alg\allowbreak Rep}.}

Fix a number field $L$, a prime $\ell$, an integer $n\geq0$, a
finite index set $V$, places $S(v)$ of $L$, a field isomorphism
\[
             \iota:\overline{\mathbb Q}_\ell\xrightarrow{\sim}
                                                              \mathbb C,
\]
and the original weight function $\lambda$ assigning a dominant
integer weight of rank $n$ to each embedding
$L\hookrightarrow\overline{\mathbb Q}_\ell$.  The isomorphism
$\iota$ is algebraic and is never assumed continuous.  Let $c$
be the original coefficient datum: its ring $c.O$ is a DVR
subalgebra of $\overline{\mathbb Q}_\ell$, finite over
$\mathbb Z_\ell$.

Put $F_v=L_{S(v)}$ and let
\[
 \mathcal E=\{(v,\allowbreak{} j):v\in V,\ j:F_v\hookrightarrow
           \overline{\mathbb Q}_\ell\text{ a continuous field map}\}.
\]
This is the literal source \texttt{Emb}.  For $e=(v,\allowbreak{} j)$ write
$\lambda_e=\lambda(j|_L)$ and $\psi_e=\iota\circ j$.
The original group and representation are
\begin{equation}\label{eq:wrp-algrep}
 G=\prod_{v\in V}\operatorname{GL}_n(F_v),\qquad
 \rho(g)=\bigotimes_{e=(v,\allowbreak{} j)\in\mathcal E}
                 \operatorname{sqRep}(\psi_e(g_v),\lambda_e)
\end{equation}
on the actual tensor product
\[
 \mathcal V_{\mathbb C}=
       \bigotimes_{e\in\mathcal E}
                  \operatorname{nablaSq}_{\mathbb C}(n,\lambda_e).
\]
All indices $i$ below run from $0$ to $n-1$, as in the source.
For a dominant weight $\lambda_e$, set $s_e=(\lambda_e)_{n-1}$
and $\nu_{e,\allowbreak{} i}=(\lambda_e)_i-s_e$ if $n>0$; for $n=0$ use
$s_e=0$ and the constant polynomial.  By HP1, its literal
polynomial factor is
\begin{equation}\label{eq:wrp-polynomial}
 \mathcal N_{\nu_e}(K)=\{f:f(Xt)=t^{\nu_e}f(X),\
                                  f(Xu)=f(X)\ (u\in N_+)\},
 \qquad \rho_e(g)f=\det(g)^{s_e}f(g^{\mathsf T}X).
\end{equation}
The covariance conditions are polynomial identities, or the
finite integer-coefficient kernel equations of HP1.  The
determinant power may be negative.

The local Iwahori datum is the original $\mathrm{iwS}(n,\allowbreak{} S)$.
We use its actual uniformizers $\varpi_v$, not new uniformizers
chosen after extending coefficients.  OST26--OST27 identify its
torus with diagonal matrices, its compact groups with integral
matrices, and its positive monoid with tuples
\[
      a_v=\operatorname{diag}(\varpi_v^{a_{v,0}},\ldots,
                             \varpi_v^{a_{v,\allowbreak{} n-1}}),
                    \quad a_{v,0}\geq\cdots\geq a_{v,\allowbreak{} n-1}.
\]
Its distinguished $z$ uses $a_{v,\allowbreak{} i}=n-1-i$ at every $v$.

\subsection{One coefficient extension for every local embedding}

% BEGIN WRP1
\paragraph{WRP1. The original embedding index is finite and its images are finite local fields.}
\label{claim:WRP1}
The set $\mathcal E$ is finite.  Every $j$ occurring in it has
$S(v)$ above $\ell$, fixes the embedded $\mathbb Q_\ell$, and
has image a finite extension of $\mathbb Q_\ell$ inside the
specified $\overline{\mathbb Q}_\ell$.  Moreover
\begin{equation}\label{eq:wrp-local-integrality}
 |j(x)|_\ell\leq1\quad(x\in\mathcal O_{F_v}),\qquad
 |j(x)|_\ell<1\quad(x\in\mathfrak m_{F_v}).
\end{equation}
\begin{proof}
Continuous maps agreeing on the dense subfield $L$ agree on its
completion.  There are only finitely many field maps
$L\to\overline{\mathbb Q}_\ell$: for a finite set of algebraic
generators each image is a root of a fixed nonzero polynomial,
so there are finitely many possible lists of images.  This and
finiteness of $V$ prove the first assertion.

If $S(v)$ has residue characteristic $p\ne\ell$, then
$p^r\to0$ in $F_v$, whereas its image under $j$ has
$\ell$-adic norm $1$ for every $r$, contradicting continuity.
Thus $S(v)$ lies above $\ell$.  The inclusion of $\mathbb Q$
extends to the usual inclusion $\mathbb Q_\ell\subset F_v$;
its topology is obtained by restricting the local valuation,
and completion supplies this inclusion.  By continuity and
density of $\mathbb Q$ in $\mathbb Q_\ell$, $j$ is the identity
on that subfield.  The completion $F_v$ is finite over
$\mathbb Q_\ell$: the $\mathbb Q_\ell$-span of a rational
basis of $L$ contains the dense subfield $L$ and is closed.
To check closedness, use the coordinate compactness argument
of GA2: a linear injection from a finite dimensional
$\mathbb Q_\ell$-space into a normed space has equivalent
coordinate and induced norms, since the coordinate unit
sphere is compact and its image avoids zero.  Hence its
image is complete and closed.  Consequently this span is
the whole $F_v$.  The $\mathbb Q_\ell$-linear injection $j$
therefore has finite dimensional image, a field.

The local integer ring $\mathcal O_{F_v}$ is compact.
Continuity bounds $|j(x)|_\ell$ on it by some constant.
For any fixed $x$ in that ring all positive powers remain
in the ring, so $|j(x)|_\ell^r$ stays bounded.  This forces
$|j(x)|_\ell\leq1$.  If $x\in\mathfrak m_{F_v}$ then
$x^r\to0$, and continuity forces $j(x)^r\to0$, hence
$|j(x)|_\ell<1$.  These arguments use no continuity of
$\iota$.
\end{proof}
% END WRP1

% BEGIN WRP2
\paragraph{WRP2. One actual coefficient DVR contains all local integer images.}
\label{claim:WRP2}
There is an original coefficient datum $c'$ with
$c.O\subset c'.O\subset\overline{\mathbb Q}_\ell$ such that
every $j$ in WRP1 restricts to a local homomorphism
\[
               \phi_e:\mathcal O_{F_v}\longrightarrow c'.O.
\]
In particular $\phi_e(\mathfrak m_{F_v}^k)
\subset\mathfrak m_{c'.O}^k$ for every $k$.
\begin{proof}
Let $E$ be the compositum inside the specified algebraic closure
of $\operatorname{Frac}(c.O)$ and the finitely many fields
$j(F_v)$.  It is finite over $\mathbb Q_\ell$.  For the first
field this follows from finite $\mathbb Z_\ell$-generation of
$c.O$: as proved in GA2, every denominator in that DVR divides
a power of $\ell$, so its fraction field is spanned by the
same generators over $\mathbb Q_\ell$.  The other fields are
finite by WRP1, and a finite compositum of finite extensions
is finite, by adjoining their finite bases successively.

Define the literal subring
\[
                       O'=\{x\in E:|x|_\ell\leq1\}.
\]
Here is a direct verification that it is an allowed coefficient
DVR.  Choose a $\mathbb Q_\ell$-basis of $E$.  The coordinate
norm and the induced norm are equivalent by the compact unit
sphere argument in WRP1.  Consequently $O'$ is contained
in $\ell^{-M}\mathbb Z_\ell^r$ and contains
$\ell^M\mathbb Z_\ell^r$ for some $M$.  It is a
$\mathbb Z_\ell$-submodule, hence finite over the Noetherian
PID $\mathbb Z_\ell$, and it spans $E$.  The finite algebra
$O'$ is Noetherian: every ideal is a submodule of its finite
$\mathbb Z_\ell$-module, and a finite list of generators
over $\mathbb Z_\ell$ also generates it as an ideal.

Its unique maximal ideal is
$\mathfrak m' =\{x:|x|_\ell<1\}$: an element of norm one
has its inverse in $O'$, while products involving an element
of norm less than one cannot equal one.  This ideal is
nonzero, since it contains $\ell$.  Choose a finite set of
its generators and let $\pi'$ be a nonzero generator of
largest norm.  Every $x\in\mathfrak m'$ is a sum of these
generators with coefficients of norm at most one, so
$|x|_\ell\leq|\pi'|_\ell<1$.  Thus $x/\pi'\in O'$,
and $\mathfrak m'=(\pi')$.  For any nonzero $x\in O'$,
choose the largest $r\geq0$ with
$|x|_\ell\leq|\pi'|_\ell^r$.  Such an $r$ exists and is
finite.  The element $x/(\pi')^r$ is in $O'$ and is a unit;
otherwise its norm would be at most $|\pi'|_\ell$, contrary
to maximality of $r$.  Hence every nonzero element is a unit
times a power of $\pi'$, and taking the smallest exponent
in any nonzero ideal proves the DVR property.  The ring is
complete: $E$ is complete by norm equivalence, $O'$ is
closed, and the powers of $\pi'$ give its norm neighborhoods.
Its residue field is finite because it is a quotient of the
finite ring $O'/\ell O'$.

Every element of $c.O$ has norm at most one.  Indeed finite
generation over $\mathbb Z_\ell$ gives a monic equation by
the determinant trick, and a root with norm greater than one
would give a uniquely largest term in that equation.  Thus
$c.O\subset O'$.  WRP1 gives $j(\mathcal O_{F_v})\subset O'$
and $j(\mathfrak m_{F_v})\subset\mathfrak m'$.  Set
$c'.O=O'$ as a subalgebra of the same algebraic closure.
The asserted maps are the literal restrictions of $j$;
they are local, and the assertion for powers follows by
expanding finite sums of products.  This constructs the
whole coefficient datum, not just an abstract isomorphic
DVR.  Its complex algebra map is always $\iota|_{O'}$.
\end{proof}
% END WRP2

\subsection{Integral weight coordinates on the unchanged polynomial factors}

% BEGIN WRP3
\paragraph{WRP3. A saturated integral lattice stable under all integral matrices.}
\label{claim:WRP3}
Write $O=c'.O$, $K=\operatorname{Frac}(O)$ and fix one dominant
weight $\lambda_e$.  Let $P_O$ be the finite free module of
polynomials of its prescribed column multidegree, and put
\[
              L_e=\mathcal N_{\nu_e}(K)\cap P_O.
\]
Then $L_e$ is a finite free $O$-lattice spanning
$\mathcal N_{\nu_e}(K)$.  It is stable under the original
$\rho_e(g)$ for $g\in\operatorname{GL}_n(O)$.  If
$g\equiv1\pmod{\mathfrak m_O^k}$, then
\begin{equation}\label{eq:wrp-congruence}
             (\rho_e(g)-1)L_e\subset\mathfrak m_O^k L_e.
\end{equation}
\begin{proof}
The submodule $L_e\subset P_O$ is finite and torsion free
over a DVR, hence free.  For clarity this module fact follows
by induction on the rank of $P_O$: project to its first
coordinate, whose nonzero image is a principal ideal,
choose a lift of its generator, and split off that free
summand; apply induction to the kernel in the remaining
coordinates.  A basis of the $K$-space has finitely many
polynomial coefficients; clearing their denominators places
scalar multiples in $L_e$, so it spans the whole space.
By the finite integer-coefficient equations in HP1, this
intersection is also the literal integral polynomial kernel:
an equation over $O$ holds if and only if it holds after
the injective map $O\to K$.

For $g\in\operatorname{GL}_n(O)$, matrix substitution
preserves integral polynomial coefficients and each column
degree.  The multiplier $\det(g)^{s_e}$ is an element of
$O^\times$, including for negative $s_e$.  Left substitution
commutes with the right covariance conditions defining
$\mathcal N_{\nu_e}$, so it preserves that subspace as well.
The same argument for $g^{-1}$ gives stability of the lattice
in both directions.

If $g\equiv1$ modulo $\mathfrak m_O^k$, reduction of matrix
substitution is the identity.  Also
$\det(g)^{s_e}\equiv1$ modulo that ideal; negative powers
cause no difficulty since the determinant is a unit and
$u^{-1}-1=-u^{-1}(u-1)$.  Thus the difference in
\eqref{eq:wrp-congruence} lies in
$\mathfrak m_O^kP_O\cap\mathcal N_{\nu_e}(K)$.
Writing $\mathfrak m_O^k=(\pi'^k)$ and dividing the
polynomial by $\pi'^k$ identifies this intersection with
$\pi'^kL_e$.  This proves the exact congruence, without
assuming an integral standard-monomial theorem.
\end{proof}
% END WRP3

% BEGIN WRP4
\paragraph{WRP4. A basis consisting of integral row-weight vectors.}
\label{claim:WRP4}
There is a finite $O$-basis $(b_{e,\allowbreak{} j})_{j\in J_e}$ of $L_e$
whose elements have diagonal weights
$\mu_{e,\allowbreak{} j}\in\mathbb Z^n$ for the action
$\rho_e$.  Under the specified embedding
$K\hookrightarrow\mathbb C$, it is a basis of the literal
$\operatorname{nablaSq}_{\mathbb C}(n,\lambda_e)$.
\begin{proof}
The ambient polynomial module decomposes as the direct sum
of its row-multidegree spaces.  Under the determinant twist,
their torus weights are the row multidegrees plus $s_e$ in
every coordinate.  Only finitely many occur.  Their
projections preserve $\mathcal N_{\nu_e}(K)$: this space
is torus-stable, and over the infinite field $K$ one may
choose a diagonal element on which the distinct occurring
Laurent characters have distinct values and interpolate
the weight projections by polynomials in its action.
Such a diagonal element exists because a nonzero finite
product of Laurent polynomials cannot vanish on the whole
torus of an infinite field.  On the ambient polynomial
module these projections just retain selected monomials,
so they preserve $P_O$ as well.  Consequently
\[
                 L_e=\bigoplus_\mu
                         (L_e\cap\mathcal N_{\nu_e}(K)_\mu).
\]
Each summand is finite free by the same DVR module argument
as WRP3.  Choose a basis in each and take their union.
HP1 identifies extension of the polynomial kernel from
$K$ to $\mathbb C$ coefficientwise for this actual embedding,
so the resulting vectors are a complex basis of the exact
source factor.  No new complex representation is introduced.
\end{proof}
% END WRP4

% BEGIN WRP5
\paragraph{WRP5. The unique lowest coordinate and the exact weight inequalities.}
\label{claim:WRP5}
The basis in WRP4 may be indexed with a distinguished $0_e$
such that
\begin{align}
 \mu_{e,0_e,\allowbreak{} i}&=(\lambda_e)_{n-1-i},\\
 \sum_i a_i\mu_{e,0_e,\allowbreak{} i}&\leq\sum_i a_i\mu_{e,\allowbreak{} j,\allowbreak{} i}
                         &&(a_0\geq\cdots\geq a_{n-1}),\\
 \sum_i(n-1-i)\mu_{e,0_e,\allowbreak{} i}
                 &<\sum_i(n-1-i)\mu_{e,\allowbreak{} j,\allowbreak{} i} &&(j\ne0_e).
 \label{eq:wrp-weights}
\end{align}
For every upper unitriangular $u$, the coefficient of
$b_{e,0_e}$ in $\rho_e(u)b_{e,\allowbreak{} j}$ is $\delta_{j,0_e}$.
Every weight has total sum $\sum_i(\lambda_e)_i$.
\begin{proof}
HP1 gives the exact highest polynomial $h_e$, fixed by the
upper unitriangular group and of weight $\lambda_e$.
Let $w_0$ be the reversing permutation matrix and put
$v_e=\rho_e(w_0)h_e$.  It is nonzero, fixed by the lower
unitriangular group, and has weight
$\lambda_e^{\mathrm{rev}}$.

The span of the upper-unitriangular orbit of $v_e$ is the
whole space.  To see this without invoking a new
highest-weight theorem, use the dense open Gaussian cell
$N_+TN_-$ of $\operatorname{GL}_n$.  Lower unipotents fix
$v_e$ and torus elements scale it.  A linear functional
vanishing on its $N_+$-orbit therefore vanishes on its orbit
under this entire open cell.  The matrix coefficient
$g\mapsto\alpha(\rho_e(g)v_e)$ is regular on
$\operatorname{GL}_n$: it is a polynomial in matrix entries
and inverse determinant.  Vanishing on the dense open cell
makes it identically zero.  Thus the span of the $N_+$-orbit
equals the span of the full group orbit, which is the whole
space by the absolute irreducibility proved in HP2.

For an upper elementary matrix $u_{ij}(t)$, $i<j$, and a
vector of weight $\mu$, polynomial expansion in $t$ gives
coefficients of weights
$\mu+r(\epsilon_i-\epsilon_j)$, $r\geq0$.  This follows
by conjugating $u_{ij}(t)$ by a diagonal matrix, or directly
from the row substitution in \eqref{eq:wrp-polynomial}.
The constant coefficient is the original vector.  Every
upper unitriangular matrix is an ordered product of such
elementary matrices.  The preceding spanning assertion
therefore proves that every occurring weight has the form
\[
       \mu=\lambda_e^{\mathrm{rev}}
                 +\sum_{i<j}r_{ij}(\epsilon_i-\epsilon_j),
                         \qquad r_{ij}\geq0.
\]
For the height with coefficients $n-1-i$, every nonzero
positive root has strictly positive height $j-i$.  The
height-minimal part of every expanded orbit vector is
therefore a scalar multiple of $v_e$.  Hence the
$\lambda_e^{\mathrm{rev}}$ weight space is one dimensional.
Its intersection with $L_e$ is a rank-one free direct
summand by WRP4; choose its generator as $b_{e,0_e}$.
Every other chosen basis vector has strictly larger height.

For any decreasing $a_i$, each positive root has
$a_i-a_j\geq0$, giving the weak inequality in
\eqref{eq:wrp-weights}; the strictly decreasing height gives
the strict one.  The same elementary expansions show that
upper unipotents can only increase weights.  They therefore
cannot create a lowest coefficient from another basis
vector, and preserve the lowest coefficient of
$b_{e,0_e}$ itself.  This proves the zeroth-row assertion.
Finally every positive root has coordinate sum zero,
giving the common total weight.  Rank zero has just the
constant polynomial and the same assertions with empty
coordinate sums.
\end{proof}
% END WRP5

% BEGIN WRP6
\paragraph{WRP6. Tensor coordinates and the specified algebraic intertwiner.}
\label{claim:WRP6}
The tensor basis
\[
                 b_J=\bigotimes_{e\in\mathcal E}b_{e,\allowbreak{} J_e},
                   \qquad J\in\prod_e J_e,
\]
is a basis of the actual $\mathcal V_{\mathbb C}$.  Its
distinguished index is $J_0=(0_e)_e$.  There is an integer
$d\geq0$ and a reindexing by $\{0,\ldots,\allowbreak{} d\}$ sending
zero to $J_0$.  In these coordinates let $\xi(g)$ be the
matrix of the unchanged representation \eqref{eq:wrp-algrep}.
Then $\xi:G\to\operatorname{GL}_{d+1}(\mathbb C)$ is a
representation, and the coordinate isomorphism
\begin{equation}\label{eq:wrp-intertwiner}
 B:\mathbb C^{d+1}\xrightarrow{\sim}\mathcal V_{\mathbb C},
              \qquad B(\xi(g)x)=\rho(g)B(x),
\end{equation}
proves exactly $\mathrm{IsAlgWt}(\iota,\lambda,\xi)$.
\begin{proof}
Tensor products of finite bases give a basis by the
multilinear expansion of elementary tensors and the
coordinate multilinear functionals.  Each factor has the
nonzero highest polynomial, so each basis index set is
nonempty, and the product index contains $J_0$.  Its finite
cardinality is therefore $d+1$ for some $d$.  Reindexing by
a finite ordinal and then a permutation places $J_0$ at
zero.  The matrix of a group representation in a basis is
invertible, with inverse the matrix at $g^{-1}$; composition
of matrices is exactly composition of operators.
The map $B$ sends standard coordinates to the displayed
basis, so the intertwining identity is its definition.
It uses precisely the maps $\psi_e=\iota\circ j$ and the
source tensor factors.  For empty $\mathcal E$ the tensor
product is $\mathbb C$, the basis has one element, and the
representation is trivial.
\end{proof}
% END WRP6

\subsection{The ten original weight-representation assertions}

% BEGIN WRP7
\paragraph{WRP7. The original lowest-weight character and both zeroth rows.}
\label{claim:WRP7}
For $t\in\mathrm{iwS}.T_f$ write its components uniquely as
$t_v=\operatorname{diag}(y_{v,0},\ldots,\allowbreak{} y_{v,\allowbreak{} n-1})$ and define
\begin{equation}\label{eq:wrp-character}
 \chi(t)=\prod_{e=(v,\allowbreak{} j)\in\mathcal E}\prod_i
                  \psi_e(y_{v,\allowbreak{} i})^{(\lambda_e)_{n-1-i}}
                  \quad\in\mathbb C^\times.
\end{equation}
This is a group character.  With the coordinates of WRP6,
\begin{align}
 \xi(u)_{0J}&=\delta_{J,\allowbreak{} J_0} &&(u\in\mathrm{iwS}.N_f),\\
 \xi(t)_{0J}&=\chi(t)\delta_{J,\allowbreak{} J_0}
                                      &&(t\in\mathrm{iwS}.T_f).
 \label{eq:wrp-rows}
\end{align}
Moreover its value is exactly the character specified by
the source endpoint:
\begin{equation}\label{eq:wrp-lowW}
 \chi(t)=\iota\left(\prod_{v\in V}\prod_i
     \operatorname{algPow}(S(v),\operatorname{lowW}(\lambda,\allowbreak{} i),
                                                       y_{v,\allowbreak{} i})\right).
\end{equation}
\begin{proof}
Every diagonal entry $y_{v,\allowbreak{} i}$ is nonzero, so the products
and their integer powers are nonzero.  Componentwise
multiplication of diagonal matrices proves multiplicativity.
Each tensor factor is diagonal in the weight basis: its
diagonal entry at $j_e$ is
$\prod_i\psi_e(y_{v,\allowbreak{} i})^{\mu_{e,\allowbreak{} j_e,\allowbreak{} i}}$.  Consequently
the whole tensor matrix is diagonal, and its entry at
$J_0$ is \eqref{eq:wrp-character}, by WRP5.  This gives
the torus row.  For an upper unitriangular tuple, the entry
in the $J_0$ row is the product of the corresponding lowest
rows of the factors, hence
$\prod_e\delta_{J_e,0_e}=\delta_{J,\allowbreak{} J_0}$ by WRP5.

The literal source definition is
\[
 \operatorname{algPow}(w,\allowbreak{} h,\allowbreak{} y)=
       \prod_{j:F_w\hookrightarrow\overline{\mathbb Q}_\ell
                         \text{ continuous}} j(y)^{h(j|_L)},
 \qquad
 \operatorname{lowW}(\lambda,\allowbreak{} i)(\tau)=\lambda(\tau)_{n-1-i}.
\]
The source writes a finite-support product; WRP1 proves that
the entire indexing set is finite, so it is the displayed
finite product.  Expand these definitions, commute the finite
products over $v$, $i$, and $j$, and apply multiplicativity
of $\iota$ and its compatibility with integer powers of
nonzero elements.  The result is exactly
\eqref{eq:wrp-character}.  Thus the character is identified
on every original torus element, with the same local
embeddings and the same field isomorphism.
\end{proof}
% END WRP7

% BEGIN WRP8
\paragraph{WRP8. Integral matrices on the three compact subgroups.}
\label{claim:WRP8}
Every entry of $\xi(g)$ is in $\iota(O)$ if
$g\in\mathrm{iwS}.N$, $g\in\mathrm{iwS}.T$, or
$g\in\mathrm{iwS}.N_m(0)$.  We may therefore take the
original field $c_0$ of $\mathrm{WtRep}$ to be zero.
\begin{proof}
OST26--OST27 identify each component of these three groups
as an element of $\operatorname{GL}_n(\mathcal O_{F_v})$.
The map $\phi_e$ of WRP2 carries both this matrix and its
inverse to integral matrices over $O$, so it gives an
element of $\operatorname{GL}_n(O)$.  WRP3 says that this
element preserves $L_e$.  Its matrix in the $O$-basis
of WRP4 consequently has all entries in $O$.  The image
of this matrix under $\iota$ is the matrix of
$\rho_e(\psi_e(g_v))$ in the complex basis, because all
substitutions and determinant powers commute with the
given field maps.  Entries of the tensor matrix are
products of the entries of these factor matrices, so
they also lie in $\iota(O)$.  This proves, separately,
the fields \texttt{int\_\allowbreak N},
\texttt{integral\_\allowbreak torus\allowbreak Character}, and \texttt{int\_\allowbreak Nm}.
The middle field concerns every matrix entry on the compact
torus, as its original definition requires.
\end{proof}
% END WRP8

% BEGIN WRP9
\paragraph{WRP9. Positive normalization and strict contraction outside column zero.}
\label{claim:WRP9}
For every $a\in\mathrm{iwS}.P$, every entry of
$\chi(a)^{-1}\xi(a)$ is in $\iota(O)$.  For the literal
distinguished $z$, every entry of this normalized matrix
outside column zero is in $\iota(\mathfrak m_O)$.
\begin{proof}
Use the actual exponents
$a_v=\operatorname{diag}(\varpi_v^{a_{v,\allowbreak{} i}})_i$ from
the setup, with each $a_{v,\allowbreak{} i}$ decreasing in $i$.
The matrix is diagonal by WRP4.  For the diagonal position
$J=(J_e)_e$ put
\[
 \Delta_e(J)=\sum_i a_{v,\allowbreak{} i}
                \bigl(\mu_{e,\allowbreak{} J_e,\allowbreak{} i}-(\lambda_e)_{n-1-i}\bigr)
                \quad(e=(v,\allowbreak{} j)).
\]
This integer is nonnegative by WRP5.  Using nonzero
$\varpi_v$ and the ordinary rules for integer exponents,
the normalized diagonal entry is the image under $\iota$
of the actual element
\begin{equation}\label{eq:wrp-normalized}
          r_J=\prod_{e=(v,\allowbreak{} j)}\phi_e(\varpi_v)^{\Delta_e(J)}
                                                        \in O.
\end{equation}
Here $\varpi_v$ is regarded as an element of the local
integer ring, and WRP2 applies.  Nonnegative exponents
are crucial: the normalization has canceled any negative
powers in the original highest weight.  Off-diagonal
entries are zero, proving \texttt{int\_\allowbreak P}.

For $z$ the coefficients are $a_{v,\allowbreak{} i}=n-1-i$.
If $J\ne J_0$, some $J_e\ne0_e$, so the strict inequality
in WRP5 gives $\Delta_e(J)>0$.  The element
$\phi_e(\varpi_v)$ lies in $\mathfrak m_O$ by WRP2,
and all other factors of \eqref{eq:wrp-normalized} lie
in $O$.  Thus $r_J\in\mathfrak m_O$.  In that column
every off-diagonal entry is zero and also lies in the
maximal ideal.  These elements are the explicit witnesses
required in \texttt{nilpotent\_\allowbreak central\_\allowbreak action}; no
passage to eigenvalues or residue-field coordinates is
being substituted for the required entrywise assertion.
If there is no nonzero column index, the assertion is
vacuous, as in the original definition.
\end{proof}
% END WRP9

% BEGIN WRP10
\paragraph{WRP10. Both integral congruence assertions at every precision.}
\label{claim:WRP10}
For every $k\geq0$ one may choose depth $c=k$ so that
\begin{align}
 \xi(y)_{IJ}&=\delta_{I,\allowbreak{} J}+\iota(x_{IJ}),
       &x_{IJ}\in\mathfrak m_O^k
                  &&(y\in\mathrm{iwS}.N_m(c)),\\
 \chi(t)&=1+\iota(x),
       &x\in\mathfrak m_O^k
                  &&(t\in\mathrm{iwS}.T_f,\ t\in\mathrm{iwS}.T_c(c)).
 \label{eq:wrp-continuity}
\end{align}
These are exactly \texttt{cont\_\allowbreak Nm} and
\texttt{continuous\_\allowbreak torus\allowbreak Character}.
\begin{proof}
By the original definitions checked in OST26--OST27,
each component of either $N_m(c)$ or $T_c(c)$ is
congruent to the identity modulo
$\mathfrak m_{F_v}^{c+1}$, and hence modulo
$\mathfrak m_{F_v}^k$ for $c=k$.  Apply $\phi_e$ and
WRP2 to obtain an integral matrix congruent to the
identity modulo $\mathfrak m_O^k$.  WRP3 gives the same
congruence for its action on $L_e$.  In the factor basis
its entries therefore have the form
$\delta_{i_e,\allowbreak{} j_e}+u_{i_ej_e}$ with
$u_{i_ej_e}\in\mathfrak m_O^k$.

Entries in the tensor basis are products of these
expressions.  On reducing modulo $\mathfrak m_O^k$,
their product is $\prod_e\delta_{I_e,\allowbreak{} J_e}=\delta_{I,\allowbreak{} J}$.
Equivalently their difference from $\delta_{I,\allowbreak{} J}$ is
an element of that ideal: expand the finite product,
and each other term contains at least one of the
$u_{i_ej_e}$.  Mapping to $\mathbb C$ by $\iota$ gives
the first formula, with witnesses in the original ideal
power rather than in a field topology.  For the torus,
the $00$ entry equals $\chi(t)$ by WRP7; take that entry
of the same congruence.  This gives the second formula.
The argument also covers $k=0$, for which the zeroth
ideal power is $O$, and an empty tensor product has
identity matrix and zero error.  No continuity assertion
about $\iota$ was used.
\end{proof}
% END WRP10

% BEGIN WRP11
\paragraph{WRP11. The exact scalar action on the original center.}
\label{claim:WRP11}
If $t\in\mathrm{iwS}.T_f$ belongs to the center of $G$,
then
\[
                 \xi(t)_{IJ}=\chi(t)\delta_{I,\allowbreak{} J}.
\]
\begin{proof}
For $n\geq2$, a diagonal matrix commuting with every
elementary matrix has all its diagonal entries equal:
compare the $(i,\allowbreak{} j)$ entry of its product with
$1+E_{ij}$ in the two orders.  Therefore every component
$t_v$ is a scalar matrix, since it commutes with the
elementary matrices supported in that component of $G$.
For $n=1$ it is already scalar; for $n=0$ all groups
and representation factors are trivial.  Write the
scalar as $s_v$ when $n>0$.

By WRP5 every weight of factor $e=(v,\allowbreak{} j)$ has the same
sum $\sum_i(\lambda_e)_i$.  Its diagonal eigenvalue
at $s_v$ is consequently
$\psi_e(s_v)^{\sum_i(\lambda_e)_i}$, independent of
the basis index.  Each factor is thus scalar, and
their product equals its own lowest-coordinate value,
which is $\chi(t)$ by WRP7.  This proves the entrywise
\texttt{central} field, including all off-diagonal
zeros, for the actual full product center.
\end{proof}
% END WRP11

% BEGIN WRP12
\paragraph{WRP12. Installation of the exact original \texttt{wt\allowbreak Rep} endpoint.}
\label{claim:WRP12}
For every original $(L,\ell,\allowbreak{} n,\allowbreak{} V,\allowbreak{} S,\iota,\lambda,\allowbreak{} c)$,
WRP1--WRP11 construct $c'\geq c$ and a representation
\[
                  R:\mathrm{WtRep}(\mathrm{iwS}(n,\allowbreak{} S),\allowbreak{} c'.O,\mathbb C)
\]
such that $\mathrm{IsAlgWt}(\iota,\lambda,\allowbreak{} R.\xi)$ and
the equality \eqref{eq:wrp-lowW} hold for every original
torus tuple.  This is the exact statement of
\texttt{Upper\allowbreak Theory.\allowbreak V.\allowbreak Hi.\allowbreak wt\allowbreak Rep}, in
\texttt{Prod\allowbreak Zpow\allowbreak Const.\allowbreak lean}, at the pinned source revision.
\begin{proof}
Use the subalgebra $c'.O$ constructed in WRP2, with
complex algebra map equal to the restriction of the
specified $\iota$.  Take $d,\xi$ from WRP6,
$\chi$ from WRP7, and $c_0=0$ from WRP8.  The ten
original assertion fields are supplied individually
by the following proved statements:
\begin{center}
\begin{tabular}{ll}
Original field & Written proof \\
\hline
\texttt{int\_\allowbreak N} & WRP8 \\
\texttt{integral\_\allowbreak torus\allowbreak Character} & WRP8 \\
\texttt{int\_\allowbreak Nm} & WRP8, with $c_0=0$ \\
\texttt{int\_\allowbreak P} & WRP9 \\
\texttt{row\_\allowbreak N} & WRP7 \\
\texttt{torus\allowbreak Action} & WRP7 \\
\texttt{nilpotent\_\allowbreak central\_\allowbreak action} & WRP9 \\
\texttt{cont\_\allowbreak Nm} & WRP10, every $k$ \\
\texttt{continuous\_\allowbreak torus\allowbreak Character} & WRP10, every $k$ \\
\texttt{central} & WRP11
\end{tabular}
\end{center}
WRP6 supplies the explicit linear intertwiner $B$ with
the literal original tensor representation, and WRP7
proves the required exact character identity.  These
are precisely the two additional conclusions of
the source existential theorem.  The construction
allows every prime $\ell$, every dominant integer
weight, every given algebraic isomorphism $\iota$,
and every original coefficient datum $c$.  It retains
the cases $n=0$, empty $V$, and places having no
continuous embedding into $\overline{\mathbb Q}_\ell$.
\end{proof}
% END WRP12

\subsection{Scope of the completed supplier}

The basis used here need not be the particular basis
selected by the source's auxiliary \texttt{sq\allowbreak Basis}
theorem over arbitrary domains.  The endpoint just
proved is existential, and its identification with
\texttt{alg\allowbreak Rep} is exactly the explicit intertwiner
\eqref{eq:wrp-intertwiner}; every source assertion holds
for this witness with the unchanged character
\eqref{eq:wrp-lowW}.  Thus this permitted simplification
does not change the representation, the coefficient
embedding, the local groups, the normalized positive
operators, or the input weight.

In conjunction with OCT1--OCT19 this supplies the
previously explicit \texttt{Wt\allowbreak Rep} hypothesis of the
ordinary control theorem for each of the actual
algebraic weights and the fixed complex embedding.
The coefficient extension is the constructed inclusion
$c.O\subset c'.O$, so subsequent use must still make
the corresponding original coefficient map explicit.
This chapter by itself does not assert the entire
\texttt{Spec\allowbreak Int\allowbreak Iff.\allowbreak spec} theorem: the finite-order torus
character chosen there, the single-map assembly over
all levels, and the arithmetic specialization contract
must be supplied where they are used.  Nor does it
assert a completed same-ideal Hida construction or
the global RH/GRH obligation.  Its completed output
is the exact quantified \texttt{wt\allowbreak Rep} supplier above.


\section{Classical ordinary control on the original function spaces}
\label{sec:ordinary-control-v42}

\status{This chapter proves the exact generic endpoint
\texttt{Ctrl\allowbreak Ord.\allowbreak ctrl\allowbreak Ord}.  It uses the completed ordinary structure
chapter and writes out the classical interpolation, lattice, spectral,
and control arguments.  The source assumes a complete Noetherian local
\emph{domain} $O$ with finite residue field, an injective map into an
algebraically closed field $\Omega$, and a supplied weight representation
with all its original fields.  These hypotheses are retained.  In
particular this theorem is distinct from the ordinary base-change
theorem, which allows arbitrary CNL rings and nonflat local maps.  The
construction of an actual algebraic weight representation and the full
arithmetic specialization theorem remain separate producers.}

\subsection{The exact control problem}

Retain the original Frame $X$, groups $G=G_a\times G_p$, subgroups
$U_b,\allowbreak{} N,\allowbreak{} N_f,\allowbreak{} T,\allowbreak{} T_f,\allowbreak{} N^-(c),\allowbreak{} T_c(c)$, positive monoid $P$, and element $z$
from the ordinary structure chapter.  Let $(O,\mathfrak m)$ be a
complete Noetherian local domain with finite residue field, and let
$\iota:O\hookrightarrow\Omega$ be an injective ring map to an
algebraically closed field.  All occurrences of $O$ in $\Omega$ below
mean this specified map.  Fix $m\geq0$ and
\[
 \tau:U_b\longrightarrow\operatorname{GL}_m(O),\qquad
 \theta,\phi:T\longrightarrow O^\times,\qquad U\leq U_b.
\]
Assume the original $\mathrm{FinLevel}(U)$ and
$\mathrm{Free}(\tau,\theta,\allowbreak{} U)$ conditions, with their stabilizer
formula from OST1--OST3.  No module freeness is assumed here.

The supplied $\mathrm{WtRep}$ consists of an integer $d\geq0$,
\[
 \xi:G_p\longrightarrow\operatorname{GL}_{d+1}(\Omega),\qquad
 \chi:T_f\longrightarrow\Omega^\times,\qquad c_0\geq0,
\]
and exactly the following conditions.  Matrices described as integral
have entries in $\iota(O)$.
\begin{enumerate}
\item $\xi(n)$ for $n\in N$, $\xi(t)$ for $t\in T$, and
$\xi(y)$ for $y\in N^-(c_0)$ are integral.
\item $\chi(a)^{-1}\xi(a)$ is integral for every $a\in P$.
\item The zeroth row of $\xi(n)$ is $(1,0,\ldots,0)$ for
$n\in N_f$; the zeroth row of $\xi(t)$ is
$(\chi(t),0,\ldots,0)$ for $t\in T_f$.
\item Every entry outside column zero of
$\chi(z)^{-1}\xi(z)$ lies in $\iota(\mathfrak m)$.
\item For every $k\geq0$ there is $c_1$ such that
$\xi(y)-1$ has entries in $\iota(\mathfrak m^k)$ for
$y\in N^-(c_1)$.
\item For every $k\geq0$ there is $c_2$ such that
$\chi(t)-1\in\iota(\mathfrak m^k)$ for
$t\in T_f\cap T_c(c_2)$.
\item For $t\in T_f\cap Z(G_p)$, one has
$\xi(t)=\chi(t)\,1$.
\end{enumerate}
This list records all ten assertion fields of the original record;
the three integrality assertions in the first item are separate fields,
as are its two zeroth-row assertions.  Fix a classical depth $c$ with
\begin{equation}\label{eq:oct-depth}
 c\geq c_0,\qquad \phi|_{T_c(c)}=1,\qquad
 \iota(\theta(t))=\iota(\phi(t))\chi(t)\quad(t\in T).
\end{equation}

Use $\tau_\Omega$ and $\phi_\Omega$ for the original coefficientwise
images.  The actual classical space $C_c$ consists of families
$F=(F_i)_{0\leq i<m}$, $F_i:G\to\Omega^{d+1}$, satisfying
\begin{align}
 F_i(\operatorname{glob}(\gamma)g)
      &=\xi(\operatorname{glob}(\gamma)_p)F_i(g),\\
 \tau_\Omega(u)F(g(u,1))&=F(g)\quad(u\in U),\\
 F(g(1,\allowbreak{} n))&=F(g)\quad(n\in N),\\
 F(g(1,\allowbreak{} y))&=F(g)\quad(y\in N^-(c)),\\
 F(g(1,\allowbreak{} t))&=\iota(\phi(t))^{-1}F(g)\quad(t\in T).
 \label{eq:oct-classical}
\end{align}
In the away equation $\tau_\Omega$ acts on the $i$ index.  This is
literally $X.\mathrm{ClO}(W,\tau,\allowbreak{} U,\allowbreak{} c,\phi)$, with no replacement by
an abstract isomorphic space.

Define the untwisted family and control map by
\begin{equation}\label{eq:oct-control}
 \widetilde F_i(g)=\xi(g_p)^{-1}F_i(g),\qquad
 \operatorname{ctrl}(F)(g)_i=\widetilde F_i(g)_0.
\end{equation}
For $a\in P$ let $R_a$ be the same quotient representatives used by
the original $U_a$, and put
\begin{equation}\label{eq:oct-T}
 T_aF=\chi(a)^{-1}\sum_{n\in R_a}r_{na}F,
       \qquad T=T_z.
\end{equation}
On the full ambient space this denotes the source's fixed quotient
choices; on $N$-fixed forms its independence is OST5.  Write
\[
 \operatorname{Lat}(k)=
 \{F:\widetilde F_i(g)_j\in\iota(\mathfrak m^k)
                  \text{ for all }i,\allowbreak{} g,\allowbreak{} j\},\qquad
 \Lambda_c=C_c\cap\operatorname{Lat}(0).
\]
More generally $\operatorname{Lat}(D)$ has untwisted values in an
$O$-submodule $D\subset\Omega$.  The classical ordinary space is
the literal source sum
\begin{equation}\label{eq:oct-clord}
 C_c^{\mathrm{ord}}=
 \sum_{\substack{\mu\in\Omega^\times\\
                  \mu,\mu^{-1}\text{ integral over }O}}
           \left(C_c\cap\bigcup_{r\geq0}\ker(T-\mu)^r\right).
\end{equation}
Let $M=\mathrm{Sp}_O(\tau,\theta,\allowbreak{} U)$ and
$\mathcal O=\mathrm{Ord}_O(\tau,\theta,\allowbreak{} U)$ be the original spaces
of $O^m$-valued functions.  The coefficient map $f\mapsto\iota f$
is applied at every $g$ and every coordinate.

\subsection{Classical integral coordinates and the control formula}

% BEGIN OCT1
\paragraph{OCT1. Extending the finite character to the Iwahori subgroup.}
\label{claim:OCT1}
There is a character $\psi_c:I(c)\to O^\times$ that is trivial
on $N$ and $N^-(c)$ and equals $\phi$ on $T$.
\begin{proof}
Use the torus coordinate $t_0:I(0)\to T$ proved in OST2 and set
$\psi_c(b)=\phi(t_0(b))$.  If $b'=nty$ with
$n\in N,\allowbreak{} t\in T,\allowbreak{} y\in N^-(c)$, OST2 gives
\[
 \phi(t_0(bb'))=\phi(t_0(b))\phi(t),\qquad
 \phi(t_0(b'))=\phi(t).
\]
The possible extra torus factor lies in $T_c(c)$ and is killed by
$\phi$ by \eqref{eq:oct-depth}.  Thus $\psi_c(bb')=
\psi_c(b)\psi_c(b')$.  Taking only one of $n,\allowbreak{} t,\allowbreak{} y$ nonidentity
proves its three required restrictions.  This uses the actual
Iwahori torus coordinate, without postulating a torus quotient.
\end{proof}
% END OCT1

% BEGIN OCT2
\paragraph{OCT2. Interpolation of every untwisted classical value.}
\label{claim:OCT2}
Choose distinct finite double-coset representatives $x_1,\ldots,\allowbreak{} x_s$
at depth $c$ as in OST1.  For every list of matrices
$v_r\in\operatorname{Mat}_{m,\allowbreak{} d+1}(\Omega)$ there exists exactly
one $F\in C_c$ with $\widetilde F(x_r)=v_r$.
\begin{proof}
Let $K=U\times I(c)$ and define its action on such matrices by
\[
 \rho(u,\allowbreak{} b)v=\iota(\psi_c(b))\,
                   \tau_\Omega(u)v\xi(b)^{\mathsf T}.
\]
It is a representation: left multiplication, right multiplication
by the transposed matrix, and the scalar commute, and
$\xi(bb')^{\mathsf T}=\xi(b')^{\mathsf T}\xi(b)^{\mathsf T}$
gives the required composition order.

If $\operatorname{glob}(\gamma)x_r=x_r(u,\allowbreak{} b)$, original Free
says that $b\in Z(G_p)\cap T$ and
$\theta(b)\tau(u)=1$.  The central field of $W$ gives
$\xi(b)=\chi(b)1$, while OCT1 and \eqref{eq:oct-depth} give
\[
 \rho(u,\allowbreak{} b)v=\iota(\phi(b))\chi(b)\tau_\Omega(u)v
                       =\iota(\theta(b))\tau_\Omega(u)v=v.
\]
Thus every stabilizer acts trivially, on every prescribed matrix.
For $g=\operatorname{glob}(\gamma)x_r(u,\allowbreak{} b)$ define
\[
                 \Phi(g)=\rho((u,\allowbreak{} b)^{-1})v_r.
\]
The representative index is unique.  Comparing two factorizations
of the same $g$ yields a stabilizer relation and hence equality of
the two displayed values.  Therefore $\Phi$ is well defined, left
$\Gamma$-invariant, and satisfies
$\Phi(g(u,\allowbreak{} b))=\rho((u,\allowbreak{} b)^{-1})\Phi(g)$.

Now define $F_i(g)=\xi(g_p)\Phi_i(g)$, or in row-matrix notation
$F(g)=\Phi(g)\xi(g_p)^{\mathsf T}$.  The cancellation of
$\xi(b)^{-1}$ against $\xi(g_pb)$ shows
$F(g(1,\allowbreak{} b))=\iota(\psi_c(b))^{-1}F(g)$; the same formula for
$(u,1)$ gives the away equation.  Left invariance of $\Phi$ gives
the required twisted left action of $F$.  The three restrictions
of OCT1 give every local equation in \eqref{eq:oct-classical}.
Finally $\widetilde F=\Phi$, so its values are the prescribed
$v_r$.  Conversely any classical $F$ has precisely this
transformation law for $\widetilde F$ and is determined by those
values.  This proves uniqueness as well as existence.
\end{proof}
% END OCT2

% BEGIN OCT3
\paragraph{OCT3. The actual integral lattice is finite free.}
\label{claim:OCT3}
The map
\[
 E_c:C_c\longrightarrow\Omega^{s m(d+1)},\qquad
       E_c(F)_{r,\allowbreak{} i,\allowbreak{} j}=\widetilde F_i(x_r)_j
\]
is an $\Omega$-linear isomorphism.  For every $O$-submodule
$D\subset\Omega$,
\begin{equation}\label{eq:oct-lattice-test}
 F\in C_c\cap\operatorname{Lat}(D)\quad\Longleftrightarrow
          E_c(F)\in D^{s m(d+1)}.
\end{equation}
Consequently $\Lambda_c\simeq O^{s m(d+1)}$ and
$C_c\cap\operatorname{Lat}(k)=\mathfrak m^k\Lambda_c$.
\begin{proof}
The isomorphism is OCT2 and linearity of its defining formulas.
The forward implication of \eqref{eq:oct-lattice-test} is immediate.
For the converse write
$g=\operatorname{glob}(\gamma)x_r(u,\allowbreak{} b)$ and
$b=nty$ with $n\in N,\allowbreak{} t\in T,\allowbreak{} y\in N^-(c)$, by OST1.
Left multiplication by $\operatorname{glob}(\gamma)$ leaves
$\widetilde F$ unchanged.  Right multiplication by $(u,1)$ acts
on its $i$ index by $\tau_\Omega(u^{-1})$, an integral matrix.
For each local factor $b_0\in\{n,\allowbreak{} t,\allowbreak{} y\}$ one has
\[
 \widetilde F_i(h(1,\allowbreak{} b_0))=
          \iota(s_0)\xi(b_0^{-1})\widetilde F_i(h),
\]
where $s_0=1$ for $n,\allowbreak{} y$ and $s_0=\phi(t)^{-1}$ for $t$.
The three matrices are integral by the first fields of $W$;
for the last one use $N^-(c)\subset N^-(c_0)$.  Multiplication
by an integral matrix preserves $D$ in every coordinate, since
$D$ is an $O$-module.  Successively removing the factors proves
the reverse implication.  Taking $D=\iota(O)$ and then
$D=\iota(\mathfrak m^k)$ proves the two final assertions using
injectivity of $\iota$.  This is an assertion about the exact
value lattice of the original functions.
\end{proof}
% END OCT3

% BEGIN OCT4
\paragraph{OCT4. Exact equivariance of the original control map.}
\label{claim:OCT4}
The map in \eqref{eq:oct-control} is $\Omega$-linear and satisfies
\begin{equation}\label{eq:oct-equivariance}
 \operatorname{ctrl}(T_aF)=U_a\operatorname{ctrl}(F),\qquad
 \operatorname{ctrl}(\mathrm{raF}(u)F)
                         =\mathrm{ra}(u)\operatorname{ctrl}(F).
\end{equation}
Here the away actions use the same $\tau_\Omega$, and the first
identity holds for every ambient family and every $a\in P$.
\begin{proof}
Linearity follows immediately from multiplication by
$\xi(g_p)^{-1}$ and extraction of coordinate zero.  For
$n\in N_f,\allowbreak{} t\in T_f$ and any vector $v$, the zeroth-row fields
of $W$ give
\[
 \bigl(\xi(g_pnt)^{-1}v\bigr)_0
      =\chi(t)^{-1}\bigl(\xi(g_p)^{-1}v\bigr)_0.
\]
Indeed the inverse matrix factors as
$\xi(t^{-1})\xi(n^{-1})\xi(g_p)^{-1}$; the first two zeroth
rows successively multiply coordinate zero by $\chi(t)^{-1}$
and leave it unchanged.  Apply this with $t=a$ and the value
$v=F_i(g(1,\allowbreak{} na))$ for each original representative $n\in R_a$.
Summing is exactly the first identity.  The away action does
not change $g_p$, and its matrix acts on the separate $i$ index;
the two finite matrix sums commute, proving the second identity.
The same separation of factors proves that away actions commute
with every local right translation.  No invariance hypothesis
was needed for these identities.
\end{proof}
% END OCT4

% BEGIN OCT5
\paragraph{OCT5. Integral classical forms control to the actual ordinary ambient space.}
\label{claim:OCT5}
The character $\theta$ is continuous in the original sense.
Every $F\in\Lambda_c$ has a unique $f\in M$ with
$\iota f=\operatorname{ctrl}(F)$.  More precisely, if $c'\geq c$
and $\xi(y)-1$ has entries in $\iota(\mathfrak m^k)$ for
$y\in N^-(c')$, then
\begin{equation}\label{eq:oct-ctrl-continuity}
                r_yf-f\in\mathrm{valIn}_O(k)
                         \quad(y\in N^-(c')).
\end{equation}
\begin{proof}
Choose $c'\geq c$ with $\chi(t)-1\in\iota(\mathfrak m^k)$
on $T_c(c')$.  There $\phi(t)=1$, so
\eqref{eq:oct-depth} and injectivity of $\iota$ give
$\theta(t)-1\in\mathfrak m^k$.  This is ContChar.

Since $F$ is integral, all its untwisted values have unique
$O$-preimages.  Define $f(g)_i$ as the preimage of
$\widetilde F_i(g)_0$.  Twisted left equivariance cancels in
$\widetilde F$, giving the left invariance of $f$.  The zeroth
row of $\xi(n^{-1})$ proves its $N$-invariance.  The zeroth
row of $\xi(t^{-1})$, together with
$\iota(\theta(t))=\iota(\phi(t))\chi(t)$, gives its
$\theta^{-1}$ torus law.  The away equation follows from OCT4
and injectivity of $\iota$.

For continuity take $y\in N^-(c')$.  Since $c'\geq c$,
$F(g(1,\allowbreak{} y))=F(g)$, and therefore
\[
 \widetilde F_i(g(1,\allowbreak{} y))=\xi(y^{-1})\widetilde F_i(g).
\]
The inverse $y^{-1}$ lies in the same subgroup.  If
$\widetilde F_i(g)_j=\iota(v_j)$ and
$\xi(y^{-1})_{0j}=\delta_{0j}+\iota(a_j)$ with
$a_j\in\mathfrak m^k$, then
\[
                  f(g(1,\allowbreak{} y))_i-f(g)_i=\sum_j a_jv_j
                                                  \in\mathfrak m^k.
\]
This proves \eqref{eq:oct-ctrl-continuity}, uniformly in $g,\allowbreak{} i$.
The continuity field of $W$ supplies such a depth for every $k$,
so $f$ satisfies every condition of the original $M$.  Its
uniqueness and the $O$-linearity of $F\mapsto f$ follow from
injectivity of $\iota$.
\end{proof}
% END OCT5

% BEGIN OCT6
\paragraph{OCT6. Integral stability and strict contraction in the control kernel.}
\label{claim:OCT6}
Every $T_a$, $a\in P$, preserves $\operatorname{Lat}(D)$ for
every $O$-submodule $D\subset\Omega$.  More specifically, if
$\mathfrak m D\subset D'$, $F\in\operatorname{Lat}(D)$, and
all control values of $F$ lie in $D'$, then
\begin{equation}\label{eq:oct-contraction}
                         TF\in\operatorname{Lat}(D').
\end{equation}
\begin{proof}
The untwisted value of one normalized summand in $T_aF$ is
\begin{equation}\label{eq:oct-summand}
 \xi(n)\bigl(\chi(a)^{-1}\xi(a)\bigr)
                         \widetilde F_i(g(1,\allowbreak{} na)).
\end{equation}
This follows by substituting
$F_i(h)=\xi(h_p)\widetilde F_i(h)$ in the defining sum and
cancelling $\xi(g_p)^{-1}\xi(g_p)$.  Both matrices in
\eqref{eq:oct-summand} are integral by the fields of $W$, so
each summand preserves $D$.

For $a=z$, every column of the second matrix except column
zero has entries in $\iota(\mathfrak m)$.  Their contributions
therefore lie in $D'$.  The contribution of column zero lies
in $D'$ because its input is the control value and its
coefficient is integral.  Multiplication by $\xi(n)$ preserves
$D'$.  The finite sum proves \eqref{eq:oct-contraction}.
In particular $T$ preserves every $\operatorname{Lat}(k)$, and
if $\operatorname{ctrl}(F)=0$ and
$F\in\operatorname{Lat}(k)$, then
$TF\in\operatorname{Lat}(k+1)$.
\end{proof}
% END OCT6

% BEGIN OCT7
\paragraph{OCT7. Classical level lowering on the original functions.}
\label{claim:OCT7}
For $c'\geq c$ one has $C_c\subset C_{c'}$.  Every $T_a$
preserves $C_c$ and $\Lambda_c$.  For integers $e,\allowbreak{} n\geq0$
with $n\geq e$,
\begin{equation}\label{eq:oct-lowering}
              T^n C_{c+e}\subset C_c,\qquad
              T^n\Lambda_{c+e}\subset\Lambda_c.
\end{equation}
\begin{proof}
The inclusion of levels follows from
$N^-(c')\subset N^-(c)$.  Twisted left equivariance is
preserved term by term under right translation, and away
equivariance commutes with it by OCT4.  For local invariance,
the permutation-of-cosets proof OST6 applies to the
representation on the classical families: $\phi(T_c(c))=1$
gives its required torus-congruence invariance.  The same
proof gives preservation of the $N$ and torus character laws
by OST5.  Multiplication by the scalar $\chi(a)^{-1}$ does
not alter any of these laws.  Thus $T_aC_c\subset C_c$;
OCT6 adds lattice preservation.

At each depth $c+r$, $\phi$ remains trivial on $T_c(c+r)$.
The strict part of OST6 gives
$T C_{c+r+1}\subset C_{c+r}$.  Apply it $e$ times and then
use stability at depth $c$ for the remaining $n-e$ iterates.
OCT6 proves the integral assertion at every step.  The
normalizing factor remains the original $\chi(z)^{-1}$.
\end{proof}
% END OCT7

\subsection{A finite lattice and its actual spectral ordinary part}

% BEGIN OCT8
\paragraph{OCT8. Factorial decomposition of a finite integral matrix.}
\label{claim:OCT8}
Let $L\simeq O^r$ be finite free and $T_0\in\operatorname{End}_O(L)$.
There is one idempotent $e_0$ commuting with $T_0$ such that
\[
 e_0=\lim_{j\to\infty}T_0^{j!},\qquad
 L=S_0\oplus Q_0,\quad S_0=e_0L,
                         \quad Q_0=(1-e_0)L.
\]
Both summands are finite free.  The restriction of $T_0$ to
$S_0$ is invertible; for every $k$ some $n$ satisfies
$T_0^nQ_0\subset\mathfrak m^k Q_0$.  Moreover
\begin{equation}\label{eq:oct-periodic}
 v\in S_0\quad\Longleftrightarrow\quad
 \forall k\ \exists q\geq1\quad T_0^qv-v\in\mathfrak m^kL.
\end{equation}
At every $k$ the factorial convergence is uniform in $v\in L$.
\begin{proof}
Each $O/\mathfrak m^k$ is finite.  Indeed
$\mathfrak m$ is finitely generated, so each
$\mathfrak m^j/\mathfrak m^{j+1}$ is finite-dimensional over
the finite residue field; induction along these quotients
gives finiteness.  Thus the endomorphism monoid of
$L/\mathfrak m^kL$ is finite.  Some $b\geq0,\allowbreak{} q\geq1$ satisfy
$T_0^{n+q}=T_0^n$ on this quotient for every $n\geq b$.
For all sufficiently large $j$, $j!\geq b$ and $q\mid j!$;
the endomorphism $T_0^{j!}$ is then one fixed idempotent $e_k$.
Idempotence follows by comparing the exponents $2j!$ and $j!$.
On its image $T_0^{j!}$ is the identity, so $T_0$ is invertible
there; on its kernel that power is zero.

These $e_k$ are compatible, since every one is given by all
sufficiently late members of the same factorial sequence.
Completeness of $O$ constructs their matrix-entry limit $e_0$;
separatedness makes it unique and gives $e_0^2=e_0$ and
$e_0T_0=T_0e_0$.  The decomposition is now literal.  Its
summands are finite projective and hence finite free over the
local ring; the elementary projective-to-free argument is
written in OST17.  Their reductions identify with the image
and kernel of $e_k$, because the decomposition is split.
On a basis of $S_0$, the determinant of $T_0$ is a unit modulo
$\mathfrak m$, hence a unit of $O$; the adjugate supplies the
inverse over $O$.  The kernel statement modulo each
$\mathfrak m^k$ gives the asserted topological nilpotence.

If the right side of \eqref{eq:oct-periodic} holds, take a
factorial divisible by its period $q$ at precision $k$.
Then $e_kv=v$ modulo $\mathfrak m^k$, for every $k$, whence
$e_0v=v$.  Conversely $T_0$ acts as an automorphism of the
finite module $S_0/\mathfrak m^kS_0$, so some positive power
fixes the given vector.  The construction of $e_k$ as an
endomorphism, rather than as a limit for one vector, proves
the asserted uniform convergence.  The case $r=0$ satisfies
the same statements with zero modules.
\end{proof}
% END OCT8

% BEGIN OCT9
\paragraph{OCT9. The finite lattice ordinary part equals the integral-unit spectrum.}
\label{claim:OCT9}
Extend the data of OCT8 to $L_\Omega=\Omega\otimes_O L$.
The $\Omega$-span of $S_0$ is exactly the sum of generalized
eigenspaces of $T_0$ whose eigenvalues $\mu$ satisfy
$\mu\ne0$ and both $\mu,\mu^{-1}$ integral over $O$.
\begin{proof}
On $S_0$, both $T_0$ and $T_0^{-1}$ are integral matrices.
Their monic characteristic polynomials show that every
eigenvalue $\mu$ and its inverse are integral over $O$;
invertibility shows $\mu\ne0$.

On $Q_0$, choose $n\geq1$ with
$T_0^nQ_0\subset\mathfrak m Q_0$.  If $q$ is its rank,
the characteristic polynomial of $T_0^n|_{Q_0}$ is
\[
             X^q+a_1X^{q-1}+\cdots+a_q,
                              \qquad a_j\in\mathfrak m^j.
\]
The coefficient condition follows by expanding principal
minors: each degree-$j$ minor is a sum of products of $j$
entries in $\mathfrak m$.  Suppose an eigenvalue $\mu$ of
$T_0|_{Q_0\otimes\Omega}$ were nonzero with $\mu^{-1}$
integral over $O$.  Substitution of $\mu^n$ and division by
$\mu^{nq}$ would give
\[
              1+a_1\mu^{-n}+\cdots+a_q\mu^{-nq}=0.
\]
The algebra $B=O[\mu^{-1}]\subset\Omega$ is a finite
$O$-module, since $\mu^{-1}$ satisfies a monic equation.
The displayed identity gives $1\in\mathfrak m B$, so
$B=\mathfrak m B$.  Nakayama gives $B=0$, contradicting
$1\ne0$ in $\Omega$.  Thus the complementary block has no
such eigenvalue.  If $q=0$ there is no complementary block.

For completeness, over an algebraically closed field the
characteristic polynomial factors as
$\prod_\mu(X-\mu)^{r_\mu}$.  The factors are pairwise
coprime; Bezout identities in $\Omega[X]$, evaluated at the
operator and combined with its characteristic equation,
decompose the vector space into the kernels
$\ker(T_0-\mu)^{r_\mu}$.  These are precisely the maximal
generalized eigenspaces.  Apply this decomposition separately
to the two $T_0$-stable summands of OCT8.  All eigenvalues on
the first block, and none on the second, meet the required
integral-unit condition.  This proves the equality.  No
valuation on $\Omega$ or normality of $O$ was used.
\end{proof}
% END OCT9

% BEGIN OCT10
\paragraph{OCT10. A finite free lattice spanning the exact classical ordinary space.}
\label{claim:OCT10}
For the original $C_c,\Lambda_c,\allowbreak{} T$, let $e_c$ be the factorial
projector on $\Lambda_c$ supplied by OCT8 and set
$S_c=e_c\Lambda_c$.  Then $S_c$ is finite free,
\begin{equation}\label{eq:oct-spectral}
                      C_c^{\mathrm{ord}}=\Omega S_c,
\end{equation}
and every $F\in S_c$ is periodic modulo every
$\operatorname{Lat}(k)$.  For every $G\in\Lambda_c$ and
every $k$,
\begin{equation}\label{eq:oct-classical-factorial}
 \exists J\ \forall j\geq J\quad
              T^{j!}G-e_cG\in\operatorname{Lat}(k).
\end{equation}
The choice of $e_c$ precedes $G,\allowbreak{} k$, and its threshold can be
chosen uniformly in $G\in\Lambda_c$ for each $k$.
\begin{proof}
By OCT3 the lattice is finite free and its extension to
$\Omega$ is exactly $C_c$ through the evaluation isomorphism.
OCT7 says that $T$ is an $O$-linear endomorphism of this
lattice.  Apply OCT8, using OCT3 to identify its intrinsic
filtration $\mathfrak m^k\Lambda_c$ with its global
untwisted-value filtration.  OCT9 gives
\eqref{eq:oct-spectral}.  Because $T$ preserves $C_c$,
the generalized eigenspace of its restriction at $\mu$ is
exactly $C_c\cap\bigcup_r\ker(T-\mu)^r$ in the original
ambient function space; the defining iterates are the same
on $C_c$.  Thus the spectral sum is exactly
\eqref{eq:oct-clord}, not a substituted definition.
\end{proof}
% END OCT10

\subsection{Integral control and its inverse}

% BEGIN OCT11
\paragraph{OCT11. Periodic integral elements in the control kernel vanish.}
\label{claim:OCT11}
Suppose an ambient family $F$ lies in $\operatorname{Lat}(0)$,
has $\operatorname{ctrl}(F)=0$, and for every $k$ has a
positive period $q$ modulo $\operatorname{Lat}(k)$:
$T^qF-F\in\operatorname{Lat}(k)$.  Then $F=0$.
\begin{proof}
Induct on $k$ to show $F\in\operatorname{Lat}(k)$.  The
initial assertion is assumed.  If it holds at $k$, OCT6 gives
$TF\in\operatorname{Lat}(k+1)$, and stability in OCT6 gives
$T^qF\in\operatorname{Lat}(k+1)$ for every $q\geq1$.
Choose a positive period at precision $k+1$ and subtract its
congruence to obtain $F\in\operatorname{Lat}(k+1)$.

For each untwisted coordinate, choose its $O$-preimage using
membership in $\operatorname{Lat}(0)$.  Injectivity of
$\iota$ and membership at every higher precision place that
same preimage in $\bigcap_k\mathfrak m^k=0$.  Therefore all
untwisted coordinates vanish.  Multiplying by the invertible
matrix $\xi(g_p)$ gives $F_i(g)=0$ for all $i,\allowbreak{} g$.
This argument does not presume that $F$ belongs to a finite
dimensional space, and uses positive, rather than zero,
periods.
\end{proof}
% END OCT11

% BEGIN OCT12
\paragraph{OCT12. The integral ordinary lattice maps injectively to the original ordinary module.}
\label{claim:OCT12}
There is a unique $O$-linear map
\[
             q_c:S_c\longrightarrow\mathcal O,
                  \qquad \iota(q_cF)=\operatorname{ctrl}(F),
\]
and it is injective.
\begin{proof}
OCT5 first supplies the unique $f\in M$.  At any precision
$k$, OCT10 gives $T^qF-F\in\operatorname{Lat}(k)$ for some
$q\geq1$.  Its zeroth untwisted coordinates, OCT4, and
coefficientwise naturality of the original finite sums give
\[
                     U_z^q f-f\in\mathrm{valIn}_O(k).
\]
Here membership follows from injectivity of $\iota$:
an $O$-value whose image lies in $\iota(\mathfrak m^k)$
already lies in $\mathfrak m^k$.  The literal periodicity
characterization OST10 therefore gives $f\in\mathcal O$.
Linearity follows by uniqueness.  If $q_cF=0$, OCT10 supplies
all the periods needed in OCT11, which then gives $F=0$.
\end{proof}
% END OCT12

% BEGIN OCT13
\paragraph{OCT13. An integral classical approximation at the required fixed depth.}
\label{claim:OCT13}
For every $f\in\mathcal O$ there are
$G\in\Lambda_c$ and $h\in M$ such that
\begin{equation}\label{eq:oct-approx}
 \iota h=\operatorname{ctrl}(G),\qquad
                         h-f\in\mathrm{valIn}_O(1).
\end{equation}
\begin{proof}
Choose a depth at which $f$ is $N^-$-invariant modulo
$\mathfrak m$, and a depth at which $\xi(y)-1$ has all
entries in $\iota(\mathfrak m)$.  Take $c'$ at least both
these depths and $c$, and write $e=c'-c$.  Triviality of
$\phi$ on $T_c(c')$ still holds.  Choose the finite distinct
representatives $x_r$ at this deeper depth.  OCT2 constructs
$F'\in C_{c'}$ with
\[
      \widetilde F'_i(x_r)_j=
          \begin{cases}\iota(f(x_r)_i),&j=0,\\0,&j\ne0.
          \end{cases}
\]
These are integral values, so OCT3 gives $F'\in\Lambda_{c'}$.
OCT5 gives its integral control $h'\in M$, with
$h'(x_r)=f(x_r)$.  Its refined continuity assertion says
that $h'$ is $N^-(c')$-invariant modulo $\mathfrak m$.
The same is true for $f$ by the choice of $c'$.  Apply
OST4 to $h'-f$, whose sampled values are zero, to obtain
$h'-f\in\mathrm{valIn}_O(1)$.

Choose a period $q\geq1$ for $f$ at this precision, and a
multiple $M_0$ of $q$ with $M_0\geq e$, for example
$M_0=q(e+1)$.  Set
\[
                         G=T^{M_0}F',\qquad h=U_z^{M_0}h'.
\]
OCT7 lowers $G$ from depth $c'$ to depth $c$, and OCT6
preserves its integrality.  OST7 gives $h\in M$.
OCT4 and coefficientwise naturality give
$\iota h=\operatorname{ctrl}(G)$.  Finally
\[
 h-f=U_z^{M_0}(h'-f)+(U_z^{M_0}f-f)
                                             \in\mathrm{valIn}_O(1).
\]
The first term belongs to the value filtration by OST5.
For the second, a positive multiple of a period is a
period: expand
$U_z^{qt}f-f=\sum_{a=0}^{t-1}U_z^{qa}(U_z^qf-f)$.
This proves \eqref{eq:oct-approx}.  The deeper depth is an
intermediate approximation choice; the final classical form
is at the original prescribed depth $c$.
\end{proof}
% END OCT13

% BEGIN OCT14
\paragraph{OCT14. Factorial projection gives surjectivity modulo the maximal ideal.}
\label{claim:OCT14}
The image of $q_c$ satisfies
\[
                         q_c(S_c)+\mathfrak m\mathcal O
                                                   =\mathcal O.
\]
\begin{proof}
For $f\in\mathcal O$, choose $G,\allowbreak{} h$ as in OCT13.  Put
$F_0=e_cG\in S_c$ and $f_0=q_cF_0\in\mathcal O$.
By OCT10, at precision one there is $J$ such that
$T^{j!}G-F_0\in\operatorname{Lat}(1)$ for $j\geq J$.
Choose a positive period $q$ for $f$ modulo
$\mathrm{valIn}_O(1)$ and then choose $j\geq\max(J,\allowbreak{} q)$.
The factorial $j!$ is divisible by $q$.  Applying control
and using OCT4 gives
\[
 U_z^{j!}h-f_0\in\mathrm{valIn}_O(1),\qquad
 U_z^{j!}f-f\in\mathrm{valIn}_O(1).
\]
Also $U_z^{j!}(h-f)$ lies in that filtration by OST5.
The identity
\[
 f-f_0=(U_z^{j!}h-f_0)-U_z^{j!}(h-f)-(U_z^{j!}f-f)
\]
therefore places $f-f_0$ in $\mathrm{valIn}_O(1)$.
Both $f$ and $f_0$ are ordinary, so the exact reverse
filtration identity OST16 gives
$f-f_0\in\mathfrak m\mathcal O$.
This step uses the established intrinsic filtration on
$\mathcal O$; a value congruence by itself would not suffice.
\end{proof}
% END OCT14

% BEGIN OCT15
\paragraph{OCT15. Integral control is an isomorphism.}
\label{claim:OCT15}
The map $q_c:S_c\to\mathcal O$ is an $O$-linear isomorphism.
\begin{proof}
OCT5 gives ContChar, so OST19 applies with the original
FinLevel and Free hypotheses and makes $\mathcal O$ finite
free.  The quotient $Q=\mathcal O/q_c(S_c)$ is therefore a
finite $O$-module.  OCT14 gives $Q=\mathfrak m Q$.
Nakayama gives $Q=0$, so $q_c$ is surjective.  Its injectivity
was proved in OCT12.  One may see the used form of Nakayama
directly: if generators $v_1,\ldots,\allowbreak{} v_r$ of $Q$ express each
$v_i$ as $\sum_j a_{ij}v_j$ with $a_{ij}\in\mathfrak m$,
then the adjugate of $1-(a_{ij})$ shows that its determinant
annihilates every generator.  This determinant lies in
$1+\mathfrak m$ and is a unit, forcing $Q=0$.
Thus the inverse lifts ordinary functions to actual classical
families and uses exactly the control formula
\eqref{eq:oct-control}.
\end{proof}
% END OCT15

% BEGIN OCT16
\paragraph{OCT16. Independent integral functions remain independent over the ambient field.}
\label{claim:OCT16}
Let $f_1,\ldots,\allowbreak{} f_r$ be $O$-linearly independent functions
from any set $D$ to $O$.  Then their coefficientwise images
are $\Omega$-linearly independent in $\Omega^D$.
\begin{proof}
The embedding of $O$ extends to an embedding of its fraction
field $K$ into $\Omega$.  Clearing finitely many denominators
in a $K$-linear relation shows that the $f_i$, now viewed as
$K$-valued functions, remain independent.  Let $V$ be their
$r$-dimensional $K$-span.  Evaluation at the points of $D$
spans $V^*$: otherwise the common kernel of these linear
functionals would contain a nonzero member of $V$, a
function zero at every point, which is impossible.  Choose
$r$ evaluations forming a basis of $V^*$.  The resulting
$r$-by-$r$ matrix $(f_i(x_j))$ has nonzero determinant in
$K$, and its entries and determinant lie in $O$.
Injectivity of $\iota$ makes that determinant nonzero in
$\Omega$.  Evaluating any $\Omega$-linear relation at those
$r$ points now forces every coefficient to vanish.  For
$r=0$ the statement is vacuous.  This is the precise point
where the source's domain hypothesis is used; it is not
silently removed from the control theorem.
\end{proof}
% END OCT16

% BEGIN OCT17
\paragraph{OCT17. The exact original classical ordinary control endpoint.}
\label{claim:OCT17}
Under exactly the standing hypotheses and
\eqref{eq:oct-depth}, the source's original map satisfies
\begin{align}
 F\in C_c^{\mathrm{ord}},\quad\operatorname{ctrl}(F)=0
                                     &\Longrightarrow F=0,\\
 \operatorname{ctrl}(C_c^{\mathrm{ord}})
                   &=\operatorname{span}_{\Omega}
                                \{\iota f:f\in\mathcal O\}.
 \label{eq:oct-endpoint}
\end{align}
These are precisely the two conclusions of
\texttt{Upper\allowbreak Theory.\allowbreak V.\allowbreak Hi.\allowbreak ctrl\allowbreak Ord}.
\begin{proof}
By OCT10, $C_c^{\mathrm{ord}}=\Omega S_c$.
Take an $O$-basis $F_1,\ldots,\allowbreak{} F_r$ of $S_c$ and put
$f_i=q_cF_i$.  By OCT15 these are an $O$-basis of the
original $\mathcal O$.  Apply OCT16 to their coordinates
as functions on $D=G\times\{0,\ldots,\allowbreak{} m-1\}$; the resulting
$\Omega$-valued functions $\iota f_i=\operatorname{ctrl}(F_i)$
are independent.  Every element of $\Omega S_c$ is a finite
$\Omega$-linear combination of the $F_i$, so zero control
forces all its coefficients to vanish.  Linearity and the
same basis show that its image is exactly the displayed
span.  This also shows that the coefficientwise tensor map
$\Omega\otimes_O\mathcal O\to\Omega^{G\times\{0,\ldots,\allowbreak{} m-1\}}$
is injective, and control identifies $C_c^{\mathrm{ord}}$
with its image.  All spaces in this argument are the
original spaces of functions from the source definitions.
\end{proof}
% END OCT17

\subsection{Compatibility and the remaining installation boundary}

% BEGIN OCT18
\paragraph{OCT18. Compatibility with original positive operators and classical depth.}
\label{claim:OCT18}
Every $T_a$, $a\in P$, preserves $S_c$, and
\[
                  q_cT_a=U_aq_c.
\]
It is bijective on $S_c$.  If $c'\geq c$ satisfies the same
standing depth requirements, then $S_c=S_{c'}$ as submodules
of the original ambient families, and their integral control
inverses are identical.  Thus the ordinary classical subspace
and the inverse control map are independent of sufficiently
deep classical level, with no new choices at each precision.
\begin{proof}
OCT7 and OCT6 show that $T_a$ preserves $\Lambda_c$.
On $N$-fixed families, the finite double-coset product rule
OST5 and multiplicativity of $\chi$ give
$T_aT=T T_a$.  Consequently $T_a$ preserves the periodic
submodule $S_c$ characterized by OCT8, or equivalently
commutes with its factorial projector.  OCT4 gives
equivariance under $q_c$.  By OST20 the original $U_a$ is
bijective on $\mathcal O$, and OCT15 therefore makes $T_a$
bijective on $S_c$.

For depths, OCT7 gives $\Lambda_c\subset\Lambda_{c'}$,
and periodicity in the global value lattices gives
$S_c\subset S_{c'}$.  The two maps $q_c,\allowbreak{} q_{c'}$ are
restrictions of the same pointwise formula and are both
isomorphisms onto $\mathcal O$.  For $F'\in S_{c'}$ choose
$F\in S_c$ with $q_cF=q_{c'}F'$.  Their difference belongs
to $S_{c'}$ and has zero control, so OCT12 gives $F'=F$.
This proves equality and identity of the inverse maps.
OCT10 then also gives equality of the classical ordinary
$\Omega$-spaces at these depths.
\end{proof}
% END OCT18

% BEGIN OCT19
\paragraph{OCT19. Compatibility with finite sums of actual away translations.}
\label{claim:OCT19}
Suppose two level/depth pairs $(U,\allowbreak{} c)$ and $(U',\allowbreak{} c')$ satisfy
the standing hypotheses for the same $O,\Omega,\allowbreak{} W,\tau,
\theta,\phi$.  Let
\[
       A_F=\sum_{u\in E}\mathrm{raF}(u),\qquad
       A_f=\sum_{u\in E}\mathrm{ra}(u),
\]
where $E$ is any finite indexed list of elements of $U_b$,
with repetitions retained.  If $A_F$ maps $C_c(U)$ into
$C_{c'}(U')$, then it maps $S_c(U)$ into $S_{c'}(U')$ and
\[
                         q_{c',\allowbreak{} U'}A_F=A_f q_{c,\allowbreak{} U}.
\]
In particular the same assertion holds for the natural
inclusion whenever $U'\subset U$ and $c'\geq c$.
\begin{proof}
An away translation acts on untwisted coordinates by the
integral matrix $\tau_\Omega(u)$ and an away translation
of $g$.  It preserves every global lattice
$\operatorname{Lat}(k)$, as does a finite sum.  It also
commutes with $T$, since the away and local actions commute
by OCT4.  Under the assumed classical-space inclusion, a
periodic element of $\Lambda_c(U)$ therefore maps to an
integral periodic element of $C_{c'}(U')$.  The criterion
OCT8 and OCT3 put its image in $S_{c'}(U')$.
OCT4 gives the displayed equality after applying $\iota$;
injectivity of $\iota$ gives the equality over $O$.
For natural inclusion use the singleton list $u=1$ and
the defining equations for the two levels.  Thus these
compatibilities concern the original translations and
actual inclusions, rather than an abstract identification
of modules of the same rank.
\end{proof}
% END OCT19

\paragraph{Exact source and proof boundary.}
The source pin is
\begin{center}
\texttt{7cbf7a0c32fca1badf05931879a37c3f8a2b03b4}.
\end{center}
The source definitions are
\texttt{Continuous\allowbreak Torus\allowbreak Character.\allowbreak lean} (the classical section)
and \texttt{Ctrl\allowbreak Defs.\allowbreak lean}; the endpoint is
\texttt{Ordinary\allowbreak Hecke\allowbreak Characters/\allowbreak Ctrl\allowbreak Ord.\allowbreak lean}, with the control
identities in \texttt{Ctrl\allowbreak Eq.\allowbreak lean}.  The interpolation,
integral-lattice, spectral, and continuity helper files and
their exact byte hashes are recorded in the accompanying
source manifest.  OCT1--OCT19 are written proofs; the source
locations identify what they reconstruct and are not used
as black-box justifications of those proofs.

The endpoint quantifies over every supplied original Frame
and every supplied $W$ satisfying all ten assertion
fields, every $m$, every $\tau,\theta,\phi$, and every
admissible $U,\allowbreak{} c$.  It proves the exact domain-valued control
theorem; it neither changes the source to arbitrary
non-domain coefficients nor discards the specified embedding
$O\hookrightarrow\Omega$.  The arithmetic use in
\texttt{Spec\allowbreak Int\allowbreak Iff.\allowbreak spec} first supplies such a $W$ over a
finite coefficient extension, a character $\phi$, and the
compatibility \eqref{eq:oct-depth}.  That construction,
the arithmetic-fiber norm dictionary, the full single-map
specialization theorem, and the same-ideal Hida/automorphy
endpoints are not proved merely by invoking OCT17.  They
remain separate obligations unless independently supplied
elsewhere in the assembled reconstruction.  No whole
\texttt{COUNT\_\allowbreak II\_\allowbreak SUPPORT} leaf is retired by this chapter.


\section{The original ordinary Hecke algebra and its residual block}
\label{sec:ordinary-hecke-algebra-v042}

\status{This chapter proves the original \texttt{hi\allowbreak Comm} and
\texttt{hi\allowbreak Alg}, using the actual ordinary modules and operators.
The base coefficient ring may have nilpotents. The domain hypothesis
is used only for the source's torsion-freeness conclusion. No
arithmetic-point separation or automorphy theorem is inferred.}

Let \(H\) be the original \texttt{HSet}, with its original
\texttt{Ax}, coefficient DVR \(O=c.O\), and an \(O\)-algebra \(A\)
of class \texttt{Is\allowbreak CO}. This means that \(A\) is a complete
Noetherian local ring, the map \(O\to A\) is local, and it induces
an isomorphism on residue fields. In particular the residue field
of \(A\) is finite, so \(A\) is of class \texttt{Is\allowbreak CNL}.
Fix the original character \(\theta\) satisfying \texttt{Is\allowbreak Wt},
and fix any original auxiliary level \(K\). Put
\[
 U=H.\mathrm{lev}(K),\qquad
 Y=H.X.\mathrm{Ord}(H.\tau_A(A),\theta,\allowbreak{} U).
\]
OST19 makes \(Y\) finite free. The operator index set is literally
\[
       H.\mathrm{Op}=(H.W'\times\{1,\ldots,\allowbreak{} n\})
                    \sqcup(H.V_\ell\times\{1,\ldots,\allowbreak{} n\}).
\]
The first kind of operator is the integral spherical double-coset
sum at \(w\), acting through \(H.\mathrm{sec}(w)\) and the original
away action. The second is \(U_{a_{v,\allowbreak{} j}}\), where \(a_{v,\allowbreak{} j}\) has
\(\operatorname{diag}(\varpi_v I_j,\allowbreak{} I_{n-j})\) at \(v\) and identity
elsewhere. In the source the displayed \(j\) is represented by
an element of \(\operatorname{Fin}(n)\) plus one.

For a prescribed function \(a:H.\mathrm{Op}\to O\), write
\(a_i^A\) for its image in \(A\), and define the actual residual
part
\begin{equation}\label{eq:oha-locpart}
 L=\left\{y\in Y:
   \ \forall i\ \forall k\ge0\ \exists r\ge0,\
               (T_i-a_i^A)^r y\in\mathfrak m_A^kY\right\}.
\end{equation}
This is the source \(H.MK(\theta,\allowbreak{} a,\allowbreak{} K)\). Its algebra \(H.TK\)
is the \(A\)-subalgebra of \(\operatorname{End}_A(L)\) generated
by the restrictions of these same \(T_i=H.\mathrm{op}(A,\allowbreak{} i)\).
The powers in \eqref{eq:oha-locpart} are allowed to depend on
\(i,\allowbreak{} k,\allowbreak{} y\); the proof below does not replace this condition by
semisimplicity.

% BEGIN OHA1
\paragraph{OHA1. Integral finite double-coset sums on arbitrary representations.}
\label{claim:OHA1}
Let \(G\) be a group, \(C\subset G\) a subgroup, and assume
\(CgC/C\) is finite for every \(g\). For any representation
\(\rho:G\to\operatorname{Aut}_A(M)\) and \(v\in M^C\), set
\[
                 [CgC]v=\sum_{xC\in CgC/C}\rho(x)v .
\]
This is independent of the representatives: replacing \(x\)
by \(xc\) fixes its summand. Left multiplication by \(C\)
permutes the summands, so the result belongs to \(M^C\).
An operator commuting with every \(\rho(x)\) commutes with
this sum.

More generally, a finitely supported bi-\(C\)-invariant integer
function on double cosets acts by
\(\sum_{xC}f(x)\rho(x)\). Expanding two finite sums and collecting
the coefficient of each \(zC\) gives
\[
 (f*h)(z)=\sum_{xC\in G/C}f(x)h(x^{-1}z).
\]
Hence this action respects convolution. Each coefficient is an
integer, and the same identity holds over every commutative \(A\)
by the map \(\mathbb Z\to A\); no averaging denominator is used.
These sums are exactly the original \texttt{dc} on fixed vectors.

If \(G\) is a topological group and \(C\) is compact and open,
the finiteness assumption holds: \(G/C\) is discrete, and the
continuous image of the compact set \(C\) under \(c\mapsto cgC\)
is finite. This argument does not require the representation
\(\rho\) to have any topology.
% END OHA1

% BEGIN OHA2
\paragraph{OHA2. Spherical commutativity over every coefficient ring.}
\label{claim:OHA2}
Let \(F\) be one of the original nonarchimedean local fields and
\(C=\operatorname{GL}_n(\mathcal O_F)\). All its spherical
double-coset sums commute on \(C\)-fixed vectors, over every
commutative coefficient ring.

Here is a direct integral proof. Smith reduction over the DVR
gives \(g=u d v\) with \(u,\allowbreak{} v\in C\) and \(d\) diagonal. To obtain
this for a matrix with fractional entries, first clear their
finitely many denominators; the integral Smith form then gives
the asserted form after dividing its diagonal by that common
denominator. Consequently
\[
       g^t=v^t d u^t
          =(v^tu^{-1})g(v^{-1}u^t)\in CgC .
\]
Thus transpose preserves every double coset and reverses
multiplication.

For explicit measure control, on the open set
\(\operatorname{GL}_n(F)\subset M_n(F)\) the measure
\(|\det g|_F^{-n}\,\allowbreak{} dg\), with \(dg\) additive matrix Haar
measure, is both left and right invariant. Indeed multiplication
by \(h\) on the left or right has additive Jacobian
\(|\det h|_F^n\), canceled by the determinant factor.
Transpose preserves this measure. Normalize it by \(C\) having
volume one. For compactly supported bi-\(C\)-invariant integer
functions its convolution is the finite sum in OHA1, so
transpose induces the anti-automorphism
\[
                       (f*h)^t=h^t*f^t .
\]
The previous double-coset calculation gives \(f^t=f\), \(h^t=h\)
and \((f*h)^t=f*h\). Hence \(f*h=h*f\), an identity of integer
functions. OHA1 transfers it to every \(A\)-representation.
In rank zero the group is trivial and the conclusion is immediate.

The diagonal with no uniformizer entries is the identity, and
the one with all \(n\) entries equal to the uniformizer is central.
Their double cosets each have one right coset. Thus the extra
source identities \(\mathrm{sphOp}_0=1\) and
\(\mathrm{sphOp}_n=\rho(\varpi I_n)\) follow as well.
This proves the entire \texttt{sph\allowbreak Comm} interface.
% END OHA2

% BEGIN OHA3
\paragraph{OHA3. The original spherical operators preserve the actual level.}
\label{claim:OHA3}
For \(w\in H.W'\), put \(s_w=H.\mathrm{sec}(w)\) and
\(C_w=\operatorname{GL}_n(\mathcal O_{L_w})\).
The original \texttt{W\_\allowbreak le} gives \(s_w(C_w)\subset U\).
Thus a \(U\)-fixed function is \(C_w\)-fixed under the
representation \(g\mapsto r_a(s_w(g))\).

For \(u\in H.UQ\), \texttt{W\_\allowbreak dec} gives the literal factorization
\[
          u=s_w(c)u',\qquad c\in C_w,\quad u'\in H.UQ,
          \qquad [u',\allowbreak{} s_w(g)]=1\quad\text{for every }g .
\]
OHA1 makes the spherical image \(C_w\)-fixed. The action of
\(u'\) commutes term by term with its finite sum and fixes the
initial function. Hence \(u\) fixes the spherical image.
Every auxiliary-level subgroup also fixes that image, since
\texttt{sections\allowbreak Commute\allowbreak At\allowbreak Split} gives the required commutation.
The stabilizer of a vector is a subgroup, so it contains their
join, exactly \(U=H.UQ\vee\bigvee_q H.\mathrm{sec}(w_q)(K_q)\).
This proves level preservation with the source order of factors.
No unproved normality of that join is used.
% END OHA3

% BEGIN OHA4
\paragraph{OHA4. Stability and pairwise commutation on the actual ordinary module.}
\label{claim:OHA4}
Every \(H.\mathrm{op}(A,\allowbreak{} i)\) preserves \(Y\), and all these
operators commute there.

For spherical stability, the source away action is
\[
       (r_a(u)f)(g)=\tau_A(u)\,\allowbreak{} f(g(u,1)).
\]
It preserves left rational invariance, upper-unipotent invariance,
the prescribed torus character and uniform lower-unipotent
continuity. For the last assertion its finitely many matrix
coefficients multiply differences already in the same ideal
\(\mathfrak m_A^k\). It also preserves every value filtration
and commutes with all \(p\)-side translations, hence with every
\(U_b\). For \(f\in Y\), OST10 supplies at every precision a
period \(r\ge1\) with \(U_z^rf-f\) in that value filtration.
Any finite sum of away translates satisfies the same period
condition. OHA3 gives the remaining level condition, proving
spherical stability in the original \(Y\).

Each \(a_{v,\allowbreak{} j}\) belongs to the positive monoid, by its decreasing
list of exponents and the original product construction.
OST7 and OST20 give its ordinary stability and invertibility.
Such operators commute with the spherical sums term by term.
Two positive operators commute by OST5 and commutativity of the
diagonal monoid.

For two spherical places, the exact \texttt{Wm\_\allowbreak mul\_\allowbreak Wm\_\allowbreak comm}
field has two alternatives. If their section images commute,
their two finite sums commute term by term. Otherwise it gives
an isomorphism \(e:\operatorname{GL}_n(L_{w'})
\simeq\operatorname{GL}_n(L_w)\) preserving the integral compact
groups, with \(s_{w'}=s_w\circ e\). It induces a bijection of
the relevant right-coset sets. OHA1 therefore identifies the
second sum, on fixed vectors, with the double-coset operator
at \(e(a_{w',\allowbreak{} j'})\) in the first group. OHA2 applies to arbitrary
double cosets, so these two sums commute. Applying OHA3 before
the second use of the coset identification checks its fixed-vector
hypothesis at both stages. This proves the full original
\(\mathrm{CommOn}(Y,\allowbreak{} H.\mathrm{op}(A))\).
% END OHA4

% BEGIN OHA5
\paragraph{OHA5. A finite commutative operator algebra over a complete local ring.}
\label{claim:OHA5}
Let \(A\) be a complete Noetherian local ring, let \(Y\) be a
finite \(A\)-module, and let \((T_i)_{i\in I}\) be any commuting
family in \(\operatorname{End}_A(Y)\). The index set may be
infinite. The algebra
\[
                      T=A[T_i:i\in I]\subset\operatorname{End}_A(Y)
\]
is finite over \(A\). If \(Y\ne0\), it has a finite product
decomposition
\[
                 T=\prod_{\mathfrak q}T_{\mathfrak q},
                 \qquad 1=\sum_{\mathfrak q}e_{\mathfrak q},
\]
where each factor is a complete Noetherian local ring, finite
over \(A\), and its maximal ideal lies over \(\mathfrak m_A\).
The factors are the localizations at all maximal ideals of \(T\).

Choose finitely many generators of \(Y\). Evaluation on them
injects \(\operatorname{End}_A(Y)\) into a finite power of \(Y\).
Noetherianity therefore makes this endomorphism module finite,
and its submodule \(T\) is finite as well. The commuting hypothesis
makes \(T\) commutative. Finiteness implies integrality, so every
maximal ideal of \(T\) contracts to \(\mathfrak m_A\).
Now apply the finite-algebra decomposition for a complete
Noetherian local base (Stacks Project, Lemma 10.160.2, Tag 0325).
Its hypotheses have just been checked on this exact \(A,\allowbreak{} T\).
Likewise every finite \(A\)-module used below is complete and
separated, by Lemma 10.97.1, Tag 00MA, since \(A=\widehat A\).

One can see explicitly why the projectors exist. The finite
dimensional algebra \(T/\mathfrak m_AT\) is a product of local
Artinian algebras. A lift \(b\) of any one of its idempotents
is successively corrected by
\[
                 b\longmapsto b-(b^2-b)(2b-1)^{-1}.
\]
The error is squared at every step and \(2b-1\) is a unit,
because its square is \(1+4(b^2-b)\). Completeness supplies
the lift, including in residue characteristic two. Lifts are
unique: if commuting idempotents \(e,\allowbreak{} f\) are congruent modulo
\(\mathfrak m_AT\), the idempotents \(e(1-f)\) and \(f(1-e)\)
are in the Jacobson radical and hence zero. Orthogonality and
the identity sum therefore lift. The resulting factors have
one maximal ideal each. This also explains the actual
idempotents in the displayed decomposition.
% END OHA5

% BEGIN OHA6
\paragraph{OHA6. The original adic definition selects exactly one idempotent block.}
\label{claim:OHA6}
In OHA5 fix any \(a_i\in A\) and use the definition of \(L\)
in \eqref{eq:oha-locpart}. There is at most one maximal ideal
\(\mathfrak q_a\) of \(T\) for which every \(T_i-a_i\) belongs
to \(\mathfrak q_a\). If it exists, then
\[
                            L=e_{\mathfrak q_a}Y .
                            \tag{OHA.1}
\]
If it does not exist, \(L=0\). Thus in both cases the projector
belongs to the original algebra \(T\).

If \(Y=0\), take the zero projector and all assertions are immediate.
Otherwise use the decomposition in OHA5. To prove uniqueness, consider
\(J=\mathfrak m_AT+(T_i-a_i:i\in I)\).
The quotient \(T/J\) is generated over \(A/\mathfrak m_A\)
by the already specified scalar images of its generators.
It is therefore either zero or the field \(A/\mathfrak m_A\).
In the nonzero case \(J\) itself is the unique maximal ideal
in question.

Decompose \(Y=\bigoplus_{\mathfrak q}Y_{\mathfrak q}\) with
\(Y_{\mathfrak q}=e_{\mathfrak q}Y\). In the selected factor,
the maximal ideal is nilpotent modulo \(\mathfrak m_AT_{\mathfrak q_a}\).
Consequently some exponent \(h\ge1\) satisfies
\(\mathfrak q_a^hT_{\mathfrak q_a}
 \subset\mathfrak m_AT_{\mathfrak q_a}\).
For every \(i,\allowbreak{} k\), \((T_i-a_i)^{hk}\) maps
\(Y_{\mathfrak q_a}\) into \(\mathfrak m_A^kY_{\mathfrak q_a}\).
This proves the inclusion from right to left in (OHA.1).

For any other factor choose \(i\) with \(T_i-a_i\notin\mathfrak q\).
This element is a unit in \(T_{\mathfrak q}\), and its inverse
preserves \(\mathfrak m_A^kY_{\mathfrak q}\).
If \(y\in L\), its \(\mathfrak q\)-component therefore belongs
to \(\mathfrak m_A^kY_{\mathfrak q}\) for every \(k\).
Separatedness of this finite \(A\)-module forces it to be zero.
There are only finitely many factors, giving (OHA.1).

This is also the literal source \texttt{loc\allowbreak Part}: for fixed
\(i,\allowbreak{} k\), the inverse images of \(\mathfrak m_A^kY\) under
successive powers of \(T_i-a_i\) are increasing, because
that operator preserves \(\mathfrak m_A^kY\). Their submodule
supremum is consequently their union. Its elementwise statement
is exactly \eqref{eq:oha-locpart}.
Any endomorphism commuting with every \(T_i\) commutes with
the projector, since it commutes with every element of \(T\).
No choice of a semisimple generic fiber has entered this proof.
% END OHA6

% BEGIN OHA7
\paragraph{OHA7. All fields of the original residual-part freeness theorem.}
\label{claim:OHA7}
The residual part \(L\) is finite, stable under the original
operators, has a stable complement in \(Y\), and satisfies
\(\mathrm{locPart}(L,\allowbreak{} T,\allowbreak{} a)=L\). If \(Y\) is projective over \(A\),
then \(L\) is finite free over \(A\).

Use \(L=eY\) and \(Z=(1-e)Y\) from OHA6. They are complementary
finite modules and are stable because \(e\in T\). For every ideal
\(I\subset A\),
\[
                       L\cap IY=IL:
\]
if \(y\) is in the left side, applying \(e\) to a finite expression
for \(y\in IY\) gives \(y=ey\in IL\); the reverse inclusion is
immediate. Powers of \(T_i-a_i\) preserve \(L\), so this equality
at \(I=\mathfrak m_A^k\) changes its residual-part condition
relative to \(Y\) into that relative to \(L\). Hence localization
is idempotent.

A direct summand of a projective module is projective. To recall
the finite local freeness step, choose lifts of a basis of
\(L/\mathfrak m_AL\). Nakayama makes the resulting map
\(A^r\to L\) surjective. Projectivity splits it. The complementary
kernel has zero reduction modulo \(\mathfrak m_A\), and is finite,
so Nakayama makes it zero. This gives a basis of \(L\).
The pairwise commutation relations restrict from \(Y\).
These statements are exactly the five conclusions of
\texttt{Loc\allowbreak Free.\allowbreak loc\allowbreak Free}.
% END OHA7

% BEGIN OHA8
\paragraph{OHA8. The whole original \texttt{hi\allowbreak Comm} endpoint.}
\label{claim:OHA8}
For every original \(K\), the original operators preserve and
commute on \(Y\), and \(Y\) is finite free. For every original
residual function \(a:H.\mathrm{Op}\to O\), they also preserve
and commute on \(H.MK(\theta,\allowbreak{} a,\allowbreak{} K)\), which is finite free.

Indeed \texttt{Is\allowbreak CO} and the finite residue field of the
coefficient DVR give \texttt{Is\allowbreak CNL} on this exact \(A\).
The \texttt{fin} field of \(H.\mathrm{Ax}\) and the
\texttt{cont/\allowbreak free} fields of \(H.\mathrm{IsWt}\) give the
hypotheses of OST19. OHA4 proves the remaining commutation
and stability. Apply OHA7 to the actual operators and the
images of the prescribed \(a_i\). This works when that block
is zero as well. It proves every quantifier of \texttt{hi\allowbreak Comm},
not just the assertion for a selected residual eigensystem.
% END OHA8

% BEGIN OHA9
\paragraph{OHA9. The actual block algebra is complete local with the prescribed residue.}
\label{claim:OHA9}
In OHA5--OHA7 assume \(L\ne0\). Let \(T_L\subset
\operatorname{End}_A(L)\) be generated by the original restricted
operators. Then \(T_L\) is finite over \(A\), complete Noetherian
local, its residue field is exactly \(A/\mathfrak m_A\), and
the map \(A\to T_L\) is local. Each \(T_i-a_i\) belongs to
its maximal ideal. If multiplication by \(b\in A\) is injective
on \(Y\), then its scalar image is a non-zero-divisor on \(T_L\).

The selected factor in OHA6 exists because \(L\ne0\).
Restriction identifies \(eT=T_{\mathfrak q_a}\) with \(T_L\).
It is onto because its images are the original generators.
It is injective because an element \(et\) killing \(eY\)
also kills \((1-e)Y\), and \(T\subset\operatorname{End}_A(Y)\)
acts faithfully. All local and completeness claims now follow
from OHA5. The residue is \(T/\mathfrak q_a=A/\mathfrak m_A\)
by the proof of uniqueness in OHA6; this also proves the
prescribed generator congruences.

For the last assertion, if \(bt=0\) as an endomorphism of \(L\),
then \(b\,\allowbreak{} t(y)=0\) for every \(y\in L\subset Y\).
Injectivity on \(Y\) gives \(t(y)=0\) for all \(y\), hence \(t=0\).
This proves the entire \texttt{Hecke\allowbreak Local.\allowbreak hecke\allowbreak Local} conclusion,
including its version for any faithfully embedded algebra whose
image is \(T_L\): transport along that specified algebra
isomorphism, keeping its scalar map.
% END OHA9

% BEGIN OHA10
\paragraph{OHA10. Inverses already belong to the finite operator algebra.}
\label{claim:OHA10}
Let \(M\) be a finite module over any commutative ring \(A\),
and let \(u\) be an \(A\)-linear automorphism. Then
\(u^{-1}\in A[u]\). In particular every subalgebra containing
\(u\) contains its inverse.

Choose finitely many generators of \(M\) and express \(u^{-1}\)
on them by a matrix over \(A\). Its monic characteristic
polynomial annihilates \(u^{-1}\) on \(M\): this follows by
applying the adjugate identity to that matrix and then the
surjection from the free module on the chosen generators.
If \(M=0\) the conclusion is immediate. Otherwise write this
annihilating polynomial as
\[
                    X^r+c_{r-1}X^{r-1}+\cdots+c_0,\qquad r\ge1 .
\]
Multiplication of its identity at \(u^{-1}\) by \(u^{r-1}\)
gives
\[
            u^{-1}=-c_{r-1}-c_{r-2}u-\cdots-c_0u^{r-1}.
\]
This proof uses no freeness, reducedness or characteristic
assumption on \(A\).

Consequently, if \(L\subset Y\) is finite and stable under
an operator invertible on finite \(Y\), its inverse preserves
\(L\), by the same polynomial formula on \(Y\).
Its restriction is invertible on \(L\), and applying the formula
on \(L\) puts the restricted inverse in the algebra generated
by that restricted operator. This is the precise unit argument
needed for the original Hecke generators.
% END OHA10

% BEGIN OHA11
\paragraph{OHA11. The full original \texttt{hi\allowbreak Alg} endpoint.}
\label{claim:OHA11}
Assume \(H.MK(\theta,\allowbreak{} a,\allowbreak{} K)\ne0\). Its original algebra
\(H.TK(\theta,\allowbreak{} a,\allowbreak{} K)\) satisfies all six conclusions of
\texttt{Hi\allowbreak Alg\allowbreak Unit.\allowbreak hi\allowbreak Alg}:
\begin{enumerate}
\item It is an \(O\)-algebra of class \texttt{Is\allowbreak CO}.
\item Its actual scalar map from \(A\) is finite.
\item If \(A\) is a domain, multiplication by any nonzero
  scalar from \(A\) is injective on this algebra.
\item Every generator differs from the original prescribed
  \(a_i\in O\) by an element of its maximal ideal.
\item Every original positive-side generator \(U_{a_{v,\allowbreak{} j}}\)
  is a unit in this algebra.
\item These very generators generate the entire algebra over \(A\).
\end{enumerate}

OHA8 supplies the finite free ordinary module and residual block.
OHA9 gives the finite complete local algebra, its local scalar
map and the congruences. The composite \(O\to A\to T_L\) is
local. The two residue maps are surjective, so every element
of \(T_L\) is congruent to the image of an element of \(O\).
Together with Noetherianity and completeness, these are
exactly the defining \texttt{Is\allowbreak CO} fields. If \(A\) is a domain,
its nonzero scalars act injectively on the free module \(Y\);
OHA9 gives the stated algebra torsion-freeness. Reducedness or
domain assumptions were not used for any earlier conclusion.

OST20 gives invertibility of every \(U_{a_{v,\allowbreak{} j}}\) on \(Y\).
OHA10 first restricts its inverse to \(L\), and then places that
inverse in \(T_L\). Hence the source \texttt{gen} itself is a
unit; this is stronger than invertibility after passing to a
fraction field. Finally \(T_L\) was defined as the algebra
generated by the restricted operators. The images of
\(\mathrm{gen}_i\) are precisely those restrictions, so the
subalgebra they generate internally is all of \(T_L\).
No field-valued eigensystem or arithmetic point was introduced.
% END OHA11

% BEGIN OHA12
\paragraph{OHA12. Installation on the same arithmetic Hecke setting at every level.}
\label{claim:OHA12}
For the original \(H=\mathrm{hset}(Q,\tau_0)\), FNL16 supplies
all its \texttt{Ax} fields and OTW12 supplies the same character
\(\theta\), fixed before all \(K\), with its full
\texttt{Is\allowbreak Wt}. Therefore OHA8 and OHA11 apply to every original
level and every original residual function \(a\), whenever the
source nonzero-block hypothesis for OHA11 holds.

In particular the original \texttt{hi\allowbreak Comm} and \texttt{hi\allowbreak Alg}
are now available after the original finite-level and ordinary
structure theorems. These statements do not require the
integral \texttt{Type\allowbreak For\allowbreak Alt} producer. They do not supply
nonvanishing of a chosen block, an arithmetic specialization
map, a Galois determinant, dense arithmetic points, or the
four Hida outputs on one ideal. Those are subsequent producer
obligations. The coefficient embedding, level, Hecke
normalizations and character have remained the same throughout.
% END OHA12

\subsection{Sources and scope of theorem reuse}

The original interfaces are \texttt{Hi\allowbreak Comm.\allowbreak lean},
\texttt{Hi\allowbreak Alg\allowbreak Unit.\allowbreak lean}, \texttt{Sph\allowbreak Comm.\allowbreak lean},
\texttt{Loc\allowbreak Idem.\allowbreak lean}, \texttt{Loc\allowbreak Free.\allowbreak lean},
\texttt{Hecke\allowbreak Local.\allowbreak lean}, and \texttt{Inv\allowbreak Mem.\allowbreak lean}, at the common
source pin. The companion provenance gives their full paths and
hashes. The only external algebra facts used as checked published
applications here are finite-module completeness and finite-algebra
local decomposition over a complete Noetherian local ring:
\url{https://stacks.math.columbia.edu/tag/00MA} and
\url{https://stacks.math.columbia.edu/tag/0325}.
They are applied to rings and modules whose Noetherianity,
completeness and finiteness were verified explicitly in OHA5.
The block identification, freeness, residue field, unit and
original operator comparisons were proved above.


\section{One ordinary specialization map for every original level}
\label{sec:specialization-map-v42}

\status{This chapter proves the exact original
\texttt{Sp.\allowbreak abs\allowbreak Spec} endpoint.  It constructs one map on the
entire ambient space of coefficient-valued functions, proves
all its level, basis, positive-operator, and away-sum assertions,
and proves both eigenvalue-integrality assertions.  The arbitrary
local coefficient map is retained, including a noninjective or
nonflat map from a nonreduced source ring.  The target ring has
the original domain hypothesis.  All input weight-representation
fields are retained; WRP1--WRP12 now produce them for the actual
algebraic weights.}

Let $X$ be an original Frame with groups
$\Gamma,\allowbreak{} G_a,\allowbreak{} G_p$, compact-side subgroup $U_b$, and local
Iwahori datum $I$.  Fix complete Noetherian local rings
$A,\allowbreak{} O$ with finite residue fields, with $O$ a domain,
and a local ring homomorphism
\[
 \alpha:A\longrightarrow O.
\]
Let $\Omega$ be algebraically closed and let
$\beta:O\hookrightarrow\Omega$ be the specified injective
homomorphism.  The specified map $A\to\Omega$ is
$\gamma=\beta\alpha$; this is exactly the source scalar-tower
hypothesis.  No injectivity or flatness of $\alpha$ is assumed.

Let $W=(d,\xi,\chi,\allowbreak{} c_0,\ldots)$ be the original
$\mathrm{WtRep}(I,\allowbreak{} O,\Omega)$ with all ten fields.  Fix
\[
 \tau:U_b\longrightarrow\operatorname{GL}_m(A),\qquad
 \theta:T\longrightarrow A^\times,\qquad
 \phi:T\longrightarrow O^\times,
\]
where $\theta$ satisfies the original $I.\mathrm{ContChar}$.
Fix $c\geq c_0$ such that $\phi$ is trivial on $T_c(c)$,
and assume, for every $t\in T$,
\begin{equation}\label{eq:smp-character}
     \gamma(\theta(t))=\beta(\phi(t))\chi(t).
\end{equation}
Put $\tau_O=\alpha\tau$, $\theta_O=\alpha\theta$,
$\tau_\Omega=\gamma\tau=\beta\tau_O$, and
$\phi_\Omega=\beta\phi$, with coefficient maps on matrices
and units understood entrywise.  The source also permits
named $\tau',\phi'$ together with equalities to these maps;
we use those equalities, not replacement representations.

Let $(U_K)_{K\in\mathcal K}$ be the original family of
subgroups of $U_b$.  It is directed toward smaller subgroups:
for $K_1,\allowbreak{} K_2$ some $K_3$ satisfies
$U_{K_3}\subset U_{K_1}\cap U_{K_2}$.
For every $K$, assume the original
$X.\mathrm{FinLevel}(U_K)$ and
$X.\mathrm{Free}(\tau,\theta,\allowbreak{} U_K)$.
The index set is not required to be nonempty.

Write
\[
 M_K=\operatorname{Ord}_A(\tau,\theta,\allowbreak{} U_K),\qquad
 N_K=\operatorname{Ord}_O(\tau_O,\theta_O,\allowbreak{} U_K).
\]
For $e\geq c$ let $C_{K,\allowbreak{} e}$ be the literal classical space
$X.\mathrm{Cl}(\tau_\Omega,\xi,\allowbreak{} U_K,\allowbreak{} e,\phi_\Omega)$,
and let $C^{\mathrm{ord}}_{K,\allowbreak{} e}$ be its literal
$\mathrm{ClOrd}$ subspace, defined by the integral-unit
generalized eigenspaces of the normalized distinguished
operator.  Classical elements are families
$F=(F_i)_{i\in\mathrm{Fin}(m)}$, with
$F_i:G_a\times G_p\to\Omega^{d+1}$.
Their actual untwisted values and control are
\begin{equation}\label{eq:smp-untwisted}
 \widetilde F_i(g)=\xi(g_p)^{-1}F_i(g),\qquad
 \operatorname{ctrl}(F)(g)_i=\widetilde F_i(g)_0.
\end{equation}
For any $O$-submodule $D\subset\Omega$,
$\operatorname{Lat}(D)$ means that every coordinate of
every untwisted value lies in $D$.
Write $\varepsilon_O$ and $\varepsilon_\Omega$ for the
original pointwise maps on functions induced by $\alpha$
and $\gamma$.  Their composites agree with the pointwise
map induced by $\beta$, by the scalar-tower equality.

\subsection{The original coefficient maps and the finite value modules}

% BEGIN SMP1
\paragraph{SMP1. Freeness and ordinary base change on the actual coefficient map.}
\label{claim:SMP1}
The original free-stabilizer condition holds for
$(\tau_O,\theta_O,\allowbreak{} U_K)$ for every $K$, and
$\theta_O$ is a continuous character in the source sense.
There is a canonical $\alpha$-semilinear map
\[
               j_K:M_K\longrightarrow N_K,
                        \qquad j_K(f)=\varepsilon_O(f),
\]
and its $O$-linear extension is an isomorphism
\begin{equation}\label{eq:smp-basechange}
                  O\otimes_A M_K\xrightarrow{\sim}N_K.
\end{equation}
Both modules are finite free over their respective rings.
\begin{proof}
The original free-stabilizer assertion has group-theoretic
centrality and membership clauses, followed by a matrix
identity involving $\tau$ and $\theta$.  The first clauses
are unchanged by a coefficient map.  Applying $\alpha$
to each entry of the identity gives exactly its version
for $\tau_O,\theta_O$, including the same scalar and
inverse characters.  Thus this is the source
\texttt{free\_\allowbreak map} on the actual functions and stabilizers.

Since $\alpha$ is local,
$\alpha(\mathfrak m_A^k)\subset\mathfrak m_O^k$ for
every $k$, by expanding finite sums of products.  Applying
$\alpha$ to the congruence expressing continuity of
$\theta$ proves continuity of $\theta_O$ at each precision.
The original finite-level hypothesis is independent of
coefficients.  OST19 and OST24 now apply with precisely
$A,\allowbreak{} O,\alpha$ and these original characters and levels.
They prove finite freeness and the isomorphism
\eqref{eq:smp-basechange}.  In OST24 its semilinear map
is exactly coefficientwise application of $\alpha$;
there is no subsequent abstract choice of a module map.
That theorem allowed every local coefficient map between
the given CNL rings, so no flatness or injectivity has
been added here.
\end{proof}
% END SMP1

% BEGIN SMP2
\paragraph{SMP2. Finite classical coordinates and the integral ordinary inverse.}
\label{claim:SMP2}
For each $K$ and $e\geq c$, $C_{K,\allowbreak{} e}$ is finite
dimensional over $\Omega$.  Its actual integral value
lattice
\[
     L_{K,\allowbreak{} e}=C_{K,\allowbreak{} e}\cap\operatorname{Lat}(\beta(O))
\]
is finite free over $O$ and spans $C_{K,\allowbreak{} e}$.
Every original normalized positive operator
$T_a=X.\mathrm{UF}(m,\allowbreak{} W,\allowbreak{} a)$ preserves both spaces.
There is a finite free integral ordinary lattice
$S_{K,\allowbreak{} e}\subset L_{K,\allowbreak{} e}$ satisfying
\begin{equation}\label{eq:smp-control-isomorphism}
 \Omega\,\allowbreak{} S_{K,\allowbreak{} e}=C^{\mathrm{ord}}_{K,\allowbreak{} e},\qquad
 q_{K,\allowbreak{} e}:S_{K,\allowbreak{} e}\xrightarrow{\sim}N_K,
       \quad \beta(q_{K,\allowbreak{} e}F)=\operatorname{ctrl}(F).
\end{equation}
In the last formula $\beta$ is applied pointwise.
\begin{proof}
Antitonicity of $T_c$ makes $\phi$ trivial on $T_c(e)$.
Also $e\geq c\geq c_0$.  Equation
\eqref{eq:smp-character} is exactly the character
balance over $O$ after using the scalar tower, and
SMP1 supplies the required free condition over $O$.
Thus every original hypothesis of OCT1--OCT19 holds
at this particular $U_K,\allowbreak{} e$.

OCT3 supplies finitely many original representatives
$x_r\in G_a\times G_p$ for which
\begin{equation}\label{eq:smp-evaluation}
 C_{K,\allowbreak{} e}\xrightarrow{\sim}\Omega^{s m(d+1)},\qquad
            F\longmapsto(\widetilde F_i(x_r)_j)_{r,\allowbreak{} i,\allowbreak{} j}.
\end{equation}
For every $O$-submodule $D\subset\Omega$, membership
in $\operatorname{Lat}(D)$ is equivalent to membership
of these finitely many coordinates in $D$.
This proves the first assertions, including the actual
integral lattice.  OCT6--OCT7 give positive-operator
stability on $L_{K,\allowbreak{} e}$ and $C_{K,\allowbreak{} e}$.  OCT10 and OCT15
give precisely \eqref{eq:smp-control-isomorphism}, with
$S_{K,\allowbreak{} e}$ the image of the integral factorial projector.
Every map displayed here is the restriction of the
original control formula \eqref{eq:smp-untwisted}.
\end{proof}
% END SMP2

% BEGIN SMP3
\paragraph{SMP3. Both original eigenvalue-integrality conclusions.}
\label{claim:SMP3}
Let $F\in C_{K,\allowbreak{} e}$ be nonzero.  If $R F=\mu F$, where
$R$ is either a positive normalized operator $T_a$ or
any finite indexed sum
\[
             A_F=\sum_{q\in Q}\operatorname{raF}_{\tau_\Omega}(u_q),
                    \qquad u_q\in U_b,
\]
then $\mu\in\Omega$ is integral over $O$.
For the away sum no preservation of $C_{K,\allowbreak{} e}$ by
$A_F$ is required.
\begin{proof}
Let $D$ be the $O$-span in $\Omega$ of the finitely
many coordinates in \eqref{eq:smp-evaluation} for $F$.
SMP2 says that all untwisted values of $F$ have their
coordinates in $D$.  The generators are themselves such
values, so $D$ is exactly the module spanned by all the
untwisted coordinates, and is finite over $O$.
It is nonzero: otherwise every untwisted value would
vanish, and invertibility of $\xi(g_p)$ would give $F=0$.

Every $T_a$ preserves $\operatorname{Lat}(D)$ by OCT6:
each normalized summand has untwisted matrix
$\xi(n)(\chi(a)^{-1}\xi(a))$, whose entries belong
to $\beta(O)$ by the original fields of $W$.
For an away translation, the untwisted coordinates are
translated in $G_a$ and multiplied in the $i$ index
by $\tau_\Omega(u_q)$, whose entries lie in $\beta(O)$.
It therefore also preserves $\operatorname{Lat}(D)$,
as does its finite sum, with every repetition retained.
The eigenvalue equation consequently implies
$\mu D\subset D$, since untwisting is $\Omega$-linear.

Here is the required determinant argument without a
freeness assumption on $D$.  Choose finite generators
$d_1,\ldots,\allowbreak{} d_r$ of this nonzero module and write
$\mu d_i=\sum_j a_{ij}d_j$ with $a_{ij}\in O$.
The column $(d_i)_i$ is killed by
$\mu I-\beta((a_{ij})_{ij})$.  Multiplying by its
adjugate shows that
$p(\mu)d_i=0$ for every $i$, where
$p(T)=\det(TI-(a_{ij})_{ij})\in O[T]$ is monic.
Some $d_i$ is nonzero in the field $\Omega$, hence
$p(\mu)=0$.  This is exactly integrality over $O$.
The proof used only the eigenvalue equation in the
ambient function space; it imposed no extra level
condition on the away sum.
\end{proof}
% END SMP3

\subsection{A single map on the entire ambient function space}

% BEGIN SMP4
\paragraph{SMP4. Original inclusions and control uniqueness at common levels.}
\label{claim:SMP4}
If $U_{K'}\subset U_K$ and $e'\geq e\geq c$, then
\[
 C^{\mathrm{ord}}_{K,\allowbreak{} e}\subset C^{\mathrm{ord}}_{K',\allowbreak{} e'}
\]
as subspaces of the same ambient families.  The control
map is injective on each of these ordinary spaces.
For a fixed $K$ the integral lattices $S_{K,\allowbreak{} e}$ and
their control inverses are identical at all $e\geq c$.
\begin{proof}
The defining classical conditions become weaker under
$U_{K'}\subset U_K$ and $N_m(e')\subset N_m(e)$.
Thus $C_{K,\allowbreak{} e}\subset C_{K',\allowbreak{} e'}$ on the actual
functions.  The normalized distinguished operator is
the same ambient operator at both levels.  Intersecting
with each of its integral-unit generalized eigenspaces,
then taking their sum, gives the asserted inclusion
of ordinary spaces.  OCT17 proves injectivity of the
same pointwise control map on each space.  Finally
OCT18 proves equality of the actual integral ordinary
lattices at the admissible depths, and equality of
their inverse control maps.  None of these statements
identifies levels merely by an abstract dimension.
\end{proof}
% END SMP4

% BEGIN SMP5
\paragraph{SMP5. Definition and uniqueness of the all-level specialization map.}
\label{claim:SMP5}
For an arbitrary $A$-valued ambient function $f$, call
a classical ambient family $F$ a good lift if
\begin{equation}\label{eq:smp-good}
 \bigl(\exists K,\exists e\geq c:
                   F\in C^{\mathrm{ord}}_{K,\allowbreak{} e}\bigr)
       \quad\text{and}\quad
             \operatorname{ctrl}(F)=\varepsilon_\Omega(f).
\end{equation}
There is at most one good lift.  Define $\mathrm{sp}(f)$
to be this unique lift if it exists, and the zero family
otherwise.  This is one map on the entire original
$X.\mathrm{Fn}(A,\allowbreak{} m)$, with the source codomain
$\mathrm{Fin}(m)\to X.\mathrm{Fn}(\Omega,\allowbreak{} d+1)$.
For every $f\in M_K$ and every $e\geq c$, it satisfies
\begin{equation}\label{eq:smp-specialization}
 \mathrm{sp}(f)=q_{K,\allowbreak{} e}^{-1}(j_K f)\in S_{K,\allowbreak{} e},\qquad
       \operatorname{ctrl}(\mathrm{sp}(f))=
                                        \varepsilon_\Omega(f).
\end{equation}
\begin{proof}
If $F_1,\allowbreak{} F_2$ are good lifts of the same $f$, choose
their witnessing $K_1,\allowbreak{} e_1,\allowbreak{} K_2,\allowbreak{} e_2$.  Directedness gives
$K_3$ with $U_{K_3}\subset U_{K_1}\cap U_{K_2}$.
Put $e_3=\max(e_1,\allowbreak{} e_2)$.  By SMP4 both families belong
to $C^{\mathrm{ord}}_{K_3,\allowbreak{} e_3}$, and their difference
has zero control.  The injectivity in SMP4 gives
$F_1=F_2$.  This proves uniqueness before making any
choice of a map.

For $f\in M_K$, use the actual base-change vector
$j_K f\in N_K$.  Its inverse under the actual isomorphism
$q_{K,\allowbreak{} e}$ exists by SMP2 and belongs to $S_{K,\allowbreak{} e}$.
Its control is $\beta(j_K f)=\varepsilon_\Omega(f)$
by SMP1 and the scalar tower.  It is therefore a good
lift, and uniqueness identifies it with the one
defined $\mathrm{sp}(f)$.  This proves
\eqref{eq:smp-specialization} simultaneously for all
admissible levels and depths.  When $\mathcal K$ is
empty there are no good lifts and the same definition
is the zero map, with all level-indexed conclusions
vacuous.  In general the definition still applies to
functions outside every $M_K$, a fact needed in SMP9.
\end{proof}
% END SMP5

% BEGIN SMP6
\paragraph{SMP6. The required additivity and semilinearity.}
\label{claim:SMP6}
For $f,\allowbreak{} f'\in M_K$ and $b\in A$,
\[
 \mathrm{sp}(f+f')=\mathrm{sp}(f)+\mathrm{sp}(f'),
 \qquad
 \mathrm{sp}(bf)=\gamma(b)\mathrm{sp}(f).
\]
The same equalities hold for ordinary functions initially
lying at different levels after passing to a common
refinement.
\begin{proof}
By SMP5 the two families on the right in the first
equality belong to $C^{\mathrm{ord}}_{K,\allowbreak{} c}$.
Their sum has control
$\varepsilon_\Omega(f)+\varepsilon_\Omega(f')
=\varepsilon_\Omega(f+f')$, by linearity of control
and of the coefficient map.  It is a good lift of
$f+f'$, so uniqueness in SMP5 proves the equality.
Likewise $\gamma(b)\mathrm{sp}(f)$ belongs to the same
ordinary $\Omega$-space and has control
$\varepsilon_\Omega(bf)$, proving semilinearity.
For distinct original levels choose a subgroup contained
in both.  The definitions of the ordinary module make
each function ordinary at this refinement: its ambient
ordinary factorial condition is unchanged and the
subgroup invariance condition is weaker.  Apply the
proved common-level equalities there.  The map itself
has not been chosen again.
\end{proof}
% END SMP6

% BEGIN SMP7
\paragraph{SMP7. The exact spanning and basis-independence assertions.}
\label{claim:SMP7}
For every $K$ and $e\geq c$,
\begin{equation}\label{eq:smp-span}
      \operatorname{span}_\Omega\mathrm{sp}(M_K)
                             =C^{\mathrm{ord}}_{K,\allowbreak{} e}.
\end{equation}
For every $A$-basis $(b_i)_{i\in J}$ of $M_K$, the
family $(\mathrm{sp}(b_i))_{i\in J}$ is linearly
independent over $\Omega$.  Thus the coefficientwise
extension of this same specialization is an isomorphism
$\Omega\otimes_A M_K\simeq C^{\mathrm{ord}}_{K,\allowbreak{} e}$.
\begin{proof}
By the exact base-change isomorphism SMP1,
$(j_K b_i)_i$ is an $O$-basis of $N_K$.  Under
$q_{K,\allowbreak{} e}^{-1}$ it becomes an $O$-basis of $S_{K,\allowbreak{} e}$;
these are the elements $\mathrm{sp}(b_i)$ by SMP5.
Since the $\Omega$-span of $S_{K,\allowbreak{} e}$ is the exact
ordinary classical space by SMP2, this proves
\eqref{eq:smp-span}.

For independence apply control.  Its values on this
family are the $O$-valued functions $j_K b_i$ followed
by the injective coefficient map $\beta$.
They are $O$-linearly independent.  OCT16 proves that
independent $O$-valued functions remain independent
over $\Omega$: for any finite subfamily, pass to
$\operatorname{Frac}(O)$, choose finitely many evaluation
coordinates giving a nonzero determinant, and use its
nonvanishing under $\beta$.  Therefore the control
values, and hence their preimages
$\mathrm{sp}(b_i)$, are $\Omega$-linearly independent.
This applies to an arbitrary basis indexing set,
since linear dependence has finite support.
Spanning and independence on this basis give the
asserted tensor isomorphism.  They do not assert that
$\mathrm{sp}$ is injective on $A$-valued functions;
the original map $\alpha$ may have a kernel.
\end{proof}
% END SMP7

\subsection{The original normalized and away operators}

% BEGIN SMP8
\paragraph{SMP8. Positive operators commute with the single map.}
\label{claim:SMP8}
For every $K$, $a\in I.P$, and $f\in M_K$,
\[
       \mathrm{sp}(U_{a,\allowbreak{} A}f)=T_a\mathrm{sp}(f),
\]
where $U_{a,\allowbreak{} A}$ is the original coefficient-valued
operator and $T_a=\chi(a)^{-1}U_{a,\Omega}$ is the
original normalized classical operator.
\begin{proof}
OCT18 says that $T_a$ preserves $S_{K,\allowbreak{} c}$ and hence
$C^{\mathrm{ord}}_{K,\allowbreak{} c}$.  OCT4 gives
\[
 \operatorname{ctrl}(T_a\mathrm{sp}(f))
     =U_{a,\Omega}\varepsilon_\Omega(f)
     =\varepsilon_\Omega(U_{a,\allowbreak{} A}f).
\]
For the last equality, use the original finite coset
sum for $U_a$ on the $N$-fixed function $f$.  Apply
$\gamma$ term by term; the coefficient map and the
same finite translations commute.  The independence
of representatives is the original coset assertion
OST5, so this also identifies the original source
operator, not a newly chosen sum.  Thus
$T_a\mathrm{sp}(f)$ is a good lift of $U_{a,\allowbreak{} A}f$.
SMP5 proves the desired equality for the one fixed
map $\mathrm{sp}$.  OST20 also ensures that
$U_{a,\allowbreak{} A}f\in M_K$, as required in the ordinary-module
interpretation of this equality.
\end{proof}
% END SMP8

% BEGIN SMP9
\paragraph{SMP9. The complete original away-sum compatibility.}
\label{claim:SMP9}
Fix $K,\allowbreak{} K'$, a finite indexed family $(u_q)_{q\in Q}$
of elements of $U_b$, and put
\[
 A_f=\sum_{q\in Q}\operatorname{ra}_{\tau}(u_q),
 \qquad A_F=\sum_{q\in Q}
                      \operatorname{raF}_{\tau_\Omega}(u_q).
\]
Assume, exactly as in the source, that $A_F$ maps
every $U_K$-fixed ambient family to a $U_{K'}$-fixed
ambient family for these away actions.  Then for
every $f\in M_K$,
\begin{equation}\label{eq:smp-away}
                   \mathrm{sp}(A_f f)=A_F\mathrm{sp}(f).
\end{equation}
There is no additional assumption that $A_f f$ is
ordinary over $A$ at any level.
\begin{proof}
An away translation preserves the twisted left
$\Gamma$-law: in its defining formula the matrix
$\tau_\Omega(u_q)$ acts on the $i$ index, while
$\xi(\operatorname{glob}(\gamma)_p)$ acts on the
separate $d+1$ coordinate index, and the two actions
commute.  It also commutes with every local right
translation, by the direct product $G_a\times G_p$.
Its finite sum therefore preserves the $N$ and
$N_m(c)$ invariances and the $\phi_\Omega$ torus law.
The assumed fixed-space assertion supplies the
remaining $U_{K'}$ invariance.  Consequently
$A_F C_{K,\allowbreak{} c}\subset C_{K',\allowbreak{} c}$.

OCT19 now gives
$A_F S_{K,\allowbreak{} c}\subset S_{K',\allowbreak{} c}$, using its integral
matrix entries and its commutation with the
distinguished normalized operator.  Thus
$A_F\mathrm{sp}(f)\in C^{\mathrm{ord}}_{K',\allowbreak{} c}$.
The control identity OCT4, summed over the same
indexed list, gives
\[
 \operatorname{ctrl}(A_F\mathrm{sp}(f))
       =A_{f,\Omega}\varepsilon_\Omega(f)
       =\varepsilon_\Omega(A_f f).
\]
The second equality follows by applying $\gamma$
to the original finite matrix products and sums
defining the away actions.  It uses the specified
$\tau_\Omega=\gamma\tau$ and retains all multiplicities.
This exhibits a good lift of the ambient function
$A_f f$, whether or not that function has an
$A$-ordinary-level description.  The definition in
SMP5 is on all ambient functions, so its uniqueness
gives \eqref{eq:smp-away} with exactly the original
domain of quantification.
\end{proof}
% END SMP9

% BEGIN SMP10
\paragraph{SMP10. Installation of every clause of the exact \texttt{abs\allowbreak Spec} endpoint.}
\label{claim:SMP10}
Under precisely the standing hypotheses, the map
$\mathrm{sp}$ of SMP5 and the classical spaces above
satisfy the original theorem
\texttt{Upper\allowbreak Theory.\allowbreak V.\allowbreak Hi.\allowbreak Sp.\allowbreak abs\allowbreak Spec}.
\begin{proof}
The conclusion before the existential map has three
parts at each $K,\allowbreak{} e\geq c$: finite dimensionality of
$C_{K,\allowbreak{} e}$, preservation by every positive normalized
operator together with integrality of its eigenvalues
on nonzero classical eigenvectors, and integrality
for eigenvalues of every finite away sum on nonzero
classical eigenvectors.  SMP2 proves finite
dimensionality and preservation; SMP3 proves both
integrality assertions, at exactly those quantifiers.

The existential map has the original domain and
codomain by SMP5.  Its six additional requirements
are supplied as follows:
\begin{center}
\begin{tabular}{ll}
Original conclusion & Written proof \\
\hline
Additivity on each ordinary module & SMP6 \\
$A$-semilinearity via the specified $A\to\Omega$ & SMP6 \\
Ordinary classical span at every $K,\allowbreak{} e\geq c$ & SMP7 \\
Independence of every original $A$-basis & SMP7 \\
Compatibility with every $a\in I.P$ & SMP8 \\
Full finite away-sum compatibility & SMP9
\end{tabular}
\end{center}
They all concern the same map from SMP5.  Its definition
preceded the choices of a level, precision, basis,
positive operator, or away sum.  The original
$\tau',\phi'$ are identified by the specified
coefficient equalities in the statement.  The
coefficient map $A\to O$ remains arbitrary local;
only $O\to\Omega$ uses the original injectivity and
domain hypotheses.  This accounts for every clause
of \texttt{Free\allowbreak Map.\allowbreak lean}'s \texttt{abs\allowbreak Spec} at the
pinned revision, with no new nonvanishing hypothesis.
\end{proof}
% END SMP10

\subsection{What this exact assembly supplies}

This completes the abstract single-map specialization
supplier on the original functions, including its
eigenvalue-integrality conclusions and its full away-sum
domain.  Its inputs are now the proved OST and OCT
packets and the stated original weight representation;
WRP supplies the latter for every original algebraic
weight after its explicit coefficient extension.
The arithmetic theorem \texttt{Spec\allowbreak Int\allowbreak Iff.\allowbreak spec} also
chooses a coefficient realization of a given arithmetic
point, identifies its finite-part torus character,
and translates the operators and classical spaces in
the original $H$-datum.  Those statements must be
checked in their own quantified arithmetic setting;
they are not inferred merely from the abstract
single-map construction.  No arithmetic density,
same-ideal Hida conclusion, or full RH/GRH endpoint
is asserted here.


\section{The original arithmetic specialization theorem}
\label{sec:arithmetic-specialization-v42}

\status{This chapter assembles the exact original
\texttt{Spec\allowbreak Int\allowbreak Iff.\allowbreak spec} endpoint and all eleven fields of
$H.\mathrm{IsSpec}$.  It keeps the given algebraic isomorphism,
arithmetic point, algebraic representation, coefficient maps,
finite-part characters, normalized Hecke operators, and one
specialization map across every level.  The necessary coefficient
point argument, elementary torus identities, and level-transport
statements are written below.  No density of arithmetic points,
nonvanishing of a Hecke block, or same-ideal Hida construction
is inferred from this specialization theorem.}

Fix the source data $L,\mathcal\Omega,\ell,\allowbreak{} n,\allowbreak{} c,\allowbreak{} H$, with
$L$ a number field, $\mathcal\Omega/L$ the specified algebraic
closure, $\ell$ prime, $c$ the specified coefficient datum,
and $H:\mathrm{HSet}(L,\mathcal\Omega,\ell,\allowbreak{} n,\allowbreak{} c)$ satisfying
all its original $\mathrm{Ax}$ fields.  In particular the
source assumes $\ell\ne2$ and $n\geq1$; these hypotheses
are retained.  Write $O_0=c.O\subset\overline{\mathbb Q}_\ell$.
Let $A$ be an $O_0$-algebra of class $\mathrm{IsCO}$:
it is complete Noetherian local, the structural map is local,
and every residue class of $A$ is represented by an element
of $O_0$.  Fix the actual character $\theta$ with
$H.\mathrm{IsWt}(\theta)$, the specified algebraic field
isomorphism
\[
             \iota:\overline{\mathbb Q}_\ell\simeq\mathbb C,
\]
an $O_0$-algebra map $x:A\to\overline{\mathbb Q}_\ell$,
and an original dominant weight function $\lambda$ such that
$H.\mathrm{Arith}(\theta,\allowbreak{} x,\lambda)$.
Thus some integer $e_0\geq0$ satisfies
\begin{equation}\label{eq:asp-arithmetic}
 \operatorname{finPart}_{v,\allowbreak{} i}(y)=1
 \quad\text{whenever}\quad y\in U_c(F_v,\allowbreak{} e_0),
 \qquad F_v=L_{H.S_l(v)},
\end{equation}
where $y$ is a local unit and the literal finite part is
\begin{equation}\label{eq:asp-finpart}
 \operatorname{finPart}_{v,\allowbreak{} i}(y)=
       \frac{x\bigl(\theta(e_{v,\allowbreak{} i}(y))\bigr)}
       {\operatorname{algPow}(H.S_l(v),
                                   \operatorname{lowW}(\lambda,\allowbreak{} i),\allowbreak{} y)}.
\end{equation}
Here $e_{v,\allowbreak{} i}(y)$ is the original diagonal element supported
at the single local coordinate $(v,\allowbreak{} i)$.  The subgroup
$U_c(F_v,\allowbreak{} r)$ consists of local units congruent to one
modulo $\mathfrak m_{F_v}^r$.  It is the source $\mathrm{Uc}$,
so the original $T_c(e)$ uses $\mathrm{Uc}(e+1)$.

The subgroup at an auxiliary level $K=(K_q)_{q\in H.Q}$
is the original join
\begin{equation}\label{eq:asp-level}
       U_K=H.\mathrm{lev}(K)
          =H.UQ\vee\bigvee_{q\in H.Q}
                            H.\mathrm{sec}(H.w_q(q))(K_q).
\end{equation}
Every $K_q$ is compact open.  All functions, representations,
and operators in this chapter use the actual Frame $H.X$.

\subsection{Capturing the given arithmetic point in one actual coefficient ring}

% BEGIN ASP1
\paragraph{ASP1. Integral elements and units under the specified field isomorphism.}
\label{claim:ASP1}
For every original coefficient datum $b$ and every
$\mu\in\mathbb C$, with the $b.O$-algebra map given by
the restriction of $\iota$,
\begin{align}
 \mu\text{ is integral over }b.O
       &\quad\Longleftrightarrow\quad|\iota^{-1}(\mu)|_\ell\leq1,
       \label{eq:asp-integral}\\
 \mathrm{IsIntUnit}(b.O,\mu)
       &\quad\Longleftrightarrow\quad|\iota^{-1}(\mu)|_\ell=1.
       \label{eq:asp-unit}
\end{align}
\begin{proof}
Every element of $b.O$ has norm at most one: it is integral
over $\mathbb Z_\ell$ by finite generation, and in a monic
equation a root of norm greater than one would give a
uniquely largest term.  The ultrametric inequality forbids
their sum from being zero.  The same largest-term argument
therefore shows that every root of a monic polynomial with
coefficients in $b.O$ has norm at most one.

Conversely, for $z\in\overline{\mathbb Q}_\ell$ with
$|z|_\ell\leq1$, the field $\mathbb Q_\ell(z)$ is finite
over $\mathbb Q_\ell$.  The norm-unit-ball proof of WRP2
shows that its integer ring is a finite
$\mathbb Z_\ell$-module.  Since it contains $z$, the
determinant trick makes $z$ integral over
$\mathbb Z_\ell$, hence over $b.O$ using the same monic
polynomial.  Transporting the polynomial by $\iota$
proves \eqref{eq:asp-integral} for the specified algebra
map.  In particular this proves the norm-integrality
test in the algebraic closure itself, without replacing
it by a completion.

The original $\mathrm{IsIntUnit}$ means
$\mu\ne0$ and integrality of both $\mu$ and $\mu^{-1}$.
For a nonzero $z$, the two norm inequalities are
$|z|_\ell\leq1$ and $|z|_\ell^{-1}\leq1$, equivalent
to $|z|_\ell=1$.  Norm one already excludes zero.
This proves \eqref{eq:asp-unit}.
\end{proof}
% END ASP1

% BEGIN ASP2
\paragraph{ASP2. The given algebra point has finite coefficients and uniform ideal-power decay.}
\label{claim:ASP2}
There is an original coefficient datum $c_1\geq c$
containing $x(A)$, and there is a real number
$0<\eta<1$ such that
\begin{equation}\label{eq:asp-point-bound}
            |x(a)|_\ell\leq\eta^k
                  \quad(a\in\mathfrak m_A^k, k\geq0).
\end{equation}
In particular $x(\mathfrak m_A)$ consists of elements
of norm strictly less than one.
\begin{proof}
First take $a\in\mathfrak m_A$.  If $|x(a)|_\ell\geq1$,
then $x(a)\ne0$ and its inverse is integral over
$\mathbb Z_\ell$ by ASP1.  Multiply a monic equation
for $x(a)^{-1}$ by the necessary power of $x(a)$ to obtain
\[
             1+b_1x(a)+\cdots+b_rx(a)^r=0,
                                  \qquad b_i\in\mathbb Z_\ell.
\]
The corresponding element $1+b_1a+\cdots+b_ra^r$ of
$A$ lies in $1+\mathfrak m_A$, so it is a unit.  A
unital ring map cannot send a unit to zero, a
contradiction.  Thus $|x(a)|_\ell<1$.

For arbitrary $a\in A$, residue surjectivity supplies
$o\in O_0$ such that $a-o\in\mathfrak m_A$.  Since
$x$ fixes the specified image of $O_0$, and those
elements have norm at most one by ASP1, the
ultrametric inequality gives $|x(a)|_\ell\leq1$.
Choose generators $t_1,\ldots,\allowbreak{} t_r$ of the Noetherian
ideal $\mathfrak m_A$, and choose
$0<\eta<1$ exceeding their finitely many image norms.
Every element of $\mathfrak m_A^k$ is a finite sum
of products of $k$ generators with coefficients in
$A$.  The bound just proved on those coefficients
gives \eqref{eq:asp-point-bound}.

Let $E\subset\overline{\mathbb Q}_\ell$ be the
compositum of $\operatorname{Frac}(O_0)$ and the fields
generated by $x(t_1),\ldots,\allowbreak{} x(t_r)$.  It is finite
over $\mathbb Q_\ell$, using GA2 for the first field.
For each $a\in A$ and $N\geq1$ there is a polynomial
$P_N\in O_0[T_1,\ldots,\allowbreak{} T_r]$ such that
$a-P_N(t_1,\ldots,\allowbreak{} t_r)\in\mathfrak m_A^N$.
Indeed start with a residue representative for $N=1$.
Given such a polynomial, express its error as a finite
sum of monomials of degree $N$ in the $t_i$, with
coefficients in $A$.  Replace each coefficient by
its representative in $O_0$ modulo $\mathfrak m_A$.
Adding the resulting polynomial improves the error
to $\mathfrak m_A^{N+1}$.

All $x(P_N(t))$ belong to $E$, and by
\eqref{eq:asp-point-bound} they converge to $x(a)$.
The finite extension $E$ is closed in the normed
algebraic closure by the coordinate compactness
argument of GA2.  Thus $x(a)\in E$, and its norm is
at most one.  The actual ring
$O_E=\{z\in E:|z|_\ell\leq1\}$ is a finite
$\mathbb Z_\ell$ DVR by the explicit proof in WRP2.
It contains $O_0$ and $x(A)$; take $c_1.O=O_E$.
This captures the specified map itself, not an
abstract map into an isomorphic coefficient field.
\end{proof}
% END ASP2

% BEGIN ASP3
\paragraph{ASP3. The original weight representation and a local coefficient map at the point.}
\label{claim:ASP3}
There are $c'\geq c_1\geq c$ and an original
$W:\mathrm{WtRep}(H.X.I,\allowbreak{} c'.O,\mathbb C)$ such that
its matrix representation $\xi$ is explicitly
intertwined with the unchanged $\operatorname{algRep}$
for $(\iota,\lambda)$ and its character satisfies
the original $\operatorname{lowW}$ formula.
Writing $O=c'.O$, the actual coefficient map
\[
                \alpha:A\longrightarrow O,
                        \qquad\alpha(a)=x(a)
\]
is local, and $O\hookrightarrow\mathbb C$ is exactly
$\beta=\iota|_O$.  Both $A$ and $O$ satisfy
the original CNL hypotheses of SMP.
\begin{proof}
Apply WRP12 with the original
$L,\ell,\allowbreak{} n,\allowbreak{} H.V,\allowbreak{} H.S_l,\iota,\lambda$ and the actual
$c_1$ from ASP2.  It constructs the coefficient
extension and all ten weight-representation fields,
with the explicit intertwiner and exact character.
The ring $O$ in that construction is the norm unit
ball of a finite extension, so its maximal ideal
consists precisely of elements of norm less than
one.  ASP2 gives
$\alpha(\mathfrak m_A)\subset\mathfrak m_O$.
Equivalently, if $\alpha(a)$ is a unit, then $a$
is a unit; thus $\alpha$ is the original required
local homomorphism.  Its composite with $\beta$
is exactly $a\mapsto\iota(x(a))$.

WRP2 proves that $O$ is complete Noetherian local
with finite residue field.  The ring $A$ is complete
Noetherian local by $\mathrm{IsCO}$.  The induced
map on residue fields from $O_0$ is injective,
as every unital map of fields is, and surjective
by the given residue representatives.  The
residue field of $O_0$ is finite: $\ell$ is a
nonunit by its norm, so this field is a quotient
of the finite ring $O_0/\ell O_0$.
Hence $A$ has finite residue field as well.
This verifies every coefficient hypothesis
needed to apply SMP to the actual point.
\end{proof}
% END ASP3

\subsection{The finite character and the unchanged classical spaces}

% BEGIN ASP4
\paragraph{ASP4. The actual finite-part torus character and its uniform depth.}
\label{claim:ASP4}
There is a character $\psi:T\to O^\times$ with
$\beta(\psi(t))=\chi(t)$ on the compact torus.
Define the original finite character
\begin{equation}\label{eq:asp-phi}
          \phi(t)=\alpha(\theta(t))\psi(t)^{-1}
                              \quad(t\in T).
\end{equation}
Then
\begin{align}
 \iota(x(\theta(t)))&=\beta(\phi(t))\chi(t),
                                       \label{eq:asp-balance}\\
 \beta(\phi(e_{v,\allowbreak{} i}(y)))&=
                             \iota(\operatorname{finPart}_{v,\allowbreak{} i}(y)).
                                       \label{eq:asp-elementary-phi}
\end{align}
For $c_* =\max(e_0,\allowbreak{} W.c_0)$, $\phi$ is trivial on
$T_c(c_*)$, and
$\operatorname{finPart}_{v,\allowbreak{} i}(y)=1$ for every
$y\in U_c(F_v,\allowbreak{} c_*+1)$.
\begin{proof}
The compact-torus integrality field of $W$ and
its zeroth-row identity give a preimage in $O$
of $\chi(t)=\xi(t)_{00}$.  Injectivity of $\beta$
makes that preimage unique.  The preimage for
$t^{-1}$ multiplies it to one, so it is a unit.
The multiplicativity and identity equations for
$\chi$ then imply the same equations in $O$,
again by injectivity.  This defines $\psi$ as
a character into the actual unit group.  Equation
\eqref{eq:asp-balance} follows from the definition
\eqref{eq:asp-phi}.

WRP7 evaluated at the single-coordinate torus
element says
\[
 \chi(e_{v,\allowbreak{} i}(y))=
       \iota(\operatorname{algPow}(H.S_l(v),
                              \operatorname{lowW}(\lambda,\allowbreak{} i),\allowbreak{} y)).
\]
All other factors in its finite product are one.
The denominator is nonzero, since every local
embedding sends the nonzero unit $y$ to a nonzero
element.  Divide \eqref{eq:asp-balance} by this
identity to obtain \eqref{eq:asp-elementary-phi},
with exactly the source finite part.

Every compact diagonal tuple $t$ is the product
over $v,\allowbreak{} i$ of $e_{v,\allowbreak{} i}(y_{v,\allowbreak{} i})$ for its actual
diagonal entries.  By OST26--OST27,
$t\in T_c(c_*)$ means
$y_{v,\allowbreak{} i}\in U_c(F_v,\allowbreak{} c_*+1)$ for every entry.
Since $c_*+1\geq e_0$, these units belong to
$U_c(F_v,\allowbreak{} e_0)$, by inclusion of ideal powers.
Equation \eqref{eq:asp-arithmetic} and
\eqref{eq:asp-elementary-phi} make $\phi$ one
on each factor, by injectivity of $\beta$.
Its product is therefore one.  The same inclusion
of unit congruence groups proves the final
finite-part assertion.  This is the precise
source depth shift by one, uniform in every
$v,\allowbreak{} i,\allowbreak{} y$.
\end{proof}
% END ASP4

% BEGIN ASP5
\paragraph{ASP5. Exact equality of the two classical-space definitions.}
\label{claim:ASP5}
For every original auxiliary level $K$ and every
depth $e$,
\begin{equation}\label{eq:asp-classical}
 H.\mathrm{ClAt}(\theta,\iota,\allowbreak{} x,\lambda,\xi,\allowbreak{} K,\allowbreak{} e)
    =H.X.\mathrm{Cl}(H.\tau_C(\iota),\xi,\allowbreak{} U_K,\allowbreak{} e,\beta\phi).
\end{equation}
The away representation on the right is exactly
the coefficientwise specialization of $H.\tau_A(A)$.
\begin{proof}
The source $\mathrm{Big}$ condition is the
twisted rational-equivariance condition together
with invariance under $H.UQ$.  Since
$H.UQ\subset U_K$, the separate $U_K$-invariance
condition already supplies this last requirement.
Thus the rational and away conditions on the two
sides are the same.  The upper-unipotent and
depth-$e$ lower-unipotent conditions are literally
the same as well.

The remaining condition on the left prescribes
the eigenvalue
$\iota(\operatorname{finPart}_{v,\allowbreak{} i}(y))^{-1}$
for each elementary compact torus element
$e_{v,\allowbreak{} i}(y)$.  By ASP4 this is
$\beta(\phi(e_{v,\allowbreak{} i}(y)))^{-1}$.
Every element of the compact torus is the finite
product of its elementary diagonal coordinates,
as in ASP4.  Applying the representation to this
product proves that the elementary eigenvalue
conditions imply the eigenvalue condition for
every $t\in T$ with character $(\beta\phi)^{-1}$.
The converse is restriction to those elements.
This proves equality of the actual subspaces.

For the coefficient identity, every entry of
$H.\tau$ belongs to $O_0$.  The map $x$ fixes
the specified $O_0$-structure, so applying
$\iota x$ after $O_0\to A$ is just applying
$\iota$ to that same entry.  This gives
$H.\tau_C(\iota)=(\operatorname{GL}(\iota x))
\circ H.\tau_A(A)$ entrywise, with the original
section maps and away action unchanged.
\end{proof}
% END ASP5

% BEGIN ASP6
\paragraph{ASP6. Exact positive normalizations and equality of ordinary spectra.}
\label{claim:ASP6}
Let $a_{v,\allowbreak{} j}$ denote the original $\mathrm{alV}(v,\allowbreak{} j)$,
with the first $j$ local diagonal entries equal to
the original uniformizer and all others equal to
one.  For every $j\geq0$,
\begin{equation}\label{eq:asp-lowval}
                  \chi(a_{v,\allowbreak{} j})=H.\mathrm{lowVal}(\iota,\lambda,\allowbreak{} v,\allowbreak{} j).
\end{equation}
The original local operator indexed by
$(v,\allowbreak{} j)\in H.V\times\mathrm{Fin}(n)$ is exactly
$T_{a_{v,\allowbreak{} j+1}}$.  The original
$H.\mathrm{uzC}(\iota,\lambda,\allowbreak{} W.d+1)$ is exactly $T_z$.
Consequently, for every $K,\allowbreak{} e$,
\begin{equation}\label{eq:asp-ordinary}
 H.\mathrm{OrdAt}(\theta,\iota,\allowbreak{} x,\lambda,\xi,\allowbreak{} K,\allowbreak{} e)
  =\mathrm{ClOrd}(W,\allowbreak{} H.\tau_C(\iota),\allowbreak{} U_K,\allowbreak{} e,\beta\phi).
\end{equation}
\begin{proof}
In WRP7's original character formula, evaluation
at $a_{v,\allowbreak{} j}$ leaves precisely the factors
$\operatorname{algPow}(H.S_l(v),
\operatorname{lowW}(\lambda,\allowbreak{} i),\varpi_v)$
with $i<j$; every other diagonal entry is one.
This is the literal definition of $H.\mathrm{lowVal}$,
giving \eqref{eq:asp-lowval}.  By definition the
source local operator uses the index $j+1$ and
normalizes its original $U$-operator by this value
inverted.  This is exactly $T_{a_{v,\allowbreak{} j+1}}$.
Each $a_{v,\allowbreak{} j}$ belongs to the original positive
monoid: its local exponent sequence is one for
$i<j$ and zero afterward, hence is decreasing;
all other components have exponent zero.  This
also covers $j\geq n$, when that local sequence
is constant one.

The diagonal identity for the distinguished element is
\begin{equation}\label{eq:asp-z-factor}
                    z=\prod_{v\in H.V}\prod_{j=0}^{n-1}a_{v,\allowbreak{} j}.
\end{equation}
Indeed, at coordinate $i$ exactly $n-1-i$ indices
$j\in\{0,\ldots,\allowbreak{} n-1\}$ satisfy $i<j$, giving
the original exponent $n-1-i$.  All factors are
diagonal and commute.  Apply the character to
\eqref{eq:asp-z-factor} and use
\eqref{eq:asp-lowval}.  The scalar normalizing
$H.\mathrm{uzC}$ is thus exactly $\chi(z)^{-1}$,
with its source product indexed from zero to
$n-1$.  This proves equality of the actual
distinguished operators.

The two definitions in \eqref{eq:asp-ordinary}
sum intersections with maximal generalized
eigenspaces of that same operator.  Their
classical spaces agree by ASP5.  The allowed
eigenvalues on the left satisfy
$|\iota^{-1}(\mu)|_\ell=1$; on the right they
satisfy $\mathrm{IsIntUnit}(O,\mu)$.
These conditions are identical by ASP1.
The sums and all their summands therefore
agree as actual subspaces, proving the claimed
ordinary-space equality.
\end{proof}
% END ASP6

\subsection{The original level family and its exact finite Hecke sums}

% BEGIN ASP7
\paragraph{ASP7. Common refinements and transport by the original auxiliary sections.}
\label{claim:ASP7}
The original family $U_K$ is directed toward smaller
subgroups.  Moreover for every $K$, $q\in H.Q$, and
$g\in\operatorname{GL}_n(L_{H.w_q(q)})$, there is
one original level $K^g$ such that, for every
representation $\rho$ of $H.X.U_b$ over any
coefficient ring,
\begin{equation}\label{eq:asp-level-transport}
 \rho\bigl(H.\mathrm{sec}(H.w_q(q))(g)\bigr)
                     (M^{U_K})\subset M^{U_{K^g}}.
\end{equation}
The choice of $K^g$ depends only on $K,\allowbreak{} q,\allowbreak{} g$, not
on the coefficients or the representation.
\begin{proof}
For $K,\allowbreak{} K'$, take the componentwise intersections
$K''_q=K_q\cap K'_q$.  They are open, and are
compact since an open subgroup is closed and
their intersection is a closed subset of a compact
subgroup.  Thus $K''$ is an original allowed
level.  Monotonicity of images and joins in
\eqref{eq:asp-level} gives
$U_{K''}\subset U_K\cap U_{K'}$.

For transport, replace the $q$ component by
$gK_qg^{-1}$ and leave the other components
unchanged.  Conjugation is a homeomorphism,
so this is again an allowed compact open level.
Put $s=H.\mathrm{sec}(H.w_q(q))(g)$ and let
$v\in M^{U_K}$.  Every element of $H.UQ$
fixes $\rho(s)v$, because it commutes with $s$
by the exact $\mathrm{Q\_\allowbreak UQ}$ field.  At a
component $q'\ne q$, the same assertion follows
from $\mathrm{Q\_\allowbreak Q}$ and invariance of $v$
under the original $K_{q'}$.  At component $q$,
write $h=gkg^{-1}$ with $k\in K_q$.  The
identity $hg=gk$ implies
\[
 \rho(H.\mathrm{sec}(H.w_q(q))(h))\rho(s)v
   =\rho(s)\rho(H.\mathrm{sec}(H.w_q(q))(k))v
   =\rho(s)v.
\]
The stabilizer of a vector under a representation
is a subgroup.  Since it contains every displayed
generator, it contains their exact join $U_{K^g}$.
This proves \eqref{eq:asp-level-transport} for every
representation at once, with the original order
of multiplication.
\end{proof}
% END ASP7

% BEGIN ASP8
\paragraph{ASP8. Original spherical operators are the same finite sums over every coefficient ring.}
\label{claim:ASP8}
For $w\in H.W_p$ and $j\geq0$, there is one
finite indexed family $(u_r)_{r\in R}\subset H.X.U_b$
such that, for every representation $\rho$ of
$H.X.U_b$,
\begin{equation}\label{eq:asp-spherical-sum}
 \mathrm{sphOp}(\rho\circ H.\mathrm{sec}(w),\allowbreak{} j)
                           =\sum_{r\in R}\rho(u_r).
\end{equation}
It is the original integral double-coset sum, with
no averaging factor.  This sum preserves
$M^{U_K}$ for every original $K$ and every $\rho$.
The same family works for both the $A$-valued
and the complex classical away actions.
\begin{proof}
Put $G_w=\operatorname{GL}_n(L_w)$ and
$C_w=\operatorname{GL}_n(\mathcal O_{L_w})$.
The original orbit $C_wa_{w,\allowbreak{} j}C_w/C_w$ is finite:
it is the image of the compact group $C_w$ in
the discrete quotient $G_w/C_w$, as proved in
OHA1.  Take this finite set as $R$, use the
original quotient representatives appearing in
the definition of $\mathrm{dc}$, and set
$u_r=H.\mathrm{sec}(w)(r)$.  Unfolding that
definition gives \eqref{eq:asp-spherical-sum}.
Its representatives and multiplicities are fixed
before $\rho$, its coefficient ring, or $K$
is selected.

For preservation, a $U_K$-fixed vector is
$C_w$-fixed because $\mathrm{W\_\allowbreak le}$ puts
$H.\mathrm{sec}(w)(C_w)$ inside $H.UQ\subset U_K$.
OHA1's finite coset permutation argument makes
its spherical image $C_w$-fixed.  Given
$u\in H.UQ$, the original $\mathrm{W\_\allowbreak dec}$
field writes
\[
 u=H.\mathrm{sec}(w)(k)u',\quad k\in C_w,
 \quad u'\in H.UQ,\quad
       [u',\allowbreak{} H.\mathrm{sec}(w)(g')]=1\ \text{for every }g'.
\]
The action of $u'$ commutes term by term with
the finite sum and fixes the input, so it fixes
the image.  The factor at $k$ fixes the image
by its $C_w$ invariance.  Each auxiliary component
also fixes the image, because the source
$\mathrm{sectionsCommuteAtSplit}$ makes its
action commute with every summand and it fixes
the input.  The stabilizer subgroup therefore
contains the join $U_K$.  This is the same
argument as OHA3, and uses only a group
representation, so it applies equally to both
required coefficient settings and to all their
ambient fixed spaces.
\end{proof}
% END ASP8

\subsection{Installation of the one arithmetic specialization map}

% BEGIN ASP9
\paragraph{ASP9. A single map for the given point and every original level.}
\label{claim:ASP9}
There is one map
\begin{equation}\label{eq:asp-map}
 \mathrm{sp}:H.X.\mathrm{Fn}(A,\allowbreak{} H.m)\longrightarrow
              \bigl(\mathrm{Fin}(H.m)\to
                              H.X.\mathrm{Fn}(\mathbb C,\allowbreak{} W.d+1)\bigr)
\end{equation}
with all the conclusions of SMP10 for the actual
$H$-data, the actual point $x$, the fixed $\iota$,
and the depth $c_*$ from ASP4.
\begin{proof}
Apply SMP10 with the original Frame $H.X$,
source ring $A$, target ring $O=c'.O$, and the
specified maps
$\alpha(a)=x(a)$ and $\beta=\iota|_O$ from ASP3.
Their composite is precisely $\iota x$.
The weight representation is the actual $W$
of ASP3, and $\tau=H.\tau_A(A)$.
The character $\theta$ has the required
continuity and free-stabilizer fields by the
given $H.\mathrm{IsWt}(\theta)$.
ASP4 supplies $\phi$, the character balance,
and its triviality on $T_c(c_*)$, with
$c_*\geq W.c_0$.

The family of subgroups is exactly
$K\mapsto H.\mathrm{lev}(K)$.  It is directed
by ASP7, is finite-level by $H.\mathrm{Ax}.\mathrm{fin}$,
and has the free condition at each $K$ by
$H.\mathrm{IsWt}.\mathrm{free}$.
Finally ASP5 identifies the complex away
representation with the required coefficient
specialization.  Thus every SMP hypothesis
has been verified on the original objects.
Take its map defined by the unique good lift.
It is defined on the whole ambient function
space and is fixed once for this arithmetic
point and weight.  In particular it is not
selected anew at a level, depth, operator,
basis, or ideal precision.
\end{proof}
% END ASP9

% BEGIN ASP10
\paragraph{ASP10. Classical finiteness, stability, and the original norm bounds.}
\label{claim:ASP10}
For every original $K$ and $e\geq c_*$, the
space $H.\mathrm{ClAt}(\theta,\iota,\allowbreak{} x,\lambda,\xi,\allowbreak{} K,\allowbreak{} e)$
is finite dimensional.  Every original operator
$H.\mathrm{opC}(\iota,\lambda,\allowbreak{} W.d+1,\allowbreak{} i)$ preserves it.
For every nonzero vector $F$ in that space and
every eigenvalue equation
\[
             H.\mathrm{opC}(i)F=\mu F,
\]
one has the original bound $|\iota^{-1}(\mu)|_\ell\leq1$.
\begin{proof}
ASP5 identifies the space with the classical
space in SMP.  Its finite dimensionality is
SMP2.  At a local index $(v,\allowbreak{} j)$ the operator
is $T_{a_{v,\allowbreak{} j+1}}$ by ASP6, and SMP2 supplies
preservation.  At an away index $(w,\allowbreak{} j)$,
ASP8 writes it as the finite sum for
$a_{w,\allowbreak{} j+1}$.  That sum preserves the ambient
$U_K$-fixed space by ASP8.  Away translations
also preserve twisted left equivariance and
commute with all local right translations,
so their sum preserves the other classical
conditions; this is the explicit argument of
SMP9.  Hence it preserves the same original
classical space.

For a nonzero eigenvector, SMP3 proves
integrality over $O$ for either its normalized
positive operator or its finite away sum.
The latter conclusion did not require
preservation of the original space, although
that property has just been checked here.
ASP1 converts the integral assertion under
the same fixed $\iota$ to exactly the required
norm bound.  This proves the original
$\mathrm{fin}$, $\mathrm{stab}$, and
$\mathrm{int}$ fields with every $K,\allowbreak{} e,\allowbreak{} i,\allowbreak{} F,\mu$
quantifier retained.
\end{proof}
% END ASP10

% BEGIN ASP11
\paragraph{ASP11. All algebraic and Hecke properties of the actual map.}
\label{claim:ASP11}
For the single map \eqref{eq:asp-map}, the
original $\mathrm{add}$, $\mathrm{smul}$,
$\mathrm{span}$, $\mathrm{indep}$, and
$\mathrm{op}$ fields of $H.\mathrm{IsSpec}$ hold.
Explicitly it is additive on every original
ordinary module and satisfies
\[
                \mathrm{sp}(bf)=\iota(x(b))\mathrm{sp}(f).
\]
Its complex span at every $K,\allowbreak{} e\geq c_*$ is
the literal $H.\mathrm{OrdAt}(K,\allowbreak{} e)$, every
original $A$-basis has complex-linearly-independent
images, and for every original Hecke index $i$,
\begin{equation}\label{eq:asp-hecke-intertwining}
              \mathrm{sp}(H.\mathrm{op}(A,\allowbreak{} i)f)
                  =H.\mathrm{opC}(\iota,\lambda,\allowbreak{} W.d+1,\allowbreak{} i)
                                                   \mathrm{sp}(f)
\end{equation}
whenever $f$ belongs to the original ordinary module
at $K$.
\begin{proof}
SMP6 gives additivity and semilinearity with scalar
$\beta(\alpha(b))=\iota(x(b))$.
SMP7 gives independence of every original
$A$-basis and the exact classical ordinary span
at every level and admissible depth.  ASP6
identifies that ordinary space with the literal
$H.\mathrm{OrdAt}$, including its original
norm-one eigenvalue condition.

For a local index $(v,\allowbreak{} j)$, the two operators
in \eqref{eq:asp-hecke-intertwining} are the
original $U_{a_{v,\allowbreak{} j+1},\allowbreak{} A}$ and
$T_{a_{v,\allowbreak{} j+1}}$ by definition and ASP6.
SMP8 therefore gives the equality.  For an
away index $(w,\allowbreak{} j)$, use ASP8's one fixed
finite list at $w,\allowbreak{} j+1$ for both coefficient
settings.  Its complex sum satisfies the
full fixed-space preservation hypothesis of
SMP9 with $K'=K$, by ASP8.  Applying SMP9
gives exactly the two sides of
\eqref{eq:asp-hecke-intertwining}, with the
original integral spherical normalization.
This verifies every operator index rather
than only the distinguished ordinary operator.
\end{proof}
% END ASP11

% BEGIN ASP12
\paragraph{ASP12. The original auxiliary-group action clause.}
\label{claim:ASP12}
For all original $K,\allowbreak{} K'$, $q\in H.Q$, and
$g\in\operatorname{GL}_n(L_{H.w_q(q)})$,
if $f$ is ordinary at $K$ and its original
away translate by $s=H.\mathrm{sec}(H.w_q(q))(g)$
is ordinary at $K'$, then
\begin{equation}\label{eq:asp-action}
       \mathrm{sp}\bigl(H.X.\operatorname{ra}_{H.\tau_A(A)}(s)f\bigr)
          =H.\operatorname{atW}(\iota,\allowbreak{} W.d+1,\allowbreak{} H.w_q(q))(g)
                                                       \mathrm{sp}(f).
\end{equation}
\begin{proof}
Take the original conjugated level $K^g$ of ASP7.
Applied to the complex ambient away representation,
\eqref{eq:asp-level-transport} says that the singleton
sum consisting of $s$ carries the whole
$U_K$-fixed space into the $U_{K^g}$-fixed space.
SMP9 therefore gives the required equality for
this singleton sum on every ordinary $f$ at $K$.
By its actual definition the complex away action
at $s$ is
$H.\operatorname{atW}(\iota,\allowbreak{} W.d+1,\allowbreak{} H.w_q(q))(g)$,
giving \eqref{eq:asp-action}.
The equality is an equality of ambient functions
for the one map already constructed.  Thus it
also holds with the additional source premise
that the $A$-valued translate is ordinary at the
specified $K'$, whether or not $K'=K^g$.
No level-dependent choice of specialization is
made in using that premise.
\end{proof}
% END ASP12

% BEGIN ASP13
\paragraph{ASP13. The full original \texttt{Spec\allowbreak Int\allowbreak Iff.\allowbreak spec} endpoint.}
\label{claim:ASP13}
For every original input in the setup there
exist $D\in\mathbb N$, a matrix representation
$\xi:\mathrm{GV}(n,\allowbreak{} H.S_l)\to\operatorname{GL}_D(\mathbb C)$,
a depth $c_*$, and the one map $\mathrm{sp}$
on the original ambient functions such that
\[
          H.\mathrm{IsSpec}(\theta,\iota,\allowbreak{} x,\lambda,
                                                \xi,\allowbreak{} c_*,\mathrm{sp})
\]
holds.  This is exactly
$\mathrm{UpperTheory.V.Hi.spec}$ in the pinned
$\mathrm{SpecIntIff.lean}$.
\begin{proof}
Take $D=W.d+1$ and $\xi=W.\xi$ from ASP3,
$c_*=\max(e_0,\allowbreak{} W.c_0)$ from ASP4, and the
single map of ASP9.  All eleven assertion
fields of the original $H.\mathrm{IsSpec}$
are furnished individually as follows:
\begin{center}
\begin{tabular}{ll}
Original field & Written proof \\
\hline
$\mathrm{alg}$ & ASP3, using the exact WRP intertwiner \\
$\mathrm{triv}$ & ASP4, with the original $c_*+1$ \\
$\mathrm{fin}$ & ASP10 \\
$\mathrm{stab}$ & ASP10 \\
$\mathrm{int}$ & ASP10 and the fixed-$\iota$ norm test \\
$\mathrm{add}$ & ASP11 \\
$\mathrm{smul}$ & ASP11, scalar $\iota(x(b))$ \\
$\mathrm{span}$ & ASP11 and ASP6 \\
$\mathrm{indep}$ & ASP11, every original $A$-basis \\
$\mathrm{op}$ & ASP11, every spherical and local index \\
$\mathrm{act}$ & ASP12, every original auxiliary section
\end{tabular}
\end{center}
The algebraic representation, finite character,
complex coefficient map, local uniformizers,
ordinary eigenvalue condition and Hecke
normalizations were identified explicitly in
ASP3--ASP6.  The coefficient-point capture
and its local map were proved for the
specified $x$, not merely for some arithmetic
point with the same residue.  The family of
levels and each finite spherical sum were
matched before applying the single-map
construction.  These observations verify the
literal consumer contract at all its
quantifiers and finish the original endpoint.
\end{proof}
% END ASP13

\subsection{Completed endpoint and remaining global work}

The original generic arithmetic specialization
theorem is now supplied in full human-readable
form, conditional only on its printed original
$H.\mathrm{Ax}$, $H.\mathrm{IsWt}$ and
$H.\mathrm{Arith}$ inputs.  For the actual
reconstruction datum, FNL16 supplies the
original $H.\mathrm{Ax}$, and the earlier
weight-character construction supplies
$H.\mathrm{IsWt}$ where invoked.  This
chapter preserves the individual arithmetic
point and original coefficient ring throughout;
the extension $c\leq c_1\leq c'$ and the
specialized map $A\to c'.O$ are explicit.

Specialization alone does not assert that
there are enough arithmetic points to separate
a Hecke algebra, that a chosen localized
ordinary block is nonzero at an arithmetic
point, or that a Hida construction uses the
same original ideal with its full automorphy
and Galois conditions.  Those are distinct
remaining statements.  No such statement,
and no whole RH or GRH obligation, is retired
by the endpoint just proved.


\section{Finite local weights for retained rows}
\label{sec:retained-rows-v42}

\status{This chapter constructs the original finite-group idempotents,
proves positivity and normalization of their original series weights,
and proves the squarefree, divisor-degree, long-prime-degree, and
root-to-twist collision parts of the local probability estimate.
It also constructs the complete
row labeling from exact additive dimensions and the prescribed good
rows.  The squarefree character identity, the comparison of the two
local characters, the root-injectivity collision bound, and the actual
automorphic row/count identifications remain supplied inputs.  No
whole global obligation is closed.}

All source references are to
\texttt{7cbf7a0c32fca1badf05931879a37c3f8a2b03b4}.
The index $k\in\{0,1\}$ used for the two auxiliary primes is kept
separate from the integer $k_0$ in the finite unipotent subgroup.
The group-theoretic arguments below work over any finite field
$\mathbb F_q$; in the original application $q=g_k$ is prime,
$N=R.r$, $k_0=R.k$, and $s=R.m/2-1$.  The original rank conditions
give $N=k_0+2s+1$ and $k_0\ge1$.

\subsection{The actual finite idempotents and their probability laws}

Put $G=\operatorname{GL}_N(\mathbb F_q)$ and let $U$ be its upper
unitriangular subgroup.  All matrix indices start at zero.  Let
$W=\{s,\ldots,\allowbreak{} s+k_0-1\}$, $c=s+k_0$, and
\[
 J=\{u\in U:u_{a,\allowbreak{} c}=0\quad(a\in W)\}.
 \tag{RRT.1}
\]
For a nontrivial additive character $\psi:\mathbb F_q\to\mathbb C^\times$,
let $\phi_U(u)=\sum_{i=0}^{N-2}u_{i,\allowbreak{} i+1}$.
The function $\phi_J$ sums the entries along each of two ordered
chains: the chain $s,\allowbreak{} s+1,\ldots,\allowbreak{} s+k_0-1$, and the chain
\[
                  0,1,\ldots,\allowbreak{} s-1,\allowbreak{} c,\allowbreak{} c+1,\ldots,\allowbreak{} N-1.
\]
When $s=0$, the second chain is the singleton $c$ and contributes
zero.  This is exactly the source's \texttt{phi\allowbreak J}, including its
truncated sums when $s=0$ or $k_0=1$.  Define in $\mathbb C[G]$
\[
 e_U=|U|^{-1}\sum_{u\in U}\psi(\phi_U(u))^{-1}[u],\qquad
 e_J=|J|^{-1}\sum_{u\in J}\psi(\phi_J(u))^{-1}[u].
 \tag{RRT.2}
\]
These are the literal \texttt{e\allowbreak U} and \texttt{e\allowbreak J}, not conjugate
surrogates.  Write $[1]a$ for the coefficient of the identity in
$a\in\mathbb C[G]$.

% BEGIN RRT1
\paragraph{RRT1. Exact idempotents and identity coefficients.}
\label{claim:RRT1}
Writing $M=N(N-1)/2$, we have
\[
 e_U^2=e_U,\quad e_J^2=e_J,\qquad
 [1]e_U=q^{-M},\quad [1]e_J=q^{k_0-M}>0.
 \tag{RRT.3}
\]
Both idempotents are self-adjoint for the involution
$a^*(g)=\overline{a(g^{-1})}$.

Indeed, for upper unitriangular $u,\allowbreak{} v$ and $a<b$,
\[
 (uv)_{a,\allowbreak{} b}=u_{a,\allowbreak{} b}+v_{a,\allowbreak{} b}
                  +\sum_{a<h<b}u_{a,\allowbreak{} h}v_{h,\allowbreak{} b}.
 \tag{RRT.4}
\]
For $a\in W,\allowbreak{} b=c$, every summand in the last sum is zero since
$h\in W$ and $v_{h,\allowbreak{} c}=0$.  Thus $J$ contains the identity and is
closed under multiplication; as a finite subset of a group, it is a
subgroup.  All entries in the two character sums are adjacent,
except possibly $(s-1,\allowbreak{} c)$.  For that entry the intervening indices
are precisely $W$, so the same calculation makes its cross terms
zero.  Consequently $\phi_U$ and $\phi_J$ are additive under
group multiplication on their respective subgroups.

For any finite subgroup $H$ and character $\lambda:H\to\mathbb C^\times$,
write $e=|H|^{-1}\sum_h\lambda(h)^{-1}[h]$.
The coefficient of $[z]$ in $e^2$ is zero outside
$H$, and for $z\in H$ it is
$|H|^{-2}\sum_{h\in H}\lambda(z)^{-1}=|H|^{-1}\lambda(z)^{-1}$.
This proves idempotence.  A finite-order complex character has
absolute value one, which proves self-adjointness.  Its coefficient
at $1$ is $|H|^{-1}$.  Finally $U$ has $M$ freely chosen entries,
and $J$ has exactly $M-k_0$ freely chosen entries: the $k_0$
forbidden entries in (RRT.1) are distinct and lie above the diagonal.
This gives (RRT.3).  No Green parametrization is used.
\hfill$\square$
% END RRT1

% BEGIN RRT2
\paragraph{RRT2. The original series weights are probabilities.}
\label{claim:RRT2}
Let $E_F$ be the actual central series idempotents
$D.\mathrm{ntypes}.\mathrm{eSer}(k,\allowbreak{} F)$ for a fixed auxiliary prime.
They are mutually orthogonal, vanish unless $F$ is monic of degree
$N$ with nonzero constant coefficient, and sum to $1$ over those
labels.  These properties are explicit fields of the original
\texttt{NTypes}.  For $a=e_U$ or $e_J$, set
\[
             w_a(F)=\frac{\operatorname{Re}[1](aE_F)}{[1]a}.
 \tag{RRT.5}
\]
Then $w_a(F)\ge0$ for every polynomial, it vanishes off the stated
labels, and its sum over all monic degree-$N$ polynomials is $1$.
These assertions apply to the original \texttt{wt\allowbreak Of}; no
\texttt{adm\allowbreak Lq} hypothesis is needed for them.

For the proof, left multiplication by $b\in\mathbb C[G]$ has
diagonal entry $[1]b$ at every vector of the group basis.  Hence
\[
                 \operatorname{tr}L_b=|G|[1]b.
 \tag{RRT.6}
\]
Since $a$ and $E_F$ commute and are idempotent, so is $aE_F$.
Any idempotent linear map is the identity on its image and zero on
its kernel.  Its trace is its rank, a nonnegative integer.
Therefore $[1](aE_F)$ is real and nonnegative.  Its denominator
is positive by RRT1.  The vanishing assertion follows from
$E_F=0$; summing $aE_F$ gives $a$, proving normalization.
In fact
\[
 w_a(F)=\frac{\dim\operatorname{im}L_{aE_F}}
                     {\dim\operatorname{im}L_a}.
 \tag{RRT.7}
\]
The monic coefficient parametrization
$F_c=X^N+\sum_{i<N}c_iX^i$ is a bijection from
$\mathbb F_q^N$ to all monic degree-$N$ polynomials.  Thus
RRT2 proves exactly the finite probability normalization used by
\texttt{Conc\allowbreak NWeight.\allowbreak sum\_\allowbreak wt\allowbreak Of}.
\hfill$\square$
% END RRT2

% BEGIN RRT3
\paragraph{RRT3. The remaining character calculations have precise multiplicity targets.}
\label{claim:RRT3}
For an irreducible complex representation $\tau$ of $G$, let
$d_\tau=\dim\tau$, let $z_\tau$ be its primitive central
idempotent, and write $m_a(\tau)=\dim\operatorname{im}\tau(a)$.
The first identity below holds for any idempotent $a$; the
weight identity is used for $a=e_U,\allowbreak{} e_J$, whose denominators
are positive:
\[
 [1](az_\tau)=\frac{d_\tau m_a(\tau)}{|G|},\qquad
 w_a(F)=\frac{\sum_{\tau\in\mathcal S_F}d_\tau m_a(\tau)}
                                      {|G|[1]a},
 \tag{RRT.8}
\]
where $\mathcal S_F$ is the set of irreducibles on which $E_F$
acts as $1$.  To see the first equality directly, the coefficient
formula
$z_\tau(g)=d_\tau|G|^{-1}\operatorname{tr}\tau(g^{-1})$
gives $[1](az_\tau)=d_\tau|G|^{-1}\operatorname{tr}\tau(a)$.
The trace of the idempotent $\tau(a)$ is $m_a(\tau)$.
Maschke's decomposition, obtained by averaging a projection onto
any invariant subspace, and Schur's lemma show that each central
idempotent is the sum of the $z_\tau$ on which it acts as $1$.
This proves the second equality.

Consequently the actual lower comparison
\[
       (1-q^{-1})w_{e_U}(F)\le w_{e_J}(F)
 \tag{RRT.9}
\]
follows if for each irreducible in that same series
\[
       (q^{k_0}-q^{k_0-1})m_{e_U}(\tau)\le m_{e_J}(\tau).
 \tag{RRT.10}
\]
Indeed, multiply (RRT.10) by $d_\tau$, sum it, and use the
ratio $[1]e_J/[1]e_U=q^{k_0}$ from RRT1.
Likewise, for a squarefree label $F=\prod P$, the two character
facts
\[
 \mathcal S_F=\{\tau_F\},\qquad
 m_{e_U}(\tau_F)=1,\qquad
 d_{\tau_F}=\frac{\prod_{j=1}^N(q^j-1)}{\prod_{P\mid F}(q^{\deg P}-1)}
 \tag{RRT.11}
\]
give the exact squarefree law
\[
                w_{e_U}(F)=\prod_{P\mid F}(q^{\deg P}-1)^{-1}.
 \tag{RRT.12}
\]
Here $|G|=q^M\prod_{j=1}^N(q^j-1)$, by counting ordered
bases, and $[1]e_U=q^{-M}$.  Statements (RRT.10)--(RRT.11)
are explicit remaining character inputs, not conclusions of
idempotence.  The source proves them using the mirabolic trace
calculation, root exchange, and the squarefree Harish--Chandra
series theorem.  The squarefree \emph{upper} comparison for $e_J$
is not used in the probability proof below.
\hfill$\square$
% END RRT3

% BEGIN RRT4
\paragraph{RRT4. One-sided domination controls every product event.}
\label{claim:RRT4}
Let $P_v,\allowbreak{} Q_v$ be probability laws on finite sets, indexed by a
set of size $f$, and suppose $Q_v\ge(1-\delta)P_v$ pointwise,
where $0\le\delta\le1$.  For every event $B$ in their product,
\[
 \left|\bigotimes_v Q_v(B)-\bigotimes_v P_v(B)\right|
       \le1-(1-\delta)^f\le f\delta.
 \tag{RRT.13}
\]
When $\delta>0$, write
$Q_v=(1-\delta)P_v+\delta R_v$; the remainder is nonnegative
and sums to $1$.  Expanding the product gives
$\bigotimes Q_v=(1-\delta)^f\bigotimes P_v+
(1-(1-\delta)^f)R$ for a probability law $R$, except when the
second coefficient is zero, in which case the two laws already
agree.  Their event probabilities differ by at most that
coefficient.  The last inequality follows by induction from
$(1-\delta)^{f+1}\ge1-(f+1)\delta$.  The cases $\delta=0$ and
$f=0$ are immediate.

For the actual local laws, (RRT.9) and RRT2 apply with
$\delta=1/q$, $P_v=w_{e_U}$, and $Q_v$ equal to either
$w_{e_U}$ or $w_{e_J}$.  Thus the cost of transferring any bad
event to the actual sharp-level product law is at most $f/q$.
This proves the source's \texttt{Comparison\allowbreak Transfer}, with a
slightly sharper first bound, from the lower comparison alone.
\hfill$\square$
% END RRT4

\subsection{The finite polynomial estimates}

For a monic squarefree polynomial $F$ with $F(0)\ne0$ put
\[
           w(F)=\prod_{P\mid F}(q^{\deg P}-1)^{-1}.
 \tag{RRT.14}
\]
The product runs over its distinct monic irreducible factors.
Let $\mathcal S_N$ denote these polynomials of degree $N$; put
$\mathcal S_0=\{1\}$.  These are finite sets, so all sums below
are ordinary finite sums.  They agree with the source's
\texttt{sq\allowbreak Set} and \texttt{sq\allowbreak Weight} under coefficient
parametrization.

% BEGIN RRT5
\paragraph{RRT5. Factor degrees are dominated by uniform permutation cycle degrees.}
\label{claim:RRT5}
For any property $\mathcal A$ of multisets of positive integers,
\[
 \sum_{\substack{F\in\mathcal S_N\\
               \mathcal A(\{\deg P:P\mid F\})}}w(F)
 \le\mathbb P_{\sigma\in S_N}
              \bigl(\mathcal A(\text{cycle lengths of }\sigma)\bigr).
 \tag{RRT.15}
\]
To prove this, let $I_d$ count the monic irreducibles of degree
$d$ with nonzero constant coefficient.  Their roots lie in
$\mathbb F_{q^d}^\times$, and each contributes $d$ distinct roots
disjoint from the roots of the other irreducibles.  Hence
$dI_d\le q^d-1$.  This root assertion follows either in the
quotient field $\mathbb F_q[X]/(P)$, which has $q^d$ elements,
or by its embedding in an algebraic closure; finite fields are
perfect and every nonzero element satisfies $x^{q^d-1}=1$.

For prescribed counts $a_d\ge0$ with $\sum_d da_d=N$, unique
factorization gives exactly $\prod_d\binom{I_d}{a_d}$ such
polynomials.  Their total weight is
\[
 \prod_d\frac{\binom{I_d}{a_d}}{(q^d-1)^{a_d}}
       \le\prod_d\frac1{d^{a_d}a_d!}.
 \tag{RRT.16}
\]
The right side is the probability of that cycle type for a
uniform permutation: arranging the $N$ symbols gives $N!$
orders, and one divides by the $d$ rotations of each $d$-cycle
and by the $a_d!$ permutations of equally long cycles.
Summing (RRT.16) over the allowed types proves (RRT.15).
In particular $\sum_{F\in\mathcal S_N}w(F)\le1$.

There is also the useful exact divisor consequence.  If $H$ is
monic squarefree, $H(0)\ne0$, and $\deg H\le N$, then
\[
                 \sum_{\substack{F\in\mathcal S_N\\H\mid F}}w(F)
                         \le w(H).
 \tag{RRT.17}
\]
Indeed $F=HR$ with $R\in\mathcal S_{N-\deg H}$ and
$\gcd(H,\allowbreak{} R)=1$, and $w(F)=w(H)w(R)$.  Dropping the coprimality
condition and using the preceding mass bound proves (RRT.17).
Thus \texttt{Type\allowbreak Weight\allowbreak Exact}, \texttt{Type\allowbreak Set\allowbreak Weight},
\texttt{Sq\allowbreak Mass\allowbreak Le\allowbreak One}, and \texttt{Divisor\allowbreak Bound} all have direct
human proofs here.
\hfill$\square$
% END RRT5

% BEGIN RRT6
\paragraph{RRT6. The squarefree exceptional mass.}
\label{claim:RRT6}
For $N\ge1$,
\[
                  \sum_{F\in\mathcal S_N}w(F)\ge1-2/q.
 \tag{RRT.18}
\]
There are $q^N$ monic degree-$N$ polynomials, and at most
$q^{N-1}$ have zero constant coefficient.  For $N\ge2$ the
number of monic squarefree degree-$N$ polynomials is
$q^N-q^{N-1}$.  Here is an elementary count: every monic
polynomial uniquely has the form $AB^2$ with $A$ squarefree,
by splitting each irreducible-factor exponent into its parity
and twice an integer.  If $S(t)$ is the formal generating series
for monic squarefree polynomials, coefficientwise finite
counting gives
\[
 \frac1{1-qt}=S(t)\frac1{1-qt^2},\qquad
                  S(t)=\frac{1-qt^2}{1-qt}.
\]
Its coefficients are $1,\allowbreak{} q,\allowbreak{} q^2-q,\allowbreak{} q^3-q^2,\ldots$.
Thus $|\mathcal S_N|\ge q^N-2q^{N-1}$; for $N=1$,
$|\mathcal S_1|=q-1$ gives the same lower bound.
Each $F\in\mathcal S_N$ has $w(F)\ge q^{-N}$, because the
sum of its factor degrees is $N$.  This proves (RRT.18).

In particular, any probability law $P$ on monic degree-$N$
polynomials that equals $w$ on $\mathcal S_N$ has
$P(\mathcal S_N^c)\le2/q$.  In $f$ independent coordinates,
the probability that some coordinate is outside $\mathcal S_N$
is at most $2f/q$.  Combining this with RRT4 accounts exactly
for the term $3f/q$ in the original $W_g$.
\hfill$\square$
% END RRT6

% BEGIN RRT7
\paragraph{RRT7. A permutation has an invariant proper size with probability at most $7/8$.}
\label{claim:RRT7}
Fix $1\le a<N$.  For a permutation $\sigma$, let $Z(\sigma)$
be its number of invariant $a$-element subsets.  Let
$Y(\sigma)$ count those invariant subsets containing a fixed
point and having a fixed point in their complement.  Uniformly
over $S_N$,
\[
                \mathbb EZ=1,\qquad \mathbb EY\ge1/4.
 \tag{RRT.19}
\]
For the first identity, any particular $a$-set is preserved by
$a!(N-a)!$ permutations.  For the second, conditional on
preserving that set, the two restrictions are independent uniform
permutations.  A uniform permutation of any nonempty finite set
has a fixed point with probability at least $1/2$.  For sets of
size one this is immediate; otherwise inclusion--exclusion gives
the derangement probability $\sum_{j=0}^n(-1)^j/j!\le1/2$:
after the first three terms, consecutive negative--positive
pairs have nonpositive sum, and a possible remaining term is
negative.  Summing over all $a$-sets proves (RRT.19).

Pointwise we have
\[
                 2\mathbf1_{Z>0}+Y\le2Z.
 \tag{RRT.20}
\]
If $Y=0$, this follows from $Z\ge1$ on the indicated event.
If $Y>0$, exchange one fixed point in a counted set with a fixed
point outside it.  This gives a different invariant $a$-set,
so $Z\ge2$; since $Y\le Z$, (RRT.20) again follows.
Taking expectations yields
$2\mathbb P(Z>0)+1/4\le2$, hence
\[
                       \mathbb P(Z>0)\le7/8.
 \tag{RRT.21}
\]
This proves the full original \texttt{Lemma\allowbreak P}, without a
large-$N$ restriction.
\hfill$\square$
% END RRT7

% BEGIN RRT8
\paragraph{RRT8. No common proper divisor degree, with the source constant.}
\label{claim:RRT8}
For $1\le a<N$, a squarefree $F$ has a divisor of degree $a$
if and only if some subset of its irreducible-factor degrees
sums to $a$.  A permutation has an invariant $a$-set if and
only if a subset of its cycle lengths sums to $a$.
RRT5 and RRT7 therefore give
\[
           \sum_{\substack{F\in\mathcal S_N\\
                      \exists H\mid F:\deg H=a}}w(F)\le7/8.
 \tag{RRT.22}
\]
This is the original \texttt{Local\allowbreak Estimate\allowbreak IV}.  For any
probability law $P$ agreeing with $w$ on $\mathcal S_N$, the
probability that all $f$ independent coordinates are squarefree
and have a divisor of this same degree is at most $(7/8)^f$.
There are $\lfloor N/2\rfloor$ possible degrees in the
source's common-degree obstruction.  The union bound gives
\[
   \mathbb P(\text{squarefree common-degree obstruction})
                          \le\frac N2(7/8)^f.
 \tag{RRT.23}
\]
No assertion about the nonsquarefree part of $P$ is needed for
this event; that part was handled separately in RRT6.
\hfill$\square$
% END RRT8

% BEGIN RRT9
\paragraph{RRT9. An elementary eventual long-prime reciprocal bound.}
\label{claim:RRT9}
There is an absolute $N_0\ge100$ such that for every integer
$N\ge N_0$,
\[
       \sum_{\substack{p\text{ prime}\\N/2<p,\ p+3\le N}}
                           \frac1p\ge\frac1{5\log N}.
 \tag{RRT.24}
\]
Here only an eventual version is needed; the source proves the
stronger statement starting at $N=100$ by additional finite
checks.  The proof below avoids those checks and keeps the same
constant $5$.

First, the primorial bound $\prod_{p\le m}p\le4^m$ follows by
strong induction.  For $m=2n$, the product over $n<p\le2n$
divides $\binom{2n}{n}\le4^n$, and one uses the inductive bound
at $n$.  For $m=2n+1$, the product over
$n+1<p\le2n+1$ divides $\binom{2n+1}{n}\le4^n$; the last
inequality follows because its two equal middle terms already
contribute twice that coefficient to the sum $2^{2n+1}$.
Use the inductive bound at $n+1$.  The initial cases are direct.

Let $b_n$ count the primes in $(n,2n]$.  Legendre's factorial
formula gives
\[
 \frac{4^n}{2n+1}\le\binom{2n}{n}
                 \le (2n)^{\sqrt{2n}}4^{2n/3}(2n)^{b_n}
 \tag{RRT.25}
\]
for all sufficiently large $n$.  To check the upper bound,
the exponent of $p$ in the middle binomial coefficient is a sum
of terms
$\lfloor2n/p^j\rfloor-2\lfloor n/p^j\rfloor\in\{0,1\}$.
Thus the full power contributed by each $p\le\sqrt{2n}$ is at
most $2n$, and there are at most $\sqrt{2n}$ such primes.
Above $\sqrt{2n}$, each exponent is at most one.  The primes up
to $2n/3$ are bounded by the primorial estimate, the primes in
$(2n/3,\allowbreak{} n]$ have exponent zero, and the remaining $b_n$ primes
contribute at most $(2n)^{b_n}$.  The lower bound in (RRT.25)
comes from the largest of the $2n+1$ binomial coefficients.

Taking logarithms gives
\[
 b_n\ge\frac{(n/3)\log4-\log(2n+1)
                           -\sqrt{2n}\log(2n)}{\log(2n)}.
 \tag{RRT.26}
\]
Set $n=\lfloor N/2\rfloor$.  Every prime in $(n,2n]$ exceeds
$N/2$, and excluding the last three possible integers leaves
at least $b_n-3$ primes satisfying $p+3\le N$.
Each contributes at least $1/N$.  The lower bound for
$(b_n-3)\log N/N$ furnished by (RRT.26) tends to
$\log4/6>1/5$.  For instance
$\log t\ge2(t-1)/(t+1)$ for $t\ge1$ follows by differentiating
the difference; applying its strict inequality at $t=2$ gives
$\log4>4/3$.  Since $\log N/\sqrt N\to0$, one absolute
threshold makes the bound at least $1/5$, proving (RRT.24).
All estimates are unconditional elementary integer counts.
\hfill$\square$
% END RRT9

% BEGIN RRT10
\paragraph{RRT10. The long-prime-degree witness appears exponentially often.}
\label{claim:RRT10}
For $N\ge N_0$ and a probability law $P$ agreeing with $w$ on
$\mathcal S_N$,
\[
 P\{F\in\mathcal S_N:\text{no irreducible factor has prime degree }p,
                         \ N/2<p,\ p+3\le N\}
                   \le1-\frac1{5\log N}.
 \tag{RRT.27}
\]
Indeed, the expected number of $p$-cycles in a uniform
permutation is $1/p$: choose the $p$ symbols, arrange them into
one of $(p-1)!$ cycles, and permute the remaining symbols.
Cycles of lengths exceeding $N/2$ cannot coexist.  Hence the
probability of at least one of the specified prime cycles is
exactly the reciprocal sum in (RRT.24).  Apply RRT5 to the
\emph{failure} property, and then RRT9.  This proves the
eventual version of the original \texttt{Local\allowbreak Estimate\allowbreak V}.
Independence across the $f$ places and $1-t\le e^{-t}$ give
\[
 P^{\otimes f}\{\text{every coordinate is squarefree and lacks
                          this witness}\}
                       \le\exp\{-f/(5\log N)\}.
 \tag{RRT.28}
\]
The event includes squarefreeness; other failures were already
assigned their separate cost in RRT6.
\hfill$\square$
% END RRT10

\subsection{The exact remaining probability inputs and the original consumer}

For a nonzero element $z$ of an algebraic closure of a finite
field, $z_{\ell'}$ denotes the unique prime-to-$\ell$ part of
its finite order.  If $\operatorname{ord}(z)=\ell^a m$ with
$(\ell,\allowbreak{} m)=1$, it is the power of $z$ whose exponent is $0$
modulo $\ell^a$ and $1$ modulo $m$.  This characterizes both
source formulas for \texttt{prime\allowbreak To\allowbreak Part}: their exponents have
these same two residues.  All roots considered below are
nonzero, and lie in a finite field, so this comparison applies.
Put
\[
 \mathcal B(q,\allowbreak{} N)=\{q\}\cup
 \{\ell\text{ prime}:\exists\,1\le d\le N,
                           \ q^d<\ell^{4v_\ell(q^d-1)}\}.
 \tag{RRT.29}
\]
For prime $q$ this is exactly either source version of
\texttt{BB}.  The formally finite range loses nothing: a prime
in the displayed second set divides some $q^d-1$, hence is at
most $q^N$.

Here are the two exact collision inequalities on this same
polynomial space.  The first remains an input.  The second
is proved completely in RRT14--RRT16 below:
\begin{align}
 \sum_{\substack{F\in\mathcal S_N\\
   \exists\ell\notin\mathcal B(q,\allowbreak{} N):\\
   z\mapsto z_{\ell'}\text{ is not injective on roots of }F}}
 w(F)&\le\frac{23\log q}{\sqrt q},                 \tag{RRT.30}\\
 \sum_{\substack{F\in\mathcal S_N\\
   \exists\ell\ne q\text{ prime},\ z\text{ root of }F:\\
                         z_{\ell'}=c_{\ell'}}}w(F)
 &\le A_q:=\frac1{4\min(q_1,\allowbreak{} d_0)}
                  +\frac{1+\log N}{q-1}\quad(c\in\mathbb F_q^\times).
                                                        \tag{RRT.31}
\end{align}
In (RRT.30) the quantified $\ell$ is prime.  The exact source
guards for these two inputs include $q\ge2^{12}$,
$q_1,\allowbreak{} d_0$ distinct odd primes, $\max(q_1,\allowbreak{} d_0)\ge16$, and
$16,\allowbreak{} q_1,\allowbreak{} d_0\mid q-1$.  The two statements are respectively
\texttt{Local\allowbreak Estimate\allowbreak II} and \texttt{Local\allowbreak Estimate\allowbreak III}.
RRT5's divisor bound is one of their now proved ingredients.

% BEGIN RRT11
\paragraph{RRT11. Assemble the original $W_g$ on a reduced local contract.}
\label{claim:RRT11}
Let $\iota$ have $f$ elements, let the twist set have $m$
elements, and prescribe $\chi_{v,\allowbreak{} j}\in\mathbb F_q^\times$.
Assume $N\ge N_0$, the guards just stated, the original
series laws (RRT.9), (RRT.12), and the root-injectivity
estimate (RRT.30).  The second collision estimate (RRT.31)
is supplied by RRT16.  Choose each actual place law $Q_v$ to be
$w_{e_J}$ or $w_{e_U}$.  Then the probability of a tuple failing
the source's \texttt{good\allowbreak At} condition is at most
\[
 \frac{3f}{q}+\frac{23f\log q}{\sqrt q}
       +2fmA_q+\frac N2(7/8)^f
                         +\exp\{-f/(5\log N)\}.
 \tag{RRT.32}
\]

To prove this, first use the common base law $P=w_{e_U}$.
RRT2 proves that it is a probability on the actual labels;
(RRT.12) identifies its squarefree weights.  A tuple satisfies
the six clauses of \texttt{Good\allowbreak At} precisely when its
polynomials are valid labels, all are squarefree, their roots
remain distinct after every allowed prime-to-$\ell$ projection,
no root collides under that projection with any
$\chi_{v,\allowbreak{} j}$ or $\chi_{v,\allowbreak{} j}^{-1}$, every degree
$1\le a\le N/2$ is absent as a divisor degree at some place,
and some place supplies the long prime-degree factor of RRT10.
The prime-to-part and exceptional-set identifications above
identify this with the conjunction of the two original
\texttt{Is\allowbreak Good} predicates defining \texttt{good\allowbreak At}.

The failure set is contained in the union of: a nonsquarefree
coordinate; a root-injectivity collision at a coordinate; one
of the two twist collisions at one of the $fm$ pairs; a common
proper divisor degree across all coordinates; and absence of
the long-prime-degree witness at every coordinate.  RRT6,
(RRT.30), (RRT.31), RRT8, and RRT10 bound these probabilities
by the respective terms in (RRT.32), with $2f/q$ in place of
$3f/q$.  Independence is needed only for the last two events;
the other three use the union bound.  RRT4 adds at most $f/q$
when passing to $\bigotimes Q_v$, proving (RRT.32).

For the actual datum take $q=g_k$, $N=R.r$, $f=R.f$,
$m=R.m$, $(q_1,\allowbreak{} d_0)=(\mathrm{host}.q,\mathrm{host}.d)$, and
the original twist table.  The cardinality field
$|\mathrm{PlG}(k)|=R.f$ makes (RRT.32) the literal $W_g$.
RRT2 supplies \emph{all} local nonnegativity assertions in
\texttt{Weight\allowbreak All\allowbreak Rows}.  Any finite set of bad tuples has
weight at most the full bad event, so (RRT.32) supplies its
second assertion as well.  Polynomial tuples that are not
valid labels contribute zero; coefficient parametrization is
injective on the others.  Repeated global row labels do not
affect this local argument.

The field \texttt{All\allowbreak Join\allowbreak A.\allowbreak level} gives $Q_v=w_{e_J}$ at
every place, with one nontrivial additive character per
auxiliary prime.  Its bundle field is not used in this local
weight argument.  Only the lower comparison and squarefree
$e_U$ formula are required from \texttt{adm\allowbreak Lq}; its
squarefree upper bound on $e_J$ is unnecessary for this
consumer.  Finally $R.f\le R.k\le R.r$ and
$x^2\le R.f$ imply $R.r\ge N_0$ after the single absolute
threshold $x\ge\sqrt{N_0}$.  Thus RRT9's eventual version
suffices for the original eventual \texttt{Link\allowbreak IRR}
application, with exactly the same $W_g$ and all choice orders
of GQT5.  Here the rank inequalities come from
\texttt{Params.\allowbreak lean}: $R.f=2R.\mathrm{ord}$,
$2R.\mathrm{ord}\le R.k$, $R.r=R.k+R.m-1$, and $R.m\ge2$.
They are explicitly used again in \texttt{sizes\allowbreak Rec\_\allowbreak at} of
\texttt{Link\allowbreak IRROf\allowbreak Join\allowbreak Ex.\allowbreak lean}.  The abstract numerical
\texttt{Sizes\allowbreak At} conditions alone do not assert $N\to\infty$;
the displayed actual-rank comparison is essential here.
\hfill$\square$
% END RRT11

\subsection{Two algebraic parts of the actual row construction}

% BEGIN RRT12
\paragraph{RRT12. Exact finite tensor cylinders.}
\label{claim:RRT12}
Let the finite local groups be $G_v$, with central orthogonal
idempotents $E_{v,\allowbreak{} t}$ summing to $1$, and nonzero idempotents
$a_v$ having positive identity coefficients.  In
$\mathbb C[\prod_vG_v]\simeq\bigotimes_v\mathbb C[G_v]$ put
\[
 a=\bigotimes_va_v,\qquad
 E_S=\sum_{(t_v)\in S}\bigotimes_vE_{v,\allowbreak{} t_v}.
 \tag{RRT.33}
\]
Then $E_S$ is central and idempotent,
$E_SE_T=E_{S\cap T}$, and $E_{\rm all}=1$.
For an event determined by a subcollection $V_0$ of places,
\[
 \frac{[1](aE_S)}{[1]a}
          =\sum_{(t_v)_{v\in V_0}\in S}
                   \prod_{v\in V_0}w_{a_v}(t_v).
 \tag{RRT.34}
\]
Indeed products of different tuple idempotents vanish and
their sums expand coordinate by coordinate.  Identity
coefficients multiply under the displayed tensor product.
The complementary coordinates sum to $a_v$, so their factors
cancel in the ratio.  RRT2 makes the remaining factors real
and nonnegative.  These steps prove every assertion directly.

For the two auxiliary-prime collections, (RRT.33) is exactly
the finite \texttt{All\allowbreak Rows\allowbreak Of.\allowbreak ES} coefficient formula.  Its
universal event and cylinder value are therefore proved on
the original finite groups.  Passing to the original adelic
function $e_T$, including its Haar factor, still requires the
actual reduction/inflation and convolution comparison
\texttt{Et\allowbreak One}; (RRT.34) does not choose or rescale that Haar
measure.  It supplies the finite algebraic identity inside
that comparison.
\hfill$\square$
% END RRT12

% BEGIN RRT13
\paragraph{RRT13. Construct labels and the retained injection without distinct-label assumptions.}
\label{claim:RRT13}
Let $L$ be the finite set of valid tuple labels, with good
subset $L^+$.  Suppose exact nonnegative integer multiplicities
$d_t$ satisfy $\sum_{t\in L}d_t=D_{\rm all}$ and that the
given retained rows $\sigma\in\{1,\ldots,\allowbreak{} D_{\rm rows}\}$ have
their original good labels $\lambda_\sigma\in L^+$, with
\[
       |\{\sigma:\lambda_\sigma=t\}|=d_t\qquad(t\in L^+).
 \tag{RRT.35}
\]
Then there is a labeling of $D_{\rm all}$ rows and an
injection of the original retained rows preserving every
label, whose image is exactly the good rows, and with exactly
$\sum_{t\in S}d_t$ rows in each set of labels $S$.

Construct the finite set
\[
 T=\{1,\ldots,\allowbreak{} D_{\rm rows}\}\ \sqcup\!
            \coprod_{t\in L\setminus L^+}\{1,\ldots,\allowbreak{} d_t\}.
 \tag{RRT.36}
\]
Keep $\lambda_\sigma$ on the first summand and assign label
$t$ to the $t$-summand on the right.  Summing (RRT.35) over
$L^+$ shows $|T|=D_{\rm all}$.  Choose any enumeration of
$T$ by the full row index set and transport this labeling.
The first inclusion gives the required injection.  Its image
is exactly the good part because every other summand has a
bad label.  Fiber sizes are (RRT.35) on good labels and the
specified $d_t$ on bad labels.  Summing those fibers over $S$
gives the assertion.  This construction includes zero
multiplicities and repeated labels, and does not require
choosing distinct polynomials for distinct rows.

In the original construction take
$d_t=\mathrm{dimX}(\{t\})$.  The hypotheses
\texttt{Rows\allowbreak Add}, \texttt{Rows\allowbreak Good}, \texttt{Et\allowbreak Univ}, and
the original \texttt{Aut\allowbreak Tie} give the exact additivity,
good-fiber equation, and total rank needed above.  Use the
original retained labels
$\lambda_\sigma(k)=D.\mathrm{nd}.\mathrm{lab}(k,
D.\mathrm{what}_2.\mathrm{ft}(\sigma))$.
RRT13 then constructs \texttt{Rows\allowbreak Lab}, with all its fields,
and hence the original \texttt{Label\allowbreak Count\allowbreak Good\allowbreak Injective} once
the same \texttt{all\allowbreak Rows\allowbreak Of} object is installed.  The
representation-theoretic producers of \texttt{Rows\allowbreak Add} and
\texttt{Rows\allowbreak Good} are separate requirements, not consequences
of finite enumeration.
\hfill$\square$
% END RRT13

\subsection{The complete root-to-twist collision estimate}

% BEGIN RRT14
\paragraph{RRT14. Two prime-power sums, with the original constants.}
\label{claim:RRT14}
For a positive integer $n$ write
$S(n)=\sum_{\ell\mid n,\ \ell\text{ prime}}\ell^{v_\ell(n)}$.
If $q\ge2$, $4\mid q-1$, and $\delta\ge1$, then
\[
 S(q^\delta-1)\le\delta(q-1)+\frac{q^\delta-1}{q-1}.
 \tag{RRT.37}
\]
If moreover $p_1,\allowbreak{} p_2$ are distinct odd primes dividing $q-1$,
$16\mid q-1$, and $\max(p_1,\allowbreak{} p_2)\ge16$, then
\[
               \frac{S(q-1)}{q-1}
                         \le\frac1{4\min(p_1,\allowbreak{} p_2)}.
 \tag{RRT.38}
\]
These are exactly \texttt{Prime\allowbreak Power\allowbreak Order\allowbreak Count} and
\texttt{Prime\allowbreak Power\allowbreak Order\allowbreak Share}.

For (RRT.37), first observe that for any finite collection of
integers at least $2$, its sum is at most its product; this
includes the empty collection, whose sum is $0$ and product
is $1$.  For a nonempty collection, induction uses
$u+v\le uv$ for $u,\allowbreak{} v\ge2$.  Put $A=q-1$ and
$B=(q^\delta-1)/A$.  For a prime $\ell\mid A$,
\[
 v_\ell(q^\delta-1)=v_\ell(A)+v_\ell(\delta).
 \tag{RRT.39}
\]
Here is the needed elementary valuation proof.  If
$u=1+\ell^h b$ with $\ell\nmid b$, and either $\ell$ is odd
and $h\ge1$ or $\ell=2$ and $h\ge2$, expansion of
$u^\ell-1$ has first term $\ell^{h+1}b$ and every other
term has valuation strictly greater than $h+1$.  Hence
raising to the power $\ell$ raises the valuation by exactly
one.  Raising to an integer coprime to $\ell$ keeps valuation
$h$, again by the first binomial term.  Splitting $\delta$
into an $\ell$-power times a coprime integer proves (RRT.39).
For $\ell=2$ the hypothesis $4\mid A$ gives $h\ge2$.

The prime-power factors of $q^\delta-1$ at primes dividing
$A$ therefore have product dividing $\delta A$.  Their sum
is at most that product, hence at most $\delta A$.  The other
prime-power factors have product dividing $B$, since these
primes do not divide $A$; their sum is at most $B$.
Adding proves (RRT.37), including empty parts.

For (RRT.38) set $m=\min(p_1,\allowbreak{} p_2)$ and isolate the
prime-power factors $X$ at $2$, $Y$ at the smaller selected
prime, and $Z$ at the larger selected prime.  Then
$X\ge16$, $Y\ge m$, $Z\ge16$, and $Z\ge m$.
Let $P$ be the product of all remaining prime-power factors,
and $T$ their sum.  Thus $P\ge1$, $T\le P$,
$q-1=XYZP$, and $S(q-1)=X+Y+Z+T$.
Each of
\[
 \frac{X}{XYZP},\qquad \frac{Y}{XYZP},\qquad
 \frac{Z}{XYZP},\qquad \frac{T}{XYZP}
\]
is at most $1/(16m)$: use respectively $YZ\ge16m$,
$XZ\ge16m$, $XY\ge16m$, and $T\le P$ together with
$XYZ\ge16m$.  Their sum proves (RRT.38).
\hfill$\square$
% END RRT14

% BEGIN RRT15
\paragraph{RRT15. Count the actual irreducible collision factors.}
\label{claim:RRT15}
Fix $c\in\mathbb F_q^\times$ and $\delta\ge1$, with $q\ge3$.
Let $\mathcal B_\delta(c)$ be the monic irreducible polynomials
of degree $\delta$ that divide at least one polynomial
\[
 X^{\ell^{v_\ell(q^\delta-1)}}
             -c^{\ell^{v_\ell(q^\delta-1)}}
       \quad(\ell\text{ prime},\ \ell\mid q^\delta-1).
 \tag{RRT.40}
\]
Write $b_\delta=|\mathcal B_\delta(c)|$.  Then
\[
                  \delta b_\delta\le S(q^\delta-1).
 \tag{RRT.41}
\]
Every monic irreducible degree-$\delta$ polynomial with a
nonzero root $z$ satisfying $z_{\ell'}=c_{\ell'}$ for some
prime $\ell$ belongs to $\mathcal B_\delta(c)$.

Indeed, $z$ and $c$ lie in $\mathbb F_{q^\delta}^\times$.
The canonical prime-to-$\ell$ map is a homomorphism on this
finite cyclic group, so the quotient $z/c$ has $\ell$-power
order.  Its order also divides $q^\delta-1$.  If that order
is greater than one, it divides
$\ell^{v_\ell(q^\delta-1)}$.  If it is one, use any prime
divisor of $q^\delta-1$, which exists since $q\ge3$.
Thus $z$ is a root of one polynomial in (RRT.40), and its
monic minimal polynomial divides that polynomial.

For each fixed $\ell$, the distinct monic irreducible factors
of degree $\delta$ in (RRT.40) are coprime.  Their product
divides its nonzero polynomial of degree
$\ell^{v_\ell(q^\delta-1)}$.  Consequently their number
times $\delta$ is at most this degree.  Summing these bounds
over $\ell$ proves (RRT.41).  Counting a factor more than
once only enlarges the bound.  Every polynomial in
$\mathcal B_\delta(c)$ has nonzero constant coefficient,
because the polynomials in (RRT.40) do.
\hfill$\square$
% END RRT15

% BEGIN RRT16
\paragraph{RRT16. The full original LocalEstimateIII.}
\label{claim:RRT16}
Assume $N\ge1$, $p_1,\allowbreak{} p_2$ are distinct odd primes,
$16,\allowbreak{} p_1,\allowbreak{} p_2\mid q-1$, and $\max(p_1,\allowbreak{} p_2)\ge16$.
Then for every $c\in\mathbb F_q^\times$,
\[
 \sum_{\substack{F\in\mathcal S_N\\
        \exists\ell\ne q\text{ prime},\ z\text{ root of }F:\\
                                      z_{\ell'}=c_{\ell'}}}w(F)
        \le\frac1{4\min(p_1,\allowbreak{} p_2)}+\frac{1+\log N}{q-1}.
 \tag{RRT.42}
\]
In fact the proof also allows the prime $\ell=q$ whenever
$q$ is prime.  The displayed weaker form is the exact
original \texttt{Local\allowbreak Estimate\allowbreak III}.

If $F$ occurs in this sum, the monic minimal polynomial of
its indicated root is a factor of $F$, has a degree
$1\le\delta\le N$, and belongs to $\mathcal B_\delta(c)$
by RRT15.  RRT5's divisor bound, followed by a union bound
over these factors, therefore bounds the left side by
\[
                  \sum_{\delta=1}^N
                                \frac{b_\delta}{q^\delta-1}.
 \tag{RRT.43}
\]
For $\delta=1$, RRT15 and (RRT.38) bound this summand by
$1/(4\min(p_1,\allowbreak{} p_2))$.  For $\delta\ge2$, RRT15 and
(RRT.37) give
\[
 \frac{b_\delta}{q^\delta-1}
       \le\frac{q-1}{q^\delta-1}+\frac1{\delta(q-1)}
       \le q^{-(\delta-1)}+\frac1{\delta(q-1)}.
 \tag{RRT.44}
\]
The last comparison follows from
$(q-1)q^{\delta-1}\le q^\delta-1$.
Finally
\[
 \sum_{\delta=2}^Nq^{-(\delta-1)}\le\frac1{q-1},
 \qquad \sum_{\delta=2}^N\frac1\delta\le\log N.
\]
The first is the finite geometric sum; the second follows
by bounding $1/\delta$ by
$\int_{\delta-1}^{\delta}dt/t$.  They also hold for $N=1$
with empty sums.  Substitution in (RRT.43) proves (RRT.42).
There is no character-theoretic hypothesis in this proof:
the weight is the explicit squarefree product (RRT.14).
Installing $p_1=q_1,\allowbreak{} p_2=d_0$ proves (RRT.31) and supplies
the entire $2fmA_q$ term in RRT11 once the base probability
law is identified by (RRT.12).
\hfill$\square$
% END RRT16

\subsection{What has advanced, and the remaining frontier}

The nonnegativity and normalization part of the actual local
weights is now proved without Green.  The source's
\texttt{Comparison\allowbreak Transfer}, \texttt{Type\allowbreak Set\allowbreak Weight},
\texttt{Sq\allowbreak Mass\allowbreak Lower}, \texttt{Sq\allowbreak Mass\allowbreak Le\allowbreak One},
\texttt{Divisor\allowbreak Bound}, \texttt{Lemma\allowbreak P}, and
\texttt{Local\allowbreak Estimate\allowbreak III} and \texttt{Local\allowbreak Estimate\allowbreak IV} have
complete finite proofs here.
The eventual \texttt{Local\allowbreak Estimate\allowbreak V} has the original constant
and is enough for the selected family.  This supplies four of
the five displayed terms of $W_g$, including their actual
probabilistic meaning, and the exact finite tensor and label
constructions required upstream of GQT5.

Three local quantitative inputs remain: the original
squarefree $e_U$ character law, the $e_J$ lower comparison,
and the root-injectivity collision inequality (RRT.30).
For the first two, (RRT.10)--(RRT.11) are explicit finite
character targets.  The source's normalized-label bridge
permutes irreducible polynomial labels by degree before
identifying the original series projector with Green's
parametrization; it must be written on the same labels, not
assumed from a name.  Its exact entry point is
\texttt{Admissible\allowbreak Label\allowbreak Weights\allowbreak Of\allowbreak Green\allowbreak Cuspidal\allowbreak Parametrisa}.
For the remaining collision input the next proofs are the
same-factor and cross-factor root-fiber counts and their
weighted divisor sums.  They can now use RRT5 directly.

The global frontier still includes the original adelic
\texttt{Hida\allowbreak Setting}, reduction/inflation and \texttt{Et\allowbreak One},
the actual \texttt{Rows\allowbreak Add}/\texttt{Rows\allowbreak Good} identifications,
and the cylinder counts.  RRT13 is the full finite enumeration
after those identifications; it does not construct an
automorphic spectrum.  GQT1--GQT5 continue from the resulting
original row object and the three sufficient counting tests.
No new proof of RH or GRH, and no completion of
\texttt{GLOBAL\_\allowbreak RETAINED\_\allowbreak DEFECT}, is claimed by this packet.


\section{What the five global estimates must prove, and two sufficient simplifications}
\label{sec:global-quantitative-v042}

\status{This chapter proves the numerical and finite-counting passage to
the original retained-row defect, on the exact stated row inputs. It
also proves that a uniform polynomial spectral error suffices for the
selected final comparison. The actual arithmetic row counts, local
weights, period estimates and spectral estimates remain input-producing
theorems. None of the five whole global obligations is declared closed.}

The source is the pinned revision
\texttt{7cbf7a0c32fca1badf05931879a37c3f8a2b03b4}.
Every statement below concerns the original family and its original
maps. In particular, the letters denoting source spaces, target spaces,
Haar measures, metrics, coefficient embeddings and auxiliary primes
do not stand for independently chosen models. The source correspondence
and full declarations are recorded in the companion
\texttt{contracts.\allowbreak json} and \texttt{source\_\allowbreak manifest.\allowbreak json}.

\subsection{The five estimates as mathematical contracts}

Fix the family, its positive small exponent $\epsilon$, and its height
exponent $C_H$. At scale $x$ let $D$ denote any actual member, and put
\[
 L=x^\epsilon,\qquad b=e^L,\qquad a=xb,\qquad
 B_0=fmbL,\qquad \Gamma=e^{C_\Gamma m^5fL}.
\]
The thresholds in this subsection are chosen before the member, word,
lift, or finite place. A statement described as closed over a class
additionally preserves the source's earlier order of choosing the
small-exponent cutoff and the parameter-bound functions.

\paragraph{Denominators.}
For a lift $j$, let $d_{Z,0}(j)$ be the actual small-level rank and set
\[
 \mathcal E_j=\sum_{\substack{w\in\mathrm{supp}\\\mu(w)\ne\mu_0}}
       \ell_w e_w\eta_j(w),\qquad
 V(D)=\sum_{p\ \mathrm{prime}}\log p
                 \sum_{1\le e<t}\operatorname{shaLen}(D,\allowbreak{} p,\allowbreak{} e).
\]
Here $\ell_w$ is exactly the source's logarithmic place weight, $e_w$
its ramification index, and the finite cohomology lengths are read from
the same arithmetic datum $D.\mathrm{nd}$. Under the original
finite-length and summability guards, the required estimate is
\begin{equation}\label{eq:gqt-denominator}
 \mathcal E_j\le d_{Z,0}(j)
   \left[C\{fmbL^2+(\log x)^2N_c\}+V(D)\right].
\end{equation}
One $C$ works for all late members and all lifts. This is
\texttt{Link\allowbreak N}, including its guards; a bound using a different
cohomology datum is not the same assertion.

\paragraph{Retained rows.}
Write $D_{\rm all}$ and $D_{\rm rows}$ for the actual full and retained
target dimensions. The stored defect satisfies
\[
 D_{\rm rows}=(1-\delta)D_{\rm all}.
\]
For $\epsilon_z=\epsilon/J>0$, the required assertion is
\begin{equation}\label{eq:gqt-defect}
                  0\le\delta\le e^{-x^{\epsilon_z/7}}.
\end{equation}
The source's theorem quantifies over every late member. The proofs
below isolate exactly the counting and probability estimates needed
for this conclusion.

\paragraph{Periods and heights.}
Let $R$ be the common retained target lattice, $N_u$ the actual shifted
source lattice, $B_u:N_u\to R^\vee$ the rational period map, and
$h_u$ its supplied squared-norm height majorant. Besides integrality,
rationality, fullness of $R$, the support cutoff and the exact Bessel
identity, \texttt{Link\allowbreak DIV} requires
\begin{align}
 s_{\rm Frame}+\mu_{\max}(R)&\le x^{C_H}b,
                                                        \label{eq:gqt-height1}\\
 h(B_u)&\le h_u\quad(B_u\ne0),                           \label{eq:gqt-height2}\\
 \mu(R^\vee)+h_u-\mathrm{Fo}
   &\le\mathrm{MAIN}_u+C n_TaL_a
       +Cm(B_0+N_c)x^{C_{\rm abs}\epsilon},               \label{eq:gqt-height3}\\
 \left|\sum_w\mathrm{sz}(w,\allowbreak{} E_{u,\allowbreak{} j,\allowbreak{} w})
                  -A_T\log(g_0g_1)\right|&\le CB_0,     \label{eq:gqt-height4}\\
 \left|(h_u+A_T\log T_u)-(h_v+A_T\log T_v)\right|
       &\le C(n_TaL_a+mB_0).                             \label{eq:gqt-height5}
\end{align}
Here $n_T=\#TT$, $A_T=\sum_{t\in TT}a_t$, and every term on
the right is the original family quantity. In particular
\[
 \begin{aligned}
 \mathrm{MAIN}_u&=
 \sum_{t\in TT}\bigl[m(\phi(a_t)-a_t\log\Omega)-a_t\log T_u\bigr],\\
 \phi(y)&=y\log y-y,\\
 \log\Omega&=\gamma+\log(\pi/v)-L'/L(1).
 \end{aligned}
\]
The Bessel identity uses one target isometry before all words, then a
nonzero scalar and a source isometry for each word. The exported
period-height implication has the precise order
\[
 \mathrm{TuSize}\Longrightarrow\ \mathrm{DGamSize}\
 \Longrightarrow\bigl(\mathrm{LinkEta}\Longrightarrow
                     \mathrm{DIVinN}\Longrightarrow\mathrm{LinkDIV}\bigr).
\]
It does not assert that those upstream inputs have already been
constructed.

\paragraph{Local gains.}
At places with residue characteristic different from $T_u$, the
original local lattices must satisfy the stated valuation inclusion,
\[
 \operatorname{len}(M/N)_{u,\allowbreak{} w}\le d_u^{\rm full}
                                      (\tau_w+\mathrm{io}_w),
\]
and $\tau_w=\mathrm{io}_w=\operatorname{len}(M/N)_{u,\allowbreak{} w}=0$
when there is no served orbit. A served orbit $t_0$ at $p$ must obey
$S_v(p,\allowbreak{} t_0)\le2mN_v(p,\allowbreak{} t_0)a/p$. At the specified signal or
inert-class primes $p\le X''$, one also needs
\begin{gather}
       0\le\mathrm{gain}_{j,\allowbreak{} w}\le4n_Ta/p,                 \label{eq:gqt-gain1}\\
 \sum_j\sum_{t\in TT_w}
  \{\mathrm{vonMangoldtClassSum}_{j,\allowbreak{} t,\allowbreak{} p}-\mathrm{loG}(C,\allowbreak{} b_o,\allowbreak{} p)\}
       -m\,\mathrm{flG}(C,\allowbreak{} p)
   \le\sum_j\mathrm{gain}_{j,\allowbreak{} w}-\tau_w.                 \label{eq:gqt-gain2}
\end{gather}
The same assertion includes served-orbit containment and
\texttt{Mean\allowbreak Kappa\allowbreak Bound}. The full formulas for the two error
envelopes are retained verbatim in the source contract record;
in particular $\mathrm{flG}(C,\allowbreak{} p)=CfN(m)^2\log_p a+
Cf(\log_p a)^2$. There must be a single function $\rho(x)\to0$
with $b_o=b^{\rho(x)}$ before the family index is varied.
The source does not posit separately chosen losses for each prime.

\paragraph{Spectrum.}
For the full period operators set $A_u=c_u^{-1/2}P_u$ using
\[
 c_u=\frac{\operatorname{vol}(N_u^{\rm an})
                 \operatorname{vol}(H_u^{\rm an})}{\dim\Pi_u}.
\]
The original proviso is
\begin{equation}\label{eq:gqt-proviso}
 C_0^{-1}\le d_u\Gamma^2/D_*\le C_0\qquad\text{for every }u.
\end{equation}
Conditional on this proviso, the desired normalized inequalities are
\begin{align}
 \|A_u^*A_v\|_{\mathrm{HS}}^2
     &\le(1+\eta_0)d_ud_v/D_*\quad(u\ne v),             \label{eq:gqt-off}\\
 \|A_u^*A_u-I\|_{\mathrm{HS}}^2&\le e_0d_u,            \label{eq:gqt-diag}\\
 D_{\rm all}&\le(1+\eta_1)D_*.                         \label{eq:gqt-dimension}
\end{align}
The full original \texttt{Link\allowbreak S} asks for $e_0\ge0$ and
$\eta_0,\eta_1,\sqrt{e_0}\le e^{-x^{\epsilon/8}}$.
The errors may depend on the member, but their displayed upper bound
and late threshold must be uniform. GQT6--GQT7 below prove a weaker
sufficient error rate for the selected final comparison. They do not
produce the period operators or prove these estimates.

\subsection{Only three counts are needed for the retained target}

% BEGIN GQT1
\paragraph{GQT1: a finite union estimate with relative counting errors.}
\label{claim:GQT1}
Let $T$ be a finite set of $n$ rows, with two subsets $B_0,\allowbreak{} B_1$ of
bad rows. Assume that the good rows are exactly
$G=T\setminus(B_0\cup B_1)$, and write $g=\#G$.
Suppose $M>0$, $0\le\alpha<1$, and $p_0,\allowbreak{} p_1\ge0$ satisfy
\[
 |n-M|\le\alpha M,\qquad
 |\#B_j-p_jM|\le\alpha p_jM\quad(j=0,1).
\]
Then $n>0$, and for any $W_j\ge p_j$,
\begin{equation}\label{eq:gqt-three-counts}
 0\le1-\frac gn
       \le\frac{1+\alpha}{1-\alpha}(p_0+p_1)
       \le\frac{1+\alpha}{1-\alpha}(W_0+W_1).
\end{equation}
Indeed $n\ge(1-\alpha)M>0$. The complement of $G$ is the union
of the two bad subsets, so $n-g\le\#B_0+\#B_1
\le(1+\alpha)(p_0+p_1)M$. Divide by the lower bound for $n$.
Finally $g\le n$ proves nonnegativity. No independence of the two
bad events is required.

In fact only the lower bound for $n$ and upper bounds for the two
bad counts were used. The opposite halves of the three absolute-value
estimates, and counts for every other subset, may be omitted.
% END GQT1

% BEGIN GQT2
\paragraph{GQT2: these are the original total and bad-cylinder counts.}
\label{claim:GQT2}
Supply one original \texttt{All\allowbreak Rows} object $\mathcal A$ for $D$,
together with \texttt{Label\allowbreak Count\allowbreak Good\allowbreak Injective}. Thus its all-row
labels $\lambda_j(\tau)$ are indexed by
$\tau\in T=\{1,\ldots,\allowbreak{} D_{\rm all}\}$; its injection of the
$D_{\rm rows}$ retained labels has image exactly the labels good
at both auxiliary places. Its dimension function counts these
labels in every cylinder. For $j=0,1$, put
\[
 B_j=\{\tau\in T:\lambda_j(\tau)\text{ is not good}\},
 \qquad S_j=\{\lambda_j(\tau):\tau\in B_j\}.
\]
These are finite sets even though the ambient polynomial-label space
is infinite. A row lies in the cylinder determined by $S_j$ exactly
when it belongs to $B_j$: every label in $S_j$ is bad, and each bad
row contributes its label. Consequently
\[
 \dim\mathcal A(\mathrm{cyl}(j,\allowbreak{} S_j))=\#B_j,
 \quad \dim\mathcal A(\mathrm{cyl}(0,\mathrm{all}))=D_{\rm all},
 \quad\#(T\setminus(B_0\cup B_1))=D_{\rm rows}.
\]
The last identity uses the supplied injection and its exact
good-image assertion, not an assumption about multiplicities of
distinct labels.

Let
\[
 p_j=\sum_{t\in S_j}\prod_v
   \frac{\operatorname{Re}(e_{j,\allowbreak{} v}e_{j,\allowbreak{} t(v)})(1)}
        {\operatorname{Re}e_{j,\allowbreak{} v}(1)}
       =\operatorname{locShare}(\mathcal A,\allowbreak{} j,\allowbreak{} S_j),
 \qquad M=D_*^{\mathcal A}(\mathcal A.B.e).
\]
The notation on the right fixes the source convention even in any
zero-denominator branch. Supply the original nonnegative local
weight assertion, so $p_j\ge0$. The source's finite cylinder
identity states
$\operatorname{Re}e_{\mathrm{cyl}(j,\allowbreak{} S_j)}(1)
=\operatorname{Re}e(1)p_j$. Its definition
$D_*^{\mathcal A}(e)=d_{\rm formal}\operatorname{vol}\,
\operatorname{Re}e(1)$ therefore gives
\[
       D_*^{\mathcal A}(e_{\mathrm{cyl}(j,\allowbreak{} S_j)})=Mp_j.
\]
No change of Haar measure or formal degree has entered this equality.

The original \texttt{Count\allowbreak All\allowbreak Rows} supplies the three relative
estimates in GQT1 and $M>0$. The weight predicate supplies
$p_j\le W_g(f,\allowbreak{} m,\allowbreak{} N,\allowbreak{} g_j,\allowbreak{} q,\allowbreak{} d)$, with $N=D.R.r$ and $g_j=D.host.g(j)$.
The stored identity $D_{\rm rows}=(1-\delta)D_{\rm all}$ now
identifies the defect in GQT1 with the literal $D.\mathrm{what}_2.
\mathrm{delp}$. Thus
\begin{equation}\label{eq:gqt-source-share}
 0\le\delta\le\frac{1+\alpha}{1-\alpha}
       \bigl(W_g(f,\allowbreak{} m,\allowbreak{} N,\allowbreak{} g_0,\allowbreak{} q,\allowbreak{} d)+W_g(f,\allowbreak{} m,\allowbreak{} N,\allowbreak{} g_1,\allowbreak{} q,\allowbreak{} d)\bigr).
\end{equation}
This proof needs the counting estimate only for the full set and the
two finite bad cylinders selected above. It needs the weight bound
only for these two sets. These weaker hypotheses therefore give a
proved alternative input contract for the original defect consumer.
They do not assert that these three arithmetic estimates are known.
% END GQT2

\subsection{The uniform numerical estimate is elementary}

The explicit source envelope used in GQT2 is
\begin{align}
 W_g(f,\allowbreak{} m,\allowbreak{} N,\allowbreak{} g,\allowbreak{} q,\allowbreak{} d)={}&\frac{3f}{g}
       +\frac{23f\log g}{\sqrt g}
       +2fm\left\{\frac1{4\min(q,\allowbreak{} d)}
                         +\frac{1+\log N}{g-1}\right\}\notag\\
       &+\frac N2(7/8)^f
         +\exp\left(-\frac f{5\log N}\right).           \label{eq:gqt-W}
\end{align}

% BEGIN GQT3
\paragraph{GQT3: one threshold for every admissible numerical tuple.}
\label{claim:GQT3}
Fix $0<\vartheta\le1$ and $C_s\in\mathbb R$. For all sufficiently
large $x$, simultaneously for every integer tuple satisfying
\begin{equation}\label{eq:gqt-sizes}
 g\ge e^{8x^\vartheta},\quad x^2\le f\le x^{C_s},\quad
 m\le x^{C_s},\quad100\le N\le x^{C_s},\quad
 \min(q,\allowbreak{} d)\ge e^{x^{\vartheta/6}},
\end{equation}
one has
\begin{equation}\label{eq:gqt-Wbound}
         W_g(f,\allowbreak{} m,\allowbreak{} N,\allowbreak{} g,\allowbreak{} q,\allowbreak{} d)
       \le x^{3C_s+3}e^{-x^{\vartheta/6}}.
\end{equation}
The threshold depends only on $\vartheta,\allowbreak{} C_s$, not the tuple.

If $C_s<2$, there are no admissible tuples for $x>1$. Assume
$C_s\ge2$, and abbreviate $t=x^\vartheta$, $v=x^{\vartheta/6}$.
For late $x$, $g\ge2$ and $\log g\ge8t\ge2$.
The function $z e^{-z/2}$ is decreasing for $z\ge2$, whence
\[
 \frac{\log g}{\sqrt g}\le8t e^{-4t},\qquad
 \frac1{g-1}\le2e^{-8t}.
\]
Use these inequalities, $\log N\le C_s\log x$,
and $(7/8)^f=e^{-f\log(8/7)}$ in \eqref{eq:gqt-W}. This gives
\begin{align*}
 W_g\le{}&3x^{C_s}e^{-8t}
       +184x^{C_s+\vartheta}e^{-4t}
       +\tfrac12x^{2C_s}e^{-v}
       +4x^{2C_s}(1+C_s\log x)e^{-8t}\\
       &+\tfrac12x^{C_s}e^{-\log(8/7)x^2}
       +\exp\left(-\frac{x^2}{5C_s\log x}\right).
\end{align*}
After multiplication by $e^v$, the first, second, fourth, fifth
and sixth terms tend to zero. For the first three of these, use
$v=o(t)$ and $\log x=o(t)$; for the fifth use $v=o(x^2)$;
for the sixth use $v\log x=o(x^2)$, valid since
$\vartheta/6\le1/6$. These statements follow, for example, by
writing $x=e^y$ and comparing a fixed power of $y$ with
$e^{cy}$ for $c>0$, using its power series. They do not conceal
constants depending on $f,\allowbreak{} m,\allowbreak{} N,\allowbreak{} g,\allowbreak{} q,\allowbreak{} d$.

Choose one threshold at which their sum is at most one. Then
$W_g\le(1+x^{2C_s}/2)e^{-v}$, which is at most
$x^{3C_s+3}e^{-v}$ for $x\ge2$. This proves
\eqref{eq:gqt-Wbound} with the source's stated exponent.
% END GQT3

% BEGIN GQT4
\paragraph{GQT4: absorb the counting error and obtain the exact defect.}
\label{claim:GQT4}
Fix also $c>0$ and set
\[
                 \alpha=e^{-cNf\log N}.
\]
Under \eqref{eq:gqt-sizes},
$\alpha\le e^{-100c\log(100)x^2}$, so $0<\alpha\le1/2$
uniformly for late $x$. Therefore
$(1+\alpha)/(1-\alpha)\le3$. Applying GQT3 at both $g_0,\allowbreak{} g_1$
in \eqref{eq:gqt-source-share} gives
\[
          0\le\delta\le6x^{3C_s+3}e^{-x^{\vartheta/6}}.
\]
Finally
\[
 x^{\vartheta/6}-x^{\vartheta/7}
                   -(3C_s+3)\log x-\log6\longrightarrow+\infty,
\]
because $x^{\vartheta/7}=o(x^{\vartheta/6})$ and
$\log x=o(x^{\vartheta/6})$. Enlarge the common threshold once.
Then
\begin{equation}\label{eq:gqt-final-defect}
                         0\le\delta\le e^{-x^{\vartheta/7}}.
\end{equation}
Every threshold in this argument depends only on $\vartheta,\allowbreak{} C_s,\allowbreak{} c$.
It precedes all row data, labels and numerical tuples.
% END GQT4

% BEGIN GQT5
\paragraph{GQT5: the exact conditional family endpoint, and its smaller input.}
\label{claim:GQT5}
Fix $c>0$, $C_s\in\mathbb R$, and the source family $\mathcal F$
with $0<\epsilon\le1$,
$J\ge1$, and $\epsilon_z=\epsilon/J$. Suppose that for every late
member there is one original $\mathcal A:\texttt{All\allowbreak Rows}(D)$ with
\texttt{Label\allowbreak Count\allowbreak Good\allowbreak Injective}, \texttt{Count\allowbreak All\allowbreak Rows} at
$\alpha=e^{-c\,\allowbreak{} D.R.r\,\allowbreak{} D.R.f\log(D.R.r)}$, and
\texttt{Weight\allowbreak All\allowbreak Rows}; also suppose \texttt{Sizes\allowbreak Rec} at
$(x,\epsilon_z,\allowbreak{} C_s)$. Then the literal \texttt{Link\allowbreak IRR} holds.

Indeed $0<\epsilon_z\le1$. Take the maximum of the supplied
family threshold and the thresholds of GQT3--GQT4 for
$\vartheta=\epsilon_z$. At any later $x$, choose the supplied
row object for each member and apply GQT2--GQT4. The resulting
\eqref{eq:gqt-final-defect} is exactly the defect and exponent
in \eqref{eq:gqt-defect}. This proves the mathematical content
of \texttt{Conc\allowbreak N.\allowbreak link\allowbreak IRR\_\allowbreak of\_\allowbreak all\allowbreak Rows}, including positivity,
the original row dimension, and the uniform quantifiers.

The same proof works if the full cylinder-count and weight hypotheses
are replaced by the three inequalities and two finite-set weight
bounds specified in GQT2. As an even weaker sufficient version, one
may take any $0\le\alpha\le1/2$ with those three inequalities:
the particular exponential formula for $\alpha$ is not required
after that uniform bound is established. The source's strict rate
for $\delta$ still follows from the same $W_g$ estimates.

This closes the numerical implication, not the global existence
and counting problem. In particular the original all-row object,
its label/good-image identification, and its local weights are
not inferred from a list of dimensions. The separate source route
\texttt{All\allowbreak Join\allowbreak A}\,$\to$\,\texttt{Weight\allowbreak All\allowbreak Rows} retains its
finite-type and character calculations.
% END GQT5

\subsection{The selected comparison needs a polynomial spectral error}

% BEGIN GQT6
\paragraph{GQT6: a weaker spectral rate gives the full frame conclusion.}
\label{claim:GQT6}
Fix $C_H>0$, $C_0\ge1$, and put $\sigma=x^{-C_H}$. Retain the actual normalized
operators and the structural hypotheses of GFA6--GFA8: the original
proviso \eqref{eq:gqt-proviso}, code cardinality
$\#W\ge C_0\Gamma^2/\sigma$, ramified source inclusions,
\texttt{Link\allowbreak IRR}, and the full \texttt{Link\allowbreak DIV}. Retain also the
full positive-rank source lattices and their original rank, height,
and retained-subspace identifications. In place of
\texttt{Link\allowbreak S}, suppose that
\eqref{eq:gqt-off}--\eqref{eq:gqt-dimension} hold at every late
member with $e_0\ge0$ and
\begin{equation}\label{eq:gqt-poly}
       \eta_0,\eta_1,\sqrt{e_0}\le\sigma/16.
\end{equation}
Negative $\eta_0,\eta_1$ may be replaced by zero. Suppose also
the actual discarded source fraction is at most $2f/b$, with
$f\le x^5$ and $b=e^{x^\epsilon}$ as in GFI15.
Then the same full \texttt{Frame\allowbreak F}$(\mathcal F,\allowbreak{} C_H)$ holds.

Here is the complete error-budget check. Choose one threshold with
\[
 \sigma\le1/20,\qquad 2x^5e^{-x^\epsilon}
                                      \le\sigma^2/256.
\]
The second follows by taking logarithms, since
$x^\epsilon-(5+2C_H)\log x-\log512\to+\infty$.
Thus we may use
\[
      d=\sigma^2/256,\qquad
      \overline\eta=83\sigma/16,
      \overline e=\sigma^2/128.                         \label{eq:gqt-frame-errors}
\]
The proviso gives $d_u>0$. The diagonal error has $e_0<1$,
so the same argument as GFA6 gives $D_{\rm all}>0$; otherwise
the period operators would be zero and the diagonal squared
defect would equal $d_u$. Likewise the nonnegative normalizing
constant $c_u$ cannot be zero. The original row defect is less
than one, so fullness and the retained-dimension identity give
$\operatorname{rank}R>0$. These checks use no exponential
spectral rate. The coefficient and Bessel-map identifications
in GFA7 therefore apply unchanged.

Now $\eta_0,\eta_1,\sqrt{e_0}\le\sigma/2$ and
$d\le1/4$. Apply the proved Gram-to-slope calculation GF8,
or equivalently its explicit specialization GFA8. For every
$H>0$ with $s-\mu_{\min}(R^\vee)\le H$, it gives the actual
frame inequality
\[
 s-\log\#W-
 \frac{4u\log(1+H/u)+(\eta+4\sqrt e)H}{1-\delta}
                         \le\mu(R^\vee),
\]
with $\eta=5\sigma+3\sqrt d$ and $e=2e_0$.
The term $u\log(1+H/u)$ is interpreted as zero at $u=0$,
as in the inherited statement. The chosen errors obey
\[
 0\le\eta\le\overline\eta,\quad
 0\le e\le\overline e\le1/4,
 \qquad
 (\overline\eta+4\sqrt{\overline e})x^{C_H}
       =\frac{83+4\sqrt2}{16}<6.
\]
These are the literal all-height inequality and uniform error
conditions required in \texttt{Frame\allowbreak F}. The errors
$\overline\eta,\overline e$ are chosen before the family member;
the original $\delta$ remains distinct from the discarded source
fraction $d$. This proves the claim.
% END GQT6

% BEGIN GQT7
\paragraph{GQT7: the weaker spectral estimate is sufficient for the final consumer.}
\label{claim:GQT7}
Keep every hypothesis of GFI15 except its original exponential
\texttt{Link\allowbreak S}. Replace that hypothesis by
\eqref{eq:gqt-off}--\eqref{eq:gqt-dimension} and
\eqref{eq:gqt-poly}, with one threshold before all members at
the fixed $C_H$. Then the same numerical conclusion appearing after
the antecedents of \texttt{Link\allowbreak GLOBAL7} follows with
$e_D=C_{\rm abs}$ and $e_L=0$:
\[
 \exists K,\allowbreak{} c_{HH},\allowbreak{} e_K,\allowbreak{} C_L,\allowbreak{} x_0\quad
 \forall x\ge x_0,\ x\in\mathcal F.\mathrm{Win}\ \Longrightarrow\
 \mathrm{ChainTied}(\mathcal F,\allowbreak{} C_{\rm abs},0,\allowbreak{} K,\allowbreak{} c_{HH},\allowbreak{} e_K,\allowbreak{} C_L,\allowbreak{} x).
\]

To check every use, GFI15 first constructs the actual ramified
dimension formula and source-lattice fullness from its supplied
local fields and proviso. Its counting, source-slope,
quotient-slope, local-loss, node and character arguments do not
use \texttt{Link\allowbreak S}. The remaining uses are GFA6's positive
target rank, GFA7's exact normalization and realization, and
GFA8--GFA9's frame inequality and error budget. GQT6 supplies
all of these under the replacement. Thus GFI6--GFI13 receive
the same nine family predicates, the same lattices, the same
height exponent, and the same bound
$(\overline\eta+4\sqrt{\overline e})x^{C_H}\le6$.
Their numerical calculation is unchanged, which proves the
displayed conclusion without assuming the stronger spectral predicate.
The subsequent already
proved window argument retains exponent $1/2+9\epsilon$.

For a closed family-class assertion, the replacement itself must
be supplied with the source's earlier parameter-choice order and
uniformity. A separately chosen bound after fixing each member
would not suffice. For the human proof at fixed $C_H$, any
uniform estimate
\[
 \eta_0,\eta_1,\sqrt{e_0}=O(x^{-C_H-\rho})\qquad(\rho>0)
\]
implies \eqref{eq:gqt-poly} by increasing the threshold.
The extra exponent need not be specified in advance if the
uniform little-oh estimate $o(x^{-C_H})$ is proved.

This is a replacement of the spectral \emph{input to the selected
consumer}. It does not prove the original stronger
\texttt{Link\allowbreak S}, does not identify every other use of that
predicate in the full Lean development, and does not eliminate
the actual spectral-geometry work. The unaltered global
obligation remains open until either the original estimates or
this sufficient alternative have been constructed on the common
family with all the stated normalizations.
% END GQT7

\subsection{Reconstruction tasks exposed by the exact contracts}

For the denominator estimate, the next substantive target is the
actual local scalar table and its common choice of bounds. The
source passes from \texttt{Norm\allowbreak Choice\allowbreak Facts}, the uniform
\texttt{Extra} bounds and the two/three-place bounds to
\texttt{Concl\allowbreak F}, then to \texttt{Adm\allowbreak F}, then to
\eqref{eq:gqt-denominator}. The equality between the exported
$\eta$-sum and the source sum is definitional. Proving that equality
again would not produce the missing estimates. The Galois cut
transport must preserve the original row rank and multiplicity
bound. The guards concerning the actual cohomology lengths remain.

For the retained target, GQT1--GQT5 remove the finite counting
and asymptotic tail from the reconstruction frontier. The
remaining work is to construct the actual labels and common
row object and prove the three required count estimates and two
bad-weight bounds. The full \texttt{Count\allowbreak All\allowbreak Rows} and
\texttt{Weight\allowbreak All\allowbreak Rows} producers remain sufficient; the weaker
contracts offer an alternative if they reduce the arithmetic work.

For periods and heights, a useful division is: actual rational
common target and exact Bessel map; all-embedding chain norms and
eigenvalue lower bounds; local divisor/arithmetic comparison;
the two finite-place and cross-word remainder estimates; and
the uniform height bound \eqref{eq:gqt-height1}. Each must refer
to the same $R,\allowbreak{} N_u,\allowbreak{} B_u$ and the same metrics. The existing
period-lattice and metric-transport arguments should be reused.
The source's \texttt{Remaining\allowbreak Division\allowbreak Bounds} is only part of
the full contract, not a synonym for \texttt{Link\allowbreak DIV}.

For local gain, the source reduces the global installation to
the same \texttt{GroupCompatibility\allowbreak{}Family} and admissible-top
condition, then actual local extremal-projector, lattice, flatness
and descent estimates. The \texttt{Loc} parameter supplies the
uniform \texttt{Second\allowbreak Theorem\allowbreak Reduced}; the local-model and
flatness constructors use further fixed formal theorems. Each
relevant result still needs a human proof or a checked published
application. The uniform subexponential
$b^{\rho(x)}$ loss and the original orbit representatives are
part of its output.

For the spectrum, retain separately the exact normalized period
expansion; the degenerate, nongeneric, loop and generic terms;
the Fourier/derivative and lattice-summation argument; and the
independent target-dimension estimate. The original source's
\texttt{Theorem\allowbreak SSelected}, \texttt{Loop\allowbreak Row\allowbreak S},
\texttt{Row\allowbreak NGP}, \texttt{Row\allowbreak Main\allowbreak Eq\allowbreak P}, \texttt{Type\allowbreak P},
\texttt{Wu\allowbreak Decay\allowbreak Of\allowbreak Gap}, \texttt{Rows\allowbreak Fin\allowbreak Dim}, and all-point count
are genuine input theorems. The earlier AG chapter already
proves the finite recurrence and a sufficient packing estimate.
GQT6--GQT7 lower the required final accuracy; they do not replace
these constructions by formal theorem names.

The new mathematics is therefore a complete conditional
retained-defect theorem, a smaller sufficient row-count contract,
and a proved weaker spectral input for the selected final
comparison. The common-family and arithmetic existence
obligations retain their previous status.

\endgroup
\end{document}
